% Generated by spec/build-formats.py from spec/xmessage-format.md. Do not edit; edit the markdown.
\documentclass[10pt,letterpaper]{article}
\usepackage[margin=0.8in]{geometry}
\usepackage{fontspec}
\setmainfont{DejaVu Serif}[Scale=0.92]
\setsansfont{DejaVu Sans}[Scale=0.92]
\setmonofont{DejaVu Sans Mono}[Scale=0.84]
\usepackage{array,longtable,booktabs}
\usepackage[shortlabels]{enumitem}
\setlist{itemsep=2pt,topsep=3pt,parsep=2pt}
\usepackage{fvextra}
\fvset{breaklines=true,breakanywhere=true,fontsize=\footnotesize,frame=leftline,framesep=6pt,
  xleftmargin=8pt,rulecolor=\color{black!35}}
\usepackage{xcolor}
\usepackage[framemethod=default]{mdframed}
\newmdenv[topline=false,bottomline=false,rightline=false,linewidth=2pt,linecolor=black!30,
  backgroundcolor=black!4,innerleftmargin=8pt,innertopmargin=4pt,innerbottommargin=4pt,
  skipabove=6pt,skipbelow=6pt]{notebox}
\usepackage[hidelinks,bookmarksnumbered=false,bookmarksopen=true,bookmarksopenlevel=1,
  pdftitle={ALAMO/ALTO X-Message Format Specification},pdfsubject={Generated from spec/xmessage-format.md}]{hyperref}
\newcommand{\code}[1]{\texttt{#1}}
\setlength{\parindent}{0pt}
\setlength{\parskip}{4pt plus 2pt}
\setcounter{secnumdepth}{0}
\setcounter{tocdepth}{3}
\setlength{\LTpre}{4pt}
\setlength{\LTpost}{6pt}
\renewcommand{\arraystretch}{1.15}
\makeatletter
\renewcommand\paragraph{\@startsection{paragraph}{4}{\z@}{2.5ex plus 1ex minus .2ex}{0.8ex}%
  {\normalfont\normalsize\bfseries\itshape}}
\makeatother
\emergencystretch=3em
\tolerance=1000
\title{ALAMO/ALTO X-Message Format Specification}
\author{XMTPS}
\date{Generated from \texttt{spec/xmessage-format.md} by \texttt{spec/build-formats.py}}
\begin{document}
\maketitle

This document specifies the binary X-message format that MRV ALAMO and ALTO profiling floats use to send their data over Iridium Short Burst Data (SBD). It also says how XMTPS decodes each message and which listing-file lines each message becomes.

Status: compiled on 2026-10-25. It brings together everything known about the format:

\begin{itemize}
\item the SOLO-II X-format specification it derives from;
\item the ALAMO X-format spreadsheet (v2.4);
\item MRV's firmware message definitions (7.3.5, 10.5.0, 10.12.0);
\item the XMTPS decoder source;
\item hand-decoded production packets.
\end{itemize}

No single earlier document was complete or current, and they disagree in places. Where they disagree, this document says so, and gives the layout that the decoder and real packets agree on.

Decoder fixes in build 3309 (2026-10-25). Writing this document turned up several decoder bugs, which build 3309 fixes. The parts below describe these bugs as found and mark each one fixed:

\begin{itemize}
\item Argo \code{n\allowbreak{}Seek} is read as two bytes.
\item \code{APS} SNROut shows the output SNR.
\item Mission \code{DRT} and \code{PPC0}/\code{PPC1} list their own values.
\item 0xEA versions 0x13/0x14 read raw samples before empty bins.
\item The 0xEA version 0x10 error heading formats.
\item Mission and BIT blobs too short for their layout are rejected.
\end{itemize}

\tableofcontents
\clearpage

\section{\texorpdfstring{How to read this document}{How to read this document}}\label{sec:how-to-read-this-document}

\textbf{Structure.} A float's data for one dive cycle is split into \emph{packets}. Each packet travels as one Iridium SBD message, a binary e-mail attachment of at most 340 bytes. A packet has the form \code{X nn mm dd p <data> \$cc>}, and its \code{<data>} is a sequence of \emph{blobs} (the SOLO-II documents call them “sensor messages”). Each blob has the form \code{ID jj <data> ;\allowbreak{}}. The 1-byte ID says what the blob holds; \code{jj} is the blob length. Part A.1 gives the details.

\textbf{Offsets} in every table count from the blob's ID byte, which is offset 0. All multi-byte fields are big-endian, including the IEEE-754 floats.

\textbf{Times} are seconds since 2000-01-01 00:00 UTC unless a table says UNIX time (seconds since 1970-01-01).

\textbf{Sensor values} sent as counts are converted with the gains and offsets in the same dive's Argo message 0xF0: \code{value =\allowbreak{} counts /\allowbreak{} gain − offset} (A.1.5).

\textbf{Source tags.} The parts mark where each statement comes from:

\begingroup\footnotesize\setlength{\tabcolsep}{3pt}
\begin{longtable}{@{}>{\raggedright\arraybackslash}p{0.1161\linewidth}>{\raggedright\arraybackslash}p{0.2232\linewidth}>{\raggedright\arraybackslash}p{0.6161\linewidth}@{}}
\toprule
\textbf{Tag (part A)} & \textbf{Tag (parts B and C)} & \textbf{Source} \\
\midrule
\endfirsthead
\toprule
\textbf{Tag (part A)} & \textbf{Tag (parts B and C)} & \textbf{Source} \\
\midrule
\endhead
\bottomrule
\endlastfoot
{}[D] & {}[decoder] & The XMTPS decoder source (\code{src/\allowbreak{}data/\allowbreak{}xmsg/\allowbreak{}}, \code{src/\allowbreak{}data/\allowbreak{}XFloat.\allowbreak{}cpp}, \code{src/\allowbreak{}data/\allowbreak{}XFloat\allowbreak{}Mail\allowbreak{}Message.\allowbreak{}*}) as of 2026-10-25. It is the reference for what is implemented. \\
{}[O] & Production & Production packets, decoded by hand from the hex dumps (\code{SHX}/\code{shx} lines) in XMTPS listings. \\
{}[S] & SOLO2 & \emph{Format of SOLOII X messages}, Argo version manual/decoder v2.0, 31 Mar 2014 (\code{SOLO2\_\allowbreak{}Xformat\_\allowbreak{}v2.\allowbreak{}0\_\allowbreak{}20Mar2014.\allowbreak{}pdf}). \\
{}[X] & XF2.4 & ALAMO X-format specification v2.4 spreadsheet (\code{xfmt-format-spec-2.\allowbreak{}4-alamo\_\allowbreak{}2018-05-15.\allowbreak{}xlsx}). \\
{}[M] & NG10.5, NG10.12, DF7.3.5 & MRV firmware definitions: \code{nextgen\_\allowbreak{}v10.\allowbreak{}5.\allowbreak{}0\_\allowbreak{}blobs.\allowbreak{}txt}, \code{nextgen-10.\allowbreak{}12.\allowbreak{}0-xmessage.\allowbreak{}txt}, \code{Data\allowbreak{}Format-v7.\allowbreak{}3.\allowbreak{}5.\allowbreak{}yaml}, plus \code{xfs.\allowbreak{}txt} (on-float storage) and \code{system-error-codes.\allowbreak{}txt}. \\
{}[I] & {}[inferred] & Not documented anywhere. Worked out from production data or decoder behavior. \\
\end{longtable}
\endgroup

The MRV documents and the SOLO-II PDF are in the Float Group's \code{mrv/\allowbreak{}doc} folder. The ALAMO spreadsheets, v2.2 to v2.4 (2015–2018), are in this repository's \code{doc/\allowbreak{}} folder, which git doesn't track.

\textbf{Firmware versions} are written as 32-bit words \code{0x\allowbreak{}MMmmppbb}: major, minor, patch, then build/product byte. For example, \code{0x0A0C0001} is firmware 10.12.0.1 and \code{0x0A0C0032} is 10.12.0.32. In the last byte, 0x31 marks WIMS (wave) floats and 0x32/0x33 mark MM2/Micro-ALTO floats (C.1.4).


\section{\texorpdfstring{Message type index}{Message type index}}\label{sec:message-type-index}

\begingroup\footnotesize\setlength{\tabcolsep}{3pt}
\begin{longtable}{@{}>{\raggedright\arraybackslash}p{0.0893\linewidth}>{\raggedright\arraybackslash}p{0.2946\linewidth}>{\raggedright\arraybackslash}p{0.0982\linewidth}>{\raggedright\arraybackslash}p{0.2411\linewidth}>{\raggedright\arraybackslash}p{0.2054\linewidth}@{}}
\toprule
\textbf{ID} & \textbf{Message} & \textbf{Section} & \textbf{Decoded by XMTPS} & \textbf{Listing lines} \\
\midrule
\endfirsthead
\toprule
\textbf{ID} & \textbf{Message} & \textbf{Section} & \textbf{Decoded by XMTPS} & \textbf{Listing lines} \\
\midrule
\endhead
\bottomrule
\endlastfoot
0x00 & GPS fix after the diagnostic dive & A.2 & yes & \code{GPS} \\
0x01 & GPS fix at the start of a dive & A.2 & yes & \code{GPS} \\
0x02 & GPS fix at the end of a dive & A.2 & yes & \code{GPS} \\
0x03 & GPS fix after a mission abort & A.2 & yes & \code{GPS} \\
0x05 & GPS fix during the built-in test & A.2 & yes & \code{GPS} \\
0x10–0x1F & Binned profile pressure & A.3 & yes & \code{PRO}/\code{pro}, \code{PRD}/\code{prd} \\
0x20–0x2F & Binned profile temperature (format 3: auxiliary temperature channel) & A.3, A.4.3 & yes & \code{PRO}/\code{pro}, \code{PRD}/\code{prd} \\
0x30–0x3F & Binned profile salinity & A.3 & yes & \code{PRO}/\code{pro}, \code{PRD}/\code{prd} \\
0x40–0x4F & Fall (descent) pressure time series & A.5.1 & yes & \code{F}/\code{f} \\
0x50–0x5F & Rise (ascent) pressure time series & A.5.1 & yes & \code{R}/\code{r} \\
0x60–0x6F & Pump operations & A.5.3 & yes & \code{PMP}/\code{pmp} \\
0x78–0x7F & Drift sample times & A.4.1 & yes & \code{D}/\code{d} \\
0x80–0x8F & Surface time series & A.5.1 & yes (formats 0–2) & \code{SFC}/\code{sfc} \\
0x90–0x97 & High-resolution profile pressure & A.3 & yes & \code{PRU}/\code{pru} \\
0x98–0x9F & Drift pressure & A.3 & yes & \code{D}/\code{d} \\
0xA0–0xA7 & High-resolution profile temperature & A.3 & yes & \code{PRU}/\code{pru} \\
0xA8–0xAF & Drift temperature & A.3 & yes & \code{D}/\code{d} \\
0xB0–0xB7 & High-resolution profile salinity & A.3 & yes & \code{PRU}/\code{pru} \\
0xB8–0xBF & Drift salinity & A.3 & yes & \code{D}/\code{d} \\
0xC1 & Ocean surface waves: auto-spectra & C.3 & yes (spectral v0–2) & \code{WVS}, \code{wvs}, \code{wvm} \\
0xC2 & Auxiliary sensor data & A.4.4 & no & — \\
0xC3 & Ocean surface waves: cross-spectra & C.3 & yes (spectral v0–2) & \code{WVS}, \code{wvs}, \code{wvm} \\
0xC5 & WIMS configuration parameters & C.7 & yes (v1, v2; inferred) & \code{WIP} \\
0xCA & Marginal-ice-zone acoustics (0xC0 before build 3249) & C.4 & subtypes 3 and 5 & \code{APS} \\
0xD0–0xD1 & Mission configuration (binary, two parts) & C.1 & yes (117 firmware versions) & \code{MIS}, \code{m}, \code{mis} \\
0xD2–0xD7 & Mission configuration, further parts / SOLO-II ASCII dump & C.1 & no & — \\
0xDC & Dynamically corrected profile salinity & A.4.2 & yes & \code{saln\allowbreak{}Corr} column of \code{pro} \\
0xDE & Mission command echo & C.2 & yes & \code{CMD} \\
0xE0 & Engineering, diagnostic dive (SOLO-II) & B.1.1 & validated only & — \\
0xE2 & Engineering, normal dive (SOLO-II) & B.1.2 & validated only & — \\
0xE3 & Engineering, mission abort (SOLO-II) & B.1.3 & validated only & — \\
0xE5 & Built-in test passed & B.2 & formats 6–11 & \code{E05} \\
0xE6 & Built-in test failed & B.2 & formats 6–11 & \code{E06} \\
0xE7 & Engineering, descent & B.3 & versions 0x10–0x11 & \code{E07} \\
0xE8 & Engineering, drift/park & B.4 & versions 0x10–0x12 & \code{E08} \\
0xE9 & Engineering, ascent & B.5 & versions 0x10–0x12 & \code{E09} \\
0xEA & Engineering, surface & B.6 & versions 0x10–0x14 & \code{E0A}, \code{e0A} \\
0xEB & Valve operations & A.5.4 & yes & \code{VLV}/\code{vlv} \\
0xEC & Engineering, ice detection & B.7 & versions 0x10–0x11 & \code{E0C}, \code{e0c} \\
0xF0 & Argo data (mission summary, gains and offsets) & A.6 & v2.1–v2.4 & \code{A} \\
0xF1 & Test pattern & — & no & — \\
0xFF & Extended records (event, reset, serial number, CTD exchange, fault dump) & C.6 & since build 3321 (format inferred) & \code{XFF} \\
\end{longtable}
\endgroup

IDs not listed are unassigned. In types 0x1n–0x8n and 0x9n–0xBn, the low nibble of the ID is the part number of a multi-part series (A.3).

\clearpage
\section{\texorpdfstring{Part A — Transport, framing and data messages}{Part A — Transport, framing and data messages}}\label{sec:part-a--transport-framing-and-data-messages}

Scope of this part: packet framing and general conventions; GPS (0x00–0x05); binned, high-resolution and drift P/T/S (0x1n/0x2n/0x3n, 0x90–0xBF); drift sample times (0x78–0x7F); dynamically corrected salinity (0xDC); auxiliary sensor data (0xC2) and the auxiliary temperature channel; fall (0x4n), rise (0x5n), pump (0x6n), surface time series (0x8n) and valve (0xEB); Argo data (0xF0).

\subsection{\texorpdfstring{A.0 Sources and conventions of this document}{A.0 Sources and conventions of this document}}\label{sec:a0-sources-and-conventions-of-this-document}

\begingroup\footnotesize\setlength{\tabcolsep}{3pt}
\begin{longtable}{@{}>{\raggedright\arraybackslash}p{0.0357\linewidth}>{\raggedright\arraybackslash}p{0.9375\linewidth}@{}}
\toprule
\textbf{Tag} & \textbf{Source} \\
\midrule
\endfirsthead
\toprule
\textbf{Tag} & \textbf{Source} \\
\midrule
\endhead
\bottomrule
\endlastfoot
\textbf{[D]} & XMTPS decoder source (\code{src/\allowbreak{}data/\allowbreak{}xmsg/\allowbreak{}*}, \code{src/\allowbreak{}data/\allowbreak{}XFloat.\allowbreak{}cpp}, \code{src/\allowbreak{}data/\allowbreak{}XFloat\allowbreak{}Mail\allowbreak{}Message.\allowbreak{}*}). Ground truth for what is implemented. \\
\textbf{[O]} & Observed in production packets (hex dumps in \code{lis9287-0325.\allowbreak{}txt} and \code{lis11186-0107.\allowbreak{}txt}, firmware 10.12.x, decoded by hand for this document). \\
\textbf{[S]} & SOLO-II X format v2.0 (\code{SOLO2\_\allowbreak{}Xformat\_\allowbreak{}v2.\allowbreak{}0\_\allowbreak{}20Mar2014.\allowbreak{}pdf}). \\
\textbf{[X]} & ALAMO X-format spec v2.4 spreadsheet (\code{xfmt-format-spec-2.\allowbreak{}4-alamo\_\allowbreak{}2018-05-15.\allowbreak{}xlsx}). \\
\textbf{[M]} & MRV documents: \code{nextgen\_\allowbreak{}v10.\allowbreak{}5.\allowbreak{}0\_\allowbreak{}blobs.\allowbreak{}txt}, \code{Data\allowbreak{}Format-v7.\allowbreak{}3.\allowbreak{}5.\allowbreak{}yaml}, \code{xfs.\allowbreak{}txt}. \\
\textbf{[I]} & Inferred (not documented anywhere; reconstructed from data or code). \\
\end{longtable}
\endgroup

\begin{itemize}
\item Offsets are byte offsets counted from the blob's ID byte = 0, unless the table says otherwise.
\item \code{u8/\allowbreak{}u16/\allowbreak{}u24/\allowbreak{}u32} are unsigned and \code{i8/\allowbreak{}i16/\allowbreak{}i24/\allowbreak{}i32} are two's-complement signed big-endian integers. \code{u4} / \code{u20} denote nibble-aligned fields that share a byte with a neighbor (see the Notes column).
\item “counts” are raw integers converted with a gain and offset from the 0xF0 Argo message (section A.1.5).
\item Listing tags (\code{GPS}, \code{PRO}/\code{pro}, \code{F}/\code{f}, …) refer to the XMTPS listing file (README, “Listing files”).
\end{itemize}

\subsection{\texorpdfstring{A.1 Packet framing and general conventions}{A.1 Packet framing and general conventions}}\label{sec:a1-packet-framing-and-general-conventions}

\subsubsection{\texorpdfstring{A.1.1 Transport: Iridium SBD and the mail gateway}{A.1.1 Transport: Iridium SBD and the mail gateway}}\label{sec:a11-transport-iridium-sbd-and-the-mail-gateway}

\begin{itemize}
\item The float stores telemetry in an on-board “XMessage file system” (XFS) of 1500 blocks [M \code{xfs.\allowbreak{}txt}]. \textbf{One XFS block = one Iridium SBD mobile-originated message = one X packet.}
\item An SBD message carries at most \textbf{340 bytes}. The X packet header (8 bytes: \code{X nn mm dd p}) and trailer (4 bytes: \code{\$cc>}) take 12, so a block carries at most \textbf{328 bytes of blobs} [M]. (SOLO-II documentation still refers to the 189-byte 9601 limit [S]; it doesn't apply to ALAMO/ALTO.)
\item A block holds one or more complete blobs, end to end with no gaps. All blobs in a block belong to the same dive. \textbf{A blob never spans two packets.} Multi-part data (profiles, time series) is split into separate blobs that carry their own part index [M][D].
\item The Iridium gateway forwards each SBD message as an e-mail [D \code{XFloat\allowbreak{}Mail\allowbreak{}Message.\allowbreak{}h}]:

\begin{itemize}
\item Sender matches \code{(\allowbreak{}sbdservice@sbd.\allowbreak{}iridium.\allowbreak{}com|\allowbreak{}service@sbd.\allowbreak{}pac.\allowbreak{}disa.\allowbreak{}mil)}. This is configurable.
\item Subject matches \code{SBD Msg From Unit:\allowbreak{} (\allowbreak{}\textbackslash{}\allowbreak{}d+)}, which gives the IMEI.
\item The binary attachment (MIME type \code{application/\allowbreak{}*}) is named \code{<IMEI>\_\allowbreak{}<MOMSN,\allowbreak{} 6 digits>.\allowbreak{}sbd} (regex \code{(\allowbreak{}\textbackslash{}\allowbreak{}d+)\_\allowbreak{}(\allowbreak{}\textbackslash{}\allowbreak{}d\{6\})\textbackslash{}\allowbreak{}.\allowbreak{}sbd}). Its bytes are the X packet.
\item The body text gives \code{MOMSN:\allowbreak{}}, \code{MTMSN:\allowbreak{}}, \code{Time of Session (\allowbreak{}UTC):\allowbreak{}}, \code{Session Status:\allowbreak{} NN - text}, \code{Message Size (\allowbreak{}bytes):\allowbreak{}}, \code{Unit Location:\allowbreak{} Lat =\allowbreak{} … Long =\allowbreak{} …} and \code{CEPradius =\allowbreak{} N} (km).
\end{itemize}
\item The decoder accepts the attachment only if it has at least \textbf{8 bytes} (\code{XMSG\_\allowbreak{}SBD\_\allowbreak{}MIN\_\allowbreak{}BYTES}) and the session status is \textbf{00 or 02} (transfer OK; 02 = OK but the location is of unacceptable quality) [D].
\item Duplicates [D]:

\begin{itemize}
\item A mail with an already-seen Message-ID is ignored.
\item A mail whose SBD bytes are identical to an already-assigned packet of the dive is a retransmission under a new MOMSN. It is kept for its Iridium position, but its blobs are not decoded again.
\item Dive data is keyed by (serial number, dive number) taken from the packet header, not from the e-mail.
\end{itemize}
\item Transmission order is not packet-ID order. In the observed dives the highest \code{p} was sent first and \code{p =\allowbreak{} 0} last [O].
\end{itemize}

\subsubsection{\texorpdfstring{A.1.2 X packet: \code{X nn mm dd p <data> \$ c c >}}{A.1.2 X packet: X nn mm dd p <data> \$ c c >}}\label{sec:a12-x-packet-x-nn-mm-dd-p-data--c-c-}

\begingroup\footnotesize\setlength{\tabcolsep}{3pt}
\begin{longtable}{@{}>{\raggedright\arraybackslash}p{0.0625\linewidth}>{\raggedright\arraybackslash}p{0.0625\linewidth}>{\raggedright\arraybackslash}p{0.0625\linewidth}>{\raggedright\arraybackslash}p{0.0804\linewidth}>{\raggedright\arraybackslash}p{0.1071\linewidth}>{\raggedright\arraybackslash}p{0.5446\linewidth}@{}}
\toprule
\textbf{Offset} & \textbf{Size} & \textbf{Type} & \textbf{Field} & \textbf{Scale/units} & \textbf{Notes} \\
\midrule
\endfirsthead
\toprule
\textbf{Offset} & \textbf{Size} & \textbf{Type} & \textbf{Field} & \textbf{Scale/units} & \textbf{Notes} \\
\midrule
\endhead
\bottomrule
\endlastfoot
0 & 1 & char & \code{X} (0x58) & — & Start marker. Not checked by the decoder [D]. \\
1 & 2 & u16 & \code{nn} & bytes & Number of bytes after \code{nn} up to, but not including, \code{\$} = 5 + len(\code{<data>}) [S]. Verified: \code{nn} = packet length − 7 on 27 packets [O]. Ignored by the decoder, which uses the attachment length [D]. \\
3 & 2 & u16 & \code{mm} & — & Float serial number (WHOI/instrument ID) [S][D]. \\
5 & 2 & i16 & \code{dd} & — & Dive (cycle) number: −1 = start-up, 0 = test/diagnostic dive, then +1 per dive [S]. Signed [D]. \\
7 & 1 & u8 & \code{p} & — & Packet ID within the dive. The data of each dive starts at \code{p =\allowbreak{} 0} [S]. Listed as \code{PID} on \code{SBD} lines. \\
8 & len−12 & — & \code{<data>} & — & Zero or more blobs (A.1.3). \\
len−4 & 1 & char & \code{\$} (0x24) & — & Checksum delimiter. \\
len−3 & 2 & 2×char & \code{cc} & — & 8-bit checksum (see below). \\
len−1 & 1 & char & \code{>} (0x3E) & — & End marker. \\
\end{longtable}
\endgroup

\textbf{Checksum.} \code{sum =\allowbreak{} (\allowbreak{}Σ bytes from 'X' (\allowbreak{}offset 0) through the byte before '\$') mod 256}. It is written as two characters, high nibble first, each nibble + 0x30. So the characters are in the range \code{0}–\code{?} (\code{0x3A}–\code{0x3F} = \code{:\allowbreak{};\allowbreak{}<=\allowbreak{}>?} stand for nibble values 10–15) [S]. All 1,351,052 X packets in the production RAW files (all floats and firmware, 2014–2026) have a correct checksum [O].

\begin{notebox}
\textbf{Decoder note [D]:} since build 3310 the decoder verifies the checksum (\code{XFloat\allowbreak{}Dive\allowbreak{}Data:\allowbreak{}:\allowbreak{}Decode\allowbreak{}XMsg\allowbreak{}Data}). A packet whose checksum does not match is discarded and logged (“X-Message checksum mismatch”); a packet without the \code{\$cc>} trailer is decoded with a warning. The blob parser reads offsets 8 … len−5.

\end{notebox}

\subsubsection{\texorpdfstring{A.1.3 Blob structure: \code{ID jj <data> ;\allowbreak{}}}{A.1.3 Blob structure: ID jj <data> ;}}\label{sec:a13-blob-structure-id-jj-data-}

\begingroup\footnotesize\setlength{\tabcolsep}{3pt}
\begin{longtable}{@{}>{\raggedright\arraybackslash}p{0.0625\linewidth}>{\raggedright\arraybackslash}p{0.0446\linewidth}>{\raggedright\arraybackslash}p{0.0446\linewidth}>{\raggedright\arraybackslash}p{0.0804\linewidth}>{\raggedright\arraybackslash}p{0.6964\linewidth}@{}}
\toprule
\textbf{Offset} & \textbf{Size} & \textbf{Type} & \textbf{Field} & \textbf{Notes} \\
\midrule
\endfirsthead
\toprule
\textbf{Offset} & \textbf{Size} & \textbf{Type} & \textbf{Field} & \textbf{Notes} \\
\midrule
\endhead
\bottomrule
\endlastfoot
0 & 1 & u8 & ID & High nibble = message class; low nibble = part index or sub-type (e.g. 0x1n = binned pressure, part n). For classes with a single type (0xD0, 0xDE, 0xE0–0xEC, 0xF0, 0xF1) the whole byte is the type. \\
1–2 & 2 & — & \code{jj} & Blob length in bytes, \textbf{including} ID, \code{jj} and the trailing \code{;\allowbreak{}}. Two layouts, see below. \\
3 & jj−4 & — & data & Type-specific. \\
jj−1 & 1 & char & \code{;\allowbreak{}} (0x3B) & Terminator. \\
\end{longtable}
\endgroup

\textbf{\code{jj} layouts:}

\begin{itemize}
\item \textbf{Encoding nibble + 12-bit size}, used by the sensor/time-series families: 0x1n–0x3n, 0x4n, 0x5n, 0x6n, 0x78–0x7F, 0x8n, 0x9n, 0xAn, 0xBn and 0xDC. Bits 15–12 of \code{jj} (high nibble of byte 1) hold the packing/encoding format; bits 11–0 hold the size (0–4095) [S][X].
\item \textbf{Full 16-bit size}, used by GPS 0x00–0x05, 0xC1, 0xD0–0xD7, 0xDE, 0xE5–0xEC (including the valve message 0xEB) and 0xF0 [X]. Here the encoding format, if any, is in byte 3 (e.g. 0xEB).
\item Exception: SOLO-legacy pump format 0 (0x6n with format nibble 0) is read with the full 16-bit size by \code{XPump} [D], which is equivalent because the high nibble is then 0.
\end{itemize}

\textbf{How the decoder splits a packet into blobs [D] (\code{XFloat\allowbreak{}Dive\allowbreak{}Data:\allowbreak{}:\allowbreak{}Decode\allowbreak{}XMsg\allowbreak{}Data}):}

\begin{enumerate}
\item Bytes 8 … len−5 are appended one at a time to a buffer.
\item At 3 buffered bytes the size is taken as \code{0x0FFF \& u16(\allowbreak{}buf,\allowbreak{} 1)}. \textbf{The 12-bit mask is applied to every type.} This is harmless because a blob is always smaller than 4096 bytes (SBD ≤ 340), so the high nibble of a 16-bit size is always 0.
\item When the buffer reaches the size, the blob is dispatched on its ID byte and the buffer is cleared. A trailing incomplete blob is silently dropped. A wrong size desynchronizes the rest of the packet.
\item Each class decoder checks the ID, a minimum length, and that \code{buf[size−1] =\allowbreak{}=\allowbreak{} ';\allowbreak{}'} (most of them via \code{XMessage:\allowbreak{}:\allowbreak{}Check\allowbreak{}XMsg\allowbreak{}Data}; the size must be at least 4). Extra bytes produce a warning; a missing terminator is an error.
\item IDs without a decoder are logged as \code{Unknown X message ID:\allowbreak{} 0x.\allowbreak{}.\allowbreak{}} and skipped, for example 0xC2 and 0x70–0x77 (and 0xFF before build 3321, C.6). 0xF1 (test pattern) is ignored silently.
\end{enumerate}

\subsubsection{\texorpdfstring{A.1.4 Byte order, bit order, epochs}{A.1.4 Byte order, bit order, epochs}}\label{sec:a14-byte-order-bit-order-epochs}

\begin{itemize}
\item \textbf{Big-endian} everywhere: multi-byte integers are sent MSB first [S][D].
\item Bit fields are numbered from the MSB of the first byte. In a byte that holds two nibbles, the \textbf{high nibble comes first}. For example, a “code + 20-bit pressure” field is \code{code =\allowbreak{} b[k] >> 4}, \code{pres =\allowbreak{} (\allowbreak{}(\allowbreak{}b[k] \& 0x0F) << 16) |\allowbreak{} b[k+1] << 8 |\allowbreak{} b[k+2]}.
\item \textbf{Time epochs:}

\begin{itemize}
\item \textbf{2000-01-01T00:00:00Z} (\code{s2000}): u32 seconds since this epoch for the 0xF0 state start times and the fall/rise/surface start/reference times. \textbf{0 means “not set”}: the decoder maps it to not-a-date-time, listed as \code{0000-00-00T00:\allowbreak{}00:\allowbreak{}00Z} [D]. The spreadsheet's “offset 36526” on these fields is the Excel day number of 2000-01-01 [X].
\item \textbf{GPS epoch 1980-01-06}: GPS week and day in the GPS message (A.2).
\item \textbf{UNIX epoch 1970-01-01} (\code{s1970}): used only by the BIT v2 engineering message (\code{XAlamo\allowbreak{}Bit2.\allowbreak{}h}, outside this part) [D].
\item Relative times: u16 seconds after a reference (fall/rise/surface), and u24/i24 seconds after the drift start (0x78).
\end{itemize}
\item Leap seconds are not applied anywhere [D].
\end{itemize}

\subsubsection{\texorpdfstring{A.1.5 Counts-to-units conversion}{A.1.5 Counts-to-units conversion}}\label{sec:a15-counts-to-units-conversion}

All CTD quantities (pressure, temperature, salinity) are sent as integer counts:

\begin{Verbatim}
value = counts / gain − offset
\end{Verbatim}

\code{gain} and \code{offset} are unsigned 16-bit integers from the 0xF0 Argo message of the \textbf{same dive} [D][X][M]. The SOLO-II document gives \code{counts · gain − offset} in one place [S]; that is a typo, and both its own 0xF0 table and the decoder use division.

\begingroup\footnotesize\setlength{\tabcolsep}{3pt}
\begin{longtable}{@{}>{\raggedright\arraybackslash}p{0.4911\linewidth}>{\raggedright\arraybackslash}p{0.4821\linewidth}@{}}
\toprule
\textbf{Data} & \textbf{Gains/offsets used [D]} \\
\midrule
\endfirsthead
\toprule
\textbf{Data} & \textbf{Gains/offsets used [D]} \\
\midrule
\endhead
\bottomrule
\endlastfoot
Binned/hi-res/drift profiles, 0xDC, aux temperature, fall, rise, pump, valve & pressure: \code{p\allowbreak{}Gain/\allowbreak{}p\allowbreak{}Off} (0xF0 bytes 40–43); temperature: \code{t\allowbreak{}Gain/\allowbreak{}t\allowbreak{}Off} (44–47); salinity: \code{s\allowbreak{}Gain/\allowbreak{}s\allowbreak{}Off} (48–51). \\
Surface time series 0x8n & surface gains/offsets (0xF0 bytes 52–63). In v2.1 these are copies of the profile gains. \\
\end{longtable}
\endgroup

Typical values [O]:

\begingroup\footnotesize\setlength{\tabcolsep}{3pt}
\begin{longtable}{@{}>{\raggedright\arraybackslash}p{0.0893\linewidth}>{\raggedright\arraybackslash}p{0.0893\linewidth}>{\raggedright\arraybackslash}p{0.0625\linewidth}>{\raggedright\arraybackslash}p{0.0982\linewidth}>{\raggedright\arraybackslash}p{0.3571\linewidth}@{}}
\toprule
\textbf{Quantity} & \textbf{Gain} & \textbf{Offset} & \textbf{Resolution} & \textbf{Note} \\
\midrule
\endfirsthead
\toprule
\textbf{Quantity} & \textbf{Gain} & \textbf{Offset} & \textbf{Resolution} & \textbf{Note} \\
\midrule
\endhead
\bottomrule
\endlastfoot
P & 25 & 10 & 0.04 dBar & 0–65535 counts covers −10 … 2611.4 dBar \\
T & 1000 & 5 & 0.001 °C &  \\
S & 1000 & 1 & 0.001 PSU &  \\
Surface P & 25 or 100 & 10 &  &  \\
\end{longtable}
\endgroup

\begin{itemize}
\item \textbf{No 0xF0 in the dive} [D]: all gains and offsets are set to 9999 (\code{IRI\_\allowbreak{}GAIN\_\allowbreak{}INVALID} / \code{IRI\_\allowbreak{}OFFS\_\allowbreak{}INVALID}), so every converted value falls out of range and is listed as the fill value. The float-level fallback is commented out. A dive's data can only be converted once its own 0xF0 has arrived.
\item \textbf{gain = 0} gives the fill value (the per-message \code{Sample\allowbreak{}Value} returns −100 internally).
\item \textbf{Signedness [D]:} composite accessors pass gain/offset as \code{int16\_\allowbreak{}t}, so a gain above 32767 would be misread. No observed float uses one.
\item \textbf{Valid ranges [D]:} values outside these limits are replaced by the fill value (\code{XDefinitions.\allowbreak{}h}):

\begingroup\footnotesize\setlength{\tabcolsep}{3pt}
\begin{longtable}{@{}>{\raggedright\arraybackslash}p{0.0804\linewidth}>{\raggedright\arraybackslash}p{0.1339\linewidth}@{}}
\toprule
\textbf{Quantity} & \textbf{Valid range} \\
\midrule
\endfirsthead
\toprule
\textbf{Quantity} & \textbf{Valid range} \\
\midrule
\endhead
\bottomrule
\endlastfoot
P & −5 … 2610 dBar \\
T & −5 … 40 °C \\
S & 2 … 41 PSU \\
\end{longtable}
\endgroup
\end{itemize}

\begin{itemize}
\item \textbf{Pressure offset correction [D]:} for RBR floats (or WMO instrument types 866/867/875) XMTPS can subtract a surface-pressure offset \code{pc}, derived from the surface time series of recent dives, from profile, drift, fall, rise, pump and valve pressures. Surface (\code{sfc}) pressures are always listed uncorrected. See README “Surface pressure offset correction”.
\end{itemize}

\subsubsection{\texorpdfstring{A.1.6 Invalid and fill values}{A.1.6 Invalid and fill values}}\label{sec:a16-invalid-and-fill-values}

\begingroup\footnotesize\setlength{\tabcolsep}{3pt}
\begin{longtable}{@{}>{\raggedright\arraybackslash}p{0.2679\linewidth}>{\raggedright\arraybackslash}p{0.2679\linewidth}>{\raggedright\arraybackslash}p{0.4196\linewidth}@{}}
\toprule
\textbf{Context} & \textbf{Encoding on the wire} & \textbf{Decoder representation / listing} \\
\midrule
\endfirsthead
\toprule
\textbf{Context} & \textbf{Encoding on the wire} & \textbf{Decoder representation / listing} \\
\midrule
\endhead
\bottomrule
\endlastfoot
Fall/rise/pump depth (SOLO) & 0xFFFF = invalid pressure [S] & No special case; converted, then range-checked [D]. \\
GPS fix & byte 3 = 0 & lat 99.99999, lon 999.99999, time unset [D]. \\
s2000 time & 0 & not-a-date-time → \code{0000-00-00T00:\allowbreak{}00:\allowbreak{}00Z}. \\
CTD values in listings & — & \code{-9.\allowbreak{}999} (pressure printed with 2 decimals) for out-of-range or missing values. \\
Pump/valve vacuum not sent & — & 65535 cinHg → \code{655.\allowbreak{}35} in \code{vlv}/\code{pmp} vac0 columns [D]. \\
Integer fields never received & — & 255 / 65535 / INT\_\allowbreak{}MIN (\code{XMSG\_\allowbreak{}BAD\_\allowbreak{}*}); dive −999, IDs −1. \\
\end{longtable}
\endgroup

\subsection{\texorpdfstring{A.2 GPS position — 0x00–0x05}{A.2 GPS position — 0x00–0x05}}\label{sec:a2-gps-position--0x000x05}

\textbf{Purpose.} A GPS fix with its quality metrics. The low nibble of the ID is the mission phase in which the fix was taken. There is one layout for all phases.

\begingroup\footnotesize\setlength{\tabcolsep}{3pt}
\begin{longtable}{@{}>{\raggedright\arraybackslash}p{0.1429\linewidth}>{\raggedright\arraybackslash}p{0.8304\linewidth}@{}}
\toprule
\textbf{Low nibble (ID)} & \textbf{Meaning [S][D]} \\
\midrule
\endfirsthead
\toprule
\textbf{Low nibble (ID)} & \textbf{Meaning [S][D]} \\
\midrule
\endhead
\bottomrule
\endlastfoot
0x00 & Diagnostic dive: fix after the first surfacing following the diagnostic (test) dive. \\
0x01 & Start of dive: just before leaving the surface. \\
0x02 & End of dive: at the surface after the ascent/profile. The best position/time for this dive's profile. \\
0x03 & After a mission abort. \\
0x05 & During the built-in self-test (BIT). \\
0x04, 0x06–0x0F & Unused. The decoder rejects IDs > 0x05. \\
\end{longtable}
\endgroup

\begin{notebox}
\textbf{Discrepancy [X]:} the spreadsheet's overview sheet labels 0x01 “end-of-dive” and 0x02 “start-of-dive”. Its GPS sheet, the SOLO-II document and the decoder all use 0x01 = start, 0x02 = end.

\end{notebox}

\begingroup\footnotesize\setlength{\tabcolsep}{3pt}
\begin{longtable}{@{}>{\raggedright\arraybackslash}p{0.0625\linewidth}>{\raggedright\arraybackslash}p{0.0446\linewidth}>{\raggedright\arraybackslash}p{0.0446\linewidth}>{\raggedright\arraybackslash}p{0.2143\linewidth}>{\raggedright\arraybackslash}p{0.1071\linewidth}>{\raggedright\arraybackslash}p{0.4464\linewidth}@{}}
\toprule
\textbf{Offset} & \textbf{Size} & \textbf{Type} & \textbf{Field} & \textbf{Scale/units} & \textbf{Notes} \\
\midrule
\endfirsthead
\toprule
\textbf{Offset} & \textbf{Size} & \textbf{Type} & \textbf{Field} & \textbf{Scale/units} & \textbf{Notes} \\
\midrule
\endhead
\bottomrule
\endlastfoot
0 & 1 & u8 & ID & — & 0x00–0x05. \\
1 & 2 & u16 & size & bytes & 24 (0x0018). Full 16-bit size. \\
3 & 1 & i8 & fix status / hemisphere & — & 0 = no/invalid fix; +2 = longitude East; −2 (0xFE) = West [S][X]. The decoder only tests for 0; the sign is also carried by the i32 longitude. \\
4 & 4 & i32 & latitude & 1e-7 deg & Positive = North. \\
8 & 4 & i32 & longitude & 1e-7 deg & Signed, −180 … +180, positive = East. \\
12 & 2 & u16 & GPS week & weeks & Full week number since 1980-01-06: “traditional” 10-bit week in the low 10 bits, rollover count in the high 6 bits [S], so the u16 is the continuous week number. \\
14 & 1 & u8 & GPS day of week & — & 0 = Sunday … 6 = Saturday. \\
15 & 1 & u8 & hour & UTC &  \\
16 & 1 & u8 & minute & UTC & Seconds are not sent; the decoder uses 0. \\
17 & 1 & u8 & time to fix & 10 s & Decoded as \code{u8 × 10} seconds (0–2550 s). \\
18 & 1 & u8 & satellites used & — & The decoder keeps only the low nibble (\code{\& 0x0F}). \\
19 & 1 & u8 & SNR min & — &  \\
20 & 1 & u8 & SNR mean & — &  \\
21 & 1 & u8 & SNR max & — & Observed as 0/0/0 on 10.12 firmware [O]. \\
22 & 1 & u8 & HDOP & 0.1 & \code{HDOP =\allowbreak{} u8 /\allowbreak{} 10}. \\
23 & 1 & char & \code{;\allowbreak{}} & — &  \\
\end{longtable}
\endgroup

\textbf{Example [O]} (float 11186, dive 107, 0x01): \code{01 00 18 02 0C D2 E2 01 46 9A 69 41 09 89 05 16 16 59 05 00 00 00 0F 3B} decodes to

\begingroup\footnotesize\setlength{\tabcolsep}{3pt}
\begin{longtable}{@{}>{\raggedright\arraybackslash}p{0.1071\linewidth}>{\raggedright\arraybackslash}p{0.1875\linewidth}@{}}
\toprule
\textbf{Field} & \textbf{Value} \\
\midrule
\endfirsthead
\toprule
\textbf{Field} & \textbf{Value} \\
\midrule
\endhead
\bottomrule
\endlastfoot
Hemisphere & East \\
Latitude & 21.51470 \\
Longitude & 118.45246 \\
Week / day & 2441 / 5 \\
Time & 2026-10-23T22:22:00Z \\
Time to fix & 890 s \\
Satellites & 5 \\
HDOP & 1.5 \\
\end{longtable}
\endgroup

\textbf{Decoding notes [D]:}

\begin{itemize}
\item Time = 1980-01-06T00:00 + (week·7 + day) days + hour:minute (UTC).
\item A fix is set invalid when:

\begin{itemize}
\item byte 3 = 0;
\item the coordinates are outside ±90/±180;
\item the time is before 2000-01-01 or in the future;
\item (when the fix is assigned) the time is later than the dive's last Iridium session time.
\end{itemize}

An invalid fix gets lat 99.99999, lon 999.99999 and no time.
\item A valid 0x05 fix advances the float's start-of-mission time. Packets sent before the start-of-mission or deployment time are ignored.
\item \textbf{Profile position selection:}

\begin{itemize}
\item A 0x02 fix of the same dive is the profile position.
\item A 0x01 fix labeled with dive N+1 is valid for dive N.
\item A 0x00 fix is used only for dive 0.
\item 0x03 and 0x05 are never used as profile positions.
\end{itemize}
\item \textbf{Listing:} one \code{GPS} line per fix: \code{GPS SN dive idx date/\allowbreak{}time position TTF n\allowbreak{}Sat SNRmin SNRavg SNRmax HDOP stat}. \code{idx} is the ID (phase) and \code{stat} is always 0.
\end{itemize}

\subsection{\texorpdfstring{A.3 Sensor series: binned profile (0x1n/0x2n/0x3n), high-resolution profile (0x90–0xB7), drift (0x98–0xBF)}{A.3 Sensor series: binned profile (0x1n/0x2n/0x3n), high-resolution profile (0x90–0xB7), drift (0x98–0xBF)}}\label{sec:a3-sensor-series-binned-profile-0x1n0x2n0x3n-high-resolution-profile-0x900xb7-drift-0x980xbf}

\textbf{Purpose.} Series of integer counts from one channel, split into numbered parts.

\begingroup\footnotesize\setlength{\tabcolsep}{3pt}
\begin{longtable}{@{}>{\raggedright\arraybackslash}p{0.2500\linewidth}>{\raggedright\arraybackslash}p{0.0893\linewidth}>{\raggedright\arraybackslash}p{0.1071\linewidth}>{\raggedright\arraybackslash}p{0.0893\linewidth}>{\raggedright\arraybackslash}p{0.1696\linewidth}>{\raggedright\arraybackslash}p{0.1250\linewidth}>{\raggedright\arraybackslash}p{0.0714\linewidth}@{}}
\toprule
\textbf{Series} & \textbf{Pressure} & \textbf{Temperature} & \textbf{Salinity} & \textbf{Part index} & \textbf{Max parts [D]} & \textbf{Listing} \\
\midrule
\endfirsthead
\toprule
\textbf{Series} & \textbf{Pressure} & \textbf{Temperature} & \textbf{Salinity} & \textbf{Part index} & \textbf{Max parts [D]} & \textbf{Listing} \\
\midrule
\endhead
\bottomrule
\endlastfoot
Binned (bin-averaged) profile & 0x10–0x1F & 0x20–0x2F & 0x30–0x3F & low nibble 0–F & 16 & \code{PRO}/\code{pro} \\
High-resolution (unbinned, ≈1 Hz) profile & 0x90–0x97 & 0xA0–0xA7 & 0xB0–0xB7 & low nibble 0–7 & 8 & \code{PRU}/\code{pru} \\
Drift/park time series & 0x98–0x9F & 0xA8–0xAF & 0xB8–0xBF & low nibble 8–F (first part = 8) & 8 & \code{D}/\code{d} \\
\end{longtable}
\endgroup

All of these use the same two encodings, selected by the high nibble of byte 1:

\begingroup\footnotesize\setlength{\tabcolsep}{3pt}
\begin{longtable}{@{}>{\raggedright\arraybackslash}p{0.1250\linewidth}>{\raggedright\arraybackslash}p{0.3661\linewidth}>{\raggedright\arraybackslash}p{0.4554\linewidth}@{}}
\toprule
\textbf{Format nibble} & \textbf{Encoding} & \textbf{Status} \\
\midrule
\endfirsthead
\toprule
\textbf{Format nibble} & \textbf{Encoding} & \textbf{Status} \\
\midrule
\endhead
\bottomrule
\endlastfoot
0 & First-difference (“legacy”) & SOLO-II heritage; ALAMO mission parameter XAS = 1. \\
1 & Curvature (second-difference) & Standard; XAS = 2. All observed ALAMO/ALTO data. \\
3 & Curvature, \textbf{auxiliary temperature channel} & Only meaningful on 0x2n; see A.4.3. \\
2, 4–15 & — & The decoder throws “Unsupported packing format”. \\
\end{longtable}
\endgroup

The same curvature body is also used by drift times (0x78, A.4.1), by corrected salinity (0xDC, after a 4-byte prefix, A.4.2) and, per the MRV YAML, by 0xC2.

\subsubsection{\texorpdfstring{A.3.1 Curvature encoding (format nibble 1 or 3)}{A.3.1 Curvature encoding (format nibble 1 or 3)}}\label{sec:a31-curvature-encoding-format-nibble-1-or-3}

\begingroup\footnotesize\setlength{\tabcolsep}{3pt}
\begin{longtable}{@{}>{\raggedright\arraybackslash}p{0.0625\linewidth}>{\raggedright\arraybackslash}p{0.1161\linewidth}>{\raggedright\arraybackslash}p{0.1250\linewidth}>{\raggedright\arraybackslash}p{0.1696\linewidth}>{\raggedright\arraybackslash}p{0.1071\linewidth}>{\raggedright\arraybackslash}p{0.3393\linewidth}@{}}
\toprule
\textbf{Offset} & \textbf{Size} & \textbf{Type} & \textbf{Field} & \textbf{Scale/units} & \textbf{Notes} \\
\midrule
\endfirsthead
\toprule
\textbf{Offset} & \textbf{Size} & \textbf{Type} & \textbf{Field} & \textbf{Scale/units} & \textbf{Notes} \\
\midrule
\endhead
\bottomrule
\endlastfoot
0 & 1 & u8 & ID & — & Class nibble + part index. \\
1 & 2 & u4+u12 & format / size & — / bytes & High nibble = 1 (or 3); low 12 bits = blob size. \\
3 & 1 & u8 & B & sub-blocks & Starting sub-block number. See the notes; \textbf{ignored by the decoder}. \\
4 & 2 & u16 & NN & values & Number of values in this part, \textbf{counting VVV} (see reassembly). \\
6 & 3 & i24 & VVV & counts & First value of the part. \\
9 & 3 & i24 & DDD & counts & First difference = v[1] − v[0]. \\
12 & 12 & 32 × u3 & p0 … p31 & — & Packing factors, one per 16-value sub-block, MSB-first: \code{p[8i+j] =\allowbreak{} (\allowbreak{}u24(\allowbreak{}12+3i) >> (\allowbreak{}21−3j)) \& 7}, i = 0..3, j = 0..7. \\
24 & ⌈Σnibbles/2⌉ & nibble stream & second differences & counts & Sub-block k holds min(16, remaining) values, each \code{p\_\allowbreak{}k} nibbles wide. \\
size−1 & 1 & char & \code{;\allowbreak{}} & — &  \\
\end{longtable}
\endgroup

\textbf{Packing-factor-to-width mapping: the truth is \emph{width = p nibbles = 4·p bits}.}

\begingroup\footnotesize\setlength{\tabcolsep}{3pt}
\begin{longtable}{@{}>{\raggedright\arraybackslash}p{0.0179\linewidth}>{\raggedright\arraybackslash}p{0.1339\linewidth}>{\raggedright\arraybackslash}p{0.4911\linewidth}>{\raggedright\arraybackslash}p{0.3036\linewidth}@{}}
\toprule
\textbf{p} & \textbf{Bits per value} & \textbf{Range (signed)} & \textbf{Bytes per full 16-value sub-block} \\
\midrule
\endfirsthead
\toprule
\textbf{p} & \textbf{Bits per value} & \textbf{Range (signed)} & \textbf{Bytes per full 16-value sub-block} \\
\midrule
\endhead
\bottomrule
\endlastfoot
0 & — & sub-block \textbf{absent}: the decoder skips it and consumes no values or nibbles & 0 \\
1 & 4 & −8 … 7 & 8 \\
2 & 8 & −128 … 127 & 16 \\
3 & 12 & −2048 … 2047 & 24 \\
4 & 16 & −32768 … 32767 & 32 \\
5 & 20 & ±2\textasciicircum{}19 & 40 \\
6 & 24 & ±2\textasciicircum{}23 & 48 \\
7 & 28 & ±2\textasciicircum{}27 & 56 \\
\end{longtable}
\endgroup

Evidence: this is \code{Abstract\allowbreak{}XSensor\allowbreak{}Data:\allowbreak{}:\allowbreak{}Unpack\allowbreak{}Sensor\allowbreak{}Data} [D]. The documents disagree with it:

\begin{itemize}
\item \textbf{[S]} says “packing factor 0 = 1 nibble”, i.e. width = p+1 nibbles.
\item \textbf{[X]} says width = \code{p\allowbreak{}N·4+4} bits (p+1 nibbles) in one sentence, then lists “p=0 → 0 nibble null, p=1 → 1 nibble, …”. Its byte counts per sub-block (16, 24, …, 64 bytes for p = 1 … 7) are 8 bytes too large.
\end{itemize}

All 24 curvature blobs in the two production dives were checked [O]. Under the decoder rule (p nibbles, p = 0 skipped) the nibble count equals the data length exactly, with at most one pad nibble. Under the p+1 rule every blob would need 30–80 \% more bytes than it has. Even all-zero second differences (constant 1800 s steps in 0x78) are sent with p = 1, never p = 0, so p = 0 is used only for unused trailing sub-blocks.

\textbf{Unpacking, exactly as the decoder does it [D]:}

\begin{enumerate}
\item Read B (ignored), NN, VVV (i24) and DDD (i24).
\item If NN = 0, throw. If NN = 1, the part's values are just [VVV]. If NN = 2, they are [VVV, VVV + DDD]. VVV is included only in the first part; see reassembly. In both cases return.
\item Otherwise, require size ≥ 24 and read the 32 packing factors.
\item Set \code{remaining =\allowbreak{} NN − 2} and k = 0, and set the nibble cursor to the high nibble of byte 24.
\item While \code{remaining > 0} and the cursor is before the terminator:

\begin{itemize}
\item If k ≥ 32, error.
\item If p\_\allowbreak{}k = 0, set k += 1 and continue (no data consumed).
\item Otherwise read n = min(16, remaining) values. For each value read p\_\allowbreak{}k nibbles, most significant first, as an unsigned integer u, then sign-extend: \code{if u ≥ 2\textasciicircum{}(\allowbreak{}4p−1) then u −=\allowbreak{} 2\textasciicircum{}(\allowbreak{}4p)}. Append it to the second differences and decrement \code{remaining}.
\item Set k += 1.
\end{itemize}

\textbf{The nibble cursor is continuous across sub-blocks}: sub-blocks are not byte-aligned. Because a full sub-block has 16·p nibbles (always even), only the last sub-block can end mid-byte. The low nibble of the last data byte is then padding.
\item First differences: \code{d[0] =\allowbreak{} DDD}, \code{d[i] =\allowbreak{} d[i−1] + s[i−1]}, i = 1 … NN−2 (cumulative sum of the second differences, seeded with DDD).
\item Values: \code{x =\allowbreak{} VVV} (not appended unless this is the first part), then for each first difference \code{x +=\allowbreak{} d[i]} and append x. That is NN − 1 values, plus VVV for the first part.
\item If the stream ends before \code{remaining} reaches 0, the decoder logs a warning and keeps the partial part (it does not throw).
\end{enumerate}

\textbf{Multi-part reassembly [D] (\code{Abstract\allowbreak{}XSensor\allowbreak{}Data\allowbreak{}Composite}):}

\begin{itemize}
\item \textbf{The first part includes VVV; later parts don't.} The “first part” is the one whose ID is 0x10, 0x20, 0x30, 0x90, 0xA0, 0xB0, 0x98, 0xA8, 0xB8 or 0x78. Part m > 0 repeats the last value of part m−1 as its VVV (an overlap of one value), and the decoder drops it, so part m contributes NN − 1 values. Total samples = NN₀ + Σ\_\allowbreak{}\{m>0\} (NN\_\allowbreak{}m − 1).

Verified [O], float 11186 dive 107, temperature 0x20/0x21/0x22 with NN = 306/290/192: 306 + 289 + 191 = 786 samples, the same count as pressure and the aux-temperature channel.
\item Parts are kept per channel, sorted by part index (the low nibble of the ID), and concatenated in index order. A duplicate index is rejected. The decoder doesn't check that the overlap VVV equals the previous part's last value.
\item \textbf{B} in the observed data is the number of 16-value sub-blocks in all preceding parts of the same channel: 0 for the first part, then 15, 19, 32, 37 … [O]. The SOLO-II document says B = 1 in the first part [S]. Not used by the decoder.
\item At most 32 sub-blocks per part, so NN ≤ 514 per blob. Observed maximum: 514 [O].
\item \textbf{Completeness} (\code{Is\allowbreak{}Data\allowbreak{}Complete}), required before any \code{pro}/\code{pru}/\code{d} rows are listed:

\begin{itemize}
\item at least one sample;
\item part indices contiguous \textbf{from the lowest received index} (a missing first part isn't detected by this test alone);
\item the same number of P and T samples;
\item aux-temperature samples (if any) equal T;
\item S samples are 0 or equal P.
\end{itemize}

Otherwise the listing prints \code{\# (\allowbreak{}Incomplete … Data)}.
\item \textbf{Missing bins:} a bin with no valid CTD data is simply not sent, so all channels skip the same bin and stay aligned [S]. The decoder has no bin numbers; \code{sample\allowbreak{}Idx} is the position in the series.
\end{itemize}

\subsubsection{\texorpdfstring{A.3.2 First-difference (“legacy”) encoding (format nibble 0)}{A.3.2 First-difference (“legacy”) encoding (format nibble 0)}}\label{sec:a32-first-difference-legacy-encoding-format-nibble-0}

\begingroup\footnotesize\setlength{\tabcolsep}{3pt}
\begin{longtable}{@{}>{\raggedright\arraybackslash}p{0.0625\linewidth}>{\raggedright\arraybackslash}p{0.0446\linewidth}>{\raggedright\arraybackslash}p{0.0625\linewidth}>{\raggedright\arraybackslash}p{0.1607\linewidth}>{\raggedright\arraybackslash}p{0.4286\linewidth}@{}}
\toprule
\textbf{Offset} & \textbf{Size} & \textbf{Type} & \textbf{Field} & \textbf{Notes} \\
\midrule
\endfirsthead
\toprule
\textbf{Offset} & \textbf{Size} & \textbf{Type} & \textbf{Field} & \textbf{Notes} \\
\midrule
\endhead
\bottomrule
\endlastfoot
0 & 1 & u8 & ID &  \\
1 & 2 & u4+u12 & format = 0 / size &  \\
3 & … & — & sub-block 0, 1, … & Each sub-block is \code{[S][v0][δ1 … δ(\allowbreak{}n−1)]}, n ≤ 25. \\
size−1 & 1 & char & \code{;\allowbreak{}} &  \\
\end{longtable}
\endgroup

Sub-block layout:

\begingroup\footnotesize\setlength{\tabcolsep}{3pt}
\begin{longtable}{@{}>{\raggedright\arraybackslash}p{0.1339\linewidth}>{\raggedright\arraybackslash}p{0.0446\linewidth}>{\raggedright\arraybackslash}p{0.6161\linewidth}@{}}
\toprule
\textbf{Sub-block byte} & \textbf{Type} & \textbf{Content} \\
\midrule
\endfirsthead
\toprule
\textbf{Sub-block byte} & \textbf{Type} & \textbf{Content} \\
\midrule
\endhead
\bottomrule
\endlastfoot
0 & u8 & scale S (1–255), chosen so that each scaled difference fits in an i8 \\
1–2 & u16 & v0 (unsigned, MSB first) \\
3 … n+1 & i8 & δk = (v[k] − v[k−1]) / S \\
\end{longtable}
\endgroup

\textbf{Decoder [D]:}

\begin{itemize}
\item \code{v[k] =\allowbreak{} v[k−1] + δk·S}. Sub-blocks hold up to 25 values (1 + 24 differences, 27 bytes); only the last sub-block may be shorter.
\item Sub-blocks are read back to back until the terminator. Parts concatenate by index \textbf{without any overlap}.
\item Quirk: the loop runs while \code{start < end − 5}, so a final sub-block with only 1 or 2 values is \textbf{silently dropped}.
\end{itemize}

\textbf{Discrepancies [S]:}

\begin{itemize}
\item The SOLO-II document gives “22 bytes” for a full 25-value sub-block, which should be 27.
\item It writes the last difference as \code{\{v[n−1] − v[n]\}/\allowbreak{}S}. That is a typo for \code{v[n−1] − v[n−2]}.
\end{itemize}

\subsubsection{\texorpdfstring{A.3.3 Listing}{A.3.3 Listing}}\label{sec:a33-listing}

\begingroup\footnotesize\setlength{\tabcolsep}{3pt}
\begin{longtable}{@{}>{\raggedright\arraybackslash}p{0.1607\linewidth}>{\raggedright\arraybackslash}p{0.8125\linewidth}@{}}
\toprule
\textbf{Series} & \textbf{Listing [D]} \\
\midrule
\endfirsthead
\toprule
\textbf{Series} & \textbf{Listing [D]} \\
\midrule
\endhead
\bottomrule
\endlastfoot
Binned & \code{PRD}/\code{prd} describes the columns, \code{PRO SN dive num\allowbreak{}Samples}, then \code{pro SN dive sample\allowbreak{}Idx pres temp saln}. With aux temperature: \code{pres temp1 temp2 saln}. With complete 0xDC data: \code{saln\allowbreak{}Corr} is appended. \\
Unbinned (hi-res) & \code{PRU SN dive num\allowbreak{}Samples}, then \code{pru SN dive sample\allowbreak{}Idx pres temp saln}. \\
Drift & Without sample times: \code{D SN dive num\allowbreak{}Samples}, then \code{d SN dive sample\allowbreak{}Idx pres temp saln}. With complete 0x78 times: \code{D SN dive n\allowbreak{}Sam drift\allowbreak{}Start} (the 0xF0 park start time), then \code{d SN dive i +sec pres temp saln}. \\
\end{longtable}
\endgroup

\subsection{\texorpdfstring{A.4 Drift sample times (0x78–0x7F), corrected salinity (0xDC), auxiliary sensor data (0xC2), auxiliary temperature channel}{A.4 Drift sample times (0x78–0x7F), corrected salinity (0xDC), auxiliary sensor data (0xC2), auxiliary temperature channel}}\label{sec:a4-drift-sample-times-0x780x7f-corrected-salinity-0xdc-auxiliary-sensor-data-0xc2-auxiliary-temperature-channel}

\subsubsection{\texorpdfstring{A.4.1 Drift sample times — 0x78–0x7F  \emph{(format inferred [I]; decoder [D] \code{XDrift\allowbreak{}Time.\allowbreak{}h})}}{A.4.1 Drift sample times — 0x78–0x7F  (format inferred [I]; decoder [D] XDriftTime.h)}}\label{sec:a41-drift-sample-times--0x780x7f--format-inferred-i-decoder-d-xdrifttimeh}

\textbf{Purpose.} The time of each drift sample (0x98/0xA8/0xB8 series), in seconds after the Argo \code{park start time} (0xF0 bytes 20–23). Not in any MRV or spreadsheet document. The format was inferred from production packets (2026-10-24).

\begin{itemize}
\item \textbf{Layout:} identical to the curvature encoding of A.3.1. ID = 0x78 + part index (0x78 = first part, which includes VVV); format nibble 1.
\item \textbf{Values:} u24 seconds, no gain or offset.
\item \textbf{Observed [O]:}

\begin{itemize}
\item VVV = 1 (the first drift sample is taken 1 s after the drift starts).
\item DDD = sample period (1800, 7200 or 14400 s; equal to the mission parameter DDT on float 11186).
\item All second differences 0 (p = 1).
\item NN equals the drift P/T/S sample count.
\end{itemize}
\item \textbf{Multi-part:} 0x79 has been observed in logs, so multi-part series do occur. Reassembly is the same as A.3.1.
\item \textbf{Decoder [D]:}

\begin{itemize}
\item Accepts \code{(\allowbreak{}ID \& 0x\allowbreak{}F8) =\allowbreak{}=\allowbreak{} 0x78}.
\item Parts are kept with the drift composite.
\item \code{Is\allowbreak{}Time\allowbreak{}Data\allowbreak{}Complete} requires contiguous indices and as many times as drift samples. Then the \code{D}/\code{d} lines switch to the time layout (A.3.3). Otherwise \code{\# (\allowbreak{}Incomplete Drift Time Data)} is printed and the old layout is used.
\item \code{xmtps-list2json.\allowbreak{}py} uses these times instead of estimating drift times.
\end{itemize}
\end{itemize}

\subsubsection{\texorpdfstring{A.4.2 Dynamically corrected salinity — 0xDC  \emph{(format inferred [I]; decoder [D] \code{XProfile\allowbreak{}Salinity\allowbreak{}Corrected.\allowbreak{}h})}}{A.4.2 Dynamically corrected salinity — 0xDC  (format inferred [I]; decoder [D] XProfileSalinityCorrected.h)}}\label{sec:a42-dynamically-corrected-salinity--0xdc--format-inferred-i-decoder-d-xprofilesalinitycorrectedh}

\textbf{Purpose.} A second binned-profile salinity channel, corrected on the float (RBR dynamic correction). It is sent only when mission parameter PPE (dynamic correction mode) = 2 [O]. Values use the salinity gain and offset, and differ from the uncorrected salinity by a few thousandths of a PSU (at most about 0.07). Not in any MRV or spreadsheet document.

\begingroup\footnotesize\setlength{\tabcolsep}{3pt}
\begin{longtable}{@{}>{\raggedright\arraybackslash}p{0.0625\linewidth}>{\raggedright\arraybackslash}p{0.0446\linewidth}>{\raggedright\arraybackslash}p{0.0714\linewidth}>{\raggedright\arraybackslash}p{0.1696\linewidth}>{\raggedright\arraybackslash}p{0.1071\linewidth}>{\raggedright\arraybackslash}p{0.4643\linewidth}@{}}
\toprule
\textbf{Offset} & \textbf{Size} & \textbf{Type} & \textbf{Field} & \textbf{Scale/units} & \textbf{Notes} \\
\midrule
\endfirsthead
\toprule
\textbf{Offset} & \textbf{Size} & \textbf{Type} & \textbf{Field} & \textbf{Scale/units} & \textbf{Notes} \\
\midrule
\endhead
\bottomrule
\endlastfoot
0 & 1 & u8 & ID & — & Always 0xDC. The part index is in byte 4, not in the ID. \\
1 & 2 & u4+u12 & format / size & — / bytes & Format 1 (curvature) observed. \\
3 & 1 & u8 & (unknown) & — & Always 0x00 [O]. The decoder warns if it is non-zero. \\
4 & 1 & u8 & part index & — & 0 … 15. The decoder rejects values > 15. \\
5 & 1 & u8 & channel & — & Always 0x01 (salinity) [O]. The decoder rejects other values. \\
6 & 1 & u8 & (unknown) & — & Always 0x00 [O]. The decoder warns if it is non-zero. \\
7 & 1 & u8 & B & sub-blocks & As A.3.1 byte 3. \\
8 & 2 & u16 & NN & values &  \\
10 & 3 & i24 & VVV & counts &  \\
13 & 3 & i24 & DDD & counts &  \\
16 & 12 & 32 × u3 & packing factors & — &  \\
28 & … & nibbles & second differences & counts &  \\
size−1 & 1 & char & \code{;\allowbreak{}} & — &  \\
\end{longtable}
\endgroup

\textbf{Decoder [D]:}

\begin{itemize}
\item Validates size (8 ≤ size ≤ received bytes), byte 5 = 1 and byte 4 ≤ 15.
\item Then rewrites the blob as a standard salinity message: ID \code{0x30 |\allowbreak{} byte4}, \code{jj} = original format nibble + (size − 4), followed by bytes 7 … size−1. The result is decoded as A.3.1, so part 0 includes VVV.
\item The total count equals the 0x3n salinity count. The parts can split differently from 0x3n [O].
\item The \code{saln\allowbreak{}Corr} column is added to \code{pro} only when the corrected channel is complete (contiguous indices and the same count as salinity). Otherwise \code{\# (\allowbreak{}Incomplete Corrected Salinity Data)} is printed.
\item In the JSON output it is \code{saln\_\allowbreak{}corr} (NaN when absent).
\end{itemize}

\subsubsection{\texorpdfstring{A.4.3 Auxiliary temperature channel (format nibble 3 on 0x2n)}{A.4.3 Auxiliary temperature channel (format nibble 3 on 0x2n)}}\label{sec:a43-auxiliary-temperature-channel-format-nibble-3-on-0x2n}

\textbf{Purpose.} A second temperature series, aligned with the binned profile. Presumably the RBR conductivity-cell thermistor, enabled by mission parameter CTC “conductivity temperature channel mode” [M] [I]. It is sent as ordinary 0x20, 0x21, … blobs whose \code{jj} high nibble is \textbf{3} instead of 1, alongside the primary temperature parts, which carry the same IDs with format 1 [O].

\begin{itemize}
\item \textbf{Encoding:} curvature, exactly as A.3.1 (\code{XSENS\_\allowbreak{}FMT\_\allowbreak{}CURVATURE\_\allowbreak{}AUX\_\allowbreak{}TEMP\_\allowbreak{}CHANNEL =\allowbreak{} 3} [D]).
\item \textbf{Units:} the temperature gain and offset.
\item \textbf{Decoder [D]:}

\begin{itemize}
\item \code{XProfile\allowbreak{}Temperature} sets \code{is\_\allowbreak{}aux\_\allowbreak{}channel(\allowbreak{})}. \code{XFloat\allowbreak{}Dive\allowbreak{}Data} then routes the part to a separate aux-temperature channel of the binned composite, so the shared IDs don't collide.
\item It must have the same count as the primary temperature.
\item Listed as the \code{temp2} column (\code{temp1} = CTD temperature), and as \code{tcnd} in the JSON.
\end{itemize}
\item \textbf{Limitations [D]:}

\begin{itemize}
\item The aux flag is honored only for binned 0x2n. A format-3 high-resolution (0xAn) or drift (0xA8+) part would be merged into the primary channel, and fail with a duplicate index.
\item The legacy \code{AXT}/\code{axt} listing section (\code{XComposite\allowbreak{}Aux\allowbreak{}Temp}) is no longer written.
\end{itemize}
\item The spreadsheet [X] and the SOLO-II document [S] don't document format 3.
\end{itemize}

\subsubsection{\texorpdfstring{A.4.4 Auxiliary sensor data — 0xC2  \emph{(not implemented by the decoder)}}{A.4.4 Auxiliary sensor data — 0xC2  (not implemented by the decoder)}}\label{sec:a44-auxiliary-sensor-data--0xc2--not-implemented-by-the-decoder}

\textbf{Status [D]:} 0xC2 is defined in \code{XMessage.\allowbreak{}h} (\code{XMSG\_\allowbreak{}TYPE\_\allowbreak{}ID\_\allowbreak{}AUX\_\allowbreak{}SENSOR\_\allowbreak{}DATA}) but has no case in the blob dispatcher, so it is logged as an unknown ID and skipped. The class \code{Auxiliary\allowbreak{}Sensor\allowbreak{}Data} (a copy of the A.3 unpacker supporting formats 0 and 1) exists but is never instantiated. No 0xC2 blob was seen in the production log sampled for this document.

\textbf{Format per MRV DataFormat-v7.3.5.yaml, blob 194 = 0xC2 [M].} Byte offsets are \textbf{inferred [I]}: the YAML gives field order and types only, and assumes the usual ID + \code{jj} header.

\begingroup\footnotesize\setlength{\tabcolsep}{3pt}
\begin{longtable}{@{}>{\raggedright\arraybackslash}p{0.0625\linewidth}>{\raggedright\arraybackslash}p{0.0446\linewidth}>{\raggedright\arraybackslash}p{0.0982\linewidth}>{\raggedright\arraybackslash}p{0.1339\linewidth}>{\raggedright\arraybackslash}p{0.1071\linewidth}>{\raggedright\arraybackslash}p{0.4732\linewidth}@{}}
\toprule
\textbf{Offset} & \textbf{Size} & \textbf{Type} & \textbf{Field} & \textbf{Scale/units} & \textbf{Notes} \\
\midrule
\endfirsthead
\toprule
\textbf{Offset} & \textbf{Size} & \textbf{Type} & \textbf{Field} & \textbf{Scale/units} & \textbf{Notes} \\
\midrule
\endhead
\bottomrule
\endlastfoot
0 & 1 & u8 & ID & — & 0xC2. \\
1 & 2 & u16? & size & bytes & Size layout not stated. Probably a full 16-bit size like the other 0xCn types [I]. \\
3 & 1 & u8 & packet index & — & Part number. \\
4 & 1 & u8 & sensor channel & — &  \\
5 & 1 & state code & state & — & Dive phase code (A.5.2). The YAML type \code{state\_\allowbreak{}code} is 1 nibble in other blobs; it is assumed here to take a whole byte [I]. \\
6 & 1 & u8 & data type & — &  \\
7 & 2 & u16 & gain & — & Applies to \code{data}: value = counts/gain − offset. \\
9 & 2 & u16 & offset & — &  \\
11 & … & packed & data & counts & “packed” = the curvature body of A.3.1 (B, NN, VVV, DDD, packing factors, nibbles) [I]. \\
size−1 & 1 & char & \code{;\allowbreak{}} & — &  \\
\end{longtable}
\endgroup

\begin{notebox}
\textbf{Discrepancy [X] vs [M]:} the v2.4 spreadsheet instead puts an auxiliary-sensor table into 0xF0 (sensor count N, then for each sensor: data type, message start index, gain, offset; 6 bytes each). MRV firmware 10.5/10.12 sends a fixed 66-byte 0xF0 without that table, and carries gain and offset in the 0xC2 blob itself.

\end{notebox}

\subsection{\texorpdfstring{A.5 Time series: fall 0x4n, rise 0x5n, pump 0x6n, surface 0x8n, valve 0xEB}{A.5 Time series: fall 0x4n, rise 0x5n, pump 0x6n, surface 0x8n, valve 0xEB}}\label{sec:a5-time-series-fall-0x4n-rise-0x5n-pump-0x6n-surface-0x8n-valve-0xeb}

\subsubsection{\texorpdfstring{A.5.1 Fall (0x40–0x4F), rise (0x50–0x5F) and surface (0x80–0x8F) time series}{A.5.1 Fall (0x40–0x4F), rise (0x50–0x5F) and surface (0x80–0x8F) time series}}\label{sec:a51-fall-0x400x4f-rise-0x500x5f-and-surface-0x800x8f-time-series}

\textbf{Purpose.}

\begin{itemize}
\item \textbf{Fall:} pressure samples from the start of descent to the start of drift. The first sample is taken when the valve opens; others at fixed pressure steps or time intervals [S].
\item \textbf{Rise:} samples from the start of the descent-to-profile, through the ascent, to the end of ascent. SOLO-II v2.0 adds a record after surface buoyancy pumping [S].
\item \textbf{Surface:} CTD samples while at the surface. The surface pressure is used for the pressure-offset correction and for the \code{SFD} statistics.
\end{itemize}

All three share one decoder (\code{Abstract\allowbreak{}XFall\allowbreak{}Rise:\allowbreak{}:\allowbreak{}Unpack\allowbreak{}Data} [D]). The low nibble of the ID is the part index (0–15). Each part carries its own reference time.

\textbf{Common header:}

\begingroup\footnotesize\setlength{\tabcolsep}{3pt}
\begin{longtable}{@{}>{\raggedright\arraybackslash}p{0.0625\linewidth}>{\raggedright\arraybackslash}p{0.0446\linewidth}>{\raggedright\arraybackslash}p{0.0625\linewidth}>{\raggedright\arraybackslash}p{0.1875\linewidth}>{\raggedright\arraybackslash}p{0.1696\linewidth}>{\raggedright\arraybackslash}p{0.3929\linewidth}@{}}
\toprule
\textbf{Offset} & \textbf{Size} & \textbf{Type} & \textbf{Field} & \textbf{Scale/units} & \textbf{Notes} \\
\midrule
\endfirsthead
\toprule
\textbf{Offset} & \textbf{Size} & \textbf{Type} & \textbf{Field} & \textbf{Scale/units} & \textbf{Notes} \\
\midrule
\endhead
\bottomrule
\endlastfoot
0 & 1 & u8 & ID & — & 0x4n fall, 0x5n rise, 0x8n surface; n = part index. \\
1 & 2 & u4+u12 & format / size & — / bytes & Format nibble selects the record layout (below). \\
3 & 4 & u32 & start/reference time & s since 2000-01-01 & Part 0 starts at the start of descent (fall), the start of ascent (rise) or surface arrival. The record times are u16 offsets, so the float starts a new part with a new reference time when an offset would exceed 65535 s [X]. \\
7 & … & — & records & — & Back to back. \\
size−1 & 1 & char & \code{;\allowbreak{}} & — &  \\
\end{longtable}
\endgroup

\textbf{Record layouts by format nibble:}

\begingroup\footnotesize\setlength{\tabcolsep}{3pt}
\begin{longtable}{@{}>{\raggedright\arraybackslash}p{0.0625\linewidth}>{\raggedright\arraybackslash}p{0.1071\linewidth}>{\raggedright\arraybackslash}p{0.3750\linewidth}>{\raggedright\arraybackslash}p{0.4018\linewidth}@{}}
\toprule
\textbf{Format} & \textbf{Record size} & \textbf{Name} & \textbf{Used by} \\
\midrule
\endfirsthead
\toprule
\textbf{Format} & \textbf{Record size} & \textbf{Name} & \textbf{Used by} \\
\midrule
\endhead
\bottomrule
\endlastfoot
0 & 4 & Legacy & SOLO-II fall/rise; legacy surface \\
1 & 5 & Standard (phase code) & fall/rise (observed [O]); “legacy” surface \\
2 & 9 & Standard with T \& S & surface (observed [O]); fall/rise also accepted by the decoder \\
3 & — & Curvature, P/T/S in separate blobs (0x80–0x83 P, 0x84–0x87 T, 0x88–0x8B S) & Spreadsheet only [X]; \textbf{not implemented} (the decoder throws “unsupported … format”) \\
\end{longtable}
\endgroup

Format 0 record (offset r = 7 + 4k):

\begingroup\footnotesize\setlength{\tabcolsep}{3pt}
\begin{longtable}{@{}>{\raggedright\arraybackslash}p{0.0625\linewidth}>{\raggedright\arraybackslash}p{0.0446\linewidth}>{\raggedright\arraybackslash}p{0.0446\linewidth}>{\raggedright\arraybackslash}p{0.0804\linewidth}>{\raggedright\arraybackslash}p{0.1607\linewidth}>{\raggedright\arraybackslash}p{0.5268\linewidth}@{}}
\toprule
\textbf{Offset} & \textbf{Size} & \textbf{Type} & \textbf{Field} & \textbf{Scale/units} & \textbf{Notes} \\
\midrule
\endfirsthead
\toprule
\textbf{Offset} & \textbf{Size} & \textbf{Type} & \textbf{Field} & \textbf{Scale/units} & \textbf{Notes} \\
\midrule
\endhead
\bottomrule
\endlastfoot
r & 2 & u16 & time & s after reference &  \\
r+2 & 2 & u16 & pressure & counts & No phase code (listed as 0 “unknown”). The spreadsheet says 20 bits here [X], but a 4-byte record only leaves 16 bits; the decoder reads 16. \\
\end{longtable}
\endgroup

Format 1 record (offset r = 7 + 5k):

\begingroup\footnotesize\setlength{\tabcolsep}{3pt}
\begin{longtable}{@{}>{\raggedright\arraybackslash}p{0.0625\linewidth}>{\raggedright\arraybackslash}p{0.0446\linewidth}>{\raggedright\arraybackslash}p{0.0446\linewidth}>{\raggedright\arraybackslash}p{0.0982\linewidth}>{\raggedright\arraybackslash}p{0.1607\linewidth}>{\raggedright\arraybackslash}p{0.3304\linewidth}@{}}
\toprule
\textbf{Offset} & \textbf{Size} & \textbf{Type} & \textbf{Field} & \textbf{Scale/units} & \textbf{Notes} \\
\midrule
\endfirsthead
\toprule
\textbf{Offset} & \textbf{Size} & \textbf{Type} & \textbf{Field} & \textbf{Scale/units} & \textbf{Notes} \\
\midrule
\endhead
\bottomrule
\endlastfoot
r & 2 & u16 & time & s after reference &  \\
r+2 & ½ & u4 & phase code & — & High nibble of byte r+2 (A.5.2). \\
r+2 & 2½ & u20 & pressure & counts & Low nibble of r+2, then r+3 and r+4. \\
\end{longtable}
\endgroup

Format 2 record (offset r = 7 + 9k):

\begingroup\footnotesize\setlength{\tabcolsep}{3pt}
\begin{longtable}{@{}>{\raggedright\arraybackslash}p{0.0625\linewidth}>{\raggedright\arraybackslash}p{0.0446\linewidth}>{\raggedright\arraybackslash}p{0.0446\linewidth}>{\raggedright\arraybackslash}p{0.1071\linewidth}>{\raggedright\arraybackslash}p{0.1607\linewidth}>{\raggedright\arraybackslash}p{0.0536\linewidth}@{}}
\toprule
\textbf{Offset} & \textbf{Size} & \textbf{Type} & \textbf{Field} & \textbf{Scale/units} & \textbf{Notes} \\
\midrule
\endfirsthead
\toprule
\textbf{Offset} & \textbf{Size} & \textbf{Type} & \textbf{Field} & \textbf{Scale/units} & \textbf{Notes} \\
\midrule
\endhead
\bottomrule
\endlastfoot
r & 2 & u16 & time & s after reference &  \\
r+2 & ½ & u4 & phase code & — &  \\
r+2 & 2½ & u20 & pressure & counts &  \\
r+5 & 2 & u16 & temperature & counts &  \\
r+7 & 2 & u16 & salinity & counts &  \\
\end{longtable}
\endgroup

Typical sizes [X] (N = index of the last record):

\begingroup\footnotesize\setlength{\tabcolsep}{3pt}
\begin{longtable}{@{}>{\raggedright\arraybackslash}p{0.0625\linewidth}>{\raggedright\arraybackslash}p{0.0714\linewidth}@{}}
\toprule
\textbf{Format} & \textbf{Size} \\
\midrule
\endfirsthead
\toprule
\textbf{Format} & \textbf{Size} \\
\midrule
\endhead
\bottomrule
\endlastfoot
0 & 4N + 12 \\
1 & 5N + 13 \\
2 & 9N + 17 \\
\end{longtable}
\endgroup

\textbf{Decoding notes [D]:}

\begin{itemize}
\item Records are read until fewer than one record plus \code{;\allowbreak{}} remains. The decoder checks \code{size ≤ received} and the terminator at size−1.
\item Formats 0 and 1 carry no T/S; those columns are listed as \code{-9.\allowbreak{}999}.
\item Fall, rise and the profile use the profile gains/offsets and the optional pressure correction. Surface uses the surface gains/offsets and is never corrected.
\item \textbf{Listing:} one group per part: \code{F|\allowbreak{}R|\allowbreak{}SFC SN dive idx n\allowbreak{}Sam start\allowbreak{}Time}, then \code{sfc SN dive idx sam +sec pres temp saln phase\allowbreak{}Code \# phase\allowbreak{}Name}, or for fall and rise \code{f|\allowbreak{}r SN dive idx sam +sec pres phase\allowbreak{}Code \# phase\allowbreak{}Name} (the temp and saln columns were dropped from \code{f} and \code{r} in build 3321; only the 9-byte packing format carries them, last used by 9035 and 9050).
\item Surface parts are also combined into the derived \code{SFD} line: number of valid samples, and the minimum, mean, maximum and variance of the surface pressure.
\end{itemize}

\textbf{Example [O]} (float 11186, 0x41, format 1): \code{41 10 3F 32 6E E1 E6 00 00 A0 79 DE …} decodes to

\begingroup\footnotesize\setlength{\tabcolsep}{3pt}
\begin{longtable}{@{}>{\raggedright\arraybackslash}p{0.1339\linewidth}>{\raggedright\arraybackslash}p{0.4375\linewidth}@{}}
\toprule
\textbf{Field} & \textbf{Value} \\
\midrule
\endfirsthead
\toprule
\textbf{Field} & \textbf{Value} \\
\midrule
\endhead
\bottomrule
\endlastfoot
Part & 1 \\
Size & 63 \\
Reference time & 0x326EE1E6 s after 2000-01-01 \\
First record & t = +0 s, code 10 (descending), P counts 0x079DE \\
\end{longtable}
\endgroup

\subsubsection{\texorpdfstring{A.5.2 Phase codes (4-bit, shared by fall/rise/surface/pump/valve)}{A.5.2 Phase codes (4-bit, shared by fall/rise/surface/pump/valve)}}\label{sec:a52-phase-codes-4-bit-shared-by-fallrisesurfacepumpvalve}

\begingroup\footnotesize\setlength{\tabcolsep}{3pt}
\begin{longtable}{@{}>{\raggedright\arraybackslash}p{0.0446\linewidth}>{\raggedright\arraybackslash}p{0.5625\linewidth}>{\raggedright\arraybackslash}p{0.2589\linewidth}@{}}
\toprule
\textbf{Code} & \textbf{ALAMO/ALTO — decoder [D] (\code{XPres\allowbreak{}Code\allowbreak{}Counts}) and spreadsheet [X]} & \textbf{SOLO-II v2.0 [S]} \\
\midrule
\endfirsthead
\toprule
\textbf{Code} & \textbf{ALAMO/ALTO — decoder [D] (\code{XPres\allowbreak{}Code\allowbreak{}Counts}) and spreadsheet [X]} & \textbf{SOLO-II v2.0 [S]} \\
\midrule
\endhead
\bottomrule
\endlastfoot
0 & unknown (legacy formats) & — \\
1 & start of descent (START\_\allowbreak{}OF\_\allowbreak{}SINK) & START\_\allowbreak{}OF\_\allowbreak{}SINK \\
2 & first descent pump (“buoyancy at 100 dbar”) & Buoyancy at 100 db \\
3 & seek to drift pressure & SEEK \\
4 & start of drift & BEGINNING\_\allowbreak{}OF\_\allowbreak{}DRIFT \\
5 & seek during drift & BEGINNING\_\allowbreak{}OF\_\allowbreak{}FALL\_\allowbreak{}TO\_\allowbreak{}PROFILE \\
6 & start of descent to profile & START\_\allowbreak{}OF\_\allowbreak{}RISE \\
7 & start of ascent & END\_\allowbreak{}OF\_\allowbreak{}RISE \\
8 & end of ascent & ICE\_\allowbreak{}TURNAROUND \\
9 & ice turn-around & — \\
10 & descending & SINKING \\
11 & drift/park & DRIFTING \\
12 & descending to profile & FALLING\_\allowbreak{}TO\_\allowbreak{}PROFILE \\
13 & ascending & RISING \\
14 & surface & SURFACE \\
15 & “invalid phase code” (accepted by the decoder) & — \\
\end{longtable}
\endgroup

\begin{notebox}
ALAMO inserted “seek during drift” at 5, which shifts SOLO-II codes 5–8 up by one. The spreadsheet lists “9: END\_\allowbreak{}OF\_\allowbreak{}RISE” and “9: ICE\_\allowbreak{}TURNAROUND”. The first is a typo for 8 (decoder and data: 8 = end of ascent [O]).

\end{notebox}

\subsubsection{\texorpdfstring{A.5.3 Pump — 0x60–0x6F}{A.5.3 Pump — 0x60–0x6F}}\label{sec:a53-pump--0x600x6f}

\textbf{Purpose.} One record per pump run: pressure, duration, battery voltage, current and case vacuum. The low nibble of the ID is the part index.

\textbf{Header:}

\begingroup\footnotesize\setlength{\tabcolsep}{3pt}
\begin{longtable}{@{}>{\raggedright\arraybackslash}p{0.0625\linewidth}>{\raggedright\arraybackslash}p{0.0446\linewidth}>{\raggedright\arraybackslash}p{0.0714\linewidth}>{\raggedright\arraybackslash}p{0.1250\linewidth}>{\raggedright\arraybackslash}p{0.4375\linewidth}@{}}
\toprule
\textbf{Offset} & \textbf{Size} & \textbf{Type} & \textbf{Field} & \textbf{Notes} \\
\midrule
\endfirsthead
\toprule
\textbf{Offset} & \textbf{Size} & \textbf{Type} & \textbf{Field} & \textbf{Notes} \\
\midrule
\endhead
\bottomrule
\endlastfoot
0 & 1 & u8 & ID &  \\
1 & 2 & u4+u12 & format / size & Format 0 reads the size as the full 16 bits [D]. \\
3 & … & records &  &  \\
size−1 & 1 & char & \code{;\allowbreak{}} &  \\
\end{longtable}
\endgroup

Format 0 record, 10 bytes, legacy SOLO-II (offset r = 3 + 10k):

\begingroup\footnotesize\setlength{\tabcolsep}{3pt}
\begin{longtable}{@{}>{\raggedright\arraybackslash}p{0.0625\linewidth}>{\raggedright\arraybackslash}p{0.0446\linewidth}>{\raggedright\arraybackslash}p{0.0446\linewidth}>{\raggedright\arraybackslash}p{0.1875\linewidth}>{\raggedright\arraybackslash}p{0.1071\linewidth}>{\raggedright\arraybackslash}p{0.3036\linewidth}@{}}
\toprule
\textbf{Offset} & \textbf{Size} & \textbf{Type} & \textbf{Field} & \textbf{Scale/units} & \textbf{Notes} \\
\midrule
\endfirsthead
\toprule
\textbf{Offset} & \textbf{Size} & \textbf{Type} & \textbf{Field} & \textbf{Scale/units} & \textbf{Notes} \\
\midrule
\endhead
\bottomrule
\endlastfoot
r & 2 & u16 & pressure & counts &  \\
r+2 & 2 & i16 & pump time & s & Signed [S]. \\
r+4 & 2 & u16 & pump battery voltage & 0.01 V &  \\
r+6 & 2 & u16 & pump current & mA &  \\
r+8 & 1 & u8 & vac0 & see note & Vacuum after the pump starts [S]. \\
r+9 & 1 & u8 & vac1 & see note & Vacuum before the pump stops [S]. \\
\end{longtable}
\endgroup

Format 1 record, 11 bytes, standard with 2 vacuums (offset r = 3 + 11k):

\begingroup\footnotesize\setlength{\tabcolsep}{3pt}
\begin{longtable}{@{}>{\raggedright\arraybackslash}p{0.0625\linewidth}>{\raggedright\arraybackslash}p{0.0446\linewidth}>{\raggedright\arraybackslash}p{0.0446\linewidth}>{\raggedright\arraybackslash}p{0.1964\linewidth}>{\raggedright\arraybackslash}p{0.1071\linewidth}>{\raggedright\arraybackslash}p{0.1964\linewidth}@{}}
\toprule
\textbf{Offset} & \textbf{Size} & \textbf{Type} & \textbf{Field} & \textbf{Scale/units} & \textbf{Notes} \\
\midrule
\endfirsthead
\toprule
\textbf{Offset} & \textbf{Size} & \textbf{Type} & \textbf{Field} & \textbf{Scale/units} & \textbf{Notes} \\
\midrule
\endhead
\bottomrule
\endlastfoot
r & ½ & u4 & phase code & — &  \\
r & 2½ & u20 & pressure & counts &  \\
r+3 & 2 & i16 & pump time & s &  \\
r+5 & 2 & u16 & voltage & 0.01 V &  \\
r+7 & 2 & u16 & current & mA &  \\
r+9 & 1 & u8 & vac0 (before pumping) & 0.1 inHg & Decoder: ×10 → cinHg. \\
r+10 & 1 & u8 & vac1 (after pumping) & 0.1 inHg &  \\
\end{longtable}
\endgroup

Format 2 record, 11 bytes, single 16-bit vacuum (offset r = 3 + 11k):

\begingroup\footnotesize\setlength{\tabcolsep}{3pt}
\begin{longtable}{@{}>{\raggedright\arraybackslash}p{0.0625\linewidth}>{\raggedright\arraybackslash}p{0.0446\linewidth}>{\raggedright\arraybackslash}p{0.0446\linewidth}>{\raggedright\arraybackslash}p{0.1071\linewidth}>{\raggedright\arraybackslash}p{0.1071\linewidth}>{\raggedright\arraybackslash}p{0.5536\linewidth}@{}}
\toprule
\textbf{Offset} & \textbf{Size} & \textbf{Type} & \textbf{Field} & \textbf{Scale/units} & \textbf{Notes} \\
\midrule
\endfirsthead
\toprule
\textbf{Offset} & \textbf{Size} & \textbf{Type} & \textbf{Field} & \textbf{Scale/units} & \textbf{Notes} \\
\midrule
\endhead
\bottomrule
\endlastfoot
r & ½ & u4 & phase code & — &  \\
r & 2½ & u20 & pressure & counts &  \\
r+3 & 2 & i16 & pump time & s &  \\
r+5 & 2 & u16 & voltage & 0.01 V &  \\
r+7 & 2 & u16 & current & mA &  \\
r+9 & 2 & u16 & case vacuum & 0.01 inHg & The decoder stores it as vac1 (“after”); vac0 is the fill value (listed 655.35). The MRV YAML calls it \code{case\_\allowbreak{}pressure} (u16, /100, inHg) [M]. \\
\end{longtable}
\endgroup

\textbf{Status of the formats:}

\begin{itemize}
\item Format 2 is what current firmware sends (observed 0x60 \code{jj} = 0x205C / 0x2122 [O]).
\item The spreadsheet documents only format 1 (“standard”) in full; it describes format 2 as “11-byte (Tentative) … with phase codes” without a layout [X]. The format-2 layout above is from the decoder and the YAML, and is confirmed by the data.
\end{itemize}

\textbf{Example [O]} (float 11186, format 2): \code{20 07 59 00 87 05 DA 00 C5 01 F0} decodes to

\begingroup\footnotesize\setlength{\tabcolsep}{3pt}
\begin{longtable}{@{}>{\raggedright\arraybackslash}p{0.1071\linewidth}>{\raggedright\arraybackslash}p{0.4375\linewidth}@{}}
\toprule
\textbf{Field} & \textbf{Value} \\
\midrule
\endfirsthead
\toprule
\textbf{Field} & \textbf{Value} \\
\midrule
\endhead
\bottomrule
\endlastfoot
Phase code & 2 (first descent pump) \\
Pressure & 0x00759 counts = 65.24 dBar (gain 25, offset 10) \\
Pump time & 135 s \\
Voltage & 14.98 V \\
Current & 0.197 A \\
Case vacuum & 4.96 inHg \\
\end{longtable}
\endgroup

\textbf{Decoding notes [D]:}

\begin{itemize}
\item Format 0 vacuum bytes are taken as cinHg without scaling, so the listing shows byte/100 inHg. The SOLO-II document calls them “vacuum counts” without units, so the format-0 vacuum units are unverified.
\item Phase codes are decoded but not listed.
\item \textbf{Listing:} \code{PMP SN dive idx n\allowbreak{}Sam}, then \code{pmp SN dive idx sam pres pump\allowbreak{}T(\allowbreak{}s) pump\allowbreak{}V(\allowbreak{}V) pump\allowbreak{}I(\allowbreak{}A) vac0(\allowbreak{}in\allowbreak{}Hg) vac1(\allowbreak{}in\allowbreak{}Hg)}.
\item \textbf{Spreadsheet typo [X]:} the terminator offset is given as 4N+14; it should be 11N+14.
\end{itemize}

\subsubsection{\texorpdfstring{A.5.4 Valve — 0xEB}{A.5.4 Valve — 0xEB}}\label{sec:a54-valve--0xeb}

\textbf{Purpose.} One record per valve opening: pressure, open time and case vacuum.

\begingroup\footnotesize\setlength{\tabcolsep}{3pt}
\begin{longtable}{@{}>{\raggedright\arraybackslash}p{0.0625\linewidth}>{\raggedright\arraybackslash}p{0.0446\linewidth}>{\raggedright\arraybackslash}p{0.0446\linewidth}>{\raggedright\arraybackslash}p{0.1429\linewidth}>{\raggedright\arraybackslash}p{0.1071\linewidth}>{\raggedright\arraybackslash}p{0.5000\linewidth}@{}}
\toprule
\textbf{Offset} & \textbf{Size} & \textbf{Type} & \textbf{Field} & \textbf{Scale/units} & \textbf{Notes} \\
\midrule
\endfirsthead
\toprule
\textbf{Offset} & \textbf{Size} & \textbf{Type} & \textbf{Field} & \textbf{Scale/units} & \textbf{Notes} \\
\midrule
\endhead
\bottomrule
\endlastfoot
0 & 1 & u8 & ID & — & 0xEB. \\
1 & 2 & u16 & size & bytes & Full 16-bit size; 7N + 12 typical. \\
3 & ½ & u4 & encoding format & — & High nibble. 0 is the only format [X]; not checked [D]. \\
3 & ½ & u4 & part index & — & Low nibble [X][D]. \\
4 + 7k & ½ & u4 & phase code & — &  \\
4 + 7k & 2½ & u20 & pressure & counts & Profile gains/offsets. \\
7 + 7k & 2 & i16 & valve open time & s &  \\
9 + 7k & 2 & u16 & case vacuum & 0.01 inHg &  \\
size−1 & 1 & char & \code{;\allowbreak{}} & — &  \\
\end{longtable}
\endgroup

\textbf{Example [O]} (float 11186): \code{EB 00 0C 00 10 00 FE 01 5B 01 DD 3B} decodes to

\begingroup\footnotesize\setlength{\tabcolsep}{3pt}
\begin{longtable}{@{}>{\raggedright\arraybackslash}p{0.1429\linewidth}>{\raggedright\arraybackslash}p{0.2054\linewidth}@{}}
\toprule
\textbf{Field} & \textbf{Value} \\
\midrule
\endfirsthead
\toprule
\textbf{Field} & \textbf{Value} \\
\midrule
\endhead
\bottomrule
\endlastfoot
Format / part & 0 / 0 \\
Phase code & 1 (start of descent) \\
Pressure & 254 counts = 0.16 dBar \\
Valve open time & 347 s \\
Case vacuum & 4.77 inHg \\
\end{longtable}
\endgroup

\textbf{Decoding notes:}

\begin{itemize}
\item \textbf{Vacuum scale [D][O]:} the spreadsheet gives gain 10 for the vacuum [X], which would put the observed values near 48 inHg, more than one atmosphere. The decoder's /100 gives 4.77 inHg, which agrees with the pump vacuums, so /100 is correct.
\item The spreadsheet says “WHEN?” for when the vacuum is measured. The decoder treats it as the vacuum after the valve operation (vac1); vac0 is listed as 655.35.
\item The MRV YAML lists the record fields but omits byte 3 [M]. The observed data has byte 3.
\item The spreadsheet overview marks 0xEB “decoder supported: no”; that is outdated.
\item \textbf{Listing:} \code{VLV SN dive idx n\allowbreak{}Sam}, then \code{vlv SN dive idx sam pres valve\allowbreak{}Time vac0 vac1}.
\end{itemize}

\subsection{\texorpdfstring{A.6 Argo data — 0xF0}{A.6 Argo data — 0xF0}}\label{sec:a6-argo-data--0xf0}

\textbf{Purpose.} Per-dive mission targets, the four state start times, and the \textbf{gains and offsets that convert all CTD counts in the dive}. It is listed as \code{A} lines with \code{midx} 0–5. Without it, nothing in the dive can be converted (A.1.5).

\textbf{Version byte [D][X]:} byte 3. \textbf{High nibble = major, low nibble = minor} (0x24 = v2.4). The SOLO-II document says the opposite (“minor in high nibble, major in low”) [S]; observed packets (\code{0x24} on firmware 10.12 [O]) and the decoder use high = major.

\textbf{Versions the decoder supports [D]:}

\begingroup\footnotesize\setlength{\tabcolsep}{3pt}
\begin{longtable}{@{}>{\raggedright\arraybackslash}p{0.1250\linewidth}>{\raggedright\arraybackslash}p{0.1875\linewidth}>{\raggedright\arraybackslash}p{0.6429\linewidth}@{}}
\toprule
\textbf{Version} & \textbf{Minimum size checked} & \textbf{Layout} \\
\midrule
\endfirsthead
\toprule
\textbf{Version} & \textbf{Minimum size checked} & \textbf{Layout} \\
\midrule
\endhead
\bottomrule
\endlastfoot
1.x (SOLO-II) & — & Not decoded; throws “unsupported X-Message format”. \\
2.1 & ≥ 40 & Short layout (below). Surface gains/offsets are copies of the profile ones. \\
2.2, 2.3 & ≥ 66 & Full layout (below). \\
2.4 & == 66 & Full layout (below). \\
\end{longtable}
\endgroup

For every version, fields beyond bytes 4–7 are decoded \textbf{only if the float-type byte = 3 (ALAMO)}. Otherwise the instrument type stays unknown and the listing prints \code{A (\allowbreak{}unknown float type)}. SOLO-II (0) and Deep Argo (1) are \code{@todo} in the decoder.

\subsubsection{\texorpdfstring{A.6.1 v2.2 / v2.3 / v2.4 — 66 bytes (MRV firmware 10.5+ “Argo packet version 0x24” [M][O])}{A.6.1 v2.2 / v2.3 / v2.4 — 66 bytes (MRV firmware 10.5+ “Argo packet version 0x24” [M][O])}}\label{sec:a61-v22--v23--v24--66-bytes-mrv-firmware-105-argo-packet-version-0x24-mo}

\begingroup\footnotesize\setlength{\tabcolsep}{3pt}
\begin{longtable}{@{}>{\raggedright\arraybackslash}p{0.0625\linewidth}>{\raggedright\arraybackslash}p{0.0446\linewidth}>{\raggedright\arraybackslash}p{0.0536\linewidth}>{\raggedright\arraybackslash}p{0.2589\linewidth}>{\raggedright\arraybackslash}p{0.1696\linewidth}>{\raggedright\arraybackslash}p{0.3304\linewidth}@{}}
\toprule
\textbf{Offset} & \textbf{Size} & \textbf{Type} & \textbf{Field} & \textbf{Scale/units} & \textbf{Notes} \\
\midrule
\endfirsthead
\toprule
\textbf{Offset} & \textbf{Size} & \textbf{Type} & \textbf{Field} & \textbf{Scale/units} & \textbf{Notes} \\
\midrule
\endhead
\bottomrule
\endlastfoot
0 & 1 & u8 & ID & — & 0xF0. \\
1 & 2 & u16 & size & bytes & 66 (0x0042). \\
3 & 1 & u4/u4 & version & major/minor & 0x22, 0x23, 0x24. \\
4 & 2 & u16 & target profile pressure & dBar & \code{prof\allowbreak{}Pres}. Observed 0 on 9287, 1000 on 11186 [O]. \\
6 & 2 & u16 & target park pressure & dBar & \code{park\allowbreak{}Pres}. \\
8 & 2 & u16 & max ascent duration & min & \code{max\allowbreak{}Asc\allowbreak{}T}. MRV: “max pump runtime during ascent (minutes)”. \\
10 & 2 & u16 & max descent-to-park duration & min & \code{max\allowbreak{}Desc2Park\allowbreak{}T}. \\
12 & 2 & u16 & max fall-to-profile duration & s & \code{max\allowbreak{}Desc2Prof\allowbreak{}T}. \\
14 & 2 & u16 & target park duration & \textbf{5 min} & MRV “divisor 0.2”; SOLO “units = 5 minutes”. The decoder lists the raw value under a \code{(\allowbreak{}min)} header: observed 288 = 24 h [O]. \\
16 & 4 & u32 & descent start time & s since 2000-01-01 & \code{start-of-descent}. \\
20 & 4 & u32 & park (drift) start time & s since 2000-01-01 & \code{start-of-drift}; the reference for 0x78 drift times. \\
24 & 4 & u32 & ascent start time & s since 2000-01-01 & \code{start-of-ascent}. \\
28 & 4 & u32 & surface start time & s since 2000-01-01 & \code{start-of-surface}. \\
32 & 1 & u8 & float type & — & 0 = SOLO-II, 1 = Deep Argo, 3 = ALAMO. Must be 3 for decoding. \\
33 & 1 & u8 & target ascent rate & cm/s & \code{tar\allowbreak{}Asc\allowbreak{}Rate}. \\
34 & 2 & u16 & max number of seeks & — & \textbf{The decoder reads only byte 34 (u8)}, so the listed \code{n\allowbreak{}Seek} is the high byte, 0 on all observed floats, where the true value is 1 [O]. \emph{(Fixed in build 3309.)} \\
36 & 2 & u16 & max surface duration & s & \code{max\allowbreak{}Surf\allowbreak{}T}. Observed raw value 1 [O]. \\
38 & 2 & u16 & seek interval & min & \code{seek\allowbreak{}Interval\allowbreak{}T}. \\
40 & 2 & u16 & pressure gain & LSB/dBar & \code{p\allowbreak{}Gain}. \\
42 & 2 & u16 & pressure offset & dBar & \code{p\allowbreak{}Off}. \\
44 & 2 & u16 & temperature gain & LSB/°C & \code{t\allowbreak{}Gain}. \\
46 & 2 & u16 & temperature offset & °C & \code{t\allowbreak{}Off}. \\
48 & 2 & u16 & salinity gain & LSB/PSU & \code{s\allowbreak{}Gain}. \\
50 & 2 & u16 & salinity offset & PSU & \code{s\allowbreak{}Off}. \\
52 & 2 & u16 & surface pressure gain & LSB/dBar & \code{p\allowbreak{}Gain\allowbreak{}Surf}. \\
54 & 2 & u16 & surface pressure offset & dBar & \code{p\allowbreak{}Off\allowbreak{}Surf}. \\
56 & 2 & u16 & surface temperature gain & LSB/°C & \code{t\allowbreak{}Gain\allowbreak{}Surf}. \\
58 & 2 & u16 & surface temperature offset & °C & \code{t\allowbreak{}Off\allowbreak{}Surf}. \\
60 & 2 & u16 & surface salinity gain & LSB/PSU & \code{s\allowbreak{}Gain\allowbreak{}Surf}. \\
62 & 2 & u16 & surface salinity offset & PSU & \code{s\allowbreak{}Off\allowbreak{}Surf}. \\
64 & 1 & u8 & mission phase number & — & \code{miss\allowbreak{}Phase} (MRV MPS 0–2). \\
65 & 1 & char & \code{;\allowbreak{}} & — &  \\
\end{longtable}
\endgroup

Example [O] (float 9287, dive 325): \code{F0 00 42 24 00 00 01 13 00 11 00 C4 0E 10 01 20 32 6F 1C A3 32 6F 4D 78 32 70 9E FB 32 70 AA 12 03 08 00 01 00 01 02 58 00 19 00 0A 03 E8 00 05 03 E8 00 01 00 19 00 0A 03 E8 00 05 03 E8 00 01 00 3B}

\begingroup\footnotesize\setlength{\tabcolsep}{3pt}
\begin{longtable}{@{}>{\raggedright\arraybackslash}p{0.1964\linewidth}>{\raggedright\arraybackslash}p{0.1875\linewidth}@{}}
\toprule
\textbf{Field} & \textbf{Value} \\
\midrule
\endfirsthead
\toprule
\textbf{Field} & \textbf{Value} \\
\midrule
\endhead
\bottomrule
\endlastfoot
Park pressure & 275 dBar \\
Max ascent & 17 min \\
Max descent to park & 196 min \\
Max fall to profile & 3600 s \\
Park duration & 288 × 5 min \\
Descent start & 2026-10-24T07:37:07Z \\
Drift start & 2026-10-24T11:05:28Z \\
Ascent start & 2026-10-25T11:05:31Z \\
Surface start & 2026-10-25T11:52:50Z \\
Float type & 3 (ALAMO) \\
Ascent rate & 8 cm/s \\
Seeks & 1 \\
Max surface duration & 1 \\
Seek interval & 600 min \\
P gain/offset & 25 / 10 \\
T gain/offset & 1000 / 5 \\
S gain/offset & 1000 / 1 \\
Surface gains/offsets & same as above \\
Mission phase & 0 \\
\end{longtable}
\endgroup

\subsubsection{\texorpdfstring{A.6.2 v2.1 — short layout (decoder only [D]; no document found)}{A.6.2 v2.1 — short layout (decoder only [D]; no document found)}}\label{sec:a62-v21--short-layout-decoder-only-d-no-document-found}

\begingroup\footnotesize\setlength{\tabcolsep}{3pt}
\begin{longtable}{@{}>{\raggedright\arraybackslash}p{0.1071\linewidth}>{\raggedright\arraybackslash}p{0.0446\linewidth}>{\raggedright\arraybackslash}p{0.0714\linewidth}>{\raggedright\arraybackslash}p{0.4018\linewidth}>{\raggedright\arraybackslash}p{0.3036\linewidth}@{}}
\toprule
\textbf{Offset} & \textbf{Size} & \textbf{Type} & \textbf{Field} & \textbf{Notes} \\
\midrule
\endfirsthead
\toprule
\textbf{Offset} & \textbf{Size} & \textbf{Type} & \textbf{Field} & \textbf{Notes} \\
\midrule
\endhead
\bottomrule
\endlastfoot
0 & 1 & u8 & ID & 0xF0 \\
1 & 2 & u16 & size & ≥ 40 required \\
3 & 1 & u4/u4 & version & 0x21 \\
4 & 2 & u16 & target profile pressure (dBar) &  \\
6 & 2 & u16 & target park pressure (dBar) &  \\
8 & 4 & u32 & descent start (s2000) &  \\
12 & 4 & u32 & park start (s2000) &  \\
16 & 4 & u32 & ascent start (s2000) &  \\
20 & 4 & u32 & surface start (s2000) &  \\
24 & 1 & u8 & float type & must be 3 \\
25 & 1 & u8 & target ascent rate (cm/s) &  \\
26–37 & 12 & 6 × u16 & P gain, P offset, T gain, T offset, S gain, S offset & Also copied to the surface gains. \\
38 … size−2 & — & — & not decoded [I] &  \\
size−1 & 1 & char & \code{;\allowbreak{}} &  \\
\end{longtable}
\endgroup

v2.1 has no durations, seeks, surface gains or mission phase. Those fields keep their “bad” values (255/65535) in the listing.

\subsubsection{\texorpdfstring{A.6.3 Other documented layouts (not decoded)}{A.6.3 Other documented layouts (not decoded)}}\label{sec:a63-other-documented-layouts-not-decoded}

\textbf{SOLO-II “v1” (37 bytes) [S].} The PDF's table is misaligned; this realignment is inferred [I]:

\begingroup\footnotesize\setlength{\tabcolsep}{3pt}
\begin{longtable}{@{}>{\raggedright\arraybackslash}p{0.0625\linewidth}>{\raggedright\arraybackslash}p{0.3482\linewidth}@{}}
\toprule
\textbf{Offset} & \textbf{Field} \\
\midrule
\endfirsthead
\toprule
\textbf{Offset} & \textbf{Field} \\
\midrule
\endhead
\bottomrule
\endlastfoot
3 & version \\
4–5 & target profile depth \\
6–7 & target park depth \\
8–9 & max rise time (min) \\
10–11 & max fall-to-park time (min) \\
12–13 & max fall-from-park-to-profile time (s) \\
14–15 & target drift time (5 min) \\
16 & float version (0) \\
17 & target ascent rate \\
18–19 & number of seeks \\
20–21 & surface time \\
22–23 & seek interval (min) \\
24–35 & P, T and S gain/offset pairs \\
36 & \code{;\allowbreak{}} \\
\end{longtable}
\endgroup

\textbf{Spreadsheet “v2.4” [X].} The SOLO-II layout plus surface gains/offsets (36–47), the mission phase (48), and an auxiliary-sensor table (49: N; then 6 bytes per sensor: type, start index, gain, offset), with size 51 + 6N. \textbf{It has no state start times.}

\begin{notebox}
\textbf{Discrepancy:} the version number 0x24 appears both in the spreadsheet and in real packets, but the layouts differ. Real 0x24 packets follow the MRV 10.5.0 layout (A.6.1, start times at 16–31, fixed 66 bytes), and so does the decoder. The spreadsheet's 0xF0 sheet does not describe what floats send.

\end{notebox}

\textbf{Listing [D]} (\code{A SN dive x\allowbreak{}Ver typ midx …}):

\begingroup\footnotesize\setlength{\tabcolsep}{3pt}
\begin{longtable}{@{}>{\raggedright\arraybackslash}p{0.0446\linewidth}>{\raggedright\arraybackslash}p{0.5714\linewidth}@{}}
\toprule
\textbf{midx} & \textbf{Fields} \\
\midrule
\endfirsthead
\toprule
\textbf{midx} & \textbf{Fields} \\
\midrule
\endhead
\bottomrule
\endlastfoot
0 & \code{miss\allowbreak{}Phase prof\allowbreak{}Pres park\allowbreak{}Pres max\allowbreak{}Asc\allowbreak{}T max\allowbreak{}Desc2Park\allowbreak{}T max\allowbreak{}Desc2Prof\allowbreak{}T} \\
1 & \code{tar\allowbreak{}Park\allowbreak{}T n\allowbreak{}Seek seek\allowbreak{}Interval\allowbreak{}T tar\allowbreak{}Asc\allowbreak{}Rate max\allowbreak{}Surf\allowbreak{}T} \\
2 & \code{start-of-descent start-of-drift} \\
3 & \code{start-of-ascent start-of-surface} \\
4 & \code{p\allowbreak{}Gain p\allowbreak{}Off t\allowbreak{}Gain t\allowbreak{}Off s\allowbreak{}Gain s\allowbreak{}Off} \\
5 & the same for the surface gains/offsets \\
\end{longtable}
\endgroup

The gains/offsets and the X version are also saved to the float record in the database.

\subsection{\texorpdfstring{A.7 Open issues found in part A}{A.7 Open issues found in part A}}\label{sec:a7-open-issues-found-in-part-a}

\begin{enumerate}
\item \textbf{Packing factor width.} Decoder and data: p nibbles, p = 0 = absent. The SOLO-II document and the spreadsheet's formula say p+1 nibbles, and the spreadsheet's per-sub-block byte counts are 8 bytes too large.
\item \textbf{B (starting sub-block).} Observed as 0 for the first part, then the cumulative sub-block count. The SOLO-II document says 1 for the first part. The decoder ignores it.
\item \textbf{0xF0 v2.4.} The spreadsheet layout (no start times, aux-sensor table) disagrees with the firmware, the MRV document and the decoder.
\item \textbf{0xF0 decoder bugs:}

\begin{itemize}
\item \code{n\allowbreak{}Seek} read as a u8 at byte 34 instead of a u16 at 34–35. \emph{(Fixed in build 3309.)}
\item \code{tar\allowbreak{}Park\allowbreak{}T} listed in 5-minute units under a “(min)” header.
\end{itemize}
\item \textbf{GPS 0x01/0x02 labels} are swapped in the spreadsheet overview.
\item \textbf{Phase codes:} SOLO-II numbering differs from ALAMO's from code 5 on. The spreadsheet has a duplicate “9”.
\item \textbf{Pump format 2} is undocumented in the spreadsheet; its single vacuum is listed by the decoder as vac1.
\item \textbf{Pump format 0 vacuum units} are unknown (SOLO-II says “counts”).
\item \textbf{Valve vacuum gain:} 10 in the spreadsheet, 100 in the decoder, and the data supports 100.
\item \textbf{Surface format 3} (per-sensor curvature blobs 0x80–0x8B) is specified [X] but not decoded. Fall/rise format 2 (9-byte with T/S) is decoded but not documented for 0x4n/0x5n.
\item \textbf{0xC2 auxiliary sensor data} is not decoded. Its byte offsets here come from the MRV YAML field order only and are unverified.
\item \textbf{0x78 and 0xDC} are undocumented. Their formats are inferred, and the meanings of 0xDC bytes 3 and 6 (always 0) and byte 5 (always 1) are unknown.
\item \textbf{Legacy first-difference decoder:} it drops a final sub-block of 1–2 values. The SOLO-II document's 22-byte sub-block size should be 27.
\item \textbf{Checksum:} documented and verified on production data. \emph{(The decoder checks it since build 3310.)}
\end{enumerate}

\clearpage
\section{\texorpdfstring{Part B — Engineering messages}{Part B — Engineering messages}}\label{sec:part-b--engineering-messages}

This part specifies the engineering message group of the ALAMO/ALTO X-message format: the SOLO-II-heritage engineering messages (0xE0, 0xE2, 0xE3), the built-in-test (BIT) messages (0xE5, 0xE6), the ALAMO dive-phase engineering messages (0xE7 descent, 0xE8 drift/park, 0xE9 ascent, 0xEA surface), the ice-detection engineering message (0xEC) and the system error codes that the 0xEA and 0xEC error lists carry. The valve message 0xEB is outside this part.

\subsection{\texorpdfstring{B.0 Conventions}{B.0 Conventions}}\label{sec:b0-conventions}

\textbf{Offsets.} Every offset in this part counts from the blob's ID byte (offset 0). Bytes 1–2 are the big-endian size \code{jj}, which counts the whole blob (ID, \code{jj}, data and the \code{;\allowbreak{}} terminator). Byte 3 of every message in this group is a format/version byte. The terminator \code{;\allowbreak{}} (0x3B) is the last byte, at offset \code{jj - 1}.

\textbf{Byte order.} All multi-byte fields are big-endian (most significant byte first). This holds for the IEEE-754 single-precision floats of BIT formats 9–11 too: the decoder assembles the four bytes big-endian into a \code{uint32} and bit-casts it to \code{float} (\code{f32} in \code{common/\allowbreak{}XDefinitions.\allowbreak{}h}).

\textbf{Types.} \code{u8}, \code{u16}, \code{u32} are unsigned integers; \code{i16} is a two's-complement signed integer; \code{f32} is an IEEE-754 binary32; \code{char[4]} is four ASCII bytes.

\textbf{Time epochs.} Two epochs are used, and each table names the one that applies:

\begingroup\footnotesize\setlength{\tabcolsep}{3pt}
\begin{longtable}{@{}>{\raggedright\arraybackslash}p{0.0536\linewidth}>{\raggedright\arraybackslash}p{0.4554\linewidth}>{\raggedright\arraybackslash}p{0.1250\linewidth}@{}}
\toprule
\textbf{Name} & \textbf{Meaning} & \textbf{Decoder macro} \\
\midrule
\endfirsthead
\toprule
\textbf{Name} & \textbf{Meaning} & \textbf{Decoder macro} \\
\midrule
\endhead
\bottomrule
\endlastfoot
\code{s2000} & \code{u32} seconds since 2000-01-01T00:00:00Z & \code{s2000(\allowbreak{}d,\allowbreak{} n)} \\
\code{s1970} & \code{u32} seconds since 1970-01-01T00:00:00Z (Unix time) & \code{s1970(\allowbreak{}d,\allowbreak{} n)} \\
\end{longtable}
\endgroup

A raw value of 0 decodes to “not a date/time” and is listed as \code{0000-00-00T00:\allowbreak{}00:\allowbreak{}00Z}.

\textbf{Sensor counts.} Fields marked \emph{CTD counts} are raw 16-bit sensor codes converted with the float's profile gains and offsets from the Argo data message (0xF0): \code{value =\allowbreak{} counts /\allowbreak{} gain − offset}. The decoder uses the profile (not surface) gains for every engineering message. A result outside the valid range (pressure −5…2610 dbar, temperature −5…40 °C, salinity 2…41 PSU) or a zero gain gives the fill value −9.999.

\textbf{Fill values.} Fields that a message version does not carry keep the decoder's fill value: \code{0x\allowbreak{}FF} (\code{u8}), \code{0x\allowbreak{}FFFF} (\code{u16}), \code{0x\allowbreak{}FFFFFFFF} (\code{u32}), \code{−32768} (\code{i16}), or not-a-date/time. Some classes leave such fields uninitialized; this is noted per message. Fill values matter only to code that reads the decoded objects: the listings print only the fields each version carries.

\textbf{Version byte.} For 0xE7–0xEC, byte 3 is the “engineering version”, written in hex (\code{0x10}, \code{0x11}, …). \code{XAlamo\allowbreak{}Engineering:\allowbreak{}:\allowbreak{}engineering\_\allowbreak{}version\_\allowbreak{}as\_\allowbreak{}float(\allowbreak{})} reads the high nibble as the major and the low nibble as the minor number (0x12 → 1.2). For 0xE5/0xE6, byte 3 is the “BIT format” (decimal 6–11). Firmware documents sometimes give the engineering version in decimal (\code{packet\allowbreak{}Version =\allowbreak{} 17} = 0x11, \code{=\allowbreak{} 20} = 0x14).

\textbf{Common validation.} Every class below first calls \code{XMessage:\allowbreak{}:\allowbreak{}Check\allowbreak{}XMsg\allowbreak{}Data(\allowbreak{})}, which:

\begin{enumerate}
\item rejects an empty blob, or one shorter than the class minimum (given per message);
\item checks the ID byte (skipped for 0xE5/0xE6, whose ID is checked later);
\item reads the size as \code{u16(\allowbreak{}d,\allowbreak{} 1) \& 0x0FFF} (only the low 12 bits are used), rejects a size below 4 or larger than the received blob, and only warns when the blob is longer than the size;
\item rejects the blob unless byte \code{size − 1} is \code{;\allowbreak{}}.
\end{enumerate}

The per-version size checks below compare the length of the received blob (\code{data.\allowbreak{}size(\allowbreak{})}), not the size field. A failed check throws \code{util:\allowbreak{}:\allowbreak{}log:\allowbreak{}:\allowbreak{}Message\allowbreak{}Exception}; the message is then not assigned to the dive and the failure is logged.

\textbf{Listing output.} Engineering lines are written in the “Engineering Data” part of the dive listing, in this order of sections: Abort, BIT failed, BIT passed, Descent, Diagnostic Dive, Park/Drift, Ascent, Ice, Standard Dive, Surface. Each data line starts with a line-type tag, the serial number (\code{SN}), the dive number and a minor index (\code{midx}); the \code{\#} lines above it are column headings and units.

\begingroup\footnotesize\setlength{\tabcolsep}{3pt}
\begin{longtable}{@{}>{\raggedright\arraybackslash}p{0.1518\linewidth}>{\raggedright\arraybackslash}p{0.3304\linewidth}>{\raggedright\arraybackslash}p{0.1964\linewidth}@{}}
\toprule
\textbf{Message} & \textbf{Listing line type(s)} & \textbf{SN formatting} \\
\midrule
\endfirsthead
\toprule
\textbf{Message} & \textbf{Listing line type(s)} & \textbf{SN formatting} \\
\midrule
\endhead
\bottomrule
\endlastfoot
0xE0, 0xE2, 0xE3 & none (an empty section heading only) & — \\
0xE5 / 0xE6 & \code{E05} / \code{E06} & \code{\%5d} (not zero-padded) \\
0xE7 & \code{E07} & \code{\%05d} \\
0xE8 & \code{E08} & \code{\%05d} \\
0xE9 & \code{E09} & \code{\%5d} \\
0xEA & \code{E0A}; error entries \code{e0A} & \code{\%05d} \\
0xEC & \code{E0C}; error entries \code{e0c} & \code{\%05d} \\
\end{longtable}
\endgroup

The Python tool \code{bin/\allowbreak{}py/\allowbreak{}xmtps-list2json.\allowbreak{}py} maps these line types to the JSON engineering sections \code{bit}, \code{descent}, \code{drift}, \code{ascent}, \code{surface} and \code{ice} (see the README).

\subsection{\texorpdfstring{B.1 SOLO-II-heritage engineering messages: 0xE0, 0xE2, 0xE3}{B.1 SOLO-II-heritage engineering messages: 0xE0, 0xE2, 0xE3}}\label{sec:b1-solo-ii-heritage-engineering-messages-0xe0-0xe2-0xe3}

\textbf{Purpose.} These are the SOLO-II engineering messages, one per mission phase: 0xE0 after the first (diagnostic) dive of a mission, 0xE2 at the end of a normal dive cycle and 0xE3 after a mission abort. The SOLO-II X-format v2.0 document (20 Mar 2014) defines them as version 5. ALAMO floats do not send them: the v2.4 format overview marks all three “ALAMO supported: no, decoder supported: no”.

\textbf{Decoder status.} \code{XEng\allowbreak{}Diagnostic} (0xE0), \code{XEng\allowbreak{}Standard} (0xE2) and \code{XEng\allowbreak{}Abort} (0xE3) are stubs. Each runs the common validation (expected ID, no minimum length beyond the size field) and decodes nothing; the listing gets the section heading (“Engineering (Diagnostic Dive)”, “Engineering (Standard Dive)”, “Engineering (Abort)”) with an empty body. The version byte is not checked.

The layouts below are transcribed from the SOLO-II v2.0 document, for reference only. In that document's text the byte column and the content column of some rows have drifted apart; the assignments marked \emph{reconstructed} were made by matching the remaining contents, in order, to the remaining byte numbers so that the total equals the stated size.

\subsubsection{\texorpdfstring{B.1.1 0xE0 — diagnostic dive (SOLO-II version 5, 80 bytes)}{B.1.1 0xE0 — diagnostic dive (SOLO-II version 5, 80 bytes)}}\label{sec:b11-0xe0--diagnostic-dive-solo-ii-version-5-80-bytes}

\begingroup\footnotesize\setlength{\tabcolsep}{3pt}
\begin{longtable}{@{}>{\raggedright\arraybackslash}p{0.0625\linewidth}>{\raggedright\arraybackslash}p{0.0446\linewidth}>{\raggedright\arraybackslash}p{0.0446\linewidth}>{\raggedright\arraybackslash}p{0.3393\linewidth}>{\raggedright\arraybackslash}p{0.1071\linewidth}>{\raggedright\arraybackslash}p{0.3214\linewidth}@{}}
\toprule
\textbf{Offset} & \textbf{Size} & \textbf{Type} & \textbf{Field} & \textbf{Scale/units} & \textbf{Notes} \\
\midrule
\endfirsthead
\toprule
\textbf{Offset} & \textbf{Size} & \textbf{Type} & \textbf{Field} & \textbf{Scale/units} & \textbf{Notes} \\
\midrule
\endhead
\bottomrule
\endlastfoot
0 & 1 & u8 & ID = 0xE0 &  &  \\
1 & 2 & u16 & size = 80 (0x50) & bytes &  \\
3 & 1 & u8 & engineering message version = 5 &  &  \\
4 & 1 & u8 & packets in current session &  &  \\
5 & 6 & — & 0 (filler) &  &  \\
11 & 2 & u16 & \code{EP->sattime} & s &  \\
13 & 2 & u16 & CPU battery voltage (\code{Vcpu}) & 0.01 V &  \\
15 & 2 & u16 & pump battery voltage at surface (\code{Vpmp}) & 0.01 V &  \\
17 & 2 & u16 & pump battery voltage at end of last pump (\code{Vple}) & 0.01 V &  \\
19 & 2 & u16 & case vacuum at beginning of BIT (\code{BTvac}) & 0.01 inHg &  \\
21 & 2 & u16 & vacuum before filling bladder at surface (\code{Air[1]}) & 0.01 inHg &  \\
23 & 2 & u16 & vacuum after filling bladder at surface (\code{Air[2]}) & 0.01 inHg &  \\
25 & 2 & u16 & ID of last interrupt (\code{ISRID}) &  &  \\
27 & 2 & u16 & average pump current at bottom (\code{HPavg\allowbreak{}I}) & 1 mA &  \\
29 & 2 & u16 & maximum pump current at bottom (\code{HPmax\allowbreak{}I}) & 1 mA &  \\
31 & 2 & u16 & total seconds pumped to surface & s &  \\
33 & 2 & u16 & seconds pumped at surface & s &  \\
35 & 2 & u16 & surface pressure before offset reset (\code{SPRX}) &  & previous dive \\
37 & 2 & u16 & pressure after offset reset (\code{SPRXL}) &  & previous dive \\
39 & 3 & u24 & pressure when “in water” sensed (\code{diag\allowbreak{}P[0]}) &  &  \\
42 & 3 & u24 & temperature when “in water” sensed (\code{diag\allowbreak{}T[0]}) &  &  \\
45 & 3 & u24 & salinity when “in water” sensed (\code{diag\allowbreak{}S[0]}) &  &  \\
48 & 2 & u16 & scans recorded by the SBE (\code{SBnscan}) &  & 0xFFFF: scan count unavailable; 0xFFFE: SBE never started \\
50 & 2 & u16 & packed SBE status: \code{(\allowbreak{}SBntry\&0x\allowbreak{}F)<<4 |\allowbreak{} (\allowbreak{}SBstrt\&0x3)<<2 |\allowbreak{} (\allowbreak{}SBstop\&0x3)} &  &  \\
52 & 3 & u24 & shallowest pressure in profile (\code{diag\allowbreak{}P[1]}) &  &  \\
55 & 3 & u24 & shallowest temperature in profile (\code{diag\allowbreak{}T[1]}) &  &  \\
58 & 3 & u24 & shallowest salinity in profile (\code{diag\allowbreak{}S[1]}) &  &  \\
61 & 2 & u16 & BIT vacuum (\code{BTvac}) & 0.01 inHg &  \\
63 & 2 & u16 & BIT motor current out (\code{BTPcur}) & 1 mA &  \\
65 & 2 & u16 & BIT pump seconds (\code{BTPsec}) & s &  \\
67 & 1 & u8 & BIT pump vacuum before pumping (\code{BTPvac[0]}) &  &  \\
68 & 1 & u8 & BIT pump vacuum after pumping (\code{BTPvac[1]}) &  &  \\
69 & 2 & u16 & BIT pump battery (\code{BTVple}) & 0.01 V &  \\
71 & 2 & u16 & BIT CPU battery (\code{BTVcpu}) & 0.01 V &  \\
73 & 2 & u16 & exception flags &  & see 0xE2 offset 89 \\
75 & 1 & u8 & vent data: tenths of seconds the vent motor ran & 0.1 s &  \\
76 & 1 & u8 & LLD status before/after vent ran &  &  \\
77 & 2 & u16 & abort code (\code{Abrt\allowbreak{}Cd}) &  & see 0xE3 offset 25 \\
79 & 1 & char & \code{;\allowbreak{}} &  &  \\
\end{longtable}
\endgroup

\subsubsection{\texorpdfstring{B.1.2 0xE2 — normal dive cycle (SOLO-II version 5, 102 bytes)}{B.1.2 0xE2 — normal dive cycle (SOLO-II version 5, 102 bytes)}}\label{sec:b12-0xe2--normal-dive-cycle-solo-ii-version-5-102-bytes}

\begingroup\footnotesize\setlength{\tabcolsep}{3pt}
\begin{longtable}{@{}>{\raggedright\arraybackslash}p{0.0625\linewidth}>{\raggedright\arraybackslash}p{0.0446\linewidth}>{\raggedright\arraybackslash}p{0.0446\linewidth}>{\raggedright\arraybackslash}p{0.3304\linewidth}>{\raggedright\arraybackslash}p{0.1071\linewidth}>{\raggedright\arraybackslash}p{0.3304\linewidth}@{}}
\toprule
\textbf{Offset} & \textbf{Size} & \textbf{Type} & \textbf{Field} & \textbf{Scale/units} & \textbf{Notes} \\
\midrule
\endfirsthead
\toprule
\textbf{Offset} & \textbf{Size} & \textbf{Type} & \textbf{Field} & \textbf{Scale/units} & \textbf{Notes} \\
\midrule
\endhead
\bottomrule
\endlastfoot
0 & 1 & u8 & ID = 0xE2 &  &  \\
1 & 2 & u16 & size = 102 (0x66) & bytes &  \\
3 & 1 & u8 & engineering message version = 5 &  &  \\
4 & 1 & u8 & packets sent in current surface session &  &  \\
5 & 2 & u16 & connection tries in previous surface session &  &  \\
7 & 2 & u16 & \code{parse\_\allowbreak{}X\_\allowbreak{}reply} &  & low byte: number of messages; high byte: error bit field \\
9 & 2 & u16 & ATSBD return status, previous session &  &  \\
11 & 2 & u16 & seconds to send all SBD messages, previous session (\code{sattime}) & s &  \\
13 & 2 & u16 & CPU battery voltage (\code{Vcpu}) & 0.01 V &  \\
15 & 2 & u16 & pump battery voltage at surface (\code{Vpmp}) & 0.01 V &  \\
17 & 2 & u16 & pump battery voltage at end of last pump (\code{Vple}) & 0.01 V &  \\
19 & 2 & u16 & case vacuum sinking at 50 dbar, oil inside (\code{Air[0]}) & 0.01 inHg &  \\
21 & 2 & u16 & case vacuum before filling oil bladder at surface (\code{Air[1]}) & 0.01 inHg &  \\
23 & 2 & u16 & case vacuum after filling bladder at surface (\code{Air[2]}) & 0.01 inHg &  \\
25 & 2 & u16 & ID of last interrupt (\code{ISRID}) &  &  \\
27 & 2 & u16 & average pump current at bottom (\code{HPavg\allowbreak{}I}) & 1 mA &  \\
29 & 2 & u16 & maximum pump current at bottom (\code{HPmax\allowbreak{}I}) & 1 mA &  \\
31 & 2 & u16 & total seconds pumped to surface & s &  \\
33 & 2 & u16 & seconds pumped at surface & s &  \\
35 & 2 & u16 & surface pressure before offset reset (\code{SPRX}) &  & previous dive \\
37 & 2 & u16 & pressure after offset reset (\code{SPRXL}) &  & previous dive \\
39 & 3 & u24 & pressure before pumping for ascent (\code{diag\allowbreak{}P[0]}) &  &  \\
42 & 3 & u24 & temperature before pumping for ascent &  &  \\
45 & 3 & u24 & salinity before pumping for ascent &  &  \\
48 & 3 & u24 & last (shallowest) pressure on ascent (\code{diag\allowbreak{}P[1]}) &  &  \\
51 & 3 & u24 & last (shallowest) temperature on ascent &  &  \\
54 & 3 & u24 & last (shallowest) salinity on ascent &  &  \\
57 & 2 & u16 & bad bins from SBE (\code{SBnbad}) &  &  \\
59 & 2 & u16 & scans recorded by SBE (\code{SBnscan}) &  & 0xFFFF / 0xFFFE as in 0xE0 \\
61 & 2 & u16 & packed SBE status (as in 0xE0 offset 50) &  &  \\
63 & 3 & u24 & average pressure, first half of drift (\code{PAVG[0]}) &  &  \\
66 & 3 & u24 & average temperature, first half of drift &  &  \\
69 & 3 & u24 & average salinity, first half of drift &  &  \\
72 & 3 & u24 & average pressure, second half of drift (\code{PAVG[1]}) &  &  \\
75 & 3 & u24 & average temperature, second half of drift &  &  \\
78 & 3 & u24 & average salinity, second half of drift &  &  \\
81 & 2 & u16 & fall time: open air valve to end of settle & s &  \\
83 & 2 & u16 & average fall rate while sinking & mm/s &  \\
85 & 2 & u16 & tenths of seconds pumped in first settle to drift (\code{Seek\allowbreak{}T}) & 0.1 s &  \\
87 & 2 & i16 & depth change in first settle (\code{Seek\allowbreak{}P}) & 0.1 dbar & signed \\
89 & 2 & u16 & exception flags &  & bit table below \\
91 & 1 & u8 & vent data: tenths of seconds the vent motor ran & 0.1 s & \emph{reconstructed} \\
92 & 1 & u8 & LLD status before and after vent ran &  & \emph{reconstructed} \\
93 & 2 & u16 & SBE pressure offset & ×800 & \emph{reconstructed} \\
95 & 2 & u16 & tenths of seconds pumped to target depth (\code{PP->Seek\allowbreak{}Sc}) & 0.1 s & \emph{reconstructed} \\
97 & 2 & u16 & packets sent in previous cycle &  & \emph{reconstructed} \\
99 & 1 & u8 & binning mode (high nibble; 2 = controller binning, else CTD binning), subcycle number (low nibble) &  & \emph{reconstructed} \\
100 & 1 & u8 & ice-detect status: 0 \code{ICE\_\allowbreak{}CHECK\_\allowbreak{}OFF}, 1 \code{ICE\_\allowbreak{}NOT\_\allowbreak{}FOUND}, 2 \code{ICE\_\allowbreak{}MIXLAYER}, 3 \code{ICE\_\allowbreak{}BREAKUP} &  & \emph{reconstructed} \\
101 & 1 & char & \code{;\allowbreak{}} &  &  \\
\end{longtable}
\endgroup

0xE2 exception flags (offset 89):

\begingroup\footnotesize\setlength{\tabcolsep}{3pt}
\begin{longtable}{@{}>{\raggedright\arraybackslash}p{0.0625\linewidth}>{\raggedright\arraybackslash}p{0.5268\linewidth}@{}}
\toprule
\textbf{Bit} & \textbf{Meaning} \\
\midrule
\endfirsthead
\toprule
\textbf{Bit} & \textbf{Meaning} \\
\midrule
\endhead
\bottomrule
\endlastfoot
0x0001 & valve failed to open \\
0x0002 & valve failed to close \\
0x0004 & questionable pressure \\
0x0008 & antenna was toggled \\
0x0010 & antenna switch failure (no satellites even after toggling) \\
0x0020 & GPS communication error \\
0x0080 & float took too long to leave the surface (valve toggled) \\
0x1000 & valve failure during sink phase \\
0x2000 & valve failure during ascent phase \\
\end{longtable}
\endgroup

\subsubsection{\texorpdfstring{B.1.3 0xE3 — after mission abort (SOLO-II version 5, 30 bytes)}{B.1.3 0xE3 — after mission abort (SOLO-II version 5, 30 bytes)}}\label{sec:b13-0xe3--after-mission-abort-solo-ii-version-5-30-bytes}

\begingroup\footnotesize\setlength{\tabcolsep}{3pt}
\begin{longtable}{@{}>{\raggedright\arraybackslash}p{0.0625\linewidth}>{\raggedright\arraybackslash}p{0.0446\linewidth}>{\raggedright\arraybackslash}p{0.0446\linewidth}>{\raggedright\arraybackslash}p{0.4196\linewidth}>{\raggedright\arraybackslash}p{0.1071\linewidth}>{\raggedright\arraybackslash}p{0.2411\linewidth}@{}}
\toprule
\textbf{Offset} & \textbf{Size} & \textbf{Type} & \textbf{Field} & \textbf{Scale/units} & \textbf{Notes} \\
\midrule
\endfirsthead
\toprule
\textbf{Offset} & \textbf{Size} & \textbf{Type} & \textbf{Field} & \textbf{Scale/units} & \textbf{Notes} \\
\midrule
\endhead
\bottomrule
\endlastfoot
0 & 1 & u8 & ID = 0xE3 &  &  \\
1 & 2 & u16 & size = 30 (0x1E) & bytes &  \\
3 & 1 & u8 & engineering message version = 5 &  &  \\
4 & 1 & u8 & packets sent in current surface session &  &  \\
5 & 2 & u16 & connection tries in last surface session &  &  \\
7 & 2 & u16 & \code{parse\_\allowbreak{}X\_\allowbreak{}reply} &  & as in 0xE2 \\
9 & 2 & u16 & ATSBD return status, last session &  &  \\
11 & 2 & u16 & seconds taken to send last SBD message & s &  \\
13 & 2 & u16 & current CPU battery voltage & 0.01 V &  \\
15 & 2 & u16 & current pump battery voltage & 0.01 V &  \\
17 & 2 & u16 & case vacuum at beginning of abort (\code{Air[1]}) & 0.01 inHg &  \\
19 & 2 & u16 & case vacuum at end of last transmission (\code{Air[0]}) & 0.01 inHg & previous cycle \\
21 & 2 & — & not described &  & gap in the source document \\
23 & 2 & u16 & ID of last interrupt (\code{ISRID}) &  &  \\
25 & 2 & u16 & abort code (\code{Abrt\allowbreak{}Cd}) &  & table below \\
27 & 2 & — & empty &  &  \\
29 & 1 & char & \code{;\allowbreak{}} &  &  \\
\end{longtable}
\endgroup

Abort codes: 0 no error; 1 current time later than \code{RTCabort}; 2 unable to wake OST; 3 unable to send dive number to SOLO-II; 4 Iridium ground station commanded abort; 5 final dive completed, mission done; 6 diagnostic dive failed (no GPS fix, pressure never exceeded \code{d\allowbreak{}Bar\allowbreak{}Go}, or no Iridium message sent); 7 pressure sensor failure.

\subsubsection{\texorpdfstring{B.1.4 SOLO-II BIT messages (0xE5/0xE6 version 5)}{B.1.4 SOLO-II BIT messages (0xE5/0xE6 version 5)}}\label{sec:b14-solo-ii-bit-messages-0xe50xe6-version-5}

SOLO-II uses 0xE5 (BIT passed, 58 bytes) and 0xE6 (BIT failed, 60 bytes), version 5. The decoder does \textbf{not} support them: \code{XEng\allowbreak{}BIT} accepts only BIT formats 6–11 and rejects version 5 with “Unsupported BIT message format: 5!”. The layouts are given for reference; 0xE6 is 0xE5 with a 2-byte BIT status register inserted at offset 5.

\begingroup\footnotesize\setlength{\tabcolsep}{3pt}
\begin{longtable}{@{}>{\raggedright\arraybackslash}p{0.1071\linewidth}>{\raggedright\arraybackslash}p{0.1071\linewidth}>{\raggedright\arraybackslash}p{0.0446\linewidth}>{\raggedright\arraybackslash}p{0.0804\linewidth}>{\raggedright\arraybackslash}p{0.4732\linewidth}>{\raggedright\arraybackslash}p{0.1071\linewidth}@{}}
\toprule
\textbf{0xE5 offset} & \textbf{0xE6 offset} & \textbf{Size} & \textbf{Type} & \textbf{Field} & \textbf{Scale/units} \\
\midrule
\endfirsthead
\toprule
\textbf{0xE5 offset} & \textbf{0xE6 offset} & \textbf{Size} & \textbf{Type} & \textbf{Field} & \textbf{Scale/units} \\
\midrule
\endhead
\bottomrule
\endlastfoot
0 & 0 & 1 & u8 & ID (0xE5 / 0xE6) &  \\
1 & 1 & 2 & u16 & size (58 / 60) & bytes \\
3 & 3 & 1 & u8 & engineering message version = 5 &  \\
4 & 4 & 1 & u8 & packets sent in this surface session &  \\
— & 5 & 2 & u16 & BIT status register &  \\
5 & 7 & 2 & u16 & SBE pressure offset & ×800 \\
7 & 9 & 2 & u16 & CPU battery voltage & 0.01 V \\
9 & 11 & 2 & u16 & no-load pump battery voltage & 0.01 V \\
11 & 13 & 2 & u16 & pump battery voltage at end of last pump & 0.01 V \\
13 & 15 & 2 & u16 & average pump current (\code{HPavg\allowbreak{}I}) & 1 mA \\
15 & 17 & 2 & u16 & seconds pumped out during test & s \\
17 & 19 & 1 & u8 & oil vacuum before filling bladder & 0.01 inHg \\
18 & 20 & 1 & u8 & oil vacuum after filling bladder & 0.01 inHg \\
19 & 21 & 2 & u16 & case vacuum at beginning of BIT, oil bladder empty (\code{Air[0]}) & 0.01 inHg \\
21 & 23 & 2 & u16 & case vacuum at end of BIT, air bladder inflated (\code{Air[1]}) & 0.01 inHg \\
23 & 25 & 1 & u8 & tries needed to open valve &  \\
24 & 26 & 1 & u8 & tries needed to close valve &  \\
25 & 27 & 2 & u16 & ID of last interrupt &  \\
27 & 29 & 30 & char[30] & string returned from the SBE \code{pt} command &  \\
57 & 59 & 1 & char & \code{;\allowbreak{}} &  \\
\end{longtable}
\endgroup

The 0xE5 rows from offset 13 on are \emph{reconstructed} (the source text is misaligned there); they agree with the 0xE6 table, which is aligned.

\subsection{\texorpdfstring{B.2 Built-in-test messages: 0xE5 (passed) and 0xE6 (failed)}{B.2 Built-in-test messages: 0xE5 (passed) and 0xE6 (failed)}}\label{sec:b2-built-in-test-messages-0xe5-passed-and-0xe6-failed}

\textbf{Purpose.} An ALAMO/ALTO float runs a built-in self-test (BIT) at deployment. It exercises the air and oil pumps and valves, the CTD (resetting the pressure offset), the case vacuum sensor, the battery, GPS and Iridium, and reports the results once, as 0xE5 when every test passed or 0xE6 when at least one failed. The two IDs share one layout; only the ID differs.

\textbf{Dispatch.} \code{XEng\allowbreak{}BITPass} and \code{XEng\allowbreak{}BITFail} both derive from \code{XEng\allowbreak{}BIT}, which runs the common validation with a 13-byte minimum (ID check skipped) and dispatches on byte 3, the BIT format:

\begingroup\footnotesize\setlength{\tabcolsep}{3pt}
\begin{longtable}{@{}>{\raggedright\arraybackslash}p{0.1786\linewidth}>{\raggedright\arraybackslash}p{0.1250\linewidth}>{\raggedright\arraybackslash}p{0.1161\linewidth}>{\raggedright\arraybackslash}p{0.4464\linewidth}@{}}
\toprule
\textbf{BIT format (byte 3)} & \textbf{Decoder class} & \textbf{Size (bytes)} & \textbf{Firmware (from documents)} \\
\midrule
\endfirsthead
\toprule
\textbf{BIT format (byte 3)} & \textbf{Decoder class} & \textbf{Size (bytes)} & \textbf{Firmware (from documents)} \\
\midrule
\endhead
\bottomrule
\endlastfoot
6 & \code{XAlamo\allowbreak{}Bit} & 13 & early ALAMO (spreadsheet history) \\
7 & \code{XAlamo\allowbreak{}Bit} & 72 & X-format v2.1–2.4 \\
8 & \code{XAlamo\allowbreak{}Bit} & 76 & X-format v2.4; firmware 7.3.5 \\
9 & \code{XAlamo\allowbreak{}Bit2} & 256 & firmware 10.5.0 \\
10 & \code{XAlamo\allowbreak{}Bit2} & 256 & firmware 10.6.x (same layout as 9) \\
11 & \code{XAlamo\allowbreak{}Bit2} & 256 or 264 & firmware 10.8 and later (observed from 10.12.0.1) \\
other & — & — & rejected: “Unsupported BIT message format” \\
\end{longtable}
\endgroup

Both classes re-check that byte 0 is 0xE5 or 0xE6 and set the passed/failed flag from it.

\subsubsection{\texorpdfstring{B.2.1 BIT format 6 (13 bytes)}{B.2.1 BIT format 6 (13 bytes)}}\label{sec:b21-bit-format-6-13-bytes}

\begingroup\footnotesize\setlength{\tabcolsep}{3pt}
\begin{longtable}{@{}>{\raggedright\arraybackslash}p{0.0625\linewidth}>{\raggedright\arraybackslash}p{0.0446\linewidth}>{\raggedright\arraybackslash}p{0.0446\linewidth}>{\raggedright\arraybackslash}p{0.2054\linewidth}>{\raggedright\arraybackslash}p{0.1071\linewidth}>{\raggedright\arraybackslash}p{0.1250\linewidth}@{}}
\toprule
\textbf{Offset} & \textbf{Size} & \textbf{Type} & \textbf{Field} & \textbf{Scale/units} & \textbf{Notes} \\
\midrule
\endfirsthead
\toprule
\textbf{Offset} & \textbf{Size} & \textbf{Type} & \textbf{Field} & \textbf{Scale/units} & \textbf{Notes} \\
\midrule
\endhead
\bottomrule
\endlastfoot
0 & 1 & u8 & ID = 0xE5 / 0xE6 &  &  \\
1 & 2 & u16 & size = 13 & bytes &  \\
3 & 1 & u8 & BIT format = 6 &  &  \\
4 & 4 & u32 & float software version &  & listed in hex \\
8 & 4 & u32 & BIT results mask &  & B.2.6 \\
12 & 1 & char & \code{;\allowbreak{}} &  &  \\
\end{longtable}
\endgroup

The blob must be exactly 13 bytes.

\subsubsection{\texorpdfstring{B.2.2 BIT format 7 (72 bytes)}{B.2.2 BIT format 7 (72 bytes)}}\label{sec:b22-bit-format-7-72-bytes}

\begingroup\footnotesize\setlength{\tabcolsep}{3pt}
\begin{longtable}{@{}>{\raggedright\arraybackslash}p{0.0625\linewidth}>{\raggedright\arraybackslash}p{0.0446\linewidth}>{\raggedright\arraybackslash}p{0.0446\linewidth}>{\raggedright\arraybackslash}p{0.2946\linewidth}>{\raggedright\arraybackslash}p{0.1786\linewidth}>{\raggedright\arraybackslash}p{0.2946\linewidth}@{}}
\toprule
\textbf{Offset} & \textbf{Size} & \textbf{Type} & \textbf{Field} & \textbf{Scale/units} & \textbf{Notes} \\
\midrule
\endfirsthead
\toprule
\textbf{Offset} & \textbf{Size} & \textbf{Type} & \textbf{Field} & \textbf{Scale/units} & \textbf{Notes} \\
\midrule
\endhead
\bottomrule
\endlastfoot
0 & 1 & u8 & ID = 0xE5 / 0xE6 &  &  \\
1 & 2 & u16 & size = 72 & bytes &  \\
3 & 1 & u8 & BIT format = 7 &  &  \\
4 & 4 & u32 & float software version &  &  \\
8 & 4 & u32 & BIT results mask &  & B.2.6 \\
12 & 2 & u16 & total BIT duration & s &  \\
14 & 2 & u16 & pressure & CTD counts &  \\
16 & 2 & u16 & temperature & CTD counts &  \\
18 & 2 & u16 & salinity & CTD counts &  \\
20 & 4 & u32 & CTD status &  & listed in hex \\
24 & 2 & u16 & case vacuum at start of BIT & ÷100 → inHg &  \\
26 & 2 & u16 & case vacuum at end of BIT & ÷100 → inHg &  \\
28 & 2 & u16 & battery voltage at start of BIT & ÷100 → V &  \\
30 & 2 & u16 & battery voltage at end of BIT & ÷100 → V &  \\
32 & 2 & u16 & air valve open position & ÷1000 → V &  \\
34 & 2 & u16 & air valve open calibration position & ÷1000 → V &  \\
36 & 2 & u16 & air valve closed position & ÷1000 → V &  \\
38 & 2 & u16 & air valve closed calibration position & ÷1000 → V &  \\
40 & 2 & u16 & air valve middle position & ÷1000 → V &  \\
42 & 2 & u16 & air valve middle calibration position & ÷1000 → V &  \\
44 & 2 & u16 & air pump test duration & s &  \\
46 & 1 & u8 & air pump cycle count &  &  \\
47 & 2 & u16 & maximum air pressure & ÷1000 → psi &  \\
49 & 2 & u16 & air pump maximum load current & ÷1000 → A &  \\
51 & 2 & u16 & air pump maximum load 12 V current & ÷1000 → A &  \\
53 & 2 & u16 & air pump maximum load battery voltage & ÷1000 → V (decoder) & spreadsheet and 7.3.5 YAML: ÷100 \\
55 & 2 & u16 & oil valve open position & ÷1000 → V &  \\
57 & 2 & u16 & oil valve open calibration position & ÷1000 → V &  \\
59 & 2 & u16 & oil valve closed position & ÷1000 → V &  \\
61 & 2 & u16 & oil valve closed calibration position & ÷1000 → V &  \\
63 & 2 & u16 & oil pump test duration & s & spreadsheet: gain 1000 \\
65 & 2 & u16 & oil pump maximum load current & ÷1000 → A &  \\
67 & 2 & u16 & oil pump maximum load 12 V current & ÷1000 → A &  \\
69 & 2 & u16 & oil pump maximum load battery voltage & ÷1000 → V & 7.3.5 YAML: ÷100 \\
71 & 1 & char & \code{;\allowbreak{}} &  &  \\
\end{longtable}
\endgroup

The blob must be exactly 72 bytes.

\subsubsection{\texorpdfstring{B.2.3 BIT format 8 (76 bytes)}{B.2.3 BIT format 8 (76 bytes)}}\label{sec:b23-bit-format-8-76-bytes}

Format 8 replaces format 7's single pressure with the pressure before the offset reset, the offset reading and the pressure after the reset; every later field moves 4 bytes on.

\begingroup\footnotesize\setlength{\tabcolsep}{3pt}
\begin{longtable}{@{}>{\raggedright\arraybackslash}p{0.0625\linewidth}>{\raggedright\arraybackslash}p{0.0446\linewidth}>{\raggedright\arraybackslash}p{0.0446\linewidth}>{\raggedright\arraybackslash}p{0.2857\linewidth}>{\raggedright\arraybackslash}p{0.1786\linewidth}>{\raggedright\arraybackslash}p{0.3036\linewidth}@{}}
\toprule
\textbf{Offset} & \textbf{Size} & \textbf{Type} & \textbf{Field} & \textbf{Scale/units} & \textbf{Notes} \\
\midrule
\endfirsthead
\toprule
\textbf{Offset} & \textbf{Size} & \textbf{Type} & \textbf{Field} & \textbf{Scale/units} & \textbf{Notes} \\
\midrule
\endhead
\bottomrule
\endlastfoot
0 & 1 & u8 & ID = 0xE5 / 0xE6 &  &  \\
1 & 2 & u16 & size = 76 & bytes & spreadsheet: “77” (its terminator offset gives 76) \\
3 & 1 & u8 & BIT format = 8 &  &  \\
4 & 4 & u32 & float software version &  & 7.3.5 YAML: four \code{u8} — major, minor, patch, unused \\
8 & 4 & u32 & BIT results mask &  & B.2.6 \\
12 & 2 & u16 & total BIT duration & s &  \\
14 & 2 & u16 & pressure before offset reset & CTD counts & 7.3.5 YAML: \code{i16}, surface gains \\
16 & 2 & u16 & pressure sensor offset reading & ÷1000 → dbar & spreadsheet (0xE6) and YAML: signed \code{i16}; decoder reads unsigned \\
18 & 2 & u16 & pressure after offset reset & CTD counts & 7.3.5 YAML: \code{i16}, surface gains \\
20 & 2 & u16 & temperature & CTD counts &  \\
22 & 2 & u16 & salinity & CTD counts &  \\
24 & 4 & u32 & CTD status &  & YAML: 4 ASCII characters \\
28 & 2 & u16 & case vacuum at start of BIT & ÷100 → inHg &  \\
30 & 2 & u16 & case vacuum at end of BIT & ÷100 → inHg &  \\
32 & 2 & u16 & battery voltage at start of BIT & ÷100 → V &  \\
34 & 2 & u16 & battery voltage at end of BIT & ÷100 → V &  \\
36 & 2 & u16 & air valve open position & ÷1000 → V &  \\
38 & 2 & u16 & air valve open calibration position & ÷1000 → V &  \\
40 & 2 & u16 & air valve closed position & ÷1000 → V &  \\
42 & 2 & u16 & air valve closed calibration position & ÷1000 → V &  \\
44 & 2 & u16 & air valve middle position & ÷1000 → V &  \\
46 & 2 & u16 & air valve middle calibration position & ÷1000 → V &  \\
48 & 2 & u16 & air pump test duration & s &  \\
50 & 1 & u8 & air pump cycle count &  &  \\
51 & 2 & u16 & maximum air pressure & ÷1000 → psi &  \\
53 & 2 & u16 & air pump maximum load current & ÷1000 → A &  \\
55 & 2 & u16 & air pump maximum load 12 V current & ÷1000 → A &  \\
57 & 2 & u16 & air pump maximum load battery voltage & ÷1000 → V (decoder) & spreadsheet and YAML: ÷100 \\
59 & 2 & u16 & oil valve open position & ÷1000 → V &  \\
61 & 2 & u16 & oil valve open calibration position & ÷1000 → V &  \\
63 & 2 & u16 & oil valve closed position & ÷1000 → V &  \\
65 & 2 & u16 & oil valve closed calibration position & ÷1000 → V &  \\
67 & 2 & u16 & oil pump test duration & s & spreadsheet: gain 1000; YAML: 1 \\
69 & 2 & u16 & oil pump maximum load current & ÷1000 → A &  \\
71 & 2 & u16 & oil pump maximum load 12 V current & ÷1000 → A &  \\
73 & 2 & u16 & oil pump maximum load battery voltage & ÷1000 → V & YAML: ÷100 \\
75 & 1 & char & \code{;\allowbreak{}} &  &  \\
\end{longtable}
\endgroup

The blob must be exactly 76 bytes. The v2.4 spreadsheet's 0xE6 format 8 sheet makes the CTD status 16 bits (offsets 24–25) and shifts every later field 2 bytes earlier; this disagrees with its own 0xE5 sheet, the 7.3.5 YAML and the decoder, which all use a 4-byte status at 24–27 for both IDs. The spreadsheet also swaps the format 7/8 sizes on the 0xE6 sheet.

\subsubsection{\texorpdfstring{B.2.4 BIT formats 9 and 10 (256 bytes)}{B.2.4 BIT formats 9 and 10 (256 bytes)}}\label{sec:b24-bit-formats-9-and-10-256-bytes}

From format 9 on, measured values are sent as IEEE-754 floats in engineering units, each subsystem gets its own error bitfield, and the test thresholds are appended.

\begingroup\footnotesize\setlength{\tabcolsep}{3pt}
\begin{longtable}{@{}>{\raggedright\arraybackslash}p{0.0625\linewidth}>{\raggedright\arraybackslash}p{0.0446\linewidth}>{\raggedright\arraybackslash}p{0.0804\linewidth}>{\raggedright\arraybackslash}p{0.3036\linewidth}>{\raggedright\arraybackslash}p{0.1250\linewidth}>{\raggedright\arraybackslash}p{0.3036\linewidth}@{}}
\toprule
\textbf{Offset} & \textbf{Size} & \textbf{Type} & \textbf{Field} & \textbf{Scale/units} & \textbf{Notes} \\
\midrule
\endfirsthead
\toprule
\textbf{Offset} & \textbf{Size} & \textbf{Type} & \textbf{Field} & \textbf{Scale/units} & \textbf{Notes} \\
\midrule
\endhead
\bottomrule
\endlastfoot
0 & 1 & u8 & ID = 0xE5 / 0xE6 &  &  \\
1 & 2 & u16 & size = 256 & bytes &  \\
3 & 1 & u8 & BIT format = 9 or 10 &  &  \\
4 & 4 & u32 & firmware version &  & bytes: major, minor, patch, product (e.g. 0x0A0C0001 = 10.12.0.1) \\
8 & 1 & u8 & BIT type &  & B.2.7; decoded, not listed \\
9 & 1 & u8 & CTD (sensor) type &  & B.2.7 \\
10 & 4 & u32 & BIT results mask &  & B.2.6 \\
14 & 4 & u32 & air pump errors &  & B.2.6 \\
18 & 4 & u32 & oil pump errors &  & B.2.6 \\
22 & 4 & u32 & air valve errors &  & decoded, not listed \\
26 & 4 & u32 & oil valve errors &  & decoded, not listed \\
30 & 4 & u32 & GPS errors &  & decoded, not listed \\
34 & 4 & u32 & CTD errors &  & B.2.6 \\
38 & 4 & u32 & satellite (Iridium) errors &  & decoded, not listed \\
42 & 4 & u32 & case pressure (vacuum) errors &  &  \\
46 & 4 & u32 & battery errors &  &  \\
50 & 4 & u32 & BIT start time & \code{s1970} &  \\
54 & 4 & u32 & BIT end time & \code{s1970} & duration = end − start \\
58 & 4 & f32 & sea pressure before calibration & dbar &  \\
62 & 4 & f32 & new sea pressure offset & dbar &  \\
66 & 4 & f32 & sea pressure after calibration & dbar &  \\
70 & 4 & f32 & temperature & °C &  \\
74 & 4 & f32 & salinity & PSU &  \\
78 & 4 & char[4] & CTD status &  & decoded as \code{u32}, listed in hex (e.g. 0x64697361 = “disa”) \\
82 & 4 & f32 & case vacuum at start & inHg &  \\
86 & 4 & f32 & case vacuum at end & inHg &  \\
90 & 4 & f32 & battery voltage at start & V &  \\
94 & 4 & f32 & battery voltage at end & V &  \\
98 & 4 & f32 & air valve open position & V &  \\
102 & 4 & f32 & air valve open calibration position & V &  \\
106 & 4 & f32 & air valve closed position & V &  \\
110 & 4 & f32 & air valve closed calibration position & V &  \\
114 & 4 & f32 & air valve middle position & V &  \\
118 & 4 & f32 & air valve middle calibration position & V &  \\
122 & 2 & u16 & air pump test duration & s &  \\
124 & 1 & u8 & air pump test cycle count &  &  \\
125 & 4 & f32 & maximum air pressure & psi (decoder) & 10.5.0 blob list: “inHg” \\
129 & 4 & f32 & maximum air pump current & A &  \\
133 & 4 & f32 & average air pump current & A &  \\
137 & 4 & f32 & maximum air pump 12 V current & A &  \\
141 & 4 & f32 & maximum battery voltage during air pump test & V &  \\
145 & 4 & f32 & oil valve open position & V &  \\
149 & 4 & f32 & oil valve open calibration position & V &  \\
153 & 4 & f32 & oil valve closed position & V &  \\
157 & 4 & f32 & oil valve closed calibration position & V &  \\
161 & 2 & u16 & oil pump test duration & s &  \\
163 & 4 & f32 & maximum oil pump current & A &  \\
167 & 4 & f32 & average oil pump current & A &  \\
171 & 4 & f32 & maximum oil pump 12 V current & A &  \\
175 & 4 & f32 & maximum battery voltage during oil pump test & V &  \\
179 & 76 & f32 × 19 & thresholds &  & B.2.8; not decoded \\
255 & 1 & char & \code{;\allowbreak{}} &  &  \\
\end{longtable}
\endgroup

The 10.5.0 blob list gives “divisor” values (1000, 100) for many of these floats; transmitted values are already in engineering units (for example battery 15.51 V, case vacuum 6.07 inHg, air valve 1.274 V in a 10.12.0.1 BIT), and the decoder uses them unscaled.

\subsubsection{\texorpdfstring{B.2.5 BIT format 11 (264 or 256 bytes)}{B.2.5 BIT format 11 (264 or 256 bytes)}}\label{sec:b25-bit-format-11-264-or-256-bytes}

Format 11 adds the traveler errors at offset 50; every later field moves 4 bytes on. The 10.12.0 definition also adds a twentieth threshold (oil pump current low warning), giving 264 bytes.

\begingroup\footnotesize\setlength{\tabcolsep}{3pt}
\begin{longtable}{@{}>{\raggedright\arraybackslash}p{0.0625\linewidth}>{\raggedright\arraybackslash}p{0.0446\linewidth}>{\raggedright\arraybackslash}p{0.0804\linewidth}>{\raggedright\arraybackslash}p{0.2857\linewidth}>{\raggedright\arraybackslash}p{0.1071\linewidth}>{\raggedright\arraybackslash}p{0.3393\linewidth}@{}}
\toprule
\textbf{Offset} & \textbf{Size} & \textbf{Type} & \textbf{Field} & \textbf{Scale/units} & \textbf{Notes} \\
\midrule
\endfirsthead
\toprule
\textbf{Offset} & \textbf{Size} & \textbf{Type} & \textbf{Field} & \textbf{Scale/units} & \textbf{Notes} \\
\midrule
\endhead
\bottomrule
\endlastfoot
0–49 &  &  & as formats 9/10 (ID through battery errors) &  &  \\
50 & 4 & u32 & traveler errors &  & added in 10.8; B.2.6; decoded, not listed \\
54 & 4 & u32 & BIT start time & \code{s1970} &  \\
58 & 4 & u32 & BIT end time & \code{s1970} &  \\
62 & 4 & f32 & sea pressure before calibration & dbar &  \\
66 & 4 & f32 & new sea pressure offset & dbar &  \\
70 & 4 & f32 & sea pressure after calibration & dbar &  \\
74 & 4 & f32 & temperature & °C &  \\
78 & 4 & f32 & salinity & PSU &  \\
82 & 4 & char[4] & CTD status &  &  \\
86 & 4 & f32 & case vacuum at start & inHg &  \\
90 & 4 & f32 & case vacuum at end & inHg &  \\
94 & 4 & f32 & battery voltage at start & V &  \\
98 & 4 & f32 & battery voltage at end & V &  \\
102 & 4 & f32 & air valve open position & V &  \\
106 & 4 & f32 & air valve open calibration position & V &  \\
110 & 4 & f32 & air valve closed position & V &  \\
114 & 4 & f32 & air valve closed calibration position & V &  \\
118 & 4 & f32 & air valve middle position & V &  \\
122 & 4 & f32 & air valve middle calibration position & V &  \\
126 & 2 & u16 & air pump test duration & s &  \\
128 & 1 & u8 & air pump test cycle count &  &  \\
129 & 4 & f32 & maximum air pressure & psi &  \\
133 & 4 & f32 & maximum air pump current & A &  \\
137 & 4 & f32 & average air pump current & A &  \\
141 & 4 & f32 & maximum air pump 12 V current & A & 10.12.0 list names this “MaxAirPumpCurrent” (see B.10) \\
145 & 4 & f32 & maximum battery voltage during air pump test & V & 10.12.0 list: “MaxAirV12Current” \\
149 & 4 & f32 & oil valve open position & V &  \\
153 & 4 & f32 & oil valve open calibration position & V &  \\
157 & 4 & f32 & oil valve closed position & V &  \\
161 & 4 & f32 & oil valve closed calibration position & V &  \\
165 & 2 & u16 & oil pump test duration & s &  \\
167 & 4 & f32 & maximum oil pump current & A &  \\
171 & 4 & f32 & average oil pump current & A &  \\
175 & 4 & f32 & maximum oil pump 12 V current & A & 10.12.0 list: “MaxOilPumpCurrent” \\
179 & 4 & f32 & maximum battery voltage during oil pump test & V & 10.12.0 list: “MaxOilV12Current” \\
183 & 80 & f32 × 20 & thresholds (264-byte form) &  & B.2.8; not decoded \\
263 & 1 & char & \code{;\allowbreak{}} &  & 264-byte form \\
\end{longtable}
\endgroup

The decoder also accepts a 256-byte format 11. No document describes it. With the fields above ending at offset 182 and the terminator at 255, it would have room for 18 thresholds (183–254), one fewer than the 19 of format 9; \emph{inferred}, unverified. The decoder reads nothing beyond offset 182 in either form.

\subsubsection{\texorpdfstring{B.2.6 BIT bitfields}{B.2.6 BIT bitfields}}\label{sec:b26-bit-bitfields}

\textbf{BIT results mask} (all formats; offset 8 in formats 6–8, offset 10 in 9–11). The spreadsheet gives bits 0x01–0x40 (\code{BIST\_\allowbreak{}RESULT\_\allowbreak{}*}); 10.5.0 and 10.12.0 add the rest.

\begingroup\footnotesize\setlength{\tabcolsep}{3pt}
\begin{longtable}{@{}>{\raggedright\arraybackslash}p{0.0982\linewidth}>{\raggedright\arraybackslash}p{0.3214\linewidth}>{\raggedright\arraybackslash}p{0.1875\linewidth}@{}}
\toprule
\textbf{Bit} & \textbf{Name (spreadsheet / 10.5.0)} & \textbf{10.12.0 name} \\
\midrule
\endfirsthead
\toprule
\textbf{Bit} & \textbf{Name (spreadsheet / 10.5.0)} & \textbf{10.12.0 name} \\
\midrule
\endhead
\bottomrule
\endlastfoot
0x00000000 & \code{BIST\_\allowbreak{}RESULT\_\allowbreak{}OK} &  \\
0x00000001 & air pump failure & AirPumpErrors \\
0x00000002 & oil pump failure & OilPumpErrors \\
0x00000004 & air valve failure & AirValveErrors \\
0x00000008 & oil valve failure & OilValveErrors \\
0x00000010 & GPS failure & GPSErrors \\
0x00000020 & CTD sensor failure & CTDSensorErrors \\
0x00000040 & satellite failure & SatModemErrors \\
0x00000080 & case pressure failure (10.5.0) & CasePressureErrors \\
0x00000100 & battery voltage failure (10.5.0) & BatteryVoltageErrors \\
0x00000200 & accelerometer failure (10.5.0 only) & — \\
0x00000400 & undocumented & — \\
\end{longtable}
\endgroup

Bit 0x400 is set in a format 11 BIT failure from firmware 10.12.0.1 (results 0x00000401); it is \emph{inferred} to flag the traveler errors added in format 11.

\textbf{Per-subsystem error bitfields} (formats 9–11):

\begingroup\footnotesize\setlength{\tabcolsep}{3pt}
\begin{longtable}{@{}>{\raggedright\arraybackslash}p{0.2500\linewidth}>{\raggedright\arraybackslash}p{0.0446\linewidth}>{\raggedright\arraybackslash}p{0.2589\linewidth}>{\raggedright\arraybackslash}p{0.2143\linewidth}@{}}
\toprule
\textbf{Field} & \textbf{Bit} & \textbf{10.5.0 meaning} & \textbf{10.12.0 meaning} \\
\midrule
\endfirsthead
\toprule
\textbf{Field} & \textbf{Bit} & \textbf{10.5.0 meaning} & \textbf{10.12.0 meaning} \\
\midrule
\endhead
\bottomrule
\endlastfoot
air pump errors & 0x01 & air pump sensor pressure low & lowPressure \\
 & 0x02 & air pump current high & hiCurrent \\
 & 0x04 & air pump cycle count & hiCycleCount \\
oil pump errors & 0x01 & oil pump current low & loCurrent \\
 & 0x02 & oil pump current high & hiCurrentWarning \\
 & 0x04 & oil pump current high & hiCurrentError \\
air valve errors & 0x01 & open fail & OpenFail \\
 & 0x02 & close fail & CloseFail \\
 & 0x04 & middle fail & MiddleFail \\
oil valve errors & 0x01 & open fail & OpenFail \\
 & 0x02 & close fail & CloseFail \\
GPS errors & 0x01 & GPS communication failure & NoDatafromGPS \\
 & 0x02 & GPS no fix & FixAttempt1Failed \\
 & 0x04 & GPS failure & FixAttempt2Failed \\
 & 0x08 & — & OtherError \\
CTD errors & 0x01 & unable to start profile & UnableToStartProfile \\
 & 0x02 & unable to communicate & NoDataFromCTD \\
 & 0x04 & unable to profile & UnabletoRetrieveProfile \\
 & 0x08 & bad data & FieldsOutsideThresholds \\
satellite errors & 0x01 & communication failure & NoDataFromModem \\
 & 0x02 & first transmit failure & TransmitFailure \\
case pressure errors & 0x01 & case pressure low & CasePressureLow \\
 & 0x02 & case pressure high & CasePressureHigh \\
battery errors & 0x01 & battery voltage low & LowVoltageWarning \\
 & 0x02 & battery voltage low & LowVoltageError \\
traveler errors (format 11) & 0x01 & — & MissingPumpType \\
 & 0x02 & — & MissingBatteryType \\
\end{longtable}
\endgroup

\subsubsection{\texorpdfstring{B.2.7 BIT enumerations (formats 9–11)}{B.2.7 BIT enumerations (formats 9–11)}}\label{sec:b27-bit-enumerations-formats-911}

\begingroup\footnotesize\setlength{\tabcolsep}{3pt}
\begin{longtable}{@{}>{\raggedright\arraybackslash}p{0.1339\linewidth}>{\raggedright\arraybackslash}p{0.0893\linewidth}>{\raggedright\arraybackslash}p{0.1071\linewidth}>{\raggedright\arraybackslash}p{0.2500\linewidth}@{}}
\toprule
\textbf{Field (offset)} & \textbf{Value} & \textbf{Meaning} & \textbf{Decoder} \\
\midrule
\endfirsthead
\toprule
\textbf{Field (offset)} & \textbf{Value} & \textbf{Meaning} & \textbf{Decoder} \\
\midrule
\endhead
\bottomrule
\endlastfoot
BIT type (8) & 0 & no sky view & \code{XAlamo\allowbreak{}BIT2Type:\allowbreak{}:\allowbreak{}NO\_\allowbreak{}SKY\_\allowbreak{}VIEW} \\
 & 1 & sky view & \code{SKY\_\allowbreak{}VIEW} \\
 & 0xFF (−1) & unknown & \code{UNKNOWN} \\
CTD type (9) & 0 & RBR & listed “RBR” \\
 & 1 & Sea-Bird & listed “SBE” \\
 & other & — & listed “unknown” \\
\end{longtable}
\endgroup

The CTD type uses the same code as the mission configuration message (0xD0 byte 9).

\subsubsection{\texorpdfstring{B.2.8 BIT thresholds (formats 9–11; not decoded)}{B.2.8 BIT thresholds (formats 9–11; not decoded)}}\label{sec:b28-bit-thresholds-formats-911-not-decoded}

Each threshold is an \code{f32}. The decoder skips them.

\begingroup\footnotesize\setlength{\tabcolsep}{3pt}
\begin{longtable}{@{}>{\raggedright\arraybackslash}p{0.1696\linewidth}>{\raggedright\arraybackslash}p{0.2054\linewidth}>{\raggedright\arraybackslash}p{0.2679\linewidth}>{\raggedright\arraybackslash}p{0.2143\linewidth}@{}}
\toprule
\textbf{Format 9/10 offset} & \textbf{Format 11 (264) offset} & \textbf{Threshold} & \textbf{Units} \\
\midrule
\endfirsthead
\toprule
\textbf{Format 9/10 offset} & \textbf{Format 11 (264) offset} & \textbf{Threshold} & \textbf{Units} \\
\midrule
\endhead
\bottomrule
\endlastfoot
179 & 183 & case vacuum minimum & inHg \\
183 & 187 & case vacuum maximum & inHg \\
187 & 191 & battery voltage low warning & V \\
191 & 195 & battery voltage low & V \\
195 & 199 & oil pump current high warning & A \\
199 & 203 & oil pump current high & A \\
— & 207 & oil pump current low warning & A \\
203 & 211 & oil pump current low & A \\
207 & 215 & air pump maximum cycle count &  \\
211 & 219 & air pump sensor pressure high & (10.5.0: “inHg”) \\
215 & 223 & air pump sensor pressure low & (10.5.0: “inHg”) \\
219 & 227 & air pump current high & A \\
223 & 231 & CTD pressure maximum & dbar \\
227 & 235 & CTD pressure minimum & dbar \\
231 & 239 & CTD temperature maximum & °C (10.5.0 says “dBar”) \\
235 & 243 & CTD temperature minimum & °C \\
239 & 247 & CTD salinity maximum & PSU \\
243 & 251 & CTD salinity minimum & PSU \\
247 & 255 & CTD pressure offset maximum & dbar \\
251 & 259 & CTD pressure offset minimum & dbar \\
\end{longtable}
\endgroup

\subsubsection{\texorpdfstring{B.2.9 Decoding notes}{B.2.9 Decoding notes}}\label{sec:b29-decoding-notes}

\begin{itemize}
\item \textbf{Size checks.} Formats 6, 7 and 8 must be exactly 13, 72 and 76 bytes. For formats 9–11 the build sets \code{XMSG\_\allowbreak{}DEBUG\_\allowbreak{}BIST} to true, so a wrong size (anything but 256 for 9/10, 256 or 264 for 11) only logs a warning and the fields are read anyway; the debug build also writes the full listing to the debug log. (The format 9/10 warning text says “Expected 260 bytes” although the check is for 256.) \emph{(Fixed in build 3309.)} In build 3309, format 9/10 messages shorter than 180 bytes and format 11 messages shorter than 184 bytes are rejected, because decoding would read past their end.
\item \textbf{CTD counts} (formats 7, 8) are converted with the profile gains/offsets from 0xF0 and range-checked; the 7.3.5 YAML says the format 8 pressures use the \emph{surface} gains.
\item \textbf{Pressure offset} (format 8, offset 16) is read unsigned and divided by 1000; a negative offset (signed per the spreadsheet and YAML) would list as about 65.5 dbar.
\item \textbf{Fill values}: formats 6–8 start every \code{u16} at 0xFFFF, the status at 0xFFFFFF and the CTD status at 0xFFFFFFFF; \code{XAlamo\allowbreak{}Bit2} does not initialize its fields (all are set by every format it accepts).
\end{itemize}

\subsubsection{\texorpdfstring{B.2.10 Listing output}{B.2.10 Listing output}}\label{sec:b210-listing-output}

\textbf{Formats 6–8 (\code{XAlamo\allowbreak{}Bit}).} Line type \code{E05} (passed) or \code{E06} (failed). \code{midx} 0 carries \code{BITVer} (decimal), \code{SWVer} (\code{0x\%08X}) and \code{BITResults} (\code{0x\%08X}). Format 6 has only \code{midx} 0.

\begingroup\footnotesize\setlength{\tabcolsep}{3pt}
\begin{longtable}{@{}>{\raggedright\arraybackslash}p{0.0446\linewidth}>{\raggedright\arraybackslash}p{0.4554\linewidth}>{\raggedright\arraybackslash}p{0.4554\linewidth}@{}}
\toprule
\textbf{midx} & \textbf{Format 7 columns} & \textbf{Format 8 columns} \\
\midrule
\endfirsthead
\toprule
\textbf{midx} & \textbf{Format 7 columns} & \textbf{Format 8 columns} \\
\midrule
\endhead
\bottomrule
\endlastfoot
0 & BITVer, SWVer, BITResults & same \\
1 & BITTime (s); pres (dbar, counts); temp (°C, counts); saln (PSU, counts) & BITTime (s), CTDStat (hex) \\
2 & CTDStat (hex); VacStart, VacEnd (inHg); BatStart, BatEnd (V) & presPostReset (dbar, counts); presPreReset (dbar, counts); presOffset (dbar) \\
3 & air valve: AVOpenPos, AVOpenPosCal, AVMidPos, AVMidPosCal, AVClosePos, AVClosePosCal (V) & temp (°C, counts); saln (PSU, counts) \\
4 & air pump: APTestTime (s), APCycleCount, APMxAPres (psi), APMxLoadCur, APMxLoadCur12V (A), APMxLoadBatV (V) & VacStart, VacEnd (inHg); BatStart, BatEnd (V) \\
5 & oil valve: OVOpenPos, OVOpenPosCal, OVClosePos, OVClosePosCal (V) & air valve (as format 7 midx 3) \\
6 & oil pump: OPTestTime (s), OPMxLoadCur, OPMxLoadCur12V (A), OPMxLoadBatV (V) & air pump (as format 7 midx 4) \\
7 & — & oil valve (as format 7 midx 5) \\
8 & — & oil pump (as format 7 midx 6) \\
\end{longtable}
\endgroup

The format 8 \code{pres\allowbreak{}Offset} column uses \code{\%15.\allowbreak{}3d} on a floating-point value; boost::format prints it with three significant digits.

\textbf{Formats 9–11 (\code{XAlamo\allowbreak{}Bit2}).} Fourteen lines, each with \code{BITVer} after \code{midx}:

\begingroup\footnotesize\setlength{\tabcolsep}{3pt}
\begin{longtable}{@{}>{\raggedright\arraybackslash}p{0.0446\linewidth}>{\raggedright\arraybackslash}p{0.7679\linewidth}@{}}
\toprule
\textbf{midx} & \textbf{Columns} \\
\midrule
\endfirsthead
\toprule
\textbf{midx} & \textbf{Columns} \\
\midrule
\endhead
\bottomrule
\endlastfoot
0 & FWVer (hex), BITResults (hex), sensorType, \code{\# sensor\allowbreak{}Type\allowbreak{}Name} \\
1 & startTime, endTime (ISO 8601), duration (s) \\
2 & CTD errors (hex), CTD status (hex), presPreCal, presOffset, presPostCal (dbar) \\
3 & temp (°C), saln (PSU) \\
4 & case pressure errors (hex), caseVacStart, caseVacEnd (inHg) \\
5 & battery errors (hex), batVoltStart, batVoltEnd (V) \\
6 & air valve positions: open, mid, closed (V) \\
7 & air valve calibration positions: open, mid, closed (V) \\
8 & air pump errors (hex), duration (s), cycles, presMax (psi) \\
9 & air pump currAvg, currMax, 12VCurrMax (A), batMax (V), powerMax (W) = 12 × 12VCurrMax \\
10 & oil valve positions: open, closed (V) \\
11 & oil valve calibration positions: open, closed (V) \\
12 & oil pump errors (hex), duration (s) \\
13 & oil pump currAvg, currMax, 12VCurrMax (A), batMax (V), powerMax (W) = 12 × 12VCurrMax \\
\end{longtable}
\endgroup

Not listed: BIT type, air/oil valve errors, GPS, satellite and traveler errors, thresholds.

\subsection{\texorpdfstring{B.3 Descent engineering: 0xE7}{B.3 Descent engineering: 0xE7}}\label{sec:b3-descent-engineering-0xe7}

\textbf{Purpose.} Sent once per dive cycle, describing the descent from the surface to the drift/park pressure: case vacuum, valve positions, and the seek (settle-to-target) results. Introduced with X-format 2.1; ALAMO/ALTO only.

\textbf{Decoder:} \code{XEng\allowbreak{}Descent}, minimum 22 bytes. Versions 0x10 and 0x11.

\subsubsection{\texorpdfstring{B.3.1 Version 0x10 (23 or 27 bytes)}{B.3.1 Version 0x10 (23 or 27 bytes)}}\label{sec:b31-version-0x10-23-or-27-bytes}

\begingroup\footnotesize\setlength{\tabcolsep}{3pt}
\begin{longtable}{@{}>{\raggedright\arraybackslash}p{0.0625\linewidth}>{\raggedright\arraybackslash}p{0.0446\linewidth}>{\raggedright\arraybackslash}p{0.0446\linewidth}>{\raggedright\arraybackslash}p{0.3214\linewidth}>{\raggedright\arraybackslash}p{0.1339\linewidth}>{\raggedright\arraybackslash}p{0.3125\linewidth}@{}}
\toprule
\textbf{Offset} & \textbf{Size} & \textbf{Type} & \textbf{Field} & \textbf{Scale/units} & \textbf{Notes} \\
\midrule
\endfirsthead
\toprule
\textbf{Offset} & \textbf{Size} & \textbf{Type} & \textbf{Field} & \textbf{Scale/units} & \textbf{Notes} \\
\midrule
\endhead
\bottomrule
\endlastfoot
0 & 1 & u8 & ID = 0xE7 &  &  \\
1 & 2 & u16 & size = 23 or 27 & bytes &  \\
3 & 1 & u8 & version = 0x10 &  &  \\
4 & 4 & u32 & descent start time & \code{s2000} &  \\
8 & 2 & u16 & cumulative pump run time during descent & s &  \\
10 & 2 & u16 & pump current & mA (÷1000 → A) &  \\
12 & 2 & u16 & system current & mA (÷1000 → A) &  \\
14 & 2 & u16 & system voltage & mV (÷1000 → V) & 27-byte form only \\
16 & 2 & u16 & case vacuum & 0.01 inHg & 27-byte form only \\
18 & 2 & i16 & first-pump run-time adjustment & s & 27-byte form only \\
20 & 6 & — & not decoded &  & 27-byte form \\
26 & 1 & char & \code{;\allowbreak{}} &  & 27-byte form (23-byte form: bytes 14–21 not decoded, \code{;\allowbreak{}} at 22) \\
\end{longtable}
\endgroup

Version 0x10 is not in the v2.4 spreadsheet; this table is the decoder's reading, and the undecoded bytes are undocumented.

\subsubsection{\texorpdfstring{B.3.2 Version 0x11 (22 bytes)}{B.3.2 Version 0x11 (22 bytes)}}\label{sec:b32-version-0x11-22-bytes}

\begingroup\footnotesize\setlength{\tabcolsep}{3pt}
\begin{longtable}{@{}>{\raggedright\arraybackslash}p{0.0625\linewidth}>{\raggedright\arraybackslash}p{0.0446\linewidth}>{\raggedright\arraybackslash}p{0.0446\linewidth}>{\raggedright\arraybackslash}p{0.3839\linewidth}>{\raggedright\arraybackslash}p{0.1071\linewidth}>{\raggedright\arraybackslash}p{0.2500\linewidth}@{}}
\toprule
\textbf{Offset} & \textbf{Size} & \textbf{Type} & \textbf{Field} & \textbf{Scale/units} & \textbf{Notes} \\
\midrule
\endfirsthead
\toprule
\textbf{Offset} & \textbf{Size} & \textbf{Type} & \textbf{Field} & \textbf{Scale/units} & \textbf{Notes} \\
\midrule
\endhead
\bottomrule
\endlastfoot
0 & 1 & u8 & ID = 0xE7 &  &  \\
1 & 2 & u16 & size = 22 & bytes &  \\
3 & 1 & u8 & version = 0x11 &  & 10.12.0: \code{packet\allowbreak{}Version =\allowbreak{} 17} \\
4 & 4 & u32 & descent start time & \code{s2000} &  \\
8 & 2 & u16 & case vacuum at start & ÷100 → inHg & 10.5.0 blob list: ÷1000 \\
10 & 2 & u16 & case vacuum at end & ÷100 → inHg & 10.5.0 blob list: ÷1000 \\
12 & 2 & u16 & oil valve open position at surface & ÷1000 → V &  \\
14 & 2 & u16 & oil valve closed position during fall pump & ÷1000 → V &  \\
16 & 2 & u16 & air valve middle position at surface & ÷1000 → V &  \\
18 & 1 & u8 & seek attempts &  &  \\
19 & 2 & i16 & seek pump time adjustment & s & 10.12.0 list: \code{U2} \\
21 & 1 & char & \code{;\allowbreak{}} &  &  \\
\end{longtable}
\endgroup

\subsubsection{\texorpdfstring{B.3.3 Decoding notes and listing}{B.3.3 Decoding notes and listing}}\label{sec:b33-decoding-notes-and-listing}

Exact sizes are required (23/27 for 0x10, 22 for 0x11); other versions are rejected. Unset fields keep the fill values. There is no error list.

\begin{itemize}
\item 0x10: one \code{E07} line (\code{midx} 0): engVer, start-of-descent, pumpT (s), pumpI, sysI (A), sysV (V), vac (inHg), pumpAdj (\code{\%+d} s).
\item 0x11: \code{midx} 0: engVer, start-of-descent, vac0, vac1 (inHg), seeks, seekPumpAdjT (\code{\%+d} s); \code{midx} 1 (“Valve Position Measurements”): OilValveOpenPosSurf, OilValveClosePosFallPump, AirValveMidPosSurf (V).
\end{itemize}

\subsection{\texorpdfstring{B.4 Drift/park engineering: 0xE8}{B.4 Drift/park engineering: 0xE8}}\label{sec:b4-driftpark-engineering-0xe8}

\textbf{Purpose.} Sent once per dive cycle for the drift (park) phase: case vacuum, the average CTD values over the drift, and the seek results. Versions before 0x12 also carried the start time and pressure of the fall to profile depth, which moved to the ascent message (0xE9 v0x12).

\textbf{Decoder:} \code{XEng\allowbreak{}Drift}, minimum 19 bytes. Versions 0x10, 0x11, 0x12.

\subsubsection{\texorpdfstring{B.4.1 Version 0x10 (19 bytes)}{B.4.1 Version 0x10 (19 bytes)}}\label{sec:b41-version-0x10-19-bytes}

\begingroup\footnotesize\setlength{\tabcolsep}{3pt}
\begin{longtable}{@{}>{\raggedright\arraybackslash}p{0.0625\linewidth}>{\raggedright\arraybackslash}p{0.0446\linewidth}>{\raggedright\arraybackslash}p{0.0446\linewidth}>{\raggedright\arraybackslash}p{0.2768\linewidth}>{\raggedright\arraybackslash}p{0.1071\linewidth}>{\raggedright\arraybackslash}p{0.0536\linewidth}@{}}
\toprule
\textbf{Offset} & \textbf{Size} & \textbf{Type} & \textbf{Field} & \textbf{Scale/units} & \textbf{Notes} \\
\midrule
\endfirsthead
\toprule
\textbf{Offset} & \textbf{Size} & \textbf{Type} & \textbf{Field} & \textbf{Scale/units} & \textbf{Notes} \\
\midrule
\endhead
\bottomrule
\endlastfoot
0 & 1 & u8 & ID = 0xE8 &  &  \\
1 & 2 & u16 & size = 19 & bytes &  \\
3 & 1 & u8 & version = 0x10 &  &  \\
4 & 4 & u32 & drift start time & \code{s2000} &  \\
8 & 2 & u16 & average drift pressure & CTD counts &  \\
10 & 2 & u16 & average drift temperature & CTD counts &  \\
12 & 4 & u32 & fall-to-profile start time & \code{s2000} &  \\
16 & 2 & u16 & fall-to-profile start pressure & CTD counts &  \\
18 & 1 & char & \code{;\allowbreak{}} &  &  \\
\end{longtable}
\endgroup

\subsubsection{\texorpdfstring{B.4.2 Version 0x11 (28 bytes)}{B.4.2 Version 0x11 (28 bytes)}}\label{sec:b42-version-0x11-28-bytes}

\begingroup\footnotesize\setlength{\tabcolsep}{3pt}
\begin{longtable}{@{}>{\raggedright\arraybackslash}p{0.0625\linewidth}>{\raggedright\arraybackslash}p{0.0446\linewidth}>{\raggedright\arraybackslash}p{0.0446\linewidth}>{\raggedright\arraybackslash}p{0.2768\linewidth}>{\raggedright\arraybackslash}p{0.1071\linewidth}>{\raggedright\arraybackslash}p{0.3839\linewidth}@{}}
\toprule
\textbf{Offset} & \textbf{Size} & \textbf{Type} & \textbf{Field} & \textbf{Scale/units} & \textbf{Notes} \\
\midrule
\endfirsthead
\toprule
\textbf{Offset} & \textbf{Size} & \textbf{Type} & \textbf{Field} & \textbf{Scale/units} & \textbf{Notes} \\
\midrule
\endhead
\bottomrule
\endlastfoot
0 & 1 & u8 & ID = 0xE8 &  &  \\
1 & 2 & u16 & size = 28 & bytes &  \\
3 & 1 & u8 & version = 0x11 &  &  \\
4 & 4 & u32 & drift start time & \code{s2000} &  \\
8 & 2 & u16 & case vacuum at start & ÷100 → inHg &  \\
10 & 2 & u16 & case vacuum at end & ÷100 → inHg &  \\
12 & 2 & u16 & average pressure & CTD counts &  \\
14 & 2 & u16 & average temperature & CTD counts &  \\
16 & 2 & u16 & average salinity & CTD counts &  \\
18 & 1 & u8 & seek attempts &  &  \\
19 & 2 & i16 & seek pump time adjustment & s & decoder reads \code{u16} into an \code{int16\_\allowbreak{}t} (same result) \\
21 & 4 & u32 & fall-to-profile start time & \code{s2000} &  \\
25 & 2 & u16 & fall-to-profile start pressure & CTD counts &  \\
27 & 1 & char & \code{;\allowbreak{}} &  &  \\
\end{longtable}
\endgroup

\subsubsection{\texorpdfstring{B.4.3 Version 0x12 (22 bytes)}{B.4.3 Version 0x12 (22 bytes)}}\label{sec:b43-version-0x12-22-bytes}

Version 0x11 without the fall-to-profile fields.

\begingroup\footnotesize\setlength{\tabcolsep}{3pt}
\begin{longtable}{@{}>{\raggedright\arraybackslash}p{0.0625\linewidth}>{\raggedright\arraybackslash}p{0.0446\linewidth}>{\raggedright\arraybackslash}p{0.0446\linewidth}>{\raggedright\arraybackslash}p{0.2321\linewidth}>{\raggedright\arraybackslash}p{0.1071\linewidth}>{\raggedright\arraybackslash}p{0.3393\linewidth}@{}}
\toprule
\textbf{Offset} & \textbf{Size} & \textbf{Type} & \textbf{Field} & \textbf{Scale/units} & \textbf{Notes} \\
\midrule
\endfirsthead
\toprule
\textbf{Offset} & \textbf{Size} & \textbf{Type} & \textbf{Field} & \textbf{Scale/units} & \textbf{Notes} \\
\midrule
\endhead
\bottomrule
\endlastfoot
0 & 1 & u8 & ID = 0xE8 &  &  \\
1 & 2 & u16 & size = 22 & bytes &  \\
3 & 1 & u8 & version = 0x12 &  & 10.12.0 list says 17 (0x11); see B.10 \\
4 & 4 & u32 & drift start time & \code{s2000} &  \\
8 & 2 & u16 & case vacuum at start & ÷100 → inHg &  \\
10 & 2 & u16 & case vacuum at end & ÷100 → inHg &  \\
12 & 2 & u16 & average pressure & CTD counts &  \\
14 & 2 & u16 & average temperature & CTD counts &  \\
16 & 2 & u16 & average salinity & CTD counts &  \\
18 & 1 & u8 & seek attempts &  &  \\
19 & 2 & i16 & seek pump time adjustment & s & 10.12.0 list: \code{U2} \\
21 & 1 & char & \code{;\allowbreak{}} &  &  \\
\end{longtable}
\endgroup

\subsubsection{\texorpdfstring{B.4.4 Decoding notes and listing}{B.4.4 Decoding notes and listing}}\label{sec:b44-decoding-notes-and-listing}

Exact sizes are required (19, 28, 22); other versions are rejected. No error list.

\begin{itemize}
\item 0x10: one \code{E08} line: engVer, start-of-drift, meanP (dbar), meanT (°C), start-of-fall, fallP (dbar).
\item 0x11 and 0x12: \code{midx} 0: engVer, start-of-drift, vac0, vac1 (inHg), seeks, seekPumpAdjT (\code{\%+d} s); \code{midx} 1 (“Drift Sensor Averages”): pres (dbar, counts), temp (°C, counts), saln (PSU, counts). The 0x11 fall-to-profile fields are decoded but not listed.
\end{itemize}

Example (README, \code{lis9287-0325.\allowbreak{}txt}):

\begin{Verbatim}
E08 09287  325    0   0x12 2026-10-24T11:05:29Z      6.65      6.69               0              +0
E08 09287  325    1          274.72         7118        9.744        14744       34.523        35523
\end{Verbatim}

\subsection{\texorpdfstring{B.5 Ascent engineering: 0xE9}{B.5 Ascent engineering: 0xE9}}\label{sec:b5-ascent-engineering-0xe9}

\textbf{Purpose.} Sent once per dive cycle for the ascent (profile) phase. From version 0x12 it also carries the start time and pressure of the fall from park to profile depth.

\textbf{Decoder:} \code{XEng\allowbreak{}Ascent}, minimum 13 bytes. Versions 0x10, 0x11, 0x12.

\subsubsection{\texorpdfstring{B.5.1 Version 0x10 (13 bytes)}{B.5.1 Version 0x10 (13 bytes)}}\label{sec:b51-version-0x10-13-bytes}

\begingroup\footnotesize\setlength{\tabcolsep}{3pt}
\begin{longtable}{@{}>{\raggedright\arraybackslash}p{0.0625\linewidth}>{\raggedright\arraybackslash}p{0.0446\linewidth}>{\raggedright\arraybackslash}p{0.0446\linewidth}>{\raggedright\arraybackslash}p{0.3482\linewidth}>{\raggedright\arraybackslash}p{0.1339\linewidth}>{\raggedright\arraybackslash}p{0.0536\linewidth}@{}}
\toprule
\textbf{Offset} & \textbf{Size} & \textbf{Type} & \textbf{Field} & \textbf{Scale/units} & \textbf{Notes} \\
\midrule
\endfirsthead
\toprule
\textbf{Offset} & \textbf{Size} & \textbf{Type} & \textbf{Field} & \textbf{Scale/units} & \textbf{Notes} \\
\midrule
\endhead
\bottomrule
\endlastfoot
0 & 1 & u8 & ID = 0xE9 &  &  \\
1 & 2 & u16 & size = 13 & bytes &  \\
3 & 1 & u8 & version = 0x10 &  &  \\
4 & 4 & u32 & ascent start time & \code{s2000} &  \\
8 & 2 & u16 & system current & mA (÷1000 → A) &  \\
10 & 2 & u16 & cumulative pump run time during ascent & s &  \\
12 & 1 & char & \code{;\allowbreak{}} &  &  \\
\end{longtable}
\endgroup

\subsubsection{\texorpdfstring{B.5.2 Version 0x11 (13 bytes)}{B.5.2 Version 0x11 (13 bytes)}}\label{sec:b52-version-0x11-13-bytes}

\begingroup\footnotesize\setlength{\tabcolsep}{3pt}
\begin{longtable}{@{}>{\raggedright\arraybackslash}p{0.0625\linewidth}>{\raggedright\arraybackslash}p{0.0446\linewidth}>{\raggedright\arraybackslash}p{0.0446\linewidth}>{\raggedright\arraybackslash}p{0.1875\linewidth}>{\raggedright\arraybackslash}p{0.1071\linewidth}>{\raggedright\arraybackslash}p{0.4732\linewidth}@{}}
\toprule
\textbf{Offset} & \textbf{Size} & \textbf{Type} & \textbf{Field} & \textbf{Scale/units} & \textbf{Notes} \\
\midrule
\endfirsthead
\toprule
\textbf{Offset} & \textbf{Size} & \textbf{Type} & \textbf{Field} & \textbf{Scale/units} & \textbf{Notes} \\
\midrule
\endhead
\bottomrule
\endlastfoot
0 & 1 & u8 & ID = 0xE9 &  &  \\
1 & 2 & u16 & size = 13 & bytes &  \\
3 & 1 & u8 & version = 0x11 &  &  \\
4 & 4 & u32 & ascent start time & \code{s2000} & spreadsheet text says “drift/park phase” (copy error) \\
8 & 2 & u16 & case vacuum at start & ÷100 → inHg &  \\
10 & 2 & u16 & case vacuum at end & ÷100 → inHg &  \\
12 & 1 & char & \code{;\allowbreak{}} &  &  \\
\end{longtable}
\endgroup

\subsubsection{\texorpdfstring{B.5.3 Version 0x12 (19 bytes)}{B.5.3 Version 0x12 (19 bytes)}}\label{sec:b53-version-0x12-19-bytes}

\begingroup\footnotesize\setlength{\tabcolsep}{3pt}
\begin{longtable}{@{}>{\raggedright\arraybackslash}p{0.0625\linewidth}>{\raggedright\arraybackslash}p{0.0446\linewidth}>{\raggedright\arraybackslash}p{0.0446\linewidth}>{\raggedright\arraybackslash}p{0.2679\linewidth}>{\raggedright\arraybackslash}p{0.1071\linewidth}>{\raggedright\arraybackslash}p{0.3393\linewidth}@{}}
\toprule
\textbf{Offset} & \textbf{Size} & \textbf{Type} & \textbf{Field} & \textbf{Scale/units} & \textbf{Notes} \\
\midrule
\endfirsthead
\toprule
\textbf{Offset} & \textbf{Size} & \textbf{Type} & \textbf{Field} & \textbf{Scale/units} & \textbf{Notes} \\
\midrule
\endhead
\bottomrule
\endlastfoot
0 & 1 & u8 & ID = 0xE9 &  &  \\
1 & 2 & u16 & size = 19 & bytes &  \\
3 & 1 & u8 & version = 0x12 &  & 10.12.0 list says 17 (0x11); see B.10 \\
4 & 4 & u32 & ascent start time & \code{s2000} &  \\
8 & 2 & u16 & case vacuum at start & ÷100 → inHg &  \\
10 & 2 & u16 & case vacuum at end & ÷100 → inHg &  \\
12 & 4 & u32 & fall-from-park start time & \code{s2000} &  \\
16 & 2 & u16 & fall-from-park start pressure & CTD counts &  \\
18 & 1 & char & \code{;\allowbreak{}} &  &  \\
\end{longtable}
\endgroup

\subsubsection{\texorpdfstring{B.5.4 Decoding notes and listing}{B.5.4 Decoding notes and listing}}\label{sec:b54-decoding-notes-and-listing}

Exact sizes are required (13, 13, 19); other versions are rejected. \code{XEng\allowbreak{}Ascent} does not initialize its fields, so fields a version does not carry are undefined rather than fill values (none of them is listed). No error list.

\begin{itemize}
\item 0x10: \code{E09} \code{midx} 0: engVer, start-of-ascent, sysI (A), pumpT (s).
\item 0x11: \code{midx} 0: engVer, start-of-ascent, vac0, vac1 (inHg).
\item 0x12: \code{midx} 0 as 0x11; \code{midx} 1: start-of-fall, fallP (dbar).
\end{itemize}

\subsection{\texorpdfstring{B.6 Surface engineering: 0xEA}{B.6 Surface engineering: 0xEA}}\label{sec:b6-surface-engineering-0xea}

\textbf{Purpose.} Sent once per dive cycle from the surface: Iridium session statistics, case vacuum, air pump and valve measurements, the surface pressure offset reading, profile bookkeeping (first/last raw CTD sample, bin counts), the GPS–RTC clock difference, and the list of system errors recorded during the dive cycle.

\textbf{Decoder:} \code{XEng\allowbreak{}Surface}, minimum 46 bytes. Versions 0x10–0x14.

\subsubsection{\texorpdfstring{B.6.1 Error list}{B.6.1 Error list}}\label{sec:b61-error-list}

Every version ends with an error count \code{N} (\code{u8}) followed by \code{N} entries and the terminator.

\begingroup\footnotesize\setlength{\tabcolsep}{3pt}
\begin{longtable}{@{}>{\raggedright\arraybackslash}p{0.0893\linewidth}>{\raggedright\arraybackslash}p{0.0982\linewidth}>{\raggedright\arraybackslash}p{0.3929\linewidth}>{\raggedright\arraybackslash}p{0.1607\linewidth}@{}}
\toprule
\textbf{Version} & \textbf{Entry size} & \textbf{Entry layout} & \textbf{First entry at} \\
\midrule
\endfirsthead
\toprule
\textbf{Version} & \textbf{Entry size} & \textbf{Entry layout} & \textbf{First entry at} \\
\midrule
\endhead
\bottomrule
\endlastfoot
0x10 & 4 & \code{u32} error code (no timestamp) & 47 \\
0x11–0x14 & 6 & \code{u32} time (\code{s2000}), then \code{u16} error code (B.9) & 45 / 49 / 67 / 75 \\
\end{longtable}
\endgroup

The message size is therefore \code{first + entry × N + 1}. The decoder reads entries while the next entry ends before the terminator, then requires the number read to equal \code{N} (\code{XMSG\_\allowbreak{}DEBUG\_\allowbreak{}ENG\_\allowbreak{}SURF} is false, so a mismatch rejects the message).

\subsubsection{\texorpdfstring{B.6.2 Version 0x10 (48 + 4N bytes)}{B.6.2 Version 0x10 (48 + 4N bytes)}}\label{sec:b62-version-0x10-48--4n-bytes}

\begingroup\footnotesize\setlength{\tabcolsep}{3pt}
\begin{longtable}{@{}>{\raggedright\arraybackslash}p{0.0625\linewidth}>{\raggedright\arraybackslash}p{0.0446\linewidth}>{\raggedright\arraybackslash}p{0.0714\linewidth}>{\raggedright\arraybackslash}p{0.3304\linewidth}>{\raggedright\arraybackslash}p{0.1339\linewidth}>{\raggedright\arraybackslash}p{0.1250\linewidth}@{}}
\toprule
\textbf{Offset} & \textbf{Size} & \textbf{Type} & \textbf{Field} & \textbf{Scale/units} & \textbf{Notes} \\
\midrule
\endfirsthead
\toprule
\textbf{Offset} & \textbf{Size} & \textbf{Type} & \textbf{Field} & \textbf{Scale/units} & \textbf{Notes} \\
\midrule
\endhead
\bottomrule
\endlastfoot
0 & 1 & u8 & ID = 0xEA &  &  \\
1 & 2 & u16 & size = 48 + 4N & bytes &  \\
3 & 1 & u8 & version = 0x10 &  &  \\
4 & 4 & u32 & surface start time & \code{s2000} &  \\
8 & 4 & u32 & CTD status &  & listed hex \\
12 & 2 & u16 & SBD packets for this dive &  &  \\
14 & 2 & u16 & transmission attempts, previous dive &  &  \\
16 & 2 & u16 & case vacuum & ÷100 → inHg &  \\
18 & 2 & u16 & battery voltage & mV (÷1000 → V) &  \\
20 & 2 & u16 & system current & mA (÷1000 → A) &  \\
22 & 24 & — & not decoded &  & undocumented \\
46 & 1 & u8 & error count N &  &  \\
47 & 4N & u32 × N & error codes &  & no timestamps \\
47+4N & 1 & char & \code{;\allowbreak{}} &  &  \\
\end{longtable}
\endgroup

Version 0x10 is not in the v2.4 spreadsheet; the table is the decoder's reading.

\subsubsection{\texorpdfstring{B.6.3 Version 0x11 (46 + 6N bytes)}{B.6.3 Version 0x11 (46 + 6N bytes)}}\label{sec:b63-version-0x11-46--6n-bytes}

\begingroup\footnotesize\setlength{\tabcolsep}{3pt}
\begin{longtable}{@{}>{\raggedright\arraybackslash}p{0.0625\linewidth}>{\raggedright\arraybackslash}p{0.0446\linewidth}>{\raggedright\arraybackslash}p{0.0446\linewidth}>{\raggedright\arraybackslash}p{0.3750\linewidth}>{\raggedright\arraybackslash}p{0.1071\linewidth}>{\raggedright\arraybackslash}p{0.0536\linewidth}@{}}
\toprule
\textbf{Offset} & \textbf{Size} & \textbf{Type} & \textbf{Field} & \textbf{Scale/units} & \textbf{Notes} \\
\midrule
\endfirsthead
\toprule
\textbf{Offset} & \textbf{Size} & \textbf{Type} & \textbf{Field} & \textbf{Scale/units} & \textbf{Notes} \\
\midrule
\endhead
\bottomrule
\endlastfoot
0 & 1 & u8 & ID = 0xEA &  &  \\
1 & 2 & u16 & size = 46 + 6N & bytes &  \\
3 & 1 & u8 & version = 0x11 &  &  \\
4 & 4 & u32 & surface start time & \code{s2000} &  \\
8 & 2 & u16 & case vacuum at start & ÷100 → inHg &  \\
10 & 2 & u16 & case vacuum at end & ÷100 → inHg &  \\
12 & 4 & u32 & CTD status &  &  \\
16 & 2 & u16 & packets to transmit this session &  &  \\
18 & 2 & u16 & transmission attempts, previous dive &  &  \\
20 & 2 & u16 & air pump current & ÷1000 → A &  \\
22 & 2 & u16 & air pump run time & s &  \\
24 & 2 & u16 & air bladder maximum pressure & ÷1000 → psi &  \\
26 & 2 & u16 & air valve open position at start & ÷1000 → V &  \\
28 & 2 & u16 & air valve middle position at end & ÷1000 → V &  \\
30 & 2 & u16 & oil valve open position during air assist & ÷1000 → V &  \\
32 & 2 & u16 & oil valve closed position at end & ÷1000 → V &  \\
34 & 2 & u16 & air valve calibration open position & ÷1000 → V &  \\
36 & 2 & u16 & air valve calibration closed position & ÷1000 → V &  \\
38 & 2 & u16 & air valve calibration middle position & ÷1000 → V &  \\
40 & 2 & u16 & oil valve calibration open position & ÷1000 → V &  \\
42 & 2 & u16 & oil valve calibration closed position & ÷1000 → V &  \\
44 & 1 & u8 & error count N &  &  \\
45 & 6N &  & error entries (B.6.1) &  &  \\
45+6N & 1 & char & \code{;\allowbreak{}} &  &  \\
\end{longtable}
\endgroup

\subsubsection{\texorpdfstring{B.6.4 Version 0x12 (50 + 6N bytes)}{B.6.4 Version 0x12 (50 + 6N bytes)}}\label{sec:b64-version-0x12-50--6n-bytes}

Version 0x11 with an unidentified 2-byte field after the case vacuum and the pressure offset reading before the error count.

\begingroup\footnotesize\setlength{\tabcolsep}{3pt}
\begin{longtable}{@{}>{\raggedright\arraybackslash}p{0.0625\linewidth}>{\raggedright\arraybackslash}p{0.0446\linewidth}>{\raggedright\arraybackslash}p{0.0446\linewidth}>{\raggedright\arraybackslash}p{0.5089\linewidth}>{\raggedright\arraybackslash}p{0.1161\linewidth}>{\raggedright\arraybackslash}p{0.1339\linewidth}@{}}
\toprule
\textbf{Offset} & \textbf{Size} & \textbf{Type} & \textbf{Field} & \textbf{Scale/units} & \textbf{Notes} \\
\midrule
\endfirsthead
\toprule
\textbf{Offset} & \textbf{Size} & \textbf{Type} & \textbf{Field} & \textbf{Scale/units} & \textbf{Notes} \\
\midrule
\endhead
\bottomrule
\endlastfoot
0–11 &  &  & as 0x11 (ID … case vacuum at end) &  & version = 0x12 \\
12 & 2 & u16 & unidentified (“???” in the spreadsheet, typical value 0) &  & not decoded \\
14 & 4 & u32 & CTD status &  &  \\
18 & 2 & u16 & packets to transmit this session &  &  \\
20 & 2 & u16 & transmission attempts, previous dive &  &  \\
22 & 2 & u16 & air pump current & ÷1000 → A &  \\
24 & 2 & u16 & air pump run time & s &  \\
26 & 2 & u16 & air bladder maximum pressure & ÷1000 → psi &  \\
28 & 2 & u16 & air valve open position at start & ÷1000 → V &  \\
30 & 2 & u16 & air valve middle position at end & ÷1000 → V &  \\
32 & 2 & u16 & oil valve open position during air assist & ÷1000 → V &  \\
34 & 2 & u16 & oil valve closed position at end & ÷1000 → V &  \\
36 & 2 & u16 & air valve calibration open position & ÷1000 → V &  \\
38 & 2 & u16 & air valve calibration closed position & ÷1000 → V &  \\
40 & 2 & u16 & air valve calibration middle position & ÷1000 → V &  \\
42 & 2 & u16 & oil valve calibration open position & ÷1000 → V &  \\
44 & 2 & u16 & oil valve calibration closed position & ÷1000 → V &  \\
46 & 2 & i16 & pressure sensor offset reading & ÷1000 → dbar &  \\
48 & 1 & u8 & error count N &  &  \\
49 & 6N &  & error entries (B.6.1) &  &  \\
49+6N & 1 & char & \code{;\allowbreak{}} &  &  \\
\end{longtable}
\endgroup

The spreadsheet gives the size as “6N+49”, but its own terminator offset (6N+49) makes it 6N+50, which is what the decoder requires as a minimum.

\subsubsection{\texorpdfstring{B.6.5 Version 0x13 (68 + 6N bytes)}{B.6.5 Version 0x13 (68 + 6N bytes)}}\label{sec:b65-version-0x13-68--6n-bytes}

\begingroup\footnotesize\setlength{\tabcolsep}{3pt}
\begin{longtable}{@{}>{\raggedright\arraybackslash}p{0.0625\linewidth}>{\raggedright\arraybackslash}p{0.0446\linewidth}>{\raggedright\arraybackslash}p{0.0446\linewidth}>{\raggedright\arraybackslash}p{0.3750\linewidth}>{\raggedright\arraybackslash}p{0.1161\linewidth}>{\raggedright\arraybackslash}p{0.2768\linewidth}@{}}
\toprule
\textbf{Offset} & \textbf{Size} & \textbf{Type} & \textbf{Field} & \textbf{Scale/units} & \textbf{Notes} \\
\midrule
\endfirsthead
\toprule
\textbf{Offset} & \textbf{Size} & \textbf{Type} & \textbf{Field} & \textbf{Scale/units} & \textbf{Notes} \\
\midrule
\endhead
\bottomrule
\endlastfoot
0 & 1 & u8 & ID = 0xEA &  &  \\
1 & 2 & u16 & size = 68 + 6N & bytes &  \\
3 & 1 & u8 & version = 0x13 &  &  \\
4 & 4 & u32 & surface start time & \code{s2000} &  \\
8 & 2 & u16 & packets to transmit this session &  &  \\
10 & 2 & u16 & transmission attempts, previous dive &  &  \\
12 & 2 & u16 & case vacuum at start & ÷100 → inHg &  \\
14 & 2 & u16 & case vacuum at end & ÷100 → inHg &  \\
16 & 2 & i16 & GPS − RTC time difference before synchronization & s &  \\
18 & 4 & u32 & CTD status &  & firmware lists: 4 characters \\
22 & 2 & u16 & first raw profile pressure & CTD counts &  \\
24 & 2 & u16 & first raw profile temperature & CTD counts &  \\
26 & 2 & u16 & first raw profile salinity & CTD counts &  \\
28 & 2 & u16 & last raw profile pressure & CTD counts &  \\
30 & 2 & u16 & last raw profile temperature & CTD counts &  \\
32 & 2 & u16 & last raw profile salinity & CTD counts &  \\
34 & 2 & u16 & number of profile bins &  &  \\
36 & 2 & u16 & spreadsheet: number of raw profile samples &  & \textbf{decoder: number of empty bins} \\
38 & 2 & u16 & spreadsheet: number of empty bins &  & \textbf{decoder: number of raw samples} \\
40 & 2 & u16 & air pump current & ÷1000 → A &  \\
42 & 2 & u16 & air pump run time & s &  \\
44 & 2 & u16 & air bladder maximum pressure & ÷1000 → psi &  \\
46 & 2 & u16 & air valve open position at start & ÷1000 → V &  \\
48 & 2 & u16 & air valve middle position at end & ÷1000 → V &  \\
50 & 2 & u16 & oil valve open position during air assist & ÷1000 → V &  \\
52 & 2 & u16 & oil valve closed position at end & ÷1000 → V &  \\
54 & 2 & u16 & air valve calibration open position & ÷1000 → V &  \\
56 & 2 & u16 & air valve calibration closed position & ÷1000 → V &  \\
58 & 2 & u16 & air valve calibration middle position & ÷1000 → V &  \\
60 & 2 & u16 & oil valve calibration open position & ÷1000 → V &  \\
62 & 2 & u16 & oil valve calibration closed position & ÷1000 → V &  \\
64 & 2 & i16 & pressure sensor offset reading & ÷1000 → dbar &  \\
66 & 1 & u8 & error count N &  &  \\
67 & 6N &  & error entries (B.6.1) &  &  \\
67+6N & 1 & char & \code{;\allowbreak{}} &  &  \\
\end{longtable}
\endgroup

\subsubsection{\texorpdfstring{B.6.6 Version 0x14 (76 + 6N bytes)}{B.6.6 Version 0x14 (76 + 6N bytes)}}\label{sec:b66-version-0x14-76--6n-bytes}

Version 0x13 with the surface data-collection (“packetize”) start and end times inserted after the start time, and the clock/packet fields reordered. Firmware 7.3.5, 10.5.0 and 10.12.0 (\code{packet\allowbreak{}Version =\allowbreak{} 20}) all use this layout.

\begingroup\footnotesize\setlength{\tabcolsep}{3pt}
\begin{longtable}{@{}>{\raggedright\arraybackslash}p{0.0625\linewidth}>{\raggedright\arraybackslash}p{0.0446\linewidth}>{\raggedright\arraybackslash}p{0.0714\linewidth}>{\raggedright\arraybackslash}p{0.3036\linewidth}>{\raggedright\arraybackslash}p{0.1161\linewidth}>{\raggedright\arraybackslash}p{0.3214\linewidth}@{}}
\toprule
\textbf{Offset} & \textbf{Size} & \textbf{Type} & \textbf{Field} & \textbf{Scale/units} & \textbf{Notes} \\
\midrule
\endfirsthead
\toprule
\textbf{Offset} & \textbf{Size} & \textbf{Type} & \textbf{Field} & \textbf{Scale/units} & \textbf{Notes} \\
\midrule
\endhead
\bottomrule
\endlastfoot
0 & 1 & u8 & ID = 0xEA &  &  \\
1 & 2 & u16 & size = 76 + 6N & bytes &  \\
3 & 1 & u8 & version = 0x14 &  &  \\
4 & 4 & u32 & surface start time & \code{s2000} &  \\
8 & 4 & u32 & surface data collection (packetize) start time & \code{s2000} &  \\
12 & 4 & u32 & surface data collection (packetize) end time & \code{s2000} &  \\
16 & 2 & u16 & packets this dive/mission &  &  \\
18 & 2 & u16 & transmission attempts, previous dive &  &  \\
20 & 2 & u16 & case vacuum at start & ÷100 → inHg &  \\
22 & 2 & u16 & case vacuum at end & ÷100 → inHg &  \\
24 & 2 & i16 & GPS − RTC time difference & s &  \\
26 & 4 & char[4] & CTD status &  & decoded as \code{u32} \\
30 & 2 & u16 & first raw profile pressure & CTD counts &  \\
32 & 2 & u16 & first raw profile temperature & CTD counts &  \\
34 & 2 & u16 & first raw profile salinity & CTD counts &  \\
36 & 2 & u16 & last raw profile pressure & CTD counts &  \\
38 & 2 & u16 & last raw profile temperature & CTD counts &  \\
40 & 2 & u16 & last raw profile salinity & CTD counts &  \\
42 & 2 & u16 & number of CTD profile bins &  &  \\
44 & 2 & u16 & firmware lists: number of CTD profile samples &  & \textbf{decoder: number of empty bins} \\
46 & 2 & u16 & firmware lists: number of empty bins &  & \textbf{decoder: number of raw samples} \\
48 & 2 & u16 & air pump current & ÷1000 → A &  \\
50 & 2 & u16 & air pump run time & s &  \\
52 & 2 & u16 & air bladder maximum pressure & ÷1000 → psi &  \\
54 & 2 & u16 & air valve open position at start & ÷1000 → V &  \\
56 & 2 & u16 & air valve middle position at end & ÷1000 → V &  \\
58 & 2 & u16 & oil valve open position during air assist & ÷1000 → V &  \\
60 & 2 & u16 & oil valve closed position at end & ÷1000 → V &  \\
62 & 2 & u16 & air valve calibration open position & ÷1000 → V &  \\
64 & 2 & u16 & air valve calibration closed position & ÷1000 → V &  \\
66 & 2 & u16 & air valve calibration middle position & ÷1000 → V &  \\
68 & 2 & u16 & oil valve calibration open position & ÷1000 → V &  \\
70 & 2 & u16 & oil valve calibration closed position & ÷1000 → V &  \\
72 & 2 & i16 & pressure sensor offset reading & ÷1000 → dbar & 10.5.0/10.12.0 lists: \code{U2}; 7.3.5 YAML: \code{i16} \\
74 & 1 & u8 & error count N &  &  \\
75 & 6N &  & error entries (B.6.1) &  & 10.5.0 lists room for 10 entries \\
75+6N & 1 & char & \code{;\allowbreak{}} &  &  \\
\end{longtable}
\endgroup

\subsubsection{\texorpdfstring{B.6.7 Decoding notes}{B.6.7 Decoding notes}}\label{sec:b67-decoding-notes}

\begin{itemize}
\item Minimum sizes: 48 (0x10), 46 (0x11), 50 (0x12), 68 (0x13), 76 (0x14); longer messages carry error entries. Unknown versions are rejected.
\item The pressure offset reading is listed with \code{XMsg\allowbreak{}Format\allowbreak{}Pressure(\allowbreak{})} (two decimals, fill −9.999 if outside −5…2610 dbar).
\item First/last raw samples are converted with the profile gains/offsets and range-checked.
\item Error codes are stored as 16-bit values (32-bit in 0x10) and listed as \code{0x\%08X}, with no name lookup; B.9 gives the names. A zero timestamp lists as \code{0000-00-00T00:\allowbreak{}00:\allowbreak{}00Z}.
\end{itemize}

\subsubsection{\texorpdfstring{B.6.8 Listing output (\code{E0A}, \code{e0A})}{B.6.8 Listing output (E0A, e0A)}}\label{sec:b68-listing-output-e0a-e0a}

\begingroup\footnotesize\setlength{\tabcolsep}{3pt}
\begin{longtable}{@{}>{\raggedright\arraybackslash}p{0.0446\linewidth}>{\raggedright\arraybackslash}p{0.1964\linewidth}>{\raggedright\arraybackslash}p{0.1964\linewidth}>{\raggedright\arraybackslash}p{0.1518\linewidth}>{\raggedright\arraybackslash}p{0.1786\linewidth}>{\raggedright\arraybackslash}p{0.1518\linewidth}@{}}
\toprule
\textbf{midx} & \textbf{0x10} & \textbf{0x11} & \textbf{0x12} & \textbf{0x13} & \textbf{0x14} \\
\midrule
\endfirsthead
\toprule
\textbf{midx} & \textbf{0x10} & \textbf{0x11} & \textbf{0x12} & \textbf{0x13} & \textbf{0x14} \\
\midrule
\endhead
\bottomrule
\endlastfoot
0 & engVer, start-of-surface, CTD-status, nSBD, nAtmpt, vac, sysV, sysI, nErr & engVer, start-of-surface, vac0, vac1, nSBD, nPrevAttempt & same as 0x11 & same & same \\
1 & — & air pump: APCur (A), APRunT (s), AirBladMxP (psi) & same & same & same \\
2 & — & air valve: AVOpenStart, AVMidEnd, AVOpenCal, AVMidCal, AVCloseCal (V) & same & same & same \\
3 & — & oil valve: OVOpenAirPmp, OVCloseEnd, OVOpenCal, OVClosCal (V) & same & same & same \\
4 & — & \code{e0A} nErr & presOffset (dbar) & presOffset & presOffset \\
5 & — & — & \code{e0A} nErr & first sample: pres (dbar, counts), temp (°C, counts), saln (PSU, counts) & same \\
6 & — & — & — & last sample (same columns) & same \\
7 & — & — & — & numBins, numEmptyBins, numRawSamples & same \\
8 & — & — & — & gpsOffset (s) & gpsOffset \\
9 & — & — & — & \code{E0A} nErr & surface sampling start, end \\
10 & — & — & — & — & \code{E0A} nErr \\
\end{longtable}
\endgroup

When there are errors, a heading and one \code{e0A} line per entry follow: \code{e0A SN dive idx date-time 0x\%08X} (\code{idx} counts from 0; version 0x10 prints \code{0000-00-00T00:\allowbreak{}00:\allowbreak{}00Z} for the missing time). The \code{n\allowbreak{}Err} line is tagged \code{e0A} in versions 0x11 and 0x12 but \code{E0A} in 0x13 and 0x14. The CTD status is listed only by version 0x10; the 0x11 fields \code{n\allowbreak{}SBD}/\code{n\allowbreak{}Prev\allowbreak{}Attempt} are printed for all later versions.

\emph{(Fixed in build 3309.)} In version 0x10 the error heading's format string had five fields but was given four values; boost::format throws in that case, so a 0x10 message with errors is expected to fail listing generation (\emph{inferred from the code}; not observed).

\subsection{\texorpdfstring{B.7 Ice-detection engineering: 0xEC}{B.7 Ice-detection engineering: 0xEC}}\label{sec:b7-ice-detection-engineering-0xec}

\textbf{Purpose.} Sent by floats with ice detection enabled, once per dive cycle. It reports the ice-avoidance mode, what happened on the way up (surfaced, ice found, temperature turnaround), how many surfacing attempts were made, how much data is waiting to be sent, and the values used by the temperature-threshold and surface-salinity tests. Version 0x11 adds an error list.

\textbf{Decoder:} \code{XEng\allowbreak{}Ice}, minimum 21 bytes. Versions 0x10 and 0x11.

\subsubsection{\texorpdfstring{B.7.1 Version 0x10 (21 bytes)}{B.7.1 Version 0x10 (21 bytes)}}\label{sec:b71-version-0x10-21-bytes}

\begingroup\footnotesize\setlength{\tabcolsep}{3pt}
\begin{longtable}{@{}>{\raggedright\arraybackslash}p{0.0625\linewidth}>{\raggedright\arraybackslash}p{0.0446\linewidth}>{\raggedright\arraybackslash}p{0.0446\linewidth}>{\raggedright\arraybackslash}p{0.3214\linewidth}>{\raggedright\arraybackslash}p{0.1071\linewidth}>{\raggedright\arraybackslash}p{0.3393\linewidth}@{}}
\toprule
\textbf{Offset} & \textbf{Size} & \textbf{Type} & \textbf{Field} & \textbf{Scale/units} & \textbf{Notes} \\
\midrule
\endfirsthead
\toprule
\textbf{Offset} & \textbf{Size} & \textbf{Type} & \textbf{Field} & \textbf{Scale/units} & \textbf{Notes} \\
\midrule
\endhead
\bottomrule
\endlastfoot
0 & 1 & u8 & ID = 0xEC &  &  \\
1 & 2 & u16 & size = 21 & bytes &  \\
3 & 1 & u8 & version = 0x10 &  &  \\
4 & 1 & u8 & operation mode &  & B.7.3 \\
5 & 1 & u8 & ice detection status &  & B.7.3 \\
6 & 1 & u8 & urgency status &  & B.7.3 \\
7 & 1 & u8 & surface attempts this dive cycle &  &  \\
8 & 2 & u16 & stored SBD messages (packets) &  & 10.5.0: “maximum stored packets” \\
10 & 4 & u32 & time since last successful transmission & s & a duration (the spreadsheet's “since 01 Jan 2000” is an error) \\
14 & 2 & u16 & minimum (shallowest) pressure during ice detection & CTD counts &  \\
16 & 2 & u16 & threshold test temperature & see notes & spreadsheet: temperature gain/offset; 7.3.5 YAML and 10.5.0: ÷1000 → °C \\
18 & 2 & u16 & surface test salinity & CTD counts & valid only where the configuration measures salinity at the surface \\
20 & 1 & char & \code{;\allowbreak{}} &  &  \\
\end{longtable}
\endgroup

\subsubsection{\texorpdfstring{B.7.2 Version 0x11 (24 + 6N bytes)}{B.7.2 Version 0x11 (24 + 6N bytes)}}\label{sec:b72-version-0x11-24--6n-bytes}

\begingroup\footnotesize\setlength{\tabcolsep}{3pt}
\begin{longtable}{@{}>{\raggedright\arraybackslash}p{0.0625\linewidth}>{\raggedright\arraybackslash}p{0.0446\linewidth}>{\raggedright\arraybackslash}p{0.0446\linewidth}>{\raggedright\arraybackslash}p{0.3125\linewidth}>{\raggedright\arraybackslash}p{0.1071\linewidth}>{\raggedright\arraybackslash}p{0.3482\linewidth}@{}}
\toprule
\textbf{Offset} & \textbf{Size} & \textbf{Type} & \textbf{Field} & \textbf{Scale/units} & \textbf{Notes} \\
\midrule
\endfirsthead
\toprule
\textbf{Offset} & \textbf{Size} & \textbf{Type} & \textbf{Field} & \textbf{Scale/units} & \textbf{Notes} \\
\midrule
\endhead
\bottomrule
\endlastfoot
0–19 &  &  & as 0x10 &  & version = 0x11 \\
20 & 2 & u16 & seek time adjustment (10.5.0 blob list) & s & decoder: undocumented, not decoded (observed 0–30) \\
22 & 1 & u8 & error count N &  &  \\
23 & 6N &  & error entries: \code{u32} time (\code{s2000}), \code{u16} error code (B.9) &  &  \\
23+6N & 1 & char & \code{;\allowbreak{}} &  &  \\
\end{longtable}
\endgroup

The v2.4 spreadsheet and the 7.3.5 YAML have no field at 20–21: error count at 20, entries from 21, size 6N + 22 (the spreadsheet says “6N+21”). The decoder follows the 10.5.0 layout, which matches transmitted messages; a message in the spreadsheet layout would be mis-decoded.

\subsubsection{\texorpdfstring{B.7.3 Enumerations}{B.7.3 Enumerations}}\label{sec:b73-enumerations}

\textbf{Operation mode} (offset 4): 0 summer, 1 winter.

\textbf{Ice detection status} (offset 5):

\begingroup\footnotesize\setlength{\tabcolsep}{3pt}
\begin{longtable}{@{}>{\raggedright\arraybackslash}p{0.0536\linewidth}>{\raggedright\arraybackslash}p{0.5446\linewidth}>{\raggedright\arraybackslash}p{0.2232\linewidth}>{\raggedright\arraybackslash}p{0.1250\linewidth}@{}}
\toprule
\textbf{Value} & \textbf{v2.4 spreadsheet} & \textbf{10.5.0 blob list} & \textbf{SOLO-II 0xE2} \\
\midrule
\endfirsthead
\toprule
\textbf{Value} & \textbf{v2.4 spreadsheet} & \textbf{10.5.0 blob list} & \textbf{SOLO-II 0xE2} \\
\midrule
\endhead
\bottomrule
\endlastfoot
0 & ice detection disabled (should never be sent) & disabled & \code{ICE\_\allowbreak{}CHECK\_\allowbreak{}OFF} \\
1 & ice not encountered during surface attempt & not found & \code{ICE\_\allowbreak{}NOT\_\allowbreak{}FOUND} \\
2 & ice encountered during surface attempt & ice found (mixed layer) & \code{ICE\_\allowbreak{}MIXLAYER} \\
3 & temperature turnaround: threshold test failed, no further surfacing attempts & ice found (breakup) & \code{ICE\_\allowbreak{}BREAKUP} \\
4 & — & ice found (surface halt) & — \\
5 & — & complete, not found & — \\
\end{longtable}
\endgroup

The meanings of 2 and 3 differ between the spreadsheet and the firmware list; value 5 occurs in transmitted messages (float 9213). The decoder lists the number only.

\textbf{Urgency status} (offset 6): 0 not triggered (disabled), 1 triggered (enabled).

\subsubsection{\texorpdfstring{B.7.4 Decoding notes and listing (\code{E0C}, \code{e0c})}{B.7.4 Decoding notes and listing (E0C, e0c)}}\label{sec:b74-decoding-notes-and-listing-e0c-e0c}

\begin{itemize}
\item Version 0x10 has no size check beyond the 21-byte minimum; 0x11 needs at least 24 bytes and an entry count equal to \code{N} (otherwise rejected). The start time is not set; the case-vacuum members exist in the class but no version fills them.
\item \code{midx} 0: engVer, mode. \code{midx} 1 (“Status”): ice-status, urgency, surfAttempts, numStoredSBD, lastXmit (s). \code{midx} 2: minPres (dbar, from the profile pressure gain/offset; there is no zero-gain guard), testTemp, testSalinity. Version 0x11 adds \code{midx} 3: nErr (\code{E0C}), then a heading and one \code{e0c SN dive idx date-time 0x\%08X} line per error.
\item \code{test\allowbreak{}Temp} is printed as the raw \code{u16} although the heading says (degC); per 7.3.5/10.5.0 it is °C × 1000.
\item \code{test\allowbreak{}Salinity} is converted with the salinity gain/offset and range-checked, but printed with \code{\%15.\allowbreak{}3d}, which boost::format renders with three significant digits: the fill value −9.999 appears as \code{-10}.
\end{itemize}

Example (float 9213, dive 8, version 0x11):

\begin{Verbatim}
E0C 09213    8    0   0x11        1
E0C 09213    8    1               5               0               3              29           96781
E0C 09213    8    2            0.16               0             -10
E0C 09213    8    3               1
e0c 09213    8    0 2026-09-29T00:56:40Z      0x00000114
\end{Verbatim}

\subsection{\texorpdfstring{B.8 Version summary}{B.8 Version summary}}\label{sec:b8-version-summary}

\begingroup\footnotesize\setlength{\tabcolsep}{3pt}
\begin{longtable}{@{}>{\raggedright\arraybackslash}p{0.1518\linewidth}>{\raggedright\arraybackslash}p{0.1964\linewidth}>{\raggedright\arraybackslash}p{0.2857\linewidth}>{\raggedright\arraybackslash}p{0.3125\linewidth}@{}}
\toprule
\textbf{ID} & \textbf{Versions decoded} & \textbf{Size(s)} & \textbf{Documented in} \\
\midrule
\endfirsthead
\toprule
\textbf{ID} & \textbf{Versions decoded} & \textbf{Size(s)} & \textbf{Documented in} \\
\midrule
\endhead
\bottomrule
\endlastfoot
0xE0, 0xE2, 0xE3 & none (validated only) & — & SOLO-II v2.0 (version 5) \\
0xE5/0xE6 & 6, 7, 8, 9, 10, 11 & 13, 72, 76, 256, 256, 256/264 & spreadsheet (7, 8); 7.3.5 (8); 10.5.0 (9); 10.12.0 (layout of 11) \\
0xE7 & 0x10, 0x11 & 23/27, 22 & spreadsheet, 7.3.5, 10.5.0, 10.12.0 (0x11) \\
0xE8 & 0x10, 0x11, 0x12 & 19, 28, 22 & spreadsheet (0x11); 7.3.5, 10.5.0, 10.12.0 (0x12 layout) \\
0xE9 & 0x10, 0x11, 0x12 & 13, 13, 19 & spreadsheet (0x11); 7.3.5, 10.5.0, 10.12.0 (0x12 layout) \\
0xEA & 0x10–0x14 & 48+4N, 46+6N, 50+6N, 68+6N, 76+6N & spreadsheet (0x11–0x13); 7.3.5, 10.5.0, 10.12.0 (0x14) \\
0xEC & 0x10, 0x11 & 21, 24+6N & spreadsheet (0x10, 0x11); 10.5.0 (0x11) \\
\end{longtable}
\endgroup

\subsection{\texorpdfstring{B.9 System error codes}{B.9 System error codes}}\label{sec:b9-system-error-codes}

The 0xEA (versions 0x10–0x14) and 0xEC (version 0x11) error lists carry ALAMO firmware system error codes. A code is 16 bits on the wire (32 bits in 0xEA v0x10; the firmware enum \code{ASSERT\_\allowbreak{}LOG\_\allowbreak{}ERROR\_\allowbreak{}*} is 32 bits wide). The \textbf{low byte} names the module or mission state and the \textbf{high byte} the specific error within it; a high byte of 0 is the generic error of that module. For example, 0x010C is module 0x0C (CTD), error 0x01 (pressure out of range), and 0x0114 is module 0x14 (ascend), error 0x01 (CTD).

The decoder does not translate codes: listings print them as \code{0x\%08X} (e.g. \code{0x00000114}), and the Python tools pass the number through. The table merges \code{system-error-codes.\allowbreak{}txt} (names with the \code{ASSERT\_\allowbreak{}LOG\_\allowbreak{}ERROR\_\allowbreak{}} prefix) and Appendix 2 of \code{nextgen\_\allowbreak{}v10.\allowbreak{}5.\allowbreak{}0\_\allowbreak{}blobs.\allowbreak{}txt} (\code{ERROR\_\allowbreak{}} prefix).

\begingroup\footnotesize\setlength{\tabcolsep}{3pt}
\begin{longtable}{@{}>{\raggedright\arraybackslash}p{0.0625\linewidth}>{\raggedright\arraybackslash}p{0.2411\linewidth}>{\raggedright\arraybackslash}p{0.2500\linewidth}@{}}
\toprule
\textbf{Code} & \textbf{Name (\code{ERROR\_\allowbreak{}…})} & \textbf{Module} \\
\midrule
\endfirsthead
\toprule
\textbf{Code} & \textbf{Name (\code{ERROR\_\allowbreak{}…})} & \textbf{Module} \\
\midrule
\endhead
\bottomrule
\endlastfoot
0x0000 & NONE & — \\
0x0001 & ACMP & analog comparator \\
0x0101 & ACMP\_\allowbreak{}XXX &  \\
0x0201 & ACMP\_\allowbreak{}YYY &  \\
0x0002 & ADC & ADC \\
0x0003 & GPIO & GPIO \\
0x0004 & I2C & I²C \\
0x0104 & I2C\_\allowbreak{}NACK &  \\
0x0204 & I2C\_\allowbreak{}HW\_\allowbreak{}FAULT &  \\
0x0005 & PWR & power \\
0x0006 & RTC & real-time clock \\
0x0007 & TIMER & timer \\
0x0008 & UART & UART \\
0x000A & ACCEL & accelerometer \\
0x000B & BLUETOOTH & Bluetooth \\
0x000C & CTD & CTD \\
0x010C & CTD\_\allowbreak{}DATA\_\allowbreak{}RANGE\_\allowbreak{}PRESSURE &  \\
0x020C & CTD\_\allowbreak{}DATA\_\allowbreak{}RANGE\_\allowbreak{}TEMPERATURE &  \\
0x030C & CTD\_\allowbreak{}DATA\_\allowbreak{}RANGE\_\allowbreak{}SALINITY &  \\
0x040C & CTD\_\allowbreak{}DATA\_\allowbreak{}RANGE\_\allowbreak{}PAR &  \\
0x050C & CTD\_\allowbreak{}DATA\_\allowbreak{}RANGE\_\allowbreak{}TEMPSAL &  \\
0x060C & CTD\_\allowbreak{}DATA\_\allowbreak{}RANGE &  \\
0x070C & CTD\_\allowbreak{}DATA\_\allowbreak{}CONVERSION &  \\
0x080C & CTD\_\allowbreak{}RETRY &  \\
0x000D & EEPROM & EEPROM \\
0x000E & GPS & GPS \\
0x000F & PUMP & pump \\
0x010F & PUMP\_\allowbreak{}MAX\_\allowbreak{}RUN\_\allowbreak{}TIME &  \\
0x0010 & SATELLITE & Iridium modem \\
0x0110 & SATELLITE\_\allowbreak{}CONFIG &  \\
0x0210 & SATELLITE\_\allowbreak{}CMD\_\allowbreak{}TIMEOUT &  \\
0x0310 & SATELLITE\_\allowbreak{}PWR &  \\
0x0410 & SATELLITE\_\allowbreak{}DATA &  \\
0x0011 & VALVE & valve \\
0x0111 & VALVE\_\allowbreak{}CALIBRATION &  \\
0x0211 & VALVE\_\allowbreak{}CONTROL &  \\
0x0311 & VALVE\_\allowbreak{}TIMEOUT &  \\
0x0012 & MISSION & mission \\
0x0112 & MISSION\_\allowbreak{}CONFIG\_\allowbreak{}DATA &  \\
0x0212 & MISSION\_\allowbreak{}PHASE\_\allowbreak{}CONFIG &  \\
0x0013 & DESCEND & descent state \\
0x0113 & DESCEND\_\allowbreak{}CTD &  \\
0x0213 & DESCEND\_\allowbreak{}AIR\_\allowbreak{}OPEN\_\allowbreak{}VALVE &  \\
0x0313 & DESCEND\_\allowbreak{}AIR\_\allowbreak{}MID\_\allowbreak{}VALVE &  \\
0x0413 & DESCEND\_\allowbreak{}AIR\_\allowbreak{}CLOSE\_\allowbreak{}VALVE &  \\
0x0513 & DESCEND\_\allowbreak{}OIL\_\allowbreak{}OPEN\_\allowbreak{}VALVE &  \\
0x0613 & DESCEND\_\allowbreak{}OIL\_\allowbreak{}CLOSE\_\allowbreak{}VALVE &  \\
0x0713 & DESCEND\_\allowbreak{}PUMP &  \\
0x0813 & DESCEND\_\allowbreak{}DATA &  \\
0x0913 & DESCEND\_\allowbreak{}MAX\_\allowbreak{}DEPTH &  \\
0x0014 & ASCEND & ascent state \\
0x0114 & ASCEND\_\allowbreak{}CTD &  \\
0x0214 & ASCEND\_\allowbreak{}CTD\_\allowbreak{}PROFILE &  \\
0x0314 & ASCEND\_\allowbreak{}AIR\_\allowbreak{}OPEN\_\allowbreak{}VALVE &  \\
0x0414 & ASCEND\_\allowbreak{}AIR\_\allowbreak{}MID\_\allowbreak{}VALVE &  \\
0x0514 & ASCEND\_\allowbreak{}AIR\_\allowbreak{}CLOSE\_\allowbreak{}VALVE &  \\
0x0614 & ASCEND\_\allowbreak{}OIL\_\allowbreak{}OPEN\_\allowbreak{}VALVE &  \\
0x0714 & ASCEND\_\allowbreak{}OIL\_\allowbreak{}CLOSE\_\allowbreak{}VALVE &  \\
0x0814 & ASCEND\_\allowbreak{}PUMP &  \\
0x0914 & ASCEND\_\allowbreak{}DATA &  \\
0x0A14 & ASCEND\_\allowbreak{}DEPTH &  \\
0x0015 & PARK & park/drift state \\
0x0115 & PARK\_\allowbreak{}CTD &  \\
0x0215 & PARK\_\allowbreak{}CTD\_\allowbreak{}INVALID\_\allowbreak{}PRESSURE &  \\
0x0515 & PARK\_\allowbreak{}DATA &  \\
0x0615 & PARK\_\allowbreak{}DEPTH\_\allowbreak{}MIN &  \\
0x0715 & PARK\_\allowbreak{}DEPTH\_\allowbreak{}MAX &  \\
0x0815 & PARK\_\allowbreak{}DATE\_\allowbreak{}TIME\_\allowbreak{}PAST &  \\
0x0016 & SURFACE & surface state \\
0x0116 & SURFACE\_\allowbreak{}CTD &  \\
0x0216 & SURFACE\_\allowbreak{}CTD\_\allowbreak{}PROFILE &  \\
0x0316 & SURFACE\_\allowbreak{}CTD\_\allowbreak{}OFFSET\_\allowbreak{}RESET &  \\
0x0416 & SURFACE\_\allowbreak{}AIR\_\allowbreak{}OPEN\_\allowbreak{}VALVE &  \\
0x0516 & SURFACE\_\allowbreak{}AIR\_\allowbreak{}MID\_\allowbreak{}VALVE &  \\
0x0616 & SURFACE\_\allowbreak{}AIR\_\allowbreak{}CLOSE\_\allowbreak{}VALVE &  \\
0x0716 & SURFACE\_\allowbreak{}OIL\_\allowbreak{}OPEN\_\allowbreak{}VALVE &  \\
0x0816 & SURFACE\_\allowbreak{}OIL\_\allowbreak{}CLOSE\_\allowbreak{}VALVE &  \\
0x0916 & SURFACE\_\allowbreak{}PUMP &  \\
0x0A16 & SURFACE\_\allowbreak{}DATA &  \\
0x0B16 & SURFACE\_\allowbreak{}DATA\_\allowbreak{}MAX\_\allowbreak{}LOOP &  \\
0x0C16 & SURFACE\_\allowbreak{}GPS &  \\
0x0D16 & SURFACE\_\allowbreak{}SATELLITE &  \\
0x0E16 & SURFACE\_\allowbreak{}CFG\_\allowbreak{}STORE &  \\
0x0F16 & SURFACE\_\allowbreak{}REMOTE\_\allowbreak{}CMD &  \\
0x1016 & SURFACE\_\allowbreak{}BEACON\_\allowbreak{}DEPTH &  \\
0x2016 & SURFACE\_\allowbreak{}CTD\_\allowbreak{}CRC & 10.5.0 appendix only \\
0x0020 & BIT & built-in test \\
0x0120 & BIT\_\allowbreak{}SATELLITE &  \\
0x0030 & MAIN & main loop \\
0x00FF & EXCEPTION & \code{system-error-codes.\allowbreak{}txt} only \\
\end{longtable}
\endgroup

No code is defined for module 0x09 or for 0x0315/0x0415 (park); a code not in the table should be reported as its module plus the raw high byte.

\subsection{\texorpdfstring{B.10 Discrepancies and open questions}{B.10 Discrepancies and open questions}}\label{sec:b10-discrepancies-and-open-questions}

Between the decoder and the documents:

\begin{enumerate}
\item \textbf{0xEA v0x13/v0x14 bin counts.} The decoder reads offset 36 (v0x13) / 44 (v0x14) as the number of \emph{empty bins} and 38 / 46 as the number of \emph{raw samples}. The v2.4 spreadsheet, the 7.3.5 YAML, the 10.5.0 list and the 10.12.0 list all give the opposite order (samples, then empty bins). Listings checked so far (floats 9213, 9216, 9287, 11186) have 0 in both fields, so data cannot settle it yet; the documents suggest the decoder's \code{num\allowbreak{}Empty\allowbreak{}Bins} and \code{num\allowbreak{}Raw\allowbreak{}Samples} columns are swapped. \emph{(Fixed in build 3309: the decoder now reads samples before empty bins, as the documents give them.)}
\item \textbf{0xEC v0x11 bytes 20–21.} Decoder and 10.5.0 list: a 2-byte field (10.5.0: “seek time adjustment”) before the error count at 22. Spreadsheet v2.4 and 7.3.5 YAML: error count at

\begin{enumerate}[start=20]
\item The decoder matches transmitted messages; whether early-firmware 0x11 messages used the spreadsheet layout is unknown.
\end{enumerate}
\item \textbf{Ice detection status values.} Spreadsheet: 0–3 (3 = temperature turnaround). 10.5.0: 0–5, with 2 = mixed layer, 3 = breakup, 4 = surface halt, 5 = complete, not found. Value 5 is seen in data. The firmware list is presumably current.
\item \textbf{0xEC threshold test temperature.} Spreadsheet: temperature counts (gain/offset). 7.3.5 and 10.5.0: °C × 1000. The decoder prints the raw number under a (degC) heading.
\item \textbf{BIT format 7/8 scales.} Air-pump maximum-load battery voltage: decoder ÷1000, spreadsheet and YAML ÷100. Oil-pump battery voltage: decoder and spreadsheet ÷1000, YAML ÷100. Oil pump test duration: spreadsheet gain 1000, decoder and YAML seconds.
\item \textbf{BIT format 8 signedness.} The pressure offset (and, per the YAML, both pressures) are signed; the decoder reads them unsigned.
\item \textbf{BIT format 8 in the spreadsheet.} The 0xE6 sheet gives a 16-bit CTD status and shifted offsets, and the 0xE5/0xE6 sheets swap the format 7/8 sizes (72/77); the 0xE5 size “77” conflicts with its own terminator at 75.
\item \textbf{BIT format 11 sizes.} The 10.12.0 list describes the format 11 layout (traveler errors, 20 thresholds, 264 bytes) but labels it \code{Packet\allowbreak{}Version =\allowbreak{} 9}; a 10.12.0.1 float sends 11. The decoder also accepts a 256-byte format 11 that no document describes.
\item \textbf{BIT format 11 field names (10.12.0 list).} Offsets 141 and 175 are named “MaxAirPumpCurrent”/”MaxOilPumpCurrent” (duplicates) and 145/179 “Max…V12Current”; by comparison with 10.5.0 they are the 12 V current maximum and the battery voltage maximum, as the decoder reads them.
\item \textbf{BIT result bit 0x400} is set in real data and undocumented; bit 0x200 (accelerometer) is in 10.5.0 only.
\item \textbf{BIT format 9 “maximum air pressure”} is “inHg” in the 10.5.0 list; the decoder lists psi.
\item \textbf{E8/E9 version in 10.12.0.} The list gives \code{packet\allowbreak{}Version =\allowbreak{} 17} (0x11) for 0xE8 and 0xE9 but the 0x12 layouts (22 and 19 bytes). If a float actually sent 0x11 with those sizes the decoder would reject it (it expects 28 and 13). A 10.12.0.1 float (11186) sends 0x12 for both, so the list is wrong on this point.
\item \textbf{Seek pump time adjustment type.} \code{I2} in 10.5.0 and 7.3.5, \code{U2} in 10.12.0; the decoder treats it as signed.
\item \textbf{0xE7 case vacuum scale.} 10.5.0 says ÷1000 for 0xE7 only (÷100 for 0xE8/0xE9 and in all other documents); the decoder uses ÷100.
\item \textbf{0xEA pressure offset reading} is \code{U2} in the 10.x lists, signed elsewhere; the decoder reads it signed.
\item \textbf{BIT format 9–11 “divisor” values} in the 10.5.0 list do not apply to the transmitted floats, which are already in engineering units.
\item \textbf{7.3.5 YAML} omits the version byte for 0xE7–0xEC (its field lists start at offset 4) and uses the surface gains for BIT pressures; the decoder uses the profile gains.
\item \textbf{Error-code tables.} \code{SURFACE\_\allowbreak{}CTD\_\allowbreak{}CRC} (0x2016) is only in the 10.5.0 appendix and \code{EXCEPTION} (0x00FF) only in \code{system-error-codes.\allowbreak{}txt}.
\item \textbf{Unexplained bytes.} 0xE7 v0x10 bytes 14–21 (23-byte form) / 20–25 (27-byte form), 0xEA v0x10 bytes 22–45 and 0xEA v0x12 bytes 12–13 are undocumented and not decoded.
\end{enumerate}

Decoder behavior worth knowing (not format discrepancies):

\begin{itemize}
\item 0xEA v0x10 with errors: the error heading has more format fields than values, which makes boost::format throw during listing generation (inferred).
\item 0xEA \code{n\allowbreak{}Err} lines are tagged \code{e0A} in v0x11/0x12 and \code{E0A} in v0x13/0x14.
\item BIT formats 9–11 are accepted at any size (warning only) because \code{XMSG\_\allowbreak{}DEBUG\_\allowbreak{}BIST} is on.
\item \code{XEng\allowbreak{}Ascent}, \code{XEng\allowbreak{}Ice} and \code{XAlamo\allowbreak{}Bit2} leave some members uninitialized; the listings do not print those members.
\end{itemize}

\clearpage
\section{\texorpdfstring{Part C — Mission configuration, command echo, waves and acoustics}{Part C — Mission configuration, command echo, waves and acoustics}}\label{sec:part-c--mission-configuration-command-echo-waves-and-acoustics}

This part covers these blob types:

\begingroup\footnotesize\setlength{\tabcolsep}{3pt}
\begin{longtable}{@{}>{\raggedright\arraybackslash}p{0.0982\linewidth}>{\raggedright\arraybackslash}p{0.2946\linewidth}>{\raggedright\arraybackslash}p{0.3036\linewidth}>{\raggedright\arraybackslash}p{0.2500\linewidth}@{}}
\toprule
\textbf{ID} & \textbf{Name} & \textbf{Decoder class} & \textbf{Status in the XMTPS decoder} \\
\midrule
\endfirsthead
\toprule
\textbf{ID} & \textbf{Name} & \textbf{Decoder class} & \textbf{Status in the XMTPS decoder} \\
\midrule
\endhead
\bottomrule
\endlastfoot
0xD0-0xD7 & Mission configuration & \code{XMfa\allowbreak{}Mission\allowbreak{}Cfg} (\code{XAlamo\allowbreak{}Mission\allowbreak{}Configuration} is the unused older version) & D0 and D1 are decoded for 117 firmware versions. D2-D7 are rejected. SOLO-II ASCII dumps are not decoded. \\
0xDE & Mission command echo & \code{XMission\allowbreak{}Cmd\allowbreak{}Echo} & Decoded since build 3310; one \code{CMD} listing line per command. \\
0xC1, 0xC3 & Ocean surface waves & \code{XWave}, \code{XWave\allowbreak{}Spectral\allowbreak{}Sample}, \code{XWave\allowbreak{}Raw\allowbreak{}Sample} & Spectral versions 0-2 are decoded. The raw IMU format is not. \\
0xCA & Marginal-ice-zone (MIZ) acoustics & \code{XAcoustic\allowbreak{}Position} & Subtypes 3 and 5 are decoded. Subtypes 1 and 2 are seen in production but ignored. \\
0xFF & Extended records & \code{XExtended\allowbreak{}Record} & Decoded since build 3321; one \code{XFF} listing line per record. The format is inferred from production messages (C.6). \\
0xC5 & WIMS configuration parameters & \code{XWims\allowbreak{}Parameters} & Decoded since build 3336; one \code{WIP} listing line per parameter. The format and most parameter names are inferred from production messages (C.7). \\
\end{longtable}
\endgroup

\subsection{\texorpdfstring{C.0 Conventions}{C.0 Conventions}}\label{sec:c0-conventions}

\begin{itemize}
\item Offsets count from the blob's ID byte, which is offset 0. The common blob header is \code{ID} (offset 0, U1), \code{jj} (offsets 1-2, U2, the blob length including ID, \code{jj}, data and terminator) and the \code{;\allowbreak{}} (0x3B) terminator at offset \code{jj - 1}.
\item Multi-byte integers are big-endian. Types: U1, U2 and U4 are unsigned 8-, 16- and 32-bit integers; I1, I2 and I4 are two's-complement signed; R4 is an IEEE-754 single-precision float sent as a big-endian 32-bit word (the decoder reads a U4 and bit-casts it, \code{f32(\allowbreak{})} in \code{XDefinitions.\allowbreak{}h}).
\item \textbf{[inferred]} marks something worked out from production messages or decoder behavior, not from a format document. \textbf{[decoder]} marks behavior of the XMTPS implementation, not of the format itself.
\item Sources:

\begin{itemize}
\item \emph{Decoder}: \code{src/\allowbreak{}data/\allowbreak{}xmsg/\allowbreak{}} in the XMTPS repository, as of 2026-10-25.
\item \emph{XF2.4}: the ALAMO X-format spec v2.4 spreadsheet (\code{xfmt-format-spec-2.\allowbreak{}4-alamo\_\allowbreak{}2018-05-15.\allowbreak{}xlsx}).
\item \emph{SOLO2}: \emph{Format of SOLOII X messages}, v2.0, 20 Mar 2014.
\item \emph{NG10.5}: \code{nextgen\_\allowbreak{}v10.\allowbreak{}5.\allowbreak{}0\_\allowbreak{}blobs.\allowbreak{}txt}.
\item \emph{NG10.12}: \code{nextgen-10.\allowbreak{}12.\allowbreak{}0-xmessage.\allowbreak{}txt}.
\item \emph{DF7.3.5}: \code{Data\allowbreak{}Format-v7.\allowbreak{}3.\allowbreak{}5.\allowbreak{}yaml}.
\item \emph{Production}: hex dumps (\code{shx} lines) in production listings under \code{Z:\allowbreak{}\textbackslash{}\allowbreak{}alamo\textbackslash{}\allowbreak{}data\textbackslash{}\allowbreak{}list}.
\end{itemize}
\item Fill values. The decoder starts every mission parameter at an “unset” value: 255 (U1), 65535 (U2), 4294967295 (U4), -32768 (I2) or \code{FLT\_\allowbreak{}MIN} (R4). A parameter that a firmware layout doesn't contain is listed with that value.
\end{itemize}

\subsection{\texorpdfstring{C.1 Mission configuration (0xD0-0xD7)}{C.1 Mission configuration (0xD0-0xD7)}}\label{sec:c1-mission-configuration-0xd0-0xd7}

\subsubsection{\texorpdfstring{C.1.1 Purpose}{C.1.1 Purpose}}\label{sec:c11-purpose}

The float sends its mission parameters: the dive-cycle, CTD, descent, park, ascent, surface, ice-detection, IMU and calibration settings stored in its non-volatile memory. ALAMO and ALTO floats send them in a binary format at each surfacing when \code{APT} = 1, as two blobs, D0 and D1. SOLO-II floats send an ASCII dump of their EEPROM instead, only in answer to a ground-station \code{P} command. The low nibble of the ID is the part index (\code{D0} = part 0, \code{D1} = part 1, and so on).

\subsubsection{\texorpdfstring{C.1.2 The two encodings}{C.1.2 The two encodings}}\label{sec:c12-the-two-encodings}

\paragraph{\texorpdfstring{SOLO-II ASCII EEPROM dump (SOLO2 p. 13)}{SOLO-II ASCII EEPROM dump (SOLO2 p. 13)}}\label{sec:solo-ii-ascii-eeprom-dump-solo2-p-13}
\mbox{}\par

\begingroup\footnotesize\setlength{\tabcolsep}{3pt}
\begin{longtable}{@{}>{\raggedleft\arraybackslash}p{0.0625\linewidth}>{\raggedleft\arraybackslash}p{0.0625\linewidth}>{\raggedright\arraybackslash}p{0.0536\linewidth}>{\raggedright\arraybackslash}p{0.1339\linewidth}>{\raggedright\arraybackslash}p{0.6161\linewidth}@{}}
\toprule
\textbf{Offset} & \textbf{Size} & \textbf{Type} & \textbf{Field} & \textbf{Notes} \\
\midrule
\endfirsthead
\toprule
\textbf{Offset} & \textbf{Size} & \textbf{Type} & \textbf{Field} & \textbf{Notes} \\
\midrule
\endhead
\bottomrule
\endlastfoot
0 & 1 & U1 & ID & 0xD0, 0xD1, 0xD2, 0xD3 (SOLO2: “possible values 0:d” \\
1 & 2 & U2 & \code{jj} & Variable. \\
3 & \code{jj} - 4 & ASCII & Parameter list & \code{NAMExx=\allowbreak{}vvvvv|\allowbreak{}} records, 6-character name, \code{=\allowbreak{}}, 5-character value and \code{|\allowbreak{}}, with no gaps between records. Example: \code{Pch\allowbreak{}Sec=\allowbreak{}   -1|\allowbreak{}Max\allowbreak{}Hrs=\allowbreak{} 1440|\allowbreak{}d\allowbreak{}Bar\allowbreak{}Go=\allowbreak{}   -1|\allowbreak{}.\allowbreak{}.\allowbreak{}.\allowbreak{}|\allowbreak{} XP2=\allowbreak{}    0|\allowbreak{} XPdly=\allowbreak{}    0|\allowbreak{}}. \\
\code{jj} - 1 & 1 & char & \code{;\allowbreak{}} & Terminator. \\
\end{longtable}
\endgroup

SOLO2 says the dump takes 4 SBD messages: D0, D1 and D2 of 328 bytes each, and D3 of 7 bytes. It also says each parameter takes 15 bytes, but the record template it gives has 13 characters (see C.5).

The two encodings are told apart by byte 3. XF2.4 calls it “message encoding format type identifier”: 0 = binary (ALAMO), 1-255 = ASCII (SOLO-II). In an ASCII dump, byte 3 is the first character of the first parameter name, so it is never 0. NG10.12 calls the same byte \code{Mission\allowbreak{}Config\allowbreak{}Version} and gives it the value 0.

{}[decoder] The only code that checks byte 3 is the unused \code{XMission\allowbreak{}Configuration} wrapper: it treats \code{d[3] =\allowbreak{}=\allowbreak{} 0} on part 0 as binary and leaves everything else undecoded. The active decoder, \code{XMfa\allowbreak{}Mission\allowbreak{}Cfg}, never reads byte 3. A SOLO-II ASCII D0 would go through the binary path and fail as an unsupported firmware version, because bytes 4-7 would be ASCII characters.

\paragraph{\texorpdfstring{ALAMO/ALTO binary format}{ALAMO/ALTO binary format}}\label{sec:alamoalto-binary-format}
\mbox{}\par

The rest of C.1 covers the binary format.

\subsubsection{\texorpdfstring{C.1.3 Decoder families}{C.1.3 Decoder families}}\label{sec:c13-decoder-families}

\begingroup\footnotesize\setlength{\tabcolsep}{3pt}
\begin{longtable}{@{}>{\raggedright\arraybackslash}p{0.2321\linewidth}>{\raggedright\arraybackslash}p{0.2321\linewidth}>{\raggedright\arraybackslash}p{0.2411\linewidth}>{\raggedright\arraybackslash}p{0.2411\linewidth}@{}}
\toprule
\textbf{Family} & \textbf{Source} & \textbf{Firmware} & \textbf{Status} \\
\midrule
\endfirsthead
\toprule
\textbf{Family} & \textbf{Source} & \textbf{Firmware} & \textbf{Status} \\
\midrule
\endhead
\bottomrule
\endlastfoot
\code{XAlamo\allowbreak{}Mission\allowbreak{}Configuration} & \code{XAlamo\allowbreak{}Mission\allowbreak{}Configuration.\allowbreak{}h} (wrapped by \code{XMission\allowbreak{}Configuration.\allowbreak{}h}) & 0x02006200 to 0x05006100. It decodes one message, using offsets equal to the MFA layouts of the same versions. & Not used: nothing creates an \code{XMission\allowbreak{}Configuration}. This is the predecessor of the MFA decoder. \\
\code{XMfa\allowbreak{}Mission\allowbreak{}Cfg} (“MFA” mission parameter message) & \code{XMfa\allowbreak{}Mission\allowbreak{}Cfg.\allowbreak{}cpp/\allowbreak{}.\allowbreak{}h} & 117 firmware words (C.1.6). Each firmware word maps to one \code{Decode\allowbreak{}XMsg\allowbreak{}Data\allowbreak{}Fw<layout>} function and one \code{Generate\allowbreak{}Listing\allowbreak{}String\allowbreak{}Fw<layout>} function. & Active. \code{XFloat.\allowbreak{}cpp} sends every 0xD0-0xD7 blob to \code{XMfa\allowbreak{}Mission\allowbreak{}Cfg:\allowbreak{}:\allowbreak{}Assign\allowbreak{}Message}. \\
\end{longtable}
\endgroup

\subsubsection{\texorpdfstring{C.1.4 Common header (binary, part 0)}{C.1.4 Common header (binary, part 0)}}\label{sec:c14-common-header-binary-part-0}

\begingroup\footnotesize\setlength{\tabcolsep}{3pt}
\begin{longtable}{@{}>{\raggedleft\arraybackslash}p{0.0625\linewidth}>{\raggedleft\arraybackslash}p{0.0446\linewidth}>{\raggedright\arraybackslash}p{0.0446\linewidth}>{\raggedright\arraybackslash}p{0.3214\linewidth}>{\raggedright\arraybackslash}p{0.4554\linewidth}@{}}
\toprule
\textbf{Offset} & \textbf{Size} & \textbf{Type} & \textbf{Field} & \textbf{Notes} \\
\midrule
\endfirsthead
\toprule
\textbf{Offset} & \textbf{Size} & \textbf{Type} & \textbf{Field} & \textbf{Notes} \\
\midrule
\endhead
\bottomrule
\endlastfoot
0 & 1 & U1 & ID & 0xD0. \\
1 & 2 & U2 & \code{jj} & 322-328 bytes depending on the firmware (C.1.6). [decoder] Not read by \code{XMfa\allowbreak{}Mission\allowbreak{}Cfg}, which checks the blob's actual length instead. \\
3 & 1 & U1 & Message config version / encoding & 0 = binary. NG10.12: \code{Mission\allowbreak{}Config\allowbreak{}Version =\allowbreak{} 0}. NG10.5: “Message version”. XF2.4: encoding format (0 = binary). [decoder] Not read. \\
4 & 1 & U1 & Firmware major & For example 0x0A = 10. \\
5 & 1 & U1 & Firmware minor & For example 0x0C = 12. \\
6 & 1 & U1 & Firmware patch &  \\
7 & 1 & U1 & Firmware build / product & NG10.5: “Float model”. NG10.12: \code{Firmware\allowbreak{}Version-Product}. The values seen are variant codes: 0x31 = WIMS, 0x32 = MM2/Micro-ALTO (10.12.0.32), 0x33 = Micro-ALTO (10.6.1.33), 0x01 or 0x04 = RBR or SBE builds. \\
4 & 4 & U4 & Firmware word & {}[decoder] Bytes 4-7 are read as one big-endian U4, \code{0x\allowbreak{}MMmmppbb}, which selects the layout. It is stored as the float's firmware version, and only a D0 updates it. Version strings in this document follow the decoder's convention: major and minor in decimal, patch and build written as their hex digits (0x0A0C0032 = 10.12.0.32, 0x02056200 = 2.5.62 \\
8 & 1 & U1 & Mission parameters (packet) version & From firmware 4.2.62 on. NG10.5 and NG10.12 both document 23 (0x17). Production 10.12.0.x floats send 0x18 (24), and 8.0.1.1 floats send 0x17. Listed as \code{Cfg\allowbreak{}Msg\allowbreak{}Ver}. In firmware before 4.2.62 the parameters start here (\code{MSN} at offset 8 \\
9 & 1 & U1 & CTD (sensor) type & From 4.2.62 on. 0 = RBR, 1 = SBE (Sea-Bird). XF2.4: high nibble = manufacturer (0 RBR, 1 Seabird), low nibble = sensor type (reserved). The decoder compares the whole byte, so 0x10 would be listed as \code{unknown}. The type selects how \code{CPN}/\code{CPP} are listed, and whether surface pressure offset correction applies (RBR corrected, SBE not \\
\end{longtable}
\endgroup

Production example (float 9213, 10.12.0.32, dive 3): \code{D0 01 47 |\allowbreak{} 00 |\allowbreak{} 0A 0C 00 32 |\allowbreak{} 18 |\allowbreak{} 00 |\allowbreak{} 00 03 .\allowbreak{}.\allowbreak{}.\allowbreak{}}, which decodes as \code{jj} = 327, binary, firmware 10.12.0.32, parameters version 24, RBR, MSN = 3.

Part 1 (D1) has no header: its parameters start at offset 3, right after \code{jj}.

\subsubsection{\texorpdfstring{C.1.5 Splitting into D0 and D1, assembly and dispatch}{C.1.5 Splitting into D0 and D1, assembly and dispatch}}\label{sec:c15-splitting-into-d0-and-d1-assembly-and-dispatch}

\begin{itemize}
\item From 4.5.63 on, the parameter list is longer than one SBD-sized blob and is cut in two. D0 holds the header and the first parameters, and D1 holds the rest from offset 3. Firmware before 4.5.63 sends D0 only, except 4.6.61: its D0 ends at \code{SAN} like 4.5.63, but the decoder doesn't decode a D1 for it.
\item Where the cut falls depends on the firmware:

\begin{itemize}
\item 5.0.61 to 10.6.x (D0 = 328 bytes): D0 ends with \code{PFV1} (offsets 325-326). D1 starts with \code{PFV2}. NG10.5 numbers the parameters as one contiguous blob; D1 offset = NG10.5 offset - 324 (for example \code{PFV2}: 327 → 3, \code{PPE}: 415 → 91).
\item 6.5.77 and 6.8.5: D0 ends with \code{PFV0}, and D1 starts with \code{PFV1}.
\item 5.0.61.1, and 10.7 and later (D0 = 327 bytes): D0 ends with \code{PSM} (offset 325). D1 starts with \code{PSD}. In 10.7 the \code{PNB}/\code{PNC}/\code{PXC} BIT thresholds (6 bytes) were dropped, so \code{PFV2}, \code{PDT} and \code{PSM} (5 bytes) now fit in D0.
\end{itemize}
\item Assembly [decoder] (\code{Assign\allowbreak{}Message}): the part index is \code{ID \& 0x0F}. Only indexes 0 and 1 are accepted (\code{XMSG\_\allowbreak{}MAX\_\allowbreak{}MFA\_\allowbreak{}MISS} = 2), so D2-D7 are rejected with “Message index exceeds”. Each part is stored per dive, and a repeated part replaces the earlier one.

\begin{itemize}
\item On D0: the firmware word is read from bytes 4-7 (the blob must be at least 8 bytes), D0 is decoded, and a stored D1 is decoded again.
\item On D1: it is decoded only if D0 is already stored. Otherwise it is kept until D0 arrives.
\end{itemize}
\item {}[decoder] The decoder never looks at the D0 parameters version (offset 8) to choose a layout. Only the firmware word does.
\end{itemize}

\subsubsection{\texorpdfstring{C.1.6 Supported firmware versions}{C.1.6 Supported firmware versions}}\label{sec:c16-supported-firmware-versions}

117 firmware words are decoded. “Decode layout” is the \code{Decode\allowbreak{}XMsg\allowbreak{}Data\allowbreak{}Fw…} function and “Listing function” is the \code{Generate\allowbreak{}Listing\allowbreak{}String\allowbreak{}Fw…} function. “D0 bytes” and “D1 bytes” are the exact blob lengths the decoder expects. A version marked “layout inferred” was mapped to an existing layout after checking production messages (README, 2026-10-25), not from an MRV document.

\begingroup\footnotesize\setlength{\tabcolsep}{3pt}
\begin{longtable}{@{}>{\raggedright\arraybackslash}p{0.0982\linewidth}>{\raggedright\arraybackslash}p{0.0982\linewidth}>{\raggedright\arraybackslash}p{0.1071\linewidth}>{\raggedleft\arraybackslash}p{0.0804\linewidth}>{\raggedleft\arraybackslash}p{0.0804\linewidth}>{\raggedright\arraybackslash}p{0.1607\linewidth}>{\raggedright\arraybackslash}p{0.1339\linewidth}>{\raggedright\arraybackslash}p{0.1339\linewidth}@{}}
\toprule
\textbf{FW word} & \textbf{Version} & \textbf{Decode layout} & \textbf{D0 bytes} & \textbf{D1 bytes} & \textbf{Listing function (style)} & \textbf{Layout inferred} & \textbf{Decoder comment} \\
\midrule
\endfirsthead
\toprule
\textbf{FW word} & \textbf{Version} & \textbf{Decode layout} & \textbf{D0 bytes} & \textbf{D1 bytes} & \textbf{Listing function (style)} & \textbf{Layout inferred} & \textbf{Decoder comment} \\
\midrule
\endhead
\bottomrule
\endlastfoot
\code{0x02006200} & 2.0.62.0 & \code{02006200} & 323 & - & \code{02006200} (\code{mis} groups) &  &  \\
\code{0x02006300} & 2.0.63.0 & \code{02006200} & 323 & - & \code{02006200} (\code{mis} groups) &  &  \\
\code{0x02006700} & 2.0.67.0 & \code{02006200} & 323 & - & \code{02006200} (\code{mis} groups) &  &  \\
\code{0x02006800} & 2.0.68.0 & \code{02006200} & 323 & - & \code{02006200} (\code{mis} groups) &  &  \\
\code{0x02006900} & 2.0.69.0 & \code{02006900} & 322 & - & \code{02006900} (\code{mis} groups) &  &  \\
\code{0x02036200} & 2.3.62.0 & \code{02006900} & 322 & - & \code{02006900} (\code{mis} groups) &  &  \\
\code{0x02046200} & 2.4.62.0 & \code{02006900} & 322 & - & \code{02006900} (\code{mis} groups) &  &  \\
\code{0x02056200} & 2.5.62.0 & \code{02006900} & 322 & - & \code{02006900} (\code{mis} groups) &  &  \\
\code{0x03006500} & 3.0.65.0 & \code{03006500} & 328 & - & \code{02006900} (\code{mis} groups) &  &  \\
\code{0x03016700} & 3.1.67.0 & \code{03016700} & 328 & - & \code{02006900} (\code{mis} groups) &  & Solo-II Hybrid \#10083 \\
\code{0x03026100} & 3.2.61.0 & \code{03026100} & 328 & - & \code{02006900} (\code{mis} groups) &  & 2016 NASCar floats \\
\code{0x04016100} & 4.1.61.0 & \code{04016100} & 328 & - & \code{02006900} (\code{mis} groups) &  & 2016 NASCar floats \\
\code{0x04026200} & 4.2.62.0 & \code{04026200} & 323 & - & \code{04026200} (\code{MIS} + \code{mis} groups) &  & PAR floats \\
\code{0x04036100} & 4.3.61.0 & \code{04036100} & 327 & - & \code{04026200} (\code{MIS} + \code{mis} groups) &  & 2017 S2h test float 10108 \\
\code{0x04046100} & 4.4.61.0 & \code{04046100} & 328 & - & \code{04026200} (\code{MIS} + \code{mis} groups) &  & 2017 Hurricane Floats \\
\code{0x04056300} & 4.5.63.0 & \code{04056300} & 328 & >= 12 & \code{04056300} (\code{MIS} + \code{m}) &  & 2017 S2h test float 10108 \\
\code{0x04066100} & 4.6.61.0 & \code{04066100} & 328 & - & \code{04026200} (\code{MIS} + \code{mis} groups) &  & 2017 S2h test float 10108 \\
\code{0x05006100} & 5.0.61.0 & \code{05006100} & 328 & 79 & \code{05006100} (\code{MIS} + \code{m}) &  & 2017 Alto test float 11011 \\
\code{0x05006101} & 5.0.61.1 & \code{05006101} & 328 & 74 & \code{05006100} (\code{MIS} + \code{m}) &  & alternate (2017 Alto test float 10109) \\
\code{0x06000000} & 6.0.0.0 & \code{06000000} & 328 & 81 & \code{06000000} (\code{MIS} + \code{m}) &  &  \\
\code{0x06000100} & 6.0.1.0 & \code{06000100} & 328 & 81 & \code{06000000} (\code{MIS} + \code{m}) &  &  \\
\code{0x06010300} & 6.1.3.0 & \code{06010300} & 328 & 81 & \code{06000000} (\code{MIS} + \code{m}) &  & 2019 ALTO RBR/O2 test float 11060 \\
\code{0x06040900} & 6.4.9.0 & \code{06040900} & 328 & 81 & \code{06000000} (\code{MIS} + \code{m}) &  & v6.4.9: K. Wood's 2019 Arctic PAR floats \\
\code{0x06050100} & 6.5.1.0 & \code{06050100} & 328 & 77 & \code{06050100} (\code{MIS} + \code{m}) &  &  \\
\code{0x06057700} & 6.5.77.0 & \code{06057700} & 328 & 79 & \code{06057700} (\code{MIS} + \code{m}) &  & v6.9.w: B. Sanabia's 2019 Hurricane floats w/ WIMS \\
\code{0x06080500} & 6.8.5.0 & \code{06080500} & 328 & 79 & \code{06080500} (\code{MIS} + \code{m}) &  & 2019 ALTO \#11011 test float \\
\code{0x07020200} & 7.2.2.0 & \code{07020200} & 328 & 76 & \code{07020200} (\code{MIS} + \code{m}) &  & 2020 ALTO \#11011 test float \\
\code{0x07020300} & 7.2.3.0 & \code{07020200} & 328 & 76 & \code{07020200} (\code{MIS} + \code{m}) &  & 2020 ALTO \#11011 test float \\
\code{0x07020400} & 7.2.4.0 & \code{07020200} & 328 & 76 & \code{07020200} (\code{MIS} + \code{m}) &  & 2020 ALTO \#11011 test float \\
\code{0x07020500} & 7.2.5.0 & \code{07020200} & 328 & 76 & \code{07020200} (\code{MIS} + \code{m}) &  & 2020 ALTO \#11011 test float \\
\code{0x07020600} & 7.2.6.0 & \code{07020200} & 328 & 76 & \code{07020200} (\code{MIS} + \code{m}) &  & 2020 ALTO \#11011 test float \\
\code{0x07020700} & 7.2.7.0 & \code{07020200} & 328 & 76 & \code{07020200} (\code{MIS} + \code{m}) &  & 2020 ALTO \#11011 test float \\
\code{0x07020800} & 7.2.8.0 & \code{07020200} & 328 & 76 & \code{07020200} (\code{MIS} + \code{m}) &  & 2020 ALTO \#11011 test float \\
\code{0x07020900} & 7.2.9.0 & \code{07020200} & 328 & 76 & \code{07020200} (\code{MIS} + \code{m}) &  & 2020 ALTO \#11011 test float \\
\code{0x07030000} & 7.3.0.0 & \code{07020200} & 328 & 76 & \code{07020200} (\code{MIS} + \code{m}) &  & 2020 ALTO RBR/ODO \#11060 test float \\
\code{0x07030400} & 7.3.4.0 & \code{07020200} & 328 & 76 & \code{07020200} (\code{MIS} + \code{m}) & yes &  \\
\code{0x07030500} & 7.3.5.0 & \code{07020200} & 328 & 76 & \code{07020200} (\code{MIS} + \code{m}) &  & 2020 ALTO RBR/ODO \#11060 test float \\
\code{0x08000000} & 8.0.0.0 & \code{08000000} & 328 & 78 & \code{08000000} (\code{MIS} + \code{m}) &  &  \\
\code{0x08000001} & 8.0.0.1 & \code{08000000} & 328 & 78 & \code{08000000} (\code{MIS} + \code{m}) &  &  \\
\code{0x08000100} & 8.0.1.0 & \code{08000000} & 328 & 78 & \code{08000000} (\code{MIS} + \code{m}) &  &  \\
\code{0x08000101} & 8.0.1.1 & \code{08000000} & 328 & 78 & \code{08000000} (\code{MIS} + \code{m}) &  & 2020 Hurricane floats \\
\code{0x08000200} & 8.0.2.0 & \code{08000000} & 328 & 78 & \code{08000000} (\code{MIS} + \code{m}) &  &  \\
\code{0x08000300} & 8.0.3.0 & \code{08000000} & 328 & 78 & \code{08000000} (\code{MIS} + \code{m}) &  &  \\
\code{0x08000400} & 8.0.4.0 & \code{08000000} & 328 & 78 & \code{08000000} (\code{MIS} + \code{m}) &  &  \\
\code{0x08000500} & 8.0.5.0 & \code{08000000} & 328 & 78 & \code{08000000} (\code{MIS} + \code{m}) &  &  \\
\code{0x08010000} & 8.1.0.0 & \code{08010000} & 328 & 76 & \code{08010000} (\code{MIS} + \code{m}) &  &  \\
\code{0x08010001} & 8.1.0.1 & \code{08000000} & 328 & 78 & \code{08000000} (\code{MIS} + \code{m}) &  & 2021 Hurricane floats \\
\code{0x08020000} & 8.2.0.0 & \code{08010000} & 328 & 76 & \code{Unsupported} (generic “unsupported” listing; see notes) &  &  \\
\code{0x09000000} & 9.0.0.0 & \code{08010000} & 328 & 76 & \code{08010000} (\code{MIS} + \code{m}) &  &  \\
\code{0x0A000000} & 10.0.0.0 & \code{08010000} & 328 & 76 & \code{08010000} (\code{MIS} + \code{m}) &  &  \\
\code{0x0A000100} & 10.0.1.0 & \code{08010000} & 328 & 76 & \code{08010000} (\code{MIS} + \code{m}) &  &  \\
\code{0x0A020000} & 10.2.0.0 & \code{08010000} & 328 & 76 & \code{08010000} (\code{MIS} + \code{m}) &  &  \\
\code{0x0A020100} & 10.2.1.0 & \code{08010000} & 328 & 76 & \code{08010000} (\code{MIS} + \code{m}) &  &  \\
\code{0x0A020200} & 10.2.2.0 & \code{08010000} & 328 & 76 & \code{08010000} (\code{MIS} + \code{m}) &  &  \\
\code{0x0A020201} & 10.2.2.1 & \code{08000000} & 328 & 78 & \code{Unsupported} (generic “unsupported” listing; see notes) &  &  \\
\code{0x0A020301} & 10.2.3.1 & \code{0A050000} & 328 & 93 & \code{0A050000} (\code{MIS} + \code{m}) & yes &  \\
\code{0x0A030000} & 10.3.0.0 & \code{08010000} & 328 & 76 & \code{08010000} (\code{MIS} + \code{m}) &  &  \\
\code{0x0A030100} & 10.3.1.0 & \code{08010000} & 328 & 76 & \code{08010000} (\code{MIS} + \code{m}) &  &  \\
\code{0x0A040000} & 10.4.0.0 & \code{0A050000} & 328 & 93 & \code{0A050000} (\code{MIS} + \code{m}) & yes &  \\
\code{0x0A040400} & 10.4.4.0 & \code{0A050000} & 328 & 93 & \code{0A050000} (\code{MIS} + \code{m}) & yes &  \\
\code{0x0A040500} & 10.4.5.0 & \code{0A050000} & 328 & 93 & \code{0A050000} (\code{MIS} + \code{m}) & yes &  \\
\code{0x0A040600} & 10.4.6.0 & \code{0A050000} & 328 & 93 & \code{0A050000} (\code{MIS} + \code{m}) & yes &  \\
\code{0x0A050000} & 10.5.0.0 & \code{0A050000} & 328 & 93 & \code{0A050000} (\code{MIS} + \code{m}) &  &  \\
\code{0x0A050001} & 10.5.0.1 & \code{0A050000} & 328 & 93 & \code{0A050000} (\code{MIS} + \code{m}) &  &  \\
\code{0x0A050002} & 10.5.0.2 & \code{0A050000} & 328 & 93 & \code{0A050000} (\code{MIS} + \code{m}) &  &  \\
\code{0x0A050003} & 10.5.0.3 & \code{0A050000} & 328 & 93 & \code{0A050000} (\code{MIS} + \code{m}) &  &  \\
\code{0x0A050004} & 10.5.0.4 & \code{0A050000} & 328 & 93 & \code{0A050000} (\code{MIS} + \code{m}) &  &  \\
\code{0x0A060000} & 10.6.0.0 & \code{0A050000} & 328 & 93 & \code{0A050000} (\code{MIS} + \code{m}) &  &  \\
\code{0x0A060001} & 10.6.0.1 & \code{0A050000} & 328 & 93 & \code{0A050000} (\code{MIS} + \code{m}) &  &  \\
\code{0x0A060002} & 10.6.0.2 & \code{0A050000} & 328 & 93 & \code{0A050000} (\code{MIS} + \code{m}) &  &  \\
\code{0x0A060003} & 10.6.0.3 & \code{0A050000} & 328 & 93 & \code{0A050000} (\code{MIS} + \code{m}) &  &  \\
\code{0x0A060004} & 10.6.0.4 & \code{0A050000} & 328 & 93 & \code{0A050000} (\code{MIS} + \code{m}) &  &  \\
\code{0x0A060011} & 10.6.0.11 & \code{0A050000} & 328 & 93 & \code{0A050000} (\code{MIS} + \code{m}) & yes &  \\
\code{0x0A060100} & 10.6.1.0 & \code{0A050000} & 328 & 93 & \code{0A050000} (\code{MIS} + \code{m}) &  &  \\
\code{0x0A060101} & 10.6.1.1 & \code{0A050000} & 328 & 93 & \code{0A050000} (\code{MIS} + \code{m}) &  &  \\
\code{0x0A060102} & 10.6.1.2 & \code{0A050000} & 328 & 93 & \code{0A050000} (\code{MIS} + \code{m}) &  &  \\
\code{0x0A060103} & 10.6.1.3 & \code{0A050000} & 328 & 93 & \code{0A050000} (\code{MIS} + \code{m}) &  &  \\
\code{0x0A060104} & 10.6.1.4 & \code{0A050000} & 328 & 93 & \code{0A050000} (\code{MIS} + \code{m}) &  &  \\
\code{0x0A060132} & 10.6.1.32 & \code{0A050000} & 328 & 93 & \code{0A050000} (\code{MIS} + \code{m}) & yes &  \\
\code{0x0A060133} & 10.6.1.33 & \code{0A050000} & 328 & 93 & \code{0A050000} (\code{MIS} + \code{m}) &  & Micro-ALTO 10.6.1 \\
\code{0x0A060200} & 10.6.2.0 & \code{0A060200} & 328 & 93 & \code{0A050000} (\code{MIS} + \code{m}) &  &  \\
\code{0x0A060201} & 10.6.2.1 & \code{0A060200} & 328 & 93 & \code{0A050000} (\code{MIS} + \code{m}) &  &  \\
\code{0x0A060202} & 10.6.2.2 & \code{0A060200} & 328 & 93 & \code{0A050000} (\code{MIS} + \code{m}) &  &  \\
\code{0x0A060203} & 10.6.2.3 & \code{0A060200} & 328 & 93 & \code{0A050000} (\code{MIS} + \code{m}) &  &  \\
\code{0x0A060204} & 10.6.2.4 & \code{0A060200} & 328 & 93 & \code{0A050000} (\code{MIS} + \code{m}) &  &  \\
\code{0x0A060231} & 10.6.2.31 & \code{0A060231} & 328 & 95 & \code{0A060231} (\code{MIS} + \code{m}) &  & WIMS \\
\code{0x0A060233} & 10.6.2.33 & \code{0A060200} & 328 & 93 & \code{0A050000} (\code{MIS} + \code{m}) &  &  \\
\code{0x0A060300} & 10.6.3.0 & \code{0A060200} & 328 & 93 & \code{0A050000} (\code{MIS} + \code{m}) &  &  \\
\code{0x0A060301} & 10.6.3.1 & \code{0A060200} & 328 & 93 & \code{0A050000} (\code{MIS} + \code{m}) &  &  \\
\code{0x0A060302} & 10.6.3.2 & \code{0A060200} & 328 & 93 & \code{0A050000} (\code{MIS} + \code{m}) &  &  \\
\code{0x0A060303} & 10.6.3.3 & \code{0A060200} & 328 & 93 & \code{0A050000} (\code{MIS} + \code{m}) &  &  \\
\code{0x0A060304} & 10.6.3.4 & \code{0A060200} & 328 & 93 & \code{0A050000} (\code{MIS} + \code{m}) &  &  \\
\code{0x0A060333} & 10.6.3.33 & \code{0A060200} & 328 & 93 & \code{0A050000} (\code{MIS} + \code{m}) &  &  \\
\code{0x0A060401} & 10.6.4.1 & \code{0A050000} & 328 & 93 & \code{0A050000} (\code{MIS} + \code{m}) & yes &  \\
\code{0x0A070001} & 10.7.0.1 & \code{0A080004} & 327 & 71 or 88 & \code{0A070101} (\code{MIS} + \code{m}) & yes &  \\
\code{0x0A070101} & 10.7.1.1 & \code{0A070101} & 327 & 88 & \code{0A070101} (\code{MIS} + \code{m}) &  &  \\
\code{0x0A070104} & 10.7.1.4 & \code{0A080004} & 327 & 71 or 88 & \code{0A070101} (\code{MIS} + \code{m}) & yes &  \\
\code{0x0A070131} & 10.7.1.31 & \code{0A080031} & 327 & 90 & \code{0A080031} (\code{MIS} + \code{m}) & yes & WIMS \\
\code{0x0A070211} & 10.7.2.11 & \code{0A080004} & 327 & 71 or 88 & \code{0A070101} (\code{MIS} + \code{m}) & yes &  \\
\code{0x0A080001} & 10.8.0.1 & \code{0A080004} & 327 & 71 or 88 & \code{0A070101} (\code{MIS} + \code{m}) & yes &  \\
\code{0x0A080004} & 10.8.0.4 & \code{0A080004} & 327 & 71 or 88 & \code{0A070101} (\code{MIS} + \code{m}) &  & SBE \\
\code{0x0A080031} & 10.8.0.31 & \code{0A080031} & 327 & 90 & \code{0A080031} (\code{MIS} + \code{m}) &  & WIMS \\
\code{0x0A090001} & 10.9.0.1 & \code{0A080004} & 327 & 71 or 88 & \code{0A070101} (\code{MIS} + \code{m}) & yes &  \\
\code{0x0A090101} & 10.9.1.1 & \code{0A080004} & 327 & 71 or 88 & \code{0A070101} (\code{MIS} + \code{m}) & yes &  \\
\code{0x0A090122} & 10.9.1.22 & \code{0A080004} & 327 & 71 or 88 & \code{0A070101} (\code{MIS} + \code{m}) & yes &  \\
\code{0x0A090131} & 10.9.1.31 & \code{0A080031} & 327 & 90 & \code{0A080031} (\code{MIS} + \code{m}) & yes & WIMS \\
\code{0x0A0A0001} & 10.10.0.1 & \code{0A080004} & 327 & 71 or 88 & \code{0A070101} (\code{MIS} + \code{m}) & yes &  \\
\code{0x0A0A0004} & 10.10.0.4 & \code{0A080004} & 327 & 71 or 88 & \code{0A070101} (\code{MIS} + \code{m}) & yes &  \\
\code{0x0A0A0022} & 10.10.0.22 & \code{0A080004} & 327 & 71 or 88 & \code{0A070101} (\code{MIS} + \code{m}) & yes &  \\
\code{0x0A0A0031} & 10.10.0.31 & \code{0A080031} & 327 & 90 & \code{0A080031} (\code{MIS} + \code{m}) &  & WIMS \\
\code{0x0A0C0001} & 10.12.0.1 & \code{0A080004} & 327 & 71 or 88 & \code{0A070101} (\code{MIS} + \code{m}) &  & RBR \\
\code{0x0A0C0004} & 10.12.0.4 & \code{0A080004} & 327 & 71 or 88 & \code{0A070101} (\code{MIS} + \code{m}) & yes &  \\
\code{0x0A0C0032} & 10.12.0.32 & \code{0A080004} & 327 & 71 or 88 & \code{0A070101} (\code{MIS} + \code{m}) &  & ALAMO/MM2 Software version 10.12.0.32 \\
\code{0x0A0C0034} & 10.12.0.34 & \code{0A080004} & 327 & 71 or 88 & \code{0A070101} (\code{MIS} + \code{m}) & yes &  \\
\code{0x0A0C0131} & 10.12.1.31 & \code{0A080031} & 327 & 90 & \code{0A080031} (\code{MIS} + \code{m}) &  & WIMS \\
\code{0x0A0D0013} & 10.13.0.13 & \code{0A080004} & 327 & 71 or 88 & \code{0A070101} (\code{MIS} + \code{m}) & yes &  \\
\code{0x0A0E0028} & 10.14.0.28 & \code{0A080004} & 327 & 71 or 88 & \code{0A070101} (\code{MIS} + \code{m}) & yes &  \\
\end{longtable}
\endgroup

Notes:

\begin{itemize}
\item Two words decode but don't list. 0x08020000 and 0x0A020201 are in the decode switch but missing from the listing switch, so \code{Generate\allowbreak{}Listing\allowbreak{}String} uses \code{Generate\allowbreak{}Listing\allowbreak{}String\allowbreak{}Unsupported}. In the build configuration (\code{XMSG\_\allowbreak{}DEBUG\_\allowbreak{}MFA\_\allowbreak{}MISS\_\allowbreak{}GFG =\allowbreak{} true}) that only logs a warning and prints a generic \code{MIS}/\code{m} block (C.1.11). That block has a fixed parameter list (it includes the \code{PNB}/\code{PNC}/\code{PXC} BIT thresholds), so parameters the layout doesn't contain show their unset values.
\item Some words share a function whose code is identical. \code{0A060200} duplicates \code{0A050000}. \code{06000100} and \code{06010300} duplicate \code{06000000}.
\item A firmware word that isn't in the table falls through to \code{Decode\allowbreak{}XMsg\allowbreak{}Data\allowbreak{}Unsupported}. It reads only the parameters version (offset 8) and CTD type (offset 9), logs “Unsupported ALAMO Mission Parameters message version: 0x…”, and gets the generic listing.
\end{itemize}

\subsubsection{\texorpdfstring{C.1.7 Current layout: 10.7 to 10.14 (\code{Fw0A080004}; the same D0 for \code{Fw0A070101} and \code{Fw0A080031})}{C.1.7 Current layout: 10.7 to 10.14 (Fw0A080004; the same D0 for Fw0A070101 and Fw0A080031)}}\label{sec:c17-current-layout-107-to-1014-fw0a080004-the-same-d0-for-fw0a070101-and-fw0a080031}

The offsets agree with NG10.12 for every parameter. The only difference is the type of \code{DRT}, which NG10.12 gives as I2 and the decoder reads as U2.

\paragraph{\texorpdfstring{D0 (327 bytes; terminator at offset 326)}{D0 (327 bytes; terminator at offset 326)}}\label{sec:d0-327-bytes-terminator-at-offset-326}
\mbox{}\par

Offsets 0-7 are the common header (C.1.4).

\begingroup\footnotesize\setlength{\tabcolsep}{3pt}
\begin{longtable}{@{}>{\raggedleft\arraybackslash}p{0.0625\linewidth}>{\raggedleft\arraybackslash}p{0.0446\linewidth}>{\raggedright\arraybackslash}p{0.0446\linewidth}>{\raggedright\arraybackslash}p{0.1250\linewidth}>{\raggedright\arraybackslash}p{0.2054\linewidth}>{\raggedright\arraybackslash}p{0.4375\linewidth}@{}}
\toprule
\textbf{Offset} & \textbf{Size} & \textbf{Type} & \textbf{Mnemonic} & \textbf{Units / scale} & \textbf{Meaning and notes} \\
\midrule
\endfirsthead
\toprule
\textbf{Offset} & \textbf{Size} & \textbf{Type} & \textbf{Mnemonic} & \textbf{Units / scale} & \textbf{Meaning and notes} \\
\midrule
\endhead
\bottomrule
\endlastfoot
8 & 1 & U1 & (params ver.) & - & Mission-parameters (packet) version, MRV \code{Mission\allowbreak{}Params\allowbreak{}Version}. Printed as \code{Cfg\allowbreak{}Msg\allowbreak{}Ver} on the \code{MIS} line. \\
9 & 1 & U1 & (CTD type) & enum & CTD type: 0 = RBR, 1 = Sea-Bird (SBE). Printed on the \code{MIS} line with its name (\code{RBR}, \code{SBE}, otherwise \code{unknown} \\
10 & 2 & U2 & MSN & - & MSN: dive cycle number. Stored as signed 16-bit (10.5.0 doc: I2; 10.12.0 doc: U2 \\
12 & 1 & U1 & MPS & - & MPS: current mission phase (0-2 \\
13 & 2 & U2 & MPC & cycles & MPC: cycle count within the current phase (0-50,000 \\
15 & 2 & U2 & MNM0 & cycles & MNM0: number of dive cycles in phase 0. \\
17 & 2 & U2 & MNM1 & cycles & MNM1: number of dive cycles in phase 1. \\
19 & 2 & U2 & MNM2 & cycles & MNM2: number of dive cycles in phase 2. \\
21 & 1 & U1 & MNS0 & - & MNS0: phase to run after phase 0. \\
22 & 1 & U1 & MNS1 & - & MNS1: phase to run after phase 1. \\
23 & 1 & U1 & MNS2 & - & MNS2: phase to run after phase 2. \\
24 & 2 & U2 & DTD0 & dBar & DTD0: target drift/park pressure, phase 0. \\
26 & 2 & U2 & DTD1 & dBar & DTD1: target drift/park pressure, phase 1. \\
28 & 2 & U2 & DTD2 & dBar & DTD2: target drift/park pressure, phase 2. \\
30 & 2 & U2 & MDS & dBar & MDS: mission abort (“death”) pressure, the maximum depth. \\
32 & 1 & U1 & APT & enum & APT: send the mission parameters at every surfacing. 0 = off, 1 = on. \\
33 & 1 & U1 & BTM & enum & BTM: Bluetooth mode. 0 = off, 1 = on. \\
34 & 1 & U1 & DDB & enum & DDB: diagnostic dive bypass. 0 = run the diagnostic dive, 1 = skip it. \\
35 & 1 & U1 & XAS & enum & XAS: profile compression. 1 = first difference, 2 = curvature (10.5.0 doc). Spreadsheet v2.4 says 0 = first difference, 1 = curvature. \\
36 & 1 & U1 & CPS & enum & CPS: CTD profiling. 0 = off, 1 = on. \\
37 & 1 & U1 & CST & enum & CST: CTD profile sampling. 1 = binned, 2 = raw (10.5.0 doc). Spreadsheet v2.4 says 0 = bin-averaged, 1 = raw. \\
38 & 2 & U2 & CTF & dBar & CTF: profile cut-off pressure. \\
40 & 2 & U2 & CTL & dBar & CTL: lower pressure limit for resetting the pressure-sensor offset. \\
42 & 2 & U2 & CTU & dBar & CTU: upper pressure limit for resetting the pressure-sensor offset. \\
44 & 1 & U1 & CTC & enum & CTC: conductivity-cell (auxiliary) temperature channel, RBR only. 0 = off, 1 = on. \\
45 & 4 & U4 & CSP & ms & CSP: CTD continuous sampling period. Listed in s. \\
49 & 2 & U2 & CPB1 & dBar & CPB1: lower boundary of profile segment 1 (bottom \\
51 & 2 & U2 & CPB2 & dBar & CPB2: lower boundary of profile segment 2 (middle \\
53 & 2 & U2 & CPB3 & dBar & CPB3: lower boundary of profile segment 3 (top \\
55 & 2 & U2 & CPN1 & RBR 0.1 dBar; SBE dBar & CPN1: segment 1 bin size. Listed as value/10 dBar for RBR and as the raw value for SBE. \\
57 & 2 & U2 & CPN2 & RBR 0.1 dBar; SBE dBar & CPN2: segment 2 bin size. \\
59 & 2 & U2 & CPN3 & RBR 0.1 dBar; SBE dBar & CPN3: segment 3 bin size. \\
61 & 2 & U2 & CPP1 & RBR ms; SBE dBar & CPP1: segment 1 sampling period (RBR) or sampling pressure interval (SBE). Listed as value/1000 s for RBR and the raw dBar value for SBE. \\
63 & 2 & U2 & CPP2 & RBR ms; SBE dBar & CPP2: segment 2 sampling period or interval. \\
65 & 2 & U2 & CPP3 & RBR ms; SBE dBar & CPP3: segment 3 sampling period or interval. \\
67 & 2 & U2 & CRD & dBar & CRD: pressure at which raw profile samples start. \\
69 & 2 & U2 & CRS & samples & CRS: maximum number of raw profile samples. 0 = off. \\
71 & 2 & U2 & MTD & dBar & MTD: start-of-mission trigger pressure. \\
73 & 4 & U4 & MTT & ms & MTT: start-of-mission pressure sample period. Listed in s. \\
77 & 4 & U4 & MTS & s & MTS: start-of-mission timeout. \\
81 & 1 & U1 & MAM & bit mask & MAM: IMU sampling phases. Spreadsheet v2.4: bit 0 surface, bit 1 descent, bit 2 drift/park, bit 3 ascent. \\
82 & 2 & U2 & MAT1 & - & MAT1: number of IMU sampling intervals. \\
84 & 2 & U2 & MAT2 & s & MAT2: length of one IMU sampling interval. \\
86 & 2 & U2 & MAP & s & MAP: IMU sample period, the time between sampling intervals. \\
88 & 2 & U2 & ACA & g & ACA: accelerometer full-scale range: 2, 4, 8 or 16 g. \\
90 & 2 & U2 & ACG & deg/s & ACG: gyroscope full-scale range: 250, 500, 1000 or 2000 deg/s. \\
92 & 2 & U2 & ACR & Hz & ACR: IMU sample rate. \\
94 & 2 & U2 & PBL & counts & PBL: calibration, battery voltage at 12 V. \\
96 & 2 & U2 & PBH & counts & PBH: calibration, battery voltage at 14 V. \\
98 & 2 & U2 & PP0 & counts & PP0: calibration, oil pump current at 0 mA. \\
100 & 2 & U2 & PP1 & counts & PP1: calibration, oil pump current at 1000 mA. \\
102 & 2 & U2 & PVL & counts & PVL: calibration, V12 signal at 12 V. \\
104 & 2 & U2 & PVH & counts & PVH: calibration, V12 signal at 14 V. \\
106 & 2 & U2 & PV0 & counts & PV0: calibration, case vacuum at 0 inHg. \\
108 & 2 & U2 & PV5 & counts & PV5: calibration, case vacuum at 5 inHg. \\
110 & 2 & U2 & PC0 & counts & PC0: calibration, V12 current at 0 mA. \\
112 & 2 & U2 & PC1 & counts & PC1: calibration, V12 current at 1000 mA. \\
114 & 1 & U1 & ICM & enum & ICM: ice detection. 0 = off, 1 = on. \\
115 & 2 & U2 & IMM & bit mask & IMM: ice-detection month mask. \\
117 & 2 & U2 & ICA0 & attempts & ICA0: maximum surfacing attempts, summer. \\
119 & 2 & U2 & ICA1 & attempts & ICA1: maximum surfacing attempts, winter. \\
121 & 2 & U2 & ICB0 & cycles & ICB0: cycles between surfacing attempts, summer. \\
123 & 2 & U2 & ICB1 & cycles & ICB1: cycles between surfacing attempts, winter. \\
125 & 2 & U2 & IMT0 & \% & IMT0: stored-message urgency threshold, summer. \\
127 & 2 & U2 & IMT1 & \% & IMT1: stored-message urgency threshold, winter. \\
129 & 2 & U2 & ILT0 & days & ILT0: time-since-last-transmission urgency threshold, summer. \\
131 & 2 & U2 & ILT1 & days & ILT1: time-since-last-transmission urgency threshold, winter. \\
133 & 2 & U2 & ICN & dBar & ICN: shallow pressure limit of the ice temperature test. \\
135 & 2 & U2 & ICX & dBar & ICX: deep pressure limit of the ice temperature test. \\
137 & 2 & I2 & ICT & 0.01 degC & ICT: ice temperature threshold, signed. Listed as value/100 degC. \\
139 & 2 & U2 & ICS & s & ICS: ice temperature test sample interval (“ice storage timeout” in the MRV docs \\
141 & 1 & U1 & ISS0 & enum & ISS0: surface search (multiple surfacing) mode, summer. \\
142 & 1 & U1 & ISS1 & enum & ISS1: surface search (multiple surfacing) mode, winter. \\
143 & 1 & U1 & ISM & bit mask & ISM: surfacing criteria mask. \\
144 & 2 & U2 & IST & 0.01 PSU & IST: surfacing salinity threshold. Listed as value/100 PSU. The header comment says 1/10 PSU. \\
146 & 2 & U2 & ISO & s & ISO: surface search timeout, the length of one surfacing attempt. \\
148 & 1 & U1 & ISP & enum & ISP: surface search profile mode. \\
149 & 1 & U1 & IUE0 & enum & IUE0: urgency-mode exit criteria, summer. \\
150 & 1 & U1 & IUE1 & enum & IUE1: urgency-mode exit criteria, winter. \\
151 & 2 & U2 & ICD0 & dBar & ICD0: turnaround hold pressure, summer. \\
153 & 2 & U2 & ICD1 & dBar & ICD1: turnaround hold pressure, winter. \\
155 & 4 & U4 & ICR0 & min (listing) & ICR0: turnaround hold duration, summer. \\
159 & 4 & U4 & ICR1 & min (listing) & ICR1: turnaround hold duration, winter. \\
163 & 1 & U1 & ICG0 & enum & ICG0: turnaround seek mode, summer. \\
164 & 1 & U1 & ICG1 & enum & ICG1: turnaround seek mode, winter. \\
165 & 1 & U1 & ISA & attempts & ISA: ice descent seek attempts. \\
166 & 2 & U2 & ISD & dBar & ISD: ice descent seek pressure delta. \\
168 & 2 & U2 & IAT & s & IAT: ice descent seek timeout. \\
170 & 2 & U2 & ISV & s & ISV: ice descent seek velocity timeout. \\
172 & 2 & U2 & ISW & s & ISW: ice descent seek wait timeout. \\
174 & 2 & U2 & ICO0 & s & ICO0: turnaround valve duration, summer. \\
176 & 2 & U2 & ICO1 & s & ICO1: turnaround valve duration, winter. \\
178 & 4 & U4 & ICF & bit mask & ICF: ice-free (summer) month mask. \\
182 & 4 & U4 & ICW & bit mask & ICW: iced-over (winter) month mask. \\
186 & 2 & U2 & DNS & - & DNS: descent-phase cycle number. \\
188 & 2 & U2 & DRA & cm/s & DRA: target descent rate. The decoder stores it in 8 bits, so values above 255 are truncated. Listed raw, labeled \code{[m/\allowbreak{}s]}. \\
190 & 2 & U2 & DSD0 & dBar & DSD0: first-pump pressure during descent, phase 0. \\
192 & 2 & U2 & DSD1 & dBar & DSD1: first-pump pressure, phase 1. \\
194 & 2 & U2 & DSD2 & dBar & DSD2: first-pump pressure, phase 2. \\
196 & 2 & U2 & DST0 & s & DST0: first-pump timeout, phase 0. \\
198 & 2 & U2 & DST1 & s & DST1: first-pump timeout, phase 1. \\
200 & 2 & U2 & DST2 & s & DST2: first-pump timeout, phase 2. \\
202 & 2 & U2 & DDF & s & DDF: diagnostic dive initial pump time. \\
204 & 2 & U2 & DDC & s & DDC: diagnostic dive calibration timeout. \\
206 & 4 & U4 & DSI & ms & DSI: descent sample interval. Listed in s. \\
210 & 2 & U2 & DDD & dBar & DDD: fall time-series storage pressure interval. \\
212 & 2 & U2 & DDT & s & DDT: fall time-series storage time interval. \\
214 & 2 & U2 & DFT0 & s & DFT0: first descent pump duration, phase 0. \\
216 & 2 & U2 & DFT1 & s & DFT1: first descent pump duration, phase 1. \\
218 & 2 & U2 & DFT2 & s & DFT2: first descent pump duration, phase 2. \\
220 & 1 & U1 & DCM & enum & DCM: descent seek/calibration mode. 0 = off, 1 = on. \\
221 & 1 & U1 & DAL & enum & DAL: descent algorithm. 0 = calibration, 1 = seek. \\
222 & 2 & U2 & DCT0 & s & DCT0: wait after the first pump (calibration timeout), phase 0. \\
224 & 2 & U2 & DCT1 & s & DCT1: same, phase 1. \\
226 & 2 & U2 & DCT2 & s & DCT2: same, phase 2. \\
228 & 1 & U1 & DCA & seeks & DCA: maximum descent seeks or calibration attempts. \\
229 & 2 & U2 & DRC & s & DRC: calculated descent calibration rate (deprecated \\
231 & 2 & U2 & DRT & s & DRT: seek rate adjustment timeout. The MRV docs say I2; the decoder reads U2. The listing's DRT line printed DRA's value (fixed in build 3309; bug \\
233 & 2 & U2 & DCX & s & DCX: seek maximum pump time. \\
235 & 2 & U2 & DCP & s & DCP: seek minimum pump time. \\
237 & 2 & U2 & DCF & s & DCF: seek minimum valve time. \\
239 & 2 & U2 & DCD & dBar & DCD: calibration pressure delta (deprecated \\
241 & 2 & U2 & DCR & - & DCR: descent calibration rate. \\
243 & 2 & U2 & DSW & s & DSW: seek interval. \\
245 & 2 & U2 & DUC & \% & DUC: seek underdamping factor. \\
247 & 2 & U2 & DPT & s & DPT: pump time per 100 dBar. \\
249 & 2 & U2 & DUT & s & DUT: valve time per 100 dBar. \\
251 & 2 & U2 & DPD & s & DPD: first seek pump duration. \\
253 & 2 & U2 & DVT & s & DVT: seek pressure sample interval (velocity measurement delay \\
255 & 2 & U2 & DVD & dBar & DVD: seek valve pressure delta. \\
257 & 2 & U2 & DVW & s & DVW: seek valve wait time. \\
259 & 2 & U2 & PFR & mL/min & PFR: pump flow rate. \\
261 & 2 & U2 & PFC & cc/m x 1000 & PFC: pump flow pressure correction (listing unit: uL/dBar \\
263 & 2 & U2 & DDP & dBar & DDP: diagnostic dive ascent trigger pressure. \\
265 & 2 & U2 & DDS & s & DDS: diagnostic dive ascent trigger timeout. \\
267 & 2 & U2 & PNS & - & PNS: park-phase cycle number. \\
269 & 4 & U4 & PTO0 & s & PTO0: park duration, phase 0. \\
273 & 4 & U4 & PTO1 & s & PTO1: park duration, phase 1. \\
277 & 4 & U4 & PTO2 & s & PTO2: park duration, phase 2. \\
281 & 4 & U4 & PTH0 & bit mask & PTH0: hour-of-day mask that ends the park phase, phase 0. \\
285 & 4 & U4 & PTH1 & bit mask & PTH1: hour mask, phase 1. \\
289 & 4 & U4 & PTH2 & bit mask & PTH2: hour mask, phase 2. \\
293 & 4 & U4 & PTD & s & PTD: park end date/time. The spreadsheet says it is set as YYYYMMDD HHMMSS and returned in seconds. \\
297 & 1 & U1 & PSA & enum & PSA: park sample mode. \\
298 & 4 & U4 & PSI & ms & PSI: park sample interval. Listed in s. \\
302 & 1 & U1 & PFM0 & enum & PFM0: fall to profile pressure after the park phase, phase 0. 0 = off, 1 = on. \\
303 & 1 & U1 & PFM1 & enum & PFM1: same, phase 1. \\
304 & 1 & U1 & PFM2 & enum & PFM2: same, phase 2. \\
305 & 2 & U2 & PFT0 & s & PFT0: fall-to-profile timeout, phase 0. \\
307 & 2 & U2 & PFT1 & s & PFT1: same, phase 1. \\
309 & 2 & U2 & PFT2 & s & PFT2: same, phase 2. \\
311 & 2 & U2 & PFD0 & dBar & PFD0: target profile pressure, phase 0. \\
313 & 2 & U2 & PFD1 & dBar & PFD1: same, phase 1. \\
315 & 2 & U2 & PFD2 & dBar & PFD2: same, phase 2. \\
317 & 2 & U2 & PFV0 & s & PFV0: fall valve-open duration, phase 0. \\
319 & 2 & U2 & PFV1 & s & PFV1: same, phase 1. \\
321 & 2 & U2 & PFV2 & s & PFV2: same, phase 2. \\
323 & 2 & U2 & PDT & s & PDT: park data storage interval. \\
325 & 1 & U1 & PSM & enum & PSM: park seek mode. 0 = off, 1 = on. \\
326 & 1 & char & \code{;\allowbreak{}} &  & Terminator. \\
\end{longtable}
\endgroup

\paragraph{\texorpdfstring{D1 (71 bytes without the CTD correction parameters, 88 bytes with them)}{D1 (71 bytes without the CTD correction parameters, 88 bytes with them)}}\label{sec:d1-71-bytes-without-the-ctd-correction-parameters-88-bytes-with-them}
\mbox{}\par

\begingroup\footnotesize\setlength{\tabcolsep}{3pt}
\begin{longtable}{@{}>{\raggedleft\arraybackslash}p{0.0804\linewidth}>{\raggedleft\arraybackslash}p{0.0446\linewidth}>{\raggedright\arraybackslash}p{0.0446\linewidth}>{\raggedright\arraybackslash}p{0.0804\linewidth}>{\raggedright\arraybackslash}p{0.1250\linewidth}>{\raggedright\arraybackslash}p{0.5446\linewidth}@{}}
\toprule
\textbf{Offset} & \textbf{Size} & \textbf{Type} & \textbf{Mnemonic} & \textbf{Units / scale} & \textbf{Meaning and notes} \\
\midrule
\endfirsthead
\toprule
\textbf{Offset} & \textbf{Size} & \textbf{Type} & \textbf{Mnemonic} & \textbf{Units / scale} & \textbf{Meaning and notes} \\
\midrule
\endhead
\bottomrule
\endlastfoot
0 & 1 & U1 & ID &  & 0xD1. \\
1 & 2 & U2 & \code{jj} &  & 71 or 88. \\
3 & 2 & U2 & PSD & dBar & PSD: park seek pressure delta. \\
5 & 2 & U2 & PST & s & PST: park seek interval (timeout \\
7 & 2 & U2 & PSV & s & PSV: park seek velocity measurement delay. \\
9 & 2 & U2 & PSW & s & PSW: park seek wait timeout. \\
11 & 1 & U1 & PSS & enum & PSS: park sampling scheme. 1 = P, 2 = T, 3 = PT, 4 = S, 7 = PTS. \\
12 & 2 & U2 & ASN & - & ASN: ascent cycle number. \\
14 & 2 & U2 & ARR & cm/s & ARR: target rise rate. Listed as value/100 m/s. \\
16 & 2 & U2 & APM & s & APM: maximum ascent pump run time. \\
18 & 2 & U2 & APS & s & APS: initial ascent pump-on duration. \\
20 & 2 & U2 & APN & s & APN: ascent pump-on duration. \\
22 & 2 & U2 & APF & s & APF: ascent pump-off duration. \\
24 & 4 & U4 & ASI & ms (decoder) & ASI: ascent sample interval. The 10.5.0 doc says seconds; the decoder lists value/1000 s. \\
28 & 2 & U2 & ASD & dBar & ASD: surface pressure threshold. \\
30 & 2 & U2 & ADD & dBar & ADD: rise time-series storage pressure interval. \\
32 & 2 & U2 & ADT & s & ADT: rise time-series storage time interval. \\
34 & 2 & U2 & ARC & ms (decoder) & ARC: rise-rate calculation adjustment. Listed as value/1000 s, labeled “sample delay correction”. \\
36 & 2 & U2 & AHR & see notes & AHR: rise rate below which the ascent halts. The 10.5.0 doc says cm/s; the decoder lists value/1000 as m/s. \\
38 & 2 & U2 & AHD & dBar & AHD: halt pressure delta. \\
40 & 2 & U2 & SFN & - & SFN: surface cycle number. \\
42 & 2 & I2 & SDC & dBar & SDC: surface pressure compensation, signed. \\
44 & 2 & U2 & SGT & min & SGT: GPS timeout. \\
46 & 2 & U2 & SET & min & SET: Iridium timeout. \\
48 & 1 & U1 & SBM & enum & SBM: beacon mode. 0 = off, 1 = on. \\
49 & 2 & U2 & SBT & s & SBT: beacon interval. \\
51 & 2 & U2 & SSA & enum & SSA: surface sample averaging mode. \\
53 & 4 & U4 & SSI & ms & SSI: surface sample interval. Listed in s. \\
57 & 2 & U2 & SDS & s & SDS: surface data storage interval. \\
59 & 1 & U1 & SAM & enum & SAM: air-assist (air pump) mode. 0 = off, 1 = on. \\
60 & 2 & U2 & SAX & see notes & SAX: air-assist maximum pressure (psi in the MRV docs). Listed as value/1000 PSI. \\
62 & 2 & U2 & SAN & see notes & SAN: air-assist minimum pressure (psi in the MRV docs). Listed as value/1000 PSI. \\
64 & 2 & U2 & SAO & s & SAO: air-assist minimum (pump) duration. \\
66 & 2 & U2 & SPT & s & SPT: surface CTD storage timeout (sampling duration \\
68 & 2 & U2 & SSP & s & SSP: surface oil pump duration. \\
70 & 4 & R4 & PPC0 & s & PPC0: RBR dynamic-correction tau (conductivity-temperature lag), IEEE-754 single. Only in the 88-byte form. \\
74 & 4 & R4 & PPC1 & - & PPC1: RBR dynamic-correction alpha (short-term thermal mass), single. Only in the 88-byte form. \\
78 & 4 & R4 & PPC2 & s & PPC2: RBR dynamic-correction tdelay (water-cell lag), single. Only in the 88-byte form. \\
82 & 4 & R4 & PPC3 & - & PPC3: RBR dynamic-correction ctcoeff (long-term thermal mass), single. Only in the 88-byte form. \\
86 & 1 & U1 & PPE & enum & PPE: RBR post-processing (dynamic correction) enable. Floats with PPE = 2 send corrected salinity (0xDC). Only in the 88-byte form. \\
70 or 87 & 1 & char & \code{;\allowbreak{}} &  & Terminator. \\
\end{longtable}
\endgroup

{}[decoder] The 88-byte form carries \code{PPC0}-\code{PPC3} and \code{PPE} (NG10.12; production 10.12.0.1 and 10.12.0.32 floats). The decoder accepts both lengths for the \code{Fw0A080004} firmware group and decodes \code{PPC}/\code{PPE} only when the blob is 88 bytes long.

\paragraph{\texorpdfstring{Variants of the current D1}{Variants of the current D1}}\label{sec:variants-of-the-current-d1}
\mbox{}\par

\begingroup\footnotesize\setlength{\tabcolsep}{3pt}
\begin{longtable}{@{}>{\raggedright\arraybackslash}p{0.4375\linewidth}>{\raggedleft\arraybackslash}p{0.0804\linewidth}>{\raggedright\arraybackslash}p{0.4375\linewidth}@{}}
\toprule
\textbf{Layout (firmware)} & \textbf{D1 bytes} & \textbf{Difference from the table above} \\
\midrule
\endfirsthead
\toprule
\textbf{Layout (firmware)} & \textbf{D1 bytes} & \textbf{Difference from the table above} \\
\midrule
\endhead
\bottomrule
\endlastfoot
\code{Fw0A070101} (10.7.1.1) & 88 & Only 88 bytes is accepted. [decoder] \code{PPC0}-\code{PPE} at offsets 70-86 are not decoded (the lines are commented out), so they keep their unset values. \\
\code{Fw0A080031} (WIMS: 10.7.1.31, 10.8.0.31, 10.9.1.31, 10.10.0.31, 10.12.1.31) & 90 & \code{WMT} (U2, s, WIMS timeout) at offset 70. \code{PPC0} at 72, \code{PPC1} at 76, \code{PPC2} at 80, \code{PPC3} at 84 (R4), \code{PPE} at 88 (U1), terminator at 89. \\
\end{longtable}
\endgroup

\subsubsection{\texorpdfstring{C.1.8 The 10.5/10.6 layout (\code{Fw0A050000}, \code{Fw0A060200}, \code{Fw0A060231})}{C.1.8 The 10.5/10.6 layout (Fw0A050000, Fw0A060200, Fw0A060231)}}\label{sec:c18-the-105106-layout-fw0a050000-fw0a060200-fw0a060231}

\begin{itemize}
\item D0 is 328 bytes and identical to the 5.0.61 to 10.3 D0 (Appendix C.A, \code{Fw05006100}/\code{Fw06000000}/\code{Fw08010000}):

\begin{itemize}
\item the header as in C.1.4;
\item parameters \code{MSN} at 10 to \code{ACR} at 92;
\item \code{PNB} 94, \code{PNC} 96, \code{PXC} 98 (U2, 0.01 V and 0.01 psi BIT thresholds);
\item calibration \code{PBL} to \code{PC1} at 100-119;
\item ice block \code{ICM} at 120 to \code{ICW} at 188;
\item descent block \code{DNS} at 192 to \code{DDS} at 271;
\item park block \code{PNS} at 273 to \code{PFV1} at 325;
\item terminator at 327.
\end{itemize}
\item D1 is 93 bytes: \code{PFV2} at 3, \code{PDT} 5, \code{PSM} 7 (U1), \code{PSD} 8, \code{PST} 10, \code{PSV} 12, \code{PSW} 14, \code{PSS} 16 (U1), \code{ASN} 17, \code{ARR} 19, \code{APM} 21, \code{APS} 23, \code{APN} 25, \code{APF} 27, \code{ASI} 29 (U4), \code{ASD} 33, \code{ADD} 35, \code{ADT} 37, \code{ARC} 39, \code{AHR} 41, \code{AHD} 43, \code{SFN} 45, \code{SDC} 47 (I2), \code{SGT} 49, \code{SET} 51, \code{SBM} 53 (U1), \code{SBT} 54, \code{SSA} 56, \code{SSI} 58 (U4), \code{SDS} 62, \code{SAM} 64 (U1), \code{SAX} 65, \code{SAN} 67, \code{SAO} 69, \code{SPT} 71, \code{SSP} 73, \code{PPC0} 75, \code{PPC1} 79, \code{PPC2} 83, \code{PPC3} 87 (R4), \code{PPE} 91 (U1), terminator at 92. Unmarked fields are U2. These offsets agree with NG10.5 for every parameter, except \code{MSN} (NG10.5: I2) and \code{DRT} (NG10.5: I2), which the decoder reads unsigned.
\item \code{Fw0A060231} (10.6.2.31, WIMS): D1 is 95 bytes, with \code{WMT} (U2) at offset 92 after \code{PPE} and the terminator at 94.
\end{itemize}

\subsubsection{\texorpdfstring{C.1.9 Layout history}{C.1.9 Layout history}}\label{sec:c19-layout-history}

Each row is the decoder's view of what changed against the layout before it. Exact offsets for every layout are in Appendix C.A.

\begingroup\footnotesize\setlength{\tabcolsep}{3pt}
\begin{longtable}{@{}>{\raggedright\arraybackslash}p{0.4018\linewidth}>{\raggedright\arraybackslash}p{0.1339\linewidth}>{\raggedright\arraybackslash}p{0.4196\linewidth}@{}}
\toprule
\textbf{Layout (firmware)} & \textbf{D0 / D1 bytes} & \textbf{Changes} \\
\midrule
\endfirsthead
\toprule
\textbf{Layout (firmware)} & \textbf{D0 / D1 bytes} & \textbf{Changes} \\
\midrule
\endhead
\bottomrule
\endlastfoot
\code{02006200} (2.0.62-2.0.68) & 323 / - & First binary layout. No parameters version or CTD-type byte: \code{MSN} is at offset 8. Single D0. Has U4 state start times (seconds since 2000-01-01): descent at 99, park at 179, ascent at 227, surface at 271. Reserved gaps at 81-98, 164-178, 216-226, 257-270 and 309-321 (XF2.4 marks them RESERVED). Single \code{DTD}, \code{DSD}, \code{DST}, \code{DFT}, \code{DCT}, \code{PTO}, \code{PTH}, \code{PFM}, \code{PFT}, \code{PFD} and \code{PFV}, all without a phase index. Has \code{MAT} (one U2), \code{PNB}/\code{PNC}/\code{PXC}, \code{SDT}, \code{SAP} and \code{SPT}. \\
\code{02006900} (2.0.69, 2.3.62-2.5.62) & 322 / - & Adds \code{CTF}, \code{CTL} and \code{CTU} after \code{CST}. Adds the ice block \code{ICM} to \code{ICW} (without \code{ISA}-\code{ISW} and \code{ISM}) in place of the first reserved gap. The other reserved gaps are removed and the start times kept. [decoder] \code{ICF}/\code{ICW} are read as U1 at 141 and 145, though they are 4-byte masks (XF2.4: U4), so only their most significant byte is decoded. \\
\code{03006500} (3.0.65) & 328 / - & A 9-byte slot at 137-145 holds \code{ISA}, \code{ISD}, \code{IAT}, \code{ISV} and \code{ISW} (XF2.4), not decoded. \code{SPT} is no longer decoded. XF2.4 gives this version 331 bytes with \code{SPT} at 326 and \code{SSP} at 328; the decoder expects 328. \code{ICF}/\code{ICW} are still read as U1. \\
\code{03016700} (3.1.67, SOLO-II hybrid 10083) & 328 / - & One extra byte after \code{ISS1} (offset 116), probably \code{ISM} [inferred], not decoded. \\
\code{03026100}, \code{04016100} (3.2.61, 4.1.61) & 328 / - & Adds \code{BTM} (Bluetooth mode) after \code{APT}. \code{SAP} is no longer decoded. \\
\code{04026200} (4.2.62, PAR floats) & 323 / - & Adds the parameters version (offset 8) and CTD type (9), so parameters start at 10. Drops the state start times. Decodes \code{ISA}-\code{ISW}. \code{ICF}/\code{ICW} become U4. Adds \code{DVD} and \code{DVW} after \code{DVT}. Decodes \code{SAP}, \code{SPT} and \code{SSP} at the end. XF2.4 differs (C.5 \\
\code{04036100} (4.3.61) & 327 / - & A 4-byte slot after \code{CPP3} (49-52, probably \code{CRD}/\code{CRS} [inferred]), not decoded. \\
\code{04046100} (4.4.61, 2017 hurricane floats) & 328 / - & A 1-byte slot after \code{CTU} (offset 27, \code{CTC} per XF2.4), not decoded. \\
\code{04056300} (4.5.63) & 328 / ≥ 12 & Decodes \code{CTC}, \code{CRD} and \code{CRS}. The 6 bytes after \code{MAM} (61-66: \code{MAT1}, \code{MAT2} and \code{MAP} per XF2.4) are not decoded. The slots for \code{ISM} (128) and \code{PTD} (246-249) are not decoded either. D0 ends with \code{SAN}. First D1: \code{SAO} 3, \code{SAP} 5, \code{SPT} 7, \code{SSP} 9 (U2 \\
\code{04066100} (4.6.61) & 328 / - & Same D0 offsets as 4.5.63, but \code{CTC}, \code{CRD} and \code{CRS} are not decoded. D1 isn't decoded. XF2.4 gives a single 336-byte D0 that goes on to \code{SSP}. \\
\code{05006100} (5.0.61, ALTO) & 328 / 79 & Restructured for mission phases. Adds \code{MPS}, \code{MPC}, \code{MNM0-2}, \code{MNS0-2}, \code{DDB}, \code{MTS}, \code{MAT1}, \code{MAT2}, \code{MAP} and \code{ISM}. Adds phase-indexed \code{DTD0-2}, \code{DSD0-2}, \code{DST0-2}, \code{DFT0-2}, \code{DCT0-2}, \code{PTO0-2}, \code{PTH0-2}, \code{PFM0-2}, \code{PFT0-2}, \code{PFD0-2} and \code{PFV0-2}, plus \code{PTD}. D0 ends with \code{PFV1}. D1: \code{PFV2} 3, \code{PDT} 5, \code{PSM} 7, \code{PSD} to \code{PSW} 8-15, ascent block from \code{ASN} 16, surface block from \code{SFN} 42 (with \code{SDT} U4 at 53), \code{SAO} 70, \code{SAP} 72, \code{SPT} 74, \code{SSP} 76. \\
\code{05006101} (5.0.61.1, ALTO 10109) & 328 / 74 & No \code{DDB}, \code{MAT1} or \code{MAP} (\code{MAT2} at 81), so later fields move back. \code{PFV2}, \code{PDT} and \code{PSM} are in D0 at 322-326. D1 starts with \code{PSD}. \\
\code{06000000}, \code{06000100}, \code{06010300} (6.0.0, 6.0.1, 6.1.3) & 328 / 81 & D0 as 5.0.61. D1 adds \code{APS} (initial ascent pump-on duration) after \code{APM}. \\
\code{06040900} (6.4.9, Arctic PAR floats) & 328 / 81 & {}[decoder] D1 is read with \code{SSA} as U1 at 59 and \code{SSI} as U2 at 60, \code{SDS} isn't decoded, and byte 62 is skipped. This is probably a decoder quirk rather than a real layout [inferred]. \\
\code{06050100} (6.5.1) & 328 / 77 & \code{SDT} removed from D1. \\
\code{06057700} (6.5.77, “6.9.w”, hurricane floats with WIMS) & 328 / 79 & D0 ends with \code{PFV0}. D1 starts with \code{PFV1} at 3. \\
\code{06080500} (6.8.5, ALTO 11011) & 328 / 79 & An unidentified 2-byte field after \code{ACR} (94-95; the decoder's comments name \code{tsp0}/\code{tsp1}/\code{tsn0}/\code{tsn1}) moves \code{PNB} and everything after it by +2. D0 ends with \code{PFV0}. D1 as 6.5.77. \\
\code{07020200} (7.2.x, 7.3.x) & 328 / 76 & D0 as 5.0.61. D1 adds \code{PSS} (park sampling scheme) after \code{PSW} and drops \code{SAP}. DF7.3.5 has the same D0/D1 parameter order, but its field lists leave out the 5-byte D0 header (offsets 3-7 \\
\code{08000000} (8.0.x, 8.1.0.1, 10.2.2.1) & 328 / 78 & D1 adds \code{WMT} (WIMS timeout) at 75. \\
\code{08010000} (8.1.0, 8.2.0, 9.0.0, 10.0.x-10.3.x) & 328 / 76 & D1 as 7.2 (no \code{WMT} \\
\code{0A050000}, \code{0A060200} (10.2.3.1, 10.4.x, 10.5.x, 10.6.0-10.6.3, 10.6.4.1) & 328 / 93 & D1 adds RBR dynamic-correction \code{PPC0}-\code{PPC3} (R4) and \code{PPE} after \code{SSP} (C.1.8 \\
\code{0A060231} (10.6.2.31, WIMS) & 328 / 95 & Adds \code{WMT} after \code{PPE}. \\
\code{0A070101}, \code{0A080004} (10.7 and later) & 327 / 71 or 88 & Removes \code{PNB}, \code{PNC} and \code{PXC}. \code{PFV2}, \code{PDT} and \code{PSM} move into D0, and D1 starts with \code{PSD} (C.1.7 \\
\code{0A080031} (10.7.1.31 and later WIMS) & 327 / 90 & \code{WMT} sits between \code{SSP} and \code{PPC0}. \\
\end{longtable}
\endgroup

\subsubsection{\texorpdfstring{C.1.10 Decoding notes}{C.1.10 Decoding notes}}\label{sec:c110-decoding-notes}

\begin{itemize}
\item Length checks [decoder]: every layout checks the exact blob length (the size of the received blob, not \code{jj}). With \code{XMSG\_\allowbreak{}DEBUG\_\allowbreak{}MFA\_\allowbreak{}MISS\_\allowbreak{}GFG =\allowbreak{} true}, which is how \code{XDefinitions.\allowbreak{}h} is set, a mismatch only logs a warning and decoding continues. Since build 3309 a blob shorter than its layout is rejected, because decoding would read past its end. A blob shorter than its layout is then read past its end. Only the 4.5.63 D1 check (at least 12 bytes) and the 8-byte minimum for D0 always throw.
\item The decoder never checks the \code{jj} field or the terminator of a mission blob (\code{Check\allowbreak{}XMsg\allowbreak{}Data} isn't called).
\item {}[decoder] \code{DRA} is held in 8 bits, so values above 255 are truncated. \code{DRT} and \code{MSN} are read unsigned and \code{MSN} is stored in an int16. \code{ICF}/\code{ICW} are read as U1 in layouts \code{02006900}-\code{04016100}.
\item Scaling [decoder]. Only the listing scales values; stored values are raw.

\begingroup\footnotesize\setlength{\tabcolsep}{3pt}
\begin{longtable}{@{}>{\raggedright\arraybackslash}p{0.4286\linewidth}>{\raggedright\arraybackslash}p{0.1518\linewidth}@{}}
\toprule
\textbf{Parameters} & \textbf{Listed as} \\
\midrule
\endfirsthead
\toprule
\textbf{Parameters} & \textbf{Listed as} \\
\midrule
\endhead
\bottomrule
\endlastfoot
\code{CSP}, \code{MTT}, \code{DSI}, \code{PSI}, \code{ASI}, \code{SSI}, \code{CPP1-3} (RBR), \code{ARC} & value / 1000 s \\
\code{CPN1-3} (RBR) & value / 10 dBar \\
\code{ARR} & value / 100 m/s \\
\code{AHR} & value / 1000 m/s \\
\code{PNB}, \code{PNC}, \code{PXC} & value / 100 \\
\code{ICT} & value / 100 degC \\
\code{IST} & value / 100 PSU \\
\code{SAN}, \code{SAX} & value / 1000 PSI \\
\end{longtable}
\endgroup

Where NG10.5 gives other units, the table rows in C.1.7 say so.
\item PPC values are IEEE floats, printed with 3 decimals. PPE is printed as an integer.
\end{itemize}

\subsubsection{\texorpdfstring{C.1.11 Listing output}{C.1.11 Listing output}}\label{sec:c111-listing-output}

The “Mission Configuration” section is printed once per dive, from the dive's D0 and D1. Three styles exist.

\begin{enumerate}
\item \textbf{\code{mis} groups}, for layouts \code{02006200} and \code{02006900} (firmware 2.0.62 to 4.1.61). A header line \code{mis SN dive idx 0x\allowbreak{}FWWORD} is followed by groups of fixed-width columns, one line per group with a minor index \code{midx}, each group with its own column-name and unit comment lines. Example group: \code{mis SN dive idx midx MSN DTD MDS APT XAS}.
\item \textbf{\code{MIS} header and \code{mis} groups}, for 4.2.62, 4.3.61, 4.4.61 and 4.6.61 (listing \code{04026200}). The header is \code{MIS SN dive idx 0x\allowbreak{}VV 0x\allowbreak{}FWWORD sensor-type \# name}, followed by \code{mis} groups as in style 1.
\item \textbf{\code{MIS} header and \code{m} lines}, for 4.5.63 and every firmware from 5.0.61 on. Format:

\begin{Verbatim}
#      SN dive        CfgMsgVer       FWVer sensor-type        sensor-type-name
MIS  9213    1             0x18  0x0A0C0032           0      # RBR
# --- General --------------------------------------------------------------------------------------
#      SN dive idx param      value      units          desc
m    9213    1   0 MSN            1      [-]          # dive cycle number
m    9213    1   3 MDS         1000      [dBar]       # maximum depth
m    9213    1  42 CPP1           1.000  [sec]        # interval 1 sample period
m    9213    1 133 ICT           -1.40   [degC]       # threshold temperature
\end{Verbatim}

\begin{itemize}
\item \code{MIS} is \code{MIS \%5d \%4d 0x\%02X 0x\%08X \%11d \# name}: serial number, dive, parameters version (offset 8), firmware word, CTD type and its name (\code{RBR}, \code{SBE} or \code{unknown}).
\item Each \code{m} line is \code{m SN dive idx NAME value [units] \# description}. \code{idx} counts the \code{m} lines from 0, in listing order, not in message order. Integers are printed with \code{\%10d}. Scaled values use 1-4 decimals.
\item Sections, in order: General; BIT (layouts before 10.7 only); Calibration Coefficients; Start-of-Mission Trigger; Diagnostic Dive; Mission Phase; Profile; Profile - Dynamic Salinity Correction (only in the \code{0A050000}, \code{0A060231} and \code{0A080031} listings); Profile - Binning Configuration (CPN/CPP depend on the CTD type: RBR scaled, SBE raw dBar, unknown raw); Profile - Raw Sampling; Profile - Surface Offset Correction; Pump Characteristics; Descent (General, Fall Time-Series, Stages 1-3); Drift/Park (General, Seek, Start of Profile Trigger, Fall to Profile Pressure); Ascent/Profile (General, Rise Time-Series, Ice-Detection, Ice-Detection Summer and Winter); Surface (General, Air Pump, Sampling, Sat./Comm., Abort/Beacon Mode); Surface - Wave-Field Measurement (\code{WST}, only in the WIMS listings \code{08000000}, \code{0A060231} and \code{0A080031}; the \code{06057700} listing prints \code{WMT} instead, which that layout never decodes, so it always shows 65535).
\item A 10.12.0.x listing has 177 \code{m} lines.
\item The Python tools keep each \code{m} line (section, index, name, value, units) in the dive JSON as \code{mission.\allowbreak{}entries}, and keep \code{cfgver}, \code{fwver} and \code{ctdtype} from the \code{MIS} line.
\end{itemize}
\end{enumerate}

\begin{itemize}
\item If no D0 has arrived, the firmware word is still 0xFFFFFFFF and the section is the single line \code{\# Incomplete mission parameters.\allowbreak{}}
\item Listing bugs [decoder]:

\begin{itemize}
\item The \code{DRT} line prints \code{DRA}'s value. \emph{(Fixed in build 3309.)}
\item In the dynamic-correction section, the \code{PPC0} line (labeled tau) prints alpha (\code{PPC1}), and the \code{PPC1} line (labeled alpha) prints tau. The README notes this as unfixed. \emph{(Fixed in build 3309.)}
\item The \code{0A070101} listing, which covers 10.7 to 10.14 RBR/SBE including 10.12.0.x, has the dynamic-correction section commented out. \code{PPC}/\code{PPE} are decoded but not listed for those floats.
\item \code{DRA} is labeled \code{[m/\allowbreak{}s]} but printed raw (cm/s per NG10.5).
\end{itemize}
\end{itemize}

\subsection{\texorpdfstring{C.2 Mission command echo (0xDE)}{C.2 Mission command echo (0xDE)}}\label{sec:c2-mission-command-echo-0xde}

\subsubsection{\texorpdfstring{Purpose}{Purpose}}\label{sec:purpose}

After it receives a ground-station command upload, the float sends back the command text it received, so that the ground station can confirm the upload (SOLO2: “ASCII listing of mission command received by float”).

\subsubsection{\texorpdfstring{Format}{Format}}\label{sec:format}

\begingroup\footnotesize\setlength{\tabcolsep}{3pt}
\begin{longtable}{@{}>{\raggedleft\arraybackslash}p{0.0625\linewidth}>{\raggedleft\arraybackslash}p{0.0625\linewidth}>{\raggedright\arraybackslash}p{0.0536\linewidth}>{\raggedright\arraybackslash}p{0.1161\linewidth}>{\raggedright\arraybackslash}p{0.6339\linewidth}@{}}
\toprule
\textbf{Offset} & \textbf{Size} & \textbf{Type} & \textbf{Field} & \textbf{Notes} \\
\midrule
\endfirsthead
\toprule
\textbf{Offset} & \textbf{Size} & \textbf{Type} & \textbf{Field} & \textbf{Notes} \\
\midrule
\endhead
\bottomrule
\endlastfoot
0 & 1 & U1 & ID & 0xDE. \\
1 & 2 & U2 & \code{jj} & Length including ID, \code{jj} and terminator. Variable. DF7.3.5 gives a text field of up to 336 characters. \\
3 & \code{jj} - 4 & ASCII & Command echo & The command string as received. \\
\code{jj} - 1 & 1 & char & \code{;\allowbreak{}} & Terminator. \\
\end{longtable}
\endgroup

XF2.4 defines only the ID, size and terminator.

Production [inferred, from all 2026 echoes on argus]. The echo is the uploaded command text exactly as sent: the commands separated by \code{;\allowbreak{}}, ending with \code{MIC 1} (which the upload itself ends with). Besides the text, the blob carries 8-9 bytes of undocumented binary data. Its position varies:

\begin{itemize}
\item \textbf{After the text:} seen in most echoes. The bytes are NULs, but the last one is sometimes not (0x01, 0x08 and 0x80 have been seen).

\begin{Verbatim}
DE 00 37 | "dft 0 100;dft 1 100;dft 2 100;duc 100;MIC 1" | 00 00 00 00 00 00 00 00 | 3B
\end{Verbatim}
\end{itemize}

\begin{itemize}
\item \textbf{Before the text:} all NULs. Seen in the dive 0 echoes of the Lowell ALAMOs (10.12.1) and of 11190 (10.12.0.1).

\begin{Verbatim}
DE 00 55 | 00 00 00 00 00 00 00 00 00 | "psi 1800000;pdt 1800;...;MIC 1" | 3B
\end{Verbatim}
\end{itemize}

The commands use the mission-parameter mnemonics, with a phase index where the parameter has one (\code{dft 0 100} sets \code{DFT0} to 100). They are echoed in whatever case they were uploaded in. A long upload may come back as two consecutive echo blobs in the same session. When several uploads are queued, each one is echoed separately. Because \code{;\allowbreak{}} also appears inside the text, the end of the blob must be found from \code{jj}, not by looking for the first \code{;\allowbreak{}}.

\textbf{Verification against uploads (2026-10-25).} The uploads sent with \code{upload-alamo-cmd} since 2026-01-01 (\code{/\allowbreak{}data1/\allowbreak{}alamo/\allowbreak{}data/\allowbreak{}sbd/\allowbreak{}outgoing}) were compared with the echoes in the RAW files:

\begin{itemize}
\item 261 uploads were echoed back exactly, some in two pieces.
\item 23 were not echoed, all to floats that have not called in since.
\item A few were not echoed although the float called in. One, to 9213 on 2026-09-28, verifiably did not take effect: its mission parameters still show the old values.
\item A handful of echoes match no recorded upload (\code{ppe}, \code{wst} commands). They were presumably sent another way.
\end{itemize}

\subsubsection{\texorpdfstring{Decoding notes [decoder]}{Decoding notes [decoder]}}\label{sec:decoding-notes-decoder}

\begin{itemize}
\item \code{XMission\allowbreak{}Cmd\allowbreak{}Echo} calls \code{Check\allowbreak{}XMsg\allowbreak{}Data(\allowbreak{}data,\allowbreak{} 0x\allowbreak{}DE,\allowbreak{} -1)}. That checks the ID, that \code{jj} (lower 12 bits) is at least 4 and no more than the data received, and that the terminator is at \code{jj - 1}.
\item Since build 3310 the text is decoded. Leading NULs are skipped, the text ends at the next NUL, and it is split at \code{;\allowbreak{}} into commands. Blank entries are dropped, surrounding spaces are trimmed, and other non-printable bytes are kept as \code{\textbackslash{}\allowbreak{}x\allowbreak{}NN}. The trailing \code{MIC 1} is listed as a command like the others. (Before build 3310 the echo was checked but not decoded, and the listing section was printed empty.)
\item Listing: under the section banner “Ground Station Upload Command Echo-Back Messages” (printed whenever a dive has at least one 0xDE blob), one line per command. The parameter name (the text before the first space) is listed in uppercase, as the mission parameter listings name them; the rest of the command is listed as received:
\end{itemize}

\begin{Verbatim}
#      SN dive eidx cidx  command
CMD  9213    3    0    0  DFT 0 100
CMD  9213    3    0    1  DFT 1 100
CMD  9213    3    0    2  DFT 2 100
CMD  9213    3    0    3  DUC 100
CMD  9213    3    0    4  MIC 1
\end{Verbatim}

\code{eidx} numbers the dive's echo-back messages, and \code{cidx} the commands within one.

\subsection{\texorpdfstring{C.3 Ocean surface waves (0xC1, 0xC3)}{C.3 Ocean surface waves (0xC1, 0xC3)}}\label{sec:c3-ocean-surface-waves-0xc1-0xc3}

\subsubsection{\texorpdfstring{C.3.1 Purpose}{C.3.1 Purpose}}\label{sec:c31-purpose}

Wave-measuring floats (ALAMO wave floats such as 9270, 9287-9306 and 9453, and WIMS floats) sample a 3-axis IMU at the surface. The float reduces the data on board to one-sided acceleration auto- and cross-spectra and sends one spectral component per blob. Shore-side processing gives wave height and direction. The format of the raw IMU time series was defined in XF2.4 but is not decoded. The decoder treats 0xC1 and 0xC3 the same way: the blob's version and type fields decide how it is decoded, not the ID.

In production [inferred], 0xC1 blobs carry the auto-spectra (\code{ZZ}, \code{YY} and \code{XX} co-spectra) and 0xC3 blobs carry the cross-spectra (\code{XY} co-spectrum, \code{ZX} and \code{ZY} quadrature spectra). One dive sends six blobs, one per component, often one per SBD packet. Floats whose segment is too long for one blob split a component over several blobs using the first-bin index.

\subsubsection{\texorpdfstring{C.3.2 Common header and version}{C.3.2 Common header and version}}\label{sec:c32-common-header-and-version}

\begingroup\footnotesize\setlength{\tabcolsep}{3pt}
\begin{longtable}{@{}>{\raggedleft\arraybackslash}p{0.0625\linewidth}>{\raggedleft\arraybackslash}p{0.0446\linewidth}>{\raggedright\arraybackslash}p{0.0446\linewidth}>{\raggedright\arraybackslash}p{0.1429\linewidth}>{\raggedright\arraybackslash}p{0.6339\linewidth}@{}}
\toprule
\textbf{Offset} & \textbf{Size} & \textbf{Type} & \textbf{Field} & \textbf{Notes} \\
\midrule
\endfirsthead
\toprule
\textbf{Offset} & \textbf{Size} & \textbf{Type} & \textbf{Field} & \textbf{Notes} \\
\midrule
\endhead
\bottomrule
\endlastfoot
0 & 1 & U1 & ID & 0xC1 or 0xC3. [decoder] The ID is not checked. \\
1 & 2 & U2 & \code{jj} & {}[decoder] Validated with \code{Check\allowbreak{}XMsg\allowbreak{}Data} (minimum 17 bytes). Only the lower 12 bits are used. \\
3 & 1 & U1 & Message version & 0 and 1 = spectral v1 layout; 2 = spectral v2 layout (with a mode byte). Any other value is ignored. The decoder has read byte 3 as the version since release 0.6.3149 (2023-08 \\
\end{longtable}
\endgroup

{}[inferred] DF7.3.5 lists the fields of blob 193 (0xC1) as \code{packet\_\allowbreak{}index}, \code{mode}, \code{spectral\_\allowbreak{}axis}, \code{spectral\_\allowbreak{}type}, \code{segment\_\allowbreak{}length\_\allowbreak{}power}, \code{base\_\allowbreak{}sample\_\allowbreak{}rate}, \code{sample\_\allowbreak{}rate\_\allowbreak{}divider}, \code{full\_\allowbreak{}scale\_\allowbreak{}range\_\allowbreak{}accelerometer}, \code{sample\_\allowbreak{}index} and \code{sample\_\allowbreak{}buffer}. These match the v1 layout below, from byte 4 on, if the version byte at offset 3 is assumed. DF7.3.5 leaves out the D0 header in the same way (C.1.9).

\subsubsection{\texorpdfstring{C.3.3 Spectral format, version 1 (and 0)}{C.3.3 Spectral format, version 1 (and 0)}}\label{sec:c33-spectral-format-version-1-and-0}

Used by the 8.x firmware (for example float 9270, 8.0.1.1, 273-byte blobs).

\begingroup\footnotesize\setlength{\tabcolsep}{3pt}
\begin{longtable}{@{}>{\raggedleft\arraybackslash}p{0.0714\linewidth}>{\raggedleft\arraybackslash}p{0.0446\linewidth}>{\raggedright\arraybackslash}p{0.0714\linewidth}>{\raggedright\arraybackslash}p{0.2589\linewidth}>{\raggedright\arraybackslash}p{0.2054\linewidth}>{\raggedright\arraybackslash}p{0.2679\linewidth}@{}}
\toprule
\textbf{Offset} & \textbf{Size} & \textbf{Type} & \textbf{Field} & \textbf{Scale / units} & \textbf{Notes} \\
\midrule
\endfirsthead
\toprule
\textbf{Offset} & \textbf{Size} & \textbf{Type} & \textbf{Field} & \textbf{Scale / units} & \textbf{Notes} \\
\midrule
\endhead
\bottomrule
\endlastfoot
0 & 1 & U1 & ID &  & 0xC1 or 0xC3. \\
1 & 2 & U2 & \code{jj} & bytes & 16 + 2·n + 1 for n bins. \\
3 & 1 & U1 & Version &  & 0 or 1. \\
4 & 1 & U1 & Packet index &  & {}[inferred] DF7.3.5 \code{packet\_\allowbreak{}index}. 0 in every production blob seen. Not decoded. \\
5 & 1 & U1 & Mode & enum & 1 = spectral (0 = raw). [decoder] Not read for v0/v1: the mode is taken to be spectral. \\
6 & 1 & U1 & Spectral axis & enum, C.3.5 &  \\
7 & 1 & U1 & Spectral type & enum, C.3.5 &  \\
8 & 1 & U1 & Segment length power p &  & Segment (FFT) length N = 2\textasciicircum{}p. The blob carries one-sided FFT bins 0 to N/2 - 1: bin 0 is the DC bin, and there is no Nyquist bin (WIMS firmware v1.x and v2.x). Production: 8 (N = 256) or 9 (N = 512 \\
9 & 2 & U2 & Base sample rate f\_\allowbreak{}base & Hz & 40 (v1 floats) or 45 (v2 floats \\
11 & 2 & U2 & Sample rate divider d &  & Effective sample rate f\_\allowbreak{}s = f\_\allowbreak{}base / (d + 1). For example 40/(9+1) = 4 Hz, 45/(17+1) = 2.5 Hz, 45/(8+1) = 5 Hz. \\
13 & 1 & U1 & Accelerometer full-scale range & g & 2, 4, 8 or 16. Selects ASSF (C.3.6 \\
14 & 2 & U2 & First bin index k0 &  & Index of the first bin in this blob (DF7.3.5 \code{sample\_\allowbreak{}index} \\
16 + 2i & 2 & U2 / I2 & Spectral value, bin k0 + i & counts, ÷ ESSF → g²/Hz & Co-spectra \code{ZZ}, \code{XX} and \code{YY} are read unsigned. \code{XY} (co) and \code{ZX}/\code{ZY} (quadrature) are read signed. i = 0 … n - 1, where n = (\code{jj} - 17)/2 and n ≤ N/2. \\
\code{jj} - 1 & 1 & char & \code{;\allowbreak{}} &  & Terminator. \\
\end{longtable}
\endgroup

\subsubsection{\texorpdfstring{C.3.4 Spectral format, version 2}{C.3.4 Spectral format, version 2}}\label{sec:c34-spectral-format-version-2}

Used by the 10.x wave and WIMS firmware (for example 9287 and 9291 with 121-byte blobs of 51 bins, and 9293-9305 and 9453 with 275-byte blobs of 128 bins). It is version 1 with an accelerometer gain inserted at offset 14 and the mode byte made explicit.

\begingroup\footnotesize\setlength{\tabcolsep}{3pt}
\begin{longtable}{@{}>{\raggedleft\arraybackslash}p{0.0714\linewidth}>{\raggedleft\arraybackslash}p{0.0446\linewidth}>{\raggedright\arraybackslash}p{0.0714\linewidth}>{\raggedright\arraybackslash}p{0.2679\linewidth}>{\raggedright\arraybackslash}p{0.2054\linewidth}>{\raggedright\arraybackslash}p{0.2589\linewidth}@{}}
\toprule
\textbf{Offset} & \textbf{Size} & \textbf{Type} & \textbf{Field} & \textbf{Scale / units} & \textbf{Notes} \\
\midrule
\endfirsthead
\toprule
\textbf{Offset} & \textbf{Size} & \textbf{Type} & \textbf{Field} & \textbf{Scale / units} & \textbf{Notes} \\
\midrule
\endhead
\bottomrule
\endlastfoot
0 & 1 & U1 & ID &  & 0xC1 or 0xC3. \\
1 & 2 & U2 & \code{jj} & bytes & 18 + 2·n + 1. [decoder] Must be at least 19. \\
3 & 1 & U1 & Version &  & 2. \\
4 & 1 & U1 & Packet index &  & {}[inferred] 0 in every production blob seen. Not decoded. \\
5 & 1 & U1 & Mode & enum & 0 = raw, 1 = spectral. Above 1 throws “Invalid sample mode!”. For mode 0 nothing is decoded. \\
6 & 1 & U1 & Spectral axis & enum, C.3.5 &  \\
7 & 1 & U1 & Spectral type & enum, C.3.5 &  \\
8 & 1 & U1 & Segment length power p &  & N = 2\textasciicircum{}p. \\
9 & 2 & U2 & Base sample rate f\_\allowbreak{}base & Hz &  \\
11 & 2 & U2 & Sample rate divider d &  & f\_\allowbreak{}s = f\_\allowbreak{}base / (d + 1 \\
13 & 1 & U1 & Accelerometer full-scale range & g & 2, 4, 8 or 16. \\
14 & 2 & U2 & Accelerometer software gain (Q3) & ÷ 8 & Gain G = value / 8. Production: 0x0020, so G = 4.0. \\
16 & 2 & U2 & First bin index k0 &  &  \\
18 + 2i & 2 & U2 / I2 & Spectral value, bin k0 + i & counts, ÷ ESSF → g²/Hz & Signedness as in v1. n = (\code{jj} - 19)/2. \\
\code{jj} - 1 & 1 & char & \code{;\allowbreak{}} &  & Terminator. \\
\end{longtable}
\endgroup

Production v2 header (float 9293, dive 306, 275 bytes): \code{C1 01 13 |\allowbreak{} 02 |\allowbreak{} 00 |\allowbreak{} 01 |\allowbreak{} 05 |\allowbreak{} 01 |\allowbreak{} 08 |\allowbreak{} 00 2D |\allowbreak{} 00 11 |\allowbreak{} 04 |\allowbreak{} 00 20 |\allowbreak{} 00 00 |\allowbreak{} 00 03 00 02 …} That is version 2, packet 0, spectral, axis ZZ, type CO, N = 256, f\_\allowbreak{}s = 45/18 = 2.5 Hz, ±4 g, G = 4.0, k0 = 0, then 128 bins.

\subsubsection{\texorpdfstring{C.3.5 Enumerations}{C.3.5 Enumerations}}\label{sec:c35-enumerations}

Spectral axis (byte 6):

\begingroup\footnotesize\setlength{\tabcolsep}{3pt}
\begin{longtable}{@{}>{\raggedleft\arraybackslash}p{0.0536\linewidth}>{\raggedright\arraybackslash}p{0.0893\linewidth}>{\raggedright\arraybackslash}p{0.0893\linewidth}>{\raggedleft\arraybackslash}p{0.0536\linewidth}>{\raggedright\arraybackslash}p{0.0446\linewidth}>{\raggedright\arraybackslash}p{0.0714\linewidth}@{}}
\toprule
\textbf{Value} & \textbf{Name} & \textbf{Decoded} & \textbf{Value} & \textbf{Name} & \textbf{Decoded} \\
\midrule
\endfirsthead
\toprule
\textbf{Value} & \textbf{Name} & \textbf{Decoded} & \textbf{Value} & \textbf{Name} & \textbf{Decoded} \\
\midrule
\endhead
\bottomrule
\endlastfoot
0 & undefined & no & 8 & X & no \\
1 & Z′ & no & 9 & XZ & no \\
2 & Y′ & no & 10 & XX & co (U2 \\
3 & X′ & no & 11 & XY & co (I2 \\
4 & Z & no & 12 & Y & no \\
5 & ZZ & co (U2) & 13 & YZ & no \\
6 & ZX & quad (I2) & 14 & YX & no \\
7 & ZY & quad (I2) & 15 & YY & co (U2 \\
\end{longtable}
\endgroup

Spectral type (byte 7): 0 undefined (throws), 1 CO (co-spectrum), 2 QU (quadrature spectrum), 3 RE, 4 IM, 5 MAG and 6 PHASE (3-6 throw “Spectral type not supported!”). [decoder] An axis/type pair other than CO with ZZ, XX, YY or XY, or QU with ZX or ZY, throws, and the blob is discarded.

Mode (v2 byte 5): 0 = raw, 1 = spectral.

\subsubsection{\texorpdfstring{C.3.6 Scale factors}{C.3.6 Scale factors}}\label{sec:c36-scale-factors}

\begingroup\footnotesize\setlength{\tabcolsep}{3pt}
\begin{longtable}{@{}>{\raggedright\arraybackslash}p{0.1518\linewidth}>{\raggedright\arraybackslash}p{0.2143\linewidth}>{\raggedright\arraybackslash}p{0.1964\linewidth}>{\raggedright\arraybackslash}p{0.1518\linewidth}>{\raggedright\arraybackslash}p{0.2143\linewidth}@{}}
\toprule
\textbf{Table} & \textbf{Index or key} & \textbf{Values} & \textbf{Units} & \textbf{Used} \\
\midrule
\endfirsthead
\toprule
\textbf{Table} & \textbf{Index or key} & \textbf{Values} & \textbf{Units} & \textbf{Used} \\
\midrule
\endhead
\bottomrule
\endlastfoot
\code{ASSF\_\allowbreak{}SPECTRAL\_\allowbreak{}V1} & full-scale range 2, 4, 8, 16 g & 8192, 2048, 512, 128 & counts per g²/Hz & v0, v1 \\
\code{ASSF\_\allowbreak{}SPECTRAL\_\allowbreak{}V2} & 2, 4, 8, 16 g & 8192, 2048, 512, 128 & counts per g²/Hz & v2. Before build 3238 (2025-10) it was \{4096, 1024, 256, 64\}, which made \code{wvs}/\code{wvm} values twice as large. \\
\code{ASSF\_\allowbreak{}RAW} & AFS\_\allowbreak{}SEL 0-3 (±2, 4, 8, 16 g) & 16384, 8192, 4096, 2048 & LSB/g & raw format (not decoded \\
\code{GSSF} & FS\_\allowbreak{}SEL 0-3 (±250, 500, 1000, 2000 °/s) & 131.072, 65.536, 32.768, 16.384 & LSB/(°/s) & raw format (not decoded \\
\code{MSSF} & 0 & 30/9 ≈ 3.333 & LSB/µT & raw format (not decoded \\
\code{FSAM} & 0 & 1 & Hz & raw format (not decoded \\
\end{longtable}
\endgroup

\begin{itemize}
\item An invalid full-scale-range byte logs an error and falls back to ASSF = 2048 (the 4 g entry).
\item Effective spectral scale (WIMS user manual rev. 3.15, “Scale Sensitivity”): ESSF = 32768 · ASSF · G² · U · f\_\allowbreak{}e² / (2 N²), in counts per g²/Hz, where:

\begin{itemize}
\item G = Q / 8 is the acceleration gain (Q the raw Q3 gain; G = 1 when Q is 0, and for v0/v1);
\item U = 1 - (5L/4 + 1) / N is the mean square of the firmware's Tukey window, with taper length L = round((N - 1) · 64 / 256) (\code{wav.\allowbreak{}twalp} = 64, which the header does not carry; U = 0.68359 for N = 256);
\item f\_\allowbreak{}e = f\_\allowbreak{}base / (d + 1) is the true effective sample rate, not the doubled rate of C.3.8.
\end{itemize}

Example: v2 at ±4 g, G = 4, N = 256, f\_\allowbreak{}e = 2.5 Hz gives ESSF = 35000; at f\_\allowbreak{}e = 5 Hz, 140000. Before 2026-10-28 (build 3332) the decoder used ESSF = ⌊ASSF · Q² / 64⌋ = ASSF · G², the formula of manual revisions up to 3.14, which ignores f\_\allowbreak{}e, N and the window (32768 in the example).
\item Acceleration spectrum: S\_\allowbreak{}aa [g²/Hz] = counts / ESSF.
\item Displacement spectrum [decoder]: S\_\allowbreak{}ηη [m²/Hz] = g² · S\_\allowbreak{}aa / (2πf)⁴, with g = 9.80665 m/s² and π = 3.14159265359. It is clamped to at most +999.9999. The cross terms \code{Cxy}, \code{Qzx} and \code{Qzy} are also clamped to at least -999.9999.
\item {}[decoder] ASSF, gain, f\_\allowbreak{}base, d and N are kept per dive and overwritten by each blob, so the last blob decoded sets the scale for all components.
\end{itemize}

\subsubsection{\texorpdfstring{C.3.7 Frequency bins}{C.3.7 Frequency bins}}\label{sec:c37-frequency-bins}

For bin k (0 … N/2 - 1), f\_\allowbreak{}k = f\_\allowbreak{}s · k / N, with f\_\allowbreak{}s = f\_\allowbreak{}base / (d + 1). Bin 0 is the DC bin; the last bin is f\_\allowbreak{}s/2 - f\_\allowbreak{}s/N. WIMS firmware up to 2.01.00614b labels its own output one bin high, (k + 1) · f\_\allowbreak{}s / N, and the decoder did the same until 2026-10-26 (see below).

\begin{itemize}
\item Build 3238 (2025-10) went back to \code{fs·(\allowbreak{}i+1)/\allowbreak{}seglen}.
\item 2026-10-26: \code{fs·(\allowbreak{}k0+i)/\allowbreak{}seglen}, with bin 0 the DC bin. The WIMS firmware copies FFT bins 0 to N/2 - 1, so \code{fs·(\allowbreak{}i+1)/\allowbreak{}seglen} was one bin too high; an intercomparison with NDBC buoy 41047 confirmed the offset.
\item The 2026-10-24 fix added k0 to k, so bins in a blob with k0 > 0 get the right frequency. Before that, the frequency ignored k0.
\item {}[decoder] The spectral array holds N/2 bins per dive. A bin that no blob filled keeps frequency -99 and fill values. For example 9287 and 9291 send only bins 0-50 of N/2 = 256.
\item Newer sensor firmware truncates the spectrum to a range of frequencies, so a blob may carry fewer than N/2 bins (n from \code{jj}, as in C.3.3 and C.3.4). [decoder] Only the n bins in the blob are decoded; before build 3315 the decoder tried all N/2 and logged “wave sample index exceeds message size” for every truncated blob.
\item {}[decoder] Bins that would reach past N/2 (k0 + i ≥ N/2) are not decoded, and a warning is logged; before build 3315 \code{vector:\allowbreak{}:\allowbreak{}at} threw and the whole blob was discarded.
\end{itemize}

\subsubsection{\texorpdfstring{C.3.8 Every-other-value correction (WIMS firmware)}{C.3.8 Every-other-value correction (WIMS firmware)}}\label{sec:c38-every-other-value-correction-wims-firmware}

{}[inferred from production; README 0.6.3235] Firmware 0x0A060231, 0x0A080031, 0x0A090131 and 0x0A0A0031 send only every other spectral value: the first half of the N/2 bins holds the spectrum at twice the bin spacing, and the second half is zeros. [decoder] \code{XFloat\allowbreak{}Dive\allowbreak{}Data:\allowbreak{}:\allowbreak{}Resolve\allowbreak{}Wave\allowbreak{}Bin\allowbreak{}Spacing(\allowbreak{})} runs after each 0xC1/0xC3 blob and after each mission configuration message, so the result does not depend on message order. It takes the firmware word from the dive's own D0 or, without one, the float's latest; if that word is one of these, \code{XWave:\allowbreak{}:\allowbreak{}set\_\allowbreak{}every\_\allowbreak{}other\_\allowbreak{}bin(\allowbreak{}true)}:

\begin{itemize}
\item sets f\_\allowbreak{}s' = 2 · f\_\allowbreak{}base / (d + 1), which also becomes the \code{WVS} \code{fsam}; the ESSF (C.3.6) still uses the true f\_\allowbreak{}e;
\item sets the frequency of every received bin, i = 0 … N/2 - 1, to f\_\allowbreak{}i = f\_\allowbreak{}s' · i / N (f\_\allowbreak{}s' · (i + 1) / N before 2026-10-26); bins not received keep -99. [inferred from the NDBC 41047 intercomparison, whose offset matches one decoded bin] The floats send the even WIMS bins (0, 2, 4, …), so bin i holds WIMS FFT bin 2i.
\end{itemize}

When the firmware word is unknown (it is not kept across decoder restarts), \code{XWave:\allowbreak{}:\allowbreak{}Looks\allowbreak{}Like\allowbreak{}Every\allowbreak{}Other\allowbreak{}Bin(\allowbreak{})} decides from the data: every bin of the upper half was received and is zero in all six components, and some bin of the lower half has a nonzero Szz. Truncated spectra leave the upper bins unreceived, not zero, so they do not match. Before 2026-10-28 (build 3332) the correction was keyed only on the float's latest known firmware word, so blobs decoded before any D0 had been received were not corrected.

The spectral values are left as they are, so the zero second half is still listed, at frequencies above f\_\allowbreak{}s/2. 10.12.1.31 (0x0A0C0131) and 10.7.1.31 (0x0A070131) are WIMS firmware too, but they are not in the correction list (C.5).

\subsubsection{\texorpdfstring{C.3.9 Raw IMU format (XF2.4 “Raw Wave Field”; not decoded)}{C.3.9 Raw IMU format (XF2.4 “Raw Wave Field”; not decoded)}}\label{sec:c39-raw-imu-format-xf24-raw-wave-field-not-decoded}

\begingroup\footnotesize\setlength{\tabcolsep}{3pt}
\begin{longtable}{@{}>{\raggedleft\arraybackslash}p{0.0804\linewidth}>{\raggedleft\arraybackslash}p{0.0446\linewidth}>{\raggedright\arraybackslash}p{0.0446\linewidth}>{\raggedright\arraybackslash}p{0.2946\linewidth}>{\raggedright\arraybackslash}p{0.1161\linewidth}>{\raggedright\arraybackslash}p{0.3393\linewidth}@{}}
\toprule
\textbf{Offset} & \textbf{Size} & \textbf{Type} & \textbf{Field} & \textbf{Scale} & \textbf{Notes} \\
\midrule
\endfirsthead
\toprule
\textbf{Offset} & \textbf{Size} & \textbf{Type} & \textbf{Field} & \textbf{Scale} & \textbf{Notes} \\
\midrule
\endhead
\bottomrule
\endlastfoot
0 & 1 & U1 & ID &  & 0xC1. \\
1 & 2 & U2 & \code{jj} &  &  \\
3 & 1 & U1 & Packet index &  & Packet order for rebuilding the time series. \\
4 & 1 & U1 & Sample frequency index & \code{FSAM} & 0 = 1 Hz. \\
5 & 1 & U1 & ASSF index & \code{ASSF\_\allowbreak{}RAW} &  \\
6 & 1 & U1 & GSSF index & \code{GSSF} &  \\
7 & 1 & U1 & MSSF index & \code{MSSF} &  \\
8 & 1 & U1 & Mode (high nibble) / phase code (low nibble) &  & Mode 0 = raw, 1 = spectral. Phase code 1-14 (14 = surface). The nibble order is [inferred] from the commented-out decoder, which takes the phase code from \code{d[8] \& 0x0F}. \\
9 + 18n & 2 & I2 & Accel x, sample n & ÷ ASSF → g & Starboard. \\
11 + 18n & 2 & I2 & Accel y & ÷ ASSF → g & Bow. \\
13 + 18n & 2 & I2 & Accel z & ÷ ASSF → g & Up. \\
15 + 18n & 2 & I2 & Gyro x & ÷ GSSF → °/s & XF2.4: bow-up (pitch). The decoder comment calls it gamma, starboard-up. \\
17 + 18n & 2 & I2 & Gyro y & ÷ GSSF → °/s & XF2.4: starboard-up (roll). The decoder comment calls it beta, bow-up. \\
19 + 18n & 2 & I2 & Gyro z & ÷ GSSF → °/s & Increasing azimuth. \\
21 + 18n & 2 & I2 & Mag x & ÷ MSSF → µT &  \\
23 + 18n & 2 & I2 & Mag y & ÷ MSSF → µT &  \\
25 + 18n & 2 & I2 & Mag z & ÷ MSSF → µT &  \\
18N + 27 & 1 & char & \code{;\allowbreak{}} &  & N = index of the last sample. \\
\end{longtable}
\endgroup

This layout has no version byte. Byte 3 is the packet index and byte 5 the ASSF index. The current decoder would therefore read a raw packet as a spectral version (packet index 0 or 1) or ignore it. In the v2 layout, byte 5 is the mode, and mode 0 (raw) is accepted but nothing is decoded. The old raw decoder and its \code{WAV}/\code{wav} listing lines are commented out in \code{XWave.\allowbreak{}cpp}.

\subsubsection{\texorpdfstring{C.3.10 Validation summary [decoder]}{C.3.10 Validation summary [decoder]}}\label{sec:c310-validation-summary-decoder}

\begin{enumerate}
\item \code{Check\allowbreak{}XMsg\allowbreak{}Data(\allowbreak{}data,\allowbreak{} -1,\allowbreak{} 17)}: the blob is at least 17 bytes, \code{jj} (lower 12 bits) is at least 4 and no more than the data received, and the byte at \code{jj - 1} is 0x3B. Extra bytes after \code{jj} only give a warning.
\item Version 0, 1 or 2; any other version is ignored without an error.
\item v2: mode 0 or 1 (otherwise it throws), and \code{jj} at least 19.
\item Each bin is read only while its two bytes lie before the terminator. For v1 the loop stops when 2i + 17 ≥ \code{jj}; for v2, when 2i + 20 ≥ \code{jj}. At most N/2 bins are read. The 2026-10-24 fix corrected the v2 bound, which used to read one bin past the terminator and broke the 51-bin floats' products.
\item Unsupported axis/type combinations throw, and the blob is dropped (\code{CATCH\_\allowbreak{}BLOCK}).
\end{enumerate}

\subsubsection{\texorpdfstring{C.3.11 Listing output (\code{WVS}, \code{wvs}, \code{wvm})}{C.3.11 Listing output (WVS, wvs, wvm)}}\label{sec:c311-listing-output-wvs-wvs-wvm}

The section “Ocean Surface Wave-Field Data” is printed when the dive has wave data with N > 0. Only versions 1 and 2 print anything: a dive whose wave blobs are all version 0 gets an empty section. In the example, the \code{wvs} and \code{wvm} values were listed before build 3332 with ESSF 32768; the current decoder lists them about 7\% lower (ESSF 35100).

\begin{Verbatim}
#      SN dive vers slen     fsam fbase fdiv     assf  again       essf
#                            (Hz)  (Hz)      (LSB/g^2)   (LSB/(g^2/Hz))
WVS  9287  325 0x02  512   5.0000   45     9     2048  4.000    35100.0
#      SN dive  idx     freq         Szz         Syy         Sxx         Cxy         Qzx         Qzy
#                       (Hz)    (g^2/Hz)    (g^2/Hz)    (g^2/Hz)    (g^2/Hz)    (g^2/Hz)    (g^2/Hz)
wvs  9287  325    0   0.0098    +0.00018    +0.00098    +0.00156    +0.00046    +0.00000    +0.00000
wvs  9287  325   51 -99.0000  -999.99999  -999.99999  -999.99999  -999.99999  -999.99999  -999.99999
#      SN dive  idx     freq         Szz         Syy         Sxx         Cxy         Qzx         Qzy
#                       (Hz)    (m^2/Hz)    (m^2/Hz)    (m^2/Hz)    (m^2/Hz)    (m^2/Hz)    (m^2/Hz)
wvm  9287  325   50   0.4980     +0.0056     +0.0159     +0.0109     +0.0003     -0.0038     -0.0024
wvm  9287  325   51 -99.0000  -1000.0000  -1000.0000  -1000.0000  -1000.0000  -1000.0000  -1000.0000
\end{Verbatim}

\begin{itemize}
\item \code{WVS} (\code{WVS \%5d \%4d 0x\%02X \%4d \%8.\allowbreak{}4f \%4d \%5d \%8d \%6.\allowbreak{}3f \%10.\allowbreak{}1f}; \code{essf} was \code{\%8d} before build 3332) has these columns:

\begin{itemize}
\item \code{vers}: the version;
\item \code{slen}: N;
\item \code{fsam}: f\_\allowbreak{}s, doubled by the C.3.8 correction;
\item \code{fbase}: f\_\allowbreak{}base;
\item \code{fdiv}: d + 1, not the raw divider;
\item \code{assf};
\item \code{again}: G;
\item \code{essf}.
\end{itemize}
\item \code{wvs} (\code{\%8.\allowbreak{}4f} frequency, \code{\%+11.\allowbreak{}5f} values) gives one line per bin, idx = 0 … N/2 - 1, with the acceleration spectra in g²/Hz. The columns are always in the order \textbf{zz, yy, xx, xy, zx, zy}: Szz, Syy, Sxx, Cxy, Qzx, Qzy.
\item \code{wvm} (\code{\%+11.\allowbreak{}4f}) gives the same bins as displacement spectra in m²/Hz. It \textbf{must use the same column order as \code{wvs}} (zz, yy, xx, xy, zx, zy) under the same header. Builds 3234-3289 (2025-09 to 2026-10-24) printed \code{wvm} as xx, xy, yy, zz, zx, zy under the zz-first header, so \code{wvm} columns in listings from those builds are mislabeled. The 2026-10-24 fix restored the \code{wvs} order. Nothing downstream parses \code{wvm}.
\item Fill: a component never received is -999.99999 in \code{wvs}. In \code{wvm} the same fill prints as \code{-1000.\allowbreak{}0000}, because \code{\%+11.\allowbreak{}4f} rounds -999.99999. Unfilled bins have frequency -99.0000.
\end{itemize}

\subsection{\texorpdfstring{C.4 Marginal-ice-zone acoustic messages (0xCA)}{C.4 Marginal-ice-zone acoustic messages (0xCA)}}\label{sec:c4-marginal-ice-zone-acoustic-messages-0xca}

\subsubsection{\texorpdfstring{C.4.1 Purpose}{C.4.1 Purpose}}\label{sec:c41-purpose}

MIZ floats (for example 9213 and 9216 with firmware 10.12.0.32) listen for acoustic navigation and communication signals (CACST, “communication and navigation” pings) in scheduled listening windows. Each listening window and each received signal is reported in its own 0xCA blob, with a subtype byte at offset 3. Shore-side processing uses the times of arrival to position the float under ice. The type ID was 0xC0 in decoder builds before 3249 (2026-05/06); the decoder now uses 0xCA (\code{XMSG\_\allowbreak{}TYPE\_\allowbreak{}ID\_\allowbreak{}MIZ\_\allowbreak{}ACOUSTIC}). XF2.4 lists 0xC0 “acoustic position” as unsupported, with an empty sheet. DF7.3.5's blob 192 (0xC0) is a different text payload (\code{sensor\_\allowbreak{}type}, \code{payload\_\allowbreak{}index}, 333-character text).

\subsubsection{\texorpdfstring{C.4.2 Common header}{C.4.2 Common header}}\label{sec:c42-common-header}

\begingroup\footnotesize\setlength{\tabcolsep}{3pt}
\begin{longtable}{@{}>{\raggedleft\arraybackslash}p{0.0625\linewidth}>{\raggedleft\arraybackslash}p{0.0446\linewidth}>{\raggedright\arraybackslash}p{0.0446\linewidth}>{\raggedright\arraybackslash}p{0.0714\linewidth}>{\raggedright\arraybackslash}p{0.7054\linewidth}@{}}
\toprule
\textbf{Offset} & \textbf{Size} & \textbf{Type} & \textbf{Field} & \textbf{Notes} \\
\midrule
\endfirsthead
\toprule
\textbf{Offset} & \textbf{Size} & \textbf{Type} & \textbf{Field} & \textbf{Notes} \\
\midrule
\endhead
\bottomrule
\endlastfoot
0 & 1 & U1 & ID & 0xCA. [decoder] Checked. \\
1 & 2 & U2 & \code{jj} & {}[decoder] \code{Check\allowbreak{}XMsg\allowbreak{}Data(\allowbreak{}data,\allowbreak{} 0x\allowbreak{}CA,\allowbreak{} 4)}: \code{jj} (lower 12 bits) at least 4 and no more than the data received, terminator at \code{jj - 1}. \\
3 & 1 & U1 & Subtype & 3 = listening window, 5 = CACST receive record. Other subtypes are ignored without an error (1 and 2 seen in production, C.4.5 \\
\end{longtable}
\endgroup

\subsubsection{\texorpdfstring{C.4.3 Subtype 3: listening window (37 bytes)}{C.4.3 Subtype 3: listening window (37 bytes)}}\label{sec:c43-subtype-3-listening-window-37-bytes}

\begingroup\footnotesize\setlength{\tabcolsep}{3pt}
\begin{longtable}{@{}>{\raggedleft\arraybackslash}p{0.0625\linewidth}>{\raggedleft\arraybackslash}p{0.0446\linewidth}>{\raggedright\arraybackslash}p{0.0446\linewidth}>{\raggedright\arraybackslash}p{0.2054\linewidth}>{\raggedright\arraybackslash}p{0.1161\linewidth}>{\raggedright\arraybackslash}p{0.4464\linewidth}@{}}
\toprule
\textbf{Offset} & \textbf{Size} & \textbf{Type} & \textbf{Field} & \textbf{Units} & \textbf{Notes} \\
\midrule
\endfirsthead
\toprule
\textbf{Offset} & \textbf{Size} & \textbf{Type} & \textbf{Field} & \textbf{Units} & \textbf{Notes} \\
\midrule
\endhead
\bottomrule
\endlastfoot
0 & 1 & U1 & ID &  & 0xCA. \\
1 & 2 & U2 & \code{jj} &  & 37 (0x0025 \\
3 & 1 & U1 & Subtype &  & 3. \\
4 & 4 & U4 & Window start time & UNIX s (UTC) &  \\
8 & 4 & U4 & “Pressure” at start & unknown & {}[decoder] Named \code{p\_\allowbreak{}start}, not listed. Production values 0x00E5-0x02DF (229-735) on floats parked at 250-300 dBar; the scale is unknown. \\
12 & 4 & U4 & “Temperature” at start & unknown & {}[decoder] Named \code{t\_\allowbreak{}start}, not listed. Always 0 in production. \\
16 & 4 & U4 & “Salinity” at start & unknown & {}[decoder] Named \code{s\_\allowbreak{}start}, not listed. In production the high 16 bits are about 0x0021-0x0022 (33-34) and the low 16 bits 0x002A-0x002D (42-45). [inferred] It may be two 16-bit quantities, not one U4. \\
20 & 4 & U4 & Window end time & UNIX s (UTC) &  \\
24 & 4 & U4 & “Pressure” at end & unknown & As offset 8. \\
28 & 4 & U4 & “Temperature” at end & unknown & As offset 12. \\
32 & 4 & U4 & “Salinity” at end & unknown & In production the low 16 bits are 0. \\
36 & 1 & char & \code{;\allowbreak{}} &  &  \\
\end{longtable}
\endgroup

Example (9213): \code{CA 00 25 03 6AB85C02 00000137 00000000 0022002D 6AB863B9 0000007F 00000000 00200000 3B} is a window from 2026-09-26T23:57:54Z to 2026-09-27T00:30:49Z.

\subsubsection{\texorpdfstring{C.4.4 Subtype 5: CACST receive record (33 bytes)}{C.4.4 Subtype 5: CACST receive record (33 bytes)}}\label{sec:c44-subtype-5-cacst-receive-record-33-bytes}

\begingroup\footnotesize\setlength{\tabcolsep}{3pt}
\begin{longtable}{@{}>{\raggedleft\arraybackslash}p{0.0625\linewidth}>{\raggedleft\arraybackslash}p{0.0446\linewidth}>{\raggedright\arraybackslash}p{0.0446\linewidth}>{\raggedright\arraybackslash}p{0.2946\linewidth}>{\raggedright\arraybackslash}p{0.1250\linewidth}>{\raggedright\arraybackslash}p{0.3482\linewidth}@{}}
\toprule
\textbf{Offset} & \textbf{Size} & \textbf{Type} & \textbf{Field} & \textbf{Units / scale} & \textbf{Notes} \\
\midrule
\endfirsthead
\toprule
\textbf{Offset} & \textbf{Size} & \textbf{Type} & \textbf{Field} & \textbf{Units / scale} & \textbf{Notes} \\
\midrule
\endhead
\bottomrule
\endlastfoot
0 & 1 & U1 & ID &  & 0xCA. \\
1 & 2 & U2 & \code{jj} &  & 33 (0x0021 \\
3 & 1 & U1 & Subtype &  & 5. \\
4 & 4 & U4 & Counter &  & {}[decoder] Named \code{counter}, not listed. In production it goes up by 1 per record, and its first value in each session looks like a UNIX time shortly before the first window (for example 0x6AB826EB = 2026-09-26T19:58:03Z) [inferred]. \\
8 & 4 & U4 & Time of arrival, seconds & UNIX s (UTC) &  \\
12 & 4 & U4 & Time of arrival, microseconds & µs & Values above 999999 make the listed time invalid. \\
16 & 4 & U4 & Matched-filter peak power & dB & Production values 41-45. \\
20 & 2 & I2 & Source ID &  & -1 (0xFFFF) in every production record. [decoder] Read as U2 and stored as int16. \\
22 & 2 & U2 & PSK error code & enum & 4 in every production record. The codes are not documented. \\
24 & 4 & I4 & Input SNR & 0.1 dB & Production: 0xFFFFFC19 = -999, so -99.9 dB, apparently a “no value” marker [inferred]. \\
28 & 4 & I4 & Output SNR (after the equalizer) & 0.1 dB & Production: 0xFFFFFF9D = -99, so -9.9 dB. \\
32 & 1 & char & \code{;\allowbreak{}} &  &  \\
\end{longtable}
\endgroup

Example (9213): \code{CA 00 21 05 6AB826ED 6AB85CEF 00033D5F 0000002D FFFF 0004 FFFFFC19 FFFFFF9D 3B} is a reception at 2026-09-27T00:01:51.212319Z with MFD power 45 dB, source -1, PSK error 4, SNR in -99.9 dB and SNR out -9.9 dB.

\subsubsection{\texorpdfstring{C.4.5 Subtypes seen in production but not decoded [inferred]}{C.4.5 Subtypes seen in production but not decoded [inferred]}}\label{sec:c45-subtypes-seen-in-production-but-not-decoded-inferred}

\begingroup\footnotesize\setlength{\tabcolsep}{3pt}
\begin{longtable}{@{}>{\raggedleft\arraybackslash}p{0.0714\linewidth}>{\raggedleft\arraybackslash}p{0.0268\linewidth}>{\raggedright\arraybackslash}p{0.4018\linewidth}>{\raggedright\arraybackslash}p{0.4464\linewidth}@{}}
\toprule
\textbf{Subtype} & \textbf{\code{jj}} & \textbf{Layout as seen} & \textbf{Notes} \\
\midrule
\endfirsthead
\toprule
\textbf{Subtype} & \textbf{\code{jj}} & \textbf{Layout as seen} & \textbf{Notes} \\
\midrule
\endhead
\bottomrule
\endlastfoot
1 & 10 & \code{CA 00 0A 01} · U1 (always 0x23) · U4 UNIX time · \code{;\allowbreak{}} & The times are on a 4-hour grid, for example 0x6AB82404 and 0x6AB85C44. Possibly a schedule or event marker. \\
2 & 14 & \code{CA 00 0E 02} · U1 (always 0x21) · U4 UNIX time · I4 · \code{;\allowbreak{}} & The times fall 120 s before the subtype 1 times. The I4 values are small and signed (-120 to +49), so possibly a clock offset or a measurement. \\
\end{longtable}
\endgroup

\subsubsection{\texorpdfstring{C.4.6 Decoding notes [decoder]}{C.4.6 Decoding notes [decoder]}}\label{sec:c46-decoding-notes-decoder}

\begin{itemize}
\item Every 0xCA blob of a dive goes into one \code{XAcoustic\allowbreak{}Position} object. A window or CACST blob must be exactly 37 or 33 bytes (the received blob length); otherwise it is rejected with an error.
\item Before listing, \code{Assign\allowbreak{}Window\allowbreak{}Indices(\allowbreak{})} numbers the windows 0, 1, … in the order they were received, not in time order. Each CACST record gets the index of the first window with start ≤ ToA seconds ≤ end, and that window's receive count goes up. A record outside every window keeps index -1.
\item Window P/T/S, the counter and the per-window receive count are decoded but not listed.
\end{itemize}

\subsubsection{\texorpdfstring{C.4.7 Listing output (\code{APS})}{C.4.7 Listing output (APS)}}\label{sec:c47-listing-output-aps}

\code{APS} lines appear under “Position Data / Acoustic Position”. That section is printed only when the dive also has GPS fixes or mail messages; acoustic data alone no longer counts as position data.

\begin{Verbatim}
# --- Acoustic Listening Windows -------------------------------------------------------------------
#      SN dive midx widx      start-date/time        end-date/time
#                                       (UTC)                (UTC)
APS  9213    1    3    0 2026-09-26T23:57:54Z 2026-09-27T00:30:49Z
# --- Acoustic Communication Cycle Receive Statistics ----------------------------------------------
#      SN dive midx widx                   date/time    src MFDPwr PSKErr  SNRIn SNROut
#                                              (UTC)          (dB)          (dB)   (dB)
APS  9213    1    5    0 2026-09-27T00:01:51.212319Z     -1     45      4  -99.9  -99.9
\end{Verbatim}

\begin{itemize}
\item Window lines (\code{APS \%5d \%4d \%4d \%4d \%20s \%20s}): \code{midx} is the subtype (3), \code{widx} the window index, then the start and end times. An unset time prints as \code{0000-00-00T00:\allowbreak{}00:\allowbreak{}00Z}.
\item Receive lines (\code{APS \%5d \%4d \%4d \%4d \%27s \%6d \%6d \%6d \%6.\allowbreak{}1f \%6.\allowbreak{}1f}): \code{midx} = 5, \code{widx} (-1 if the record is in no window), ToA with microseconds (an invalid time prints as \code{0000-00-00T00:\allowbreak{}00:\allowbreak{}00.\allowbreak{}000000Z}), source ID, MFD power, PSK error code, SNRIn and SNROut in dB.
\item \textbf{Bug:} the SNROut column prints SNRIn again. The example above should show -9.9 for SNROut. The README notes this as still present. \emph{(Fixed in build 3309.)}
\end{itemize}

\subsection{\texorpdfstring{C.5 Discrepancies and open questions}{C.5 Discrepancies and open questions}}\label{sec:c5-discrepancies-and-open-questions}

Mission configuration:

\begin{enumerate}
\item \textbf{Parameters version.} NG10.5 and NG10.12 document \code{Mission\allowbreak{}Params\allowbreak{}Version} = 23 (0x17). Production 10.12.0.x floats send 0x18. The decoder doesn't use the field.
\item \textbf{Byte 3 of D0.} XF2.4 calls it the encoding type (0 = binary, otherwise ASCII). NG10.5 calls it “Message version” and NG10.12 \code{Mission\allowbreak{}Config\allowbreak{}Version} = 0. The active decoder ignores it, so a SOLO-II ASCII dump would be misread (C.1.2).
\item \textbf{Byte 7 of the firmware word.} The task calls it “build”. NG10.5 calls it “Float model” and NG10.12 “Product”. Values such as 0x31 (WIMS) and 0x32/0x33 (MM2/Micro-ALTO) suggest a product or variant code. The decoder uses all 32 bits as the key.
\item \textbf{\code{DRT} type.} NG10.5 and NG10.12 give I2; the decoder reads U2. \textbf{\code{MSN}}: NG10.5 says I2, NG10.12 U2.
\item \textbf{XF2.4 against the decoder for old firmware:}

\begin{itemize}
\item 2.0.69 (XF2.4: 323 bytes with \code{CTF}/\code{CTL}/\code{CTU} but no ice block; decoder: 322 bytes with ice, the same as XF2.4's 2.5.62 layout).
\item 3.0.65 (XF2.4: 331 bytes ending \code{SPT}, \code{SSP}; decoder: 328 bytes, no \code{SPT}/\code{SSP}).
\item 4.2.62 (XF2.4: 327 bytes with no gap at 119 and a 5-byte “empty?” field at 161; decoder: 323 bytes, \code{IST} at 120, \code{DNS} at 162).
\item 4.6.61 (XF2.4: one 336-byte D0 ending \code{SSP}; decoder: 328 bytes ending \code{SAN}, and D1 not decoded).
\end{itemize}

The decoder's sizes presumably come from real messages; the spreadsheet looks like intent.
\item \textbf{\code{ICF}/\code{ICW} as U1} in the 2.0.69 to 4.1.61 decoders: only the most significant byte of each 4-byte mask is decoded.
\item \textbf{10.7.1.1 D1} is 88 bytes and probably carries \code{PPC}/\code{PPE}, but the decoder doesn't decode them.
\item \textbf{0x08020000 and 0x0A020201} decode but have no listing function, so they get the generic “unsupported” listing.
\item \textbf{Listing bugs:}

\begin{itemize}
\item \code{DRT} prints \code{DRA}. \emph{(Fixed in build 3309.)}
\item The \code{PPC0}/\code{PPC1} values are swapped against their labels. \emph{(Fixed in build 3309.)}
\item \code{PPC}/\code{PPE} are not listed for 10.7-10.14 non-WIMS floats.
\item The 6.5.77 listing shows \code{WMT} = 65535 because that layout doesn't decode it.
\end{itemize}
\item \textbf{Units:}

\begin{itemize}
\item \code{AHR} (NG10.5: cm/s; listed value/1000 m/s, so 1000 lists as 1.000 m/s).
\item \code{ASI} (NG10.5: s; listed value/1000 s).
\item \code{IST} (header comment: 1/10 PSU; listed value/100).
\item \code{SAN}/\code{SAX} (NG10.5: psi; listed value/1000).
\item \code{DRA} labeled m/s, printed raw.
\item \code{XAS} and \code{CST} codes differ between NG10.5 and XF2.4.
\end{itemize}

These need confirming with MRV.
\item \textbf{SOLO-II ASCII record length.} SOLO2 says 15 bytes per parameter, but its template \code{NAMExx=\allowbreak{}vvvvv|\allowbreak{}} is 13 characters. And 324 bytes of text per 328-byte blob is not a whole number of either.
\item \textbf{D2-D7.} SOLO-II used D0-D3 for the ASCII dump. The binary decoder rejects anything above D1.
\end{enumerate}

Command echo:

\begin{enumerate}[start=13]
\item \textbf{0xDE.} The meaning of the trailing \code{MIC 1} and the 8 NUL bytes is unknown. Neither the decoder nor the listing shows the echo text.
\end{enumerate}

Waves:

\begin{enumerate}[start=14]
\item \textbf{0xC1 against 0xC3.} No document gives the difference. In production, C1 carries the auto-spectra and C3 the cross-spectra; the decoder treats them the same.
\item \textbf{Byte 4 of the wave header} (packet index per DF7.3.5, after alignment) is always 0 in production and not decoded. DF7.3.5 leaves out the version byte.
\item \textbf{Raw IMU format (XF2.4)} has no version byte and conflicts with the version/mode bytes of the spectral formats. It is not decoded.
\item \textbf{The every-other-value correction} covers 0x0A060231, 0x0A080031, 0x0A090131 and 0x0A0A0031. The other WIMS words, 0x0A070131 and 0x0A0C0131, are not covered; whether they need it is unverified.
\item \textbf{ASSF\_\allowbreak{}SPECTRAL} values are an empirical table (counts per g²/Hz). The v2 table was changed in build 3238. The derivation from the IMU LSB/g is not documented.
\item \textbf{Raw gyro axis names:} XF2.4 and the decoder comments disagree on pitch and roll.
\end{enumerate}

Acoustics:

\begin{enumerate}[start=20]
\item \textbf{0xCA P/T/S fields} don't look like plain pressure, temperature and salinity in production (temperature always 0, salinity seemingly two 16-bit halves). Their scaling is unknown.
\item \textbf{The CACST “counter”} looks like a UNIX-time seed plus a running count.
\item \textbf{PSK error codes and source ID} semantics are undocumented.
\item \textbf{Subtypes 1 and 2} are not decoded.
\item \textbf{The SNROut listing column} repeats SNRIn. \emph{(Fixed in build 3309.)}
\end{enumerate}

\subsection{\texorpdfstring{C.6 Extended records (0xFF)}{C.6 Extended records (0xFF)}}\label{sec:c6-extended-records-0xff}

No format document describes ID 0xFF; the decoder's ID list calls it \code{UNKNOWN}. Everything below is \textbf{[inferred]} from the 6,262 0xFF blobs in the production RAW files (survey of 2026-10-27, 160 floats, all firmware from 10.2 to 10.14, WIMS and not). Before build 3321 the decoder logged each one as \code{Unknown X message ID:\allowbreak{} 0xff} and dropped it (several hundred warnings a day).

\textbf{Layout.} The usual ID and 12-bit size (bytes 0-2, A.1.3) and \code{;\allowbreak{}} terminator, with a record type at byte 3 and, except for 0xDB, its version at byte 4. Several record types are also X-message IDs (0xCA, 0xD1, 0xDB, 0xEC), but the records do not match those messages: the type is a separate code space.

\begingroup\footnotesize\setlength{\tabcolsep}{3pt}
\begin{longtable}{@{}>{\raggedright\arraybackslash}p{0.0446\linewidth}>{\raggedright\arraybackslash}p{0.1071\linewidth}>{\raggedright\arraybackslash}p{0.0714\linewidth}>{\raggedright\arraybackslash}p{0.0536\linewidth}>{\raggedright\arraybackslash}p{0.0625\linewidth}>{\raggedright\arraybackslash}p{0.2411\linewidth}>{\raggedright\arraybackslash}p{0.3214\linewidth}@{}}
\toprule
\textbf{Type} & \textbf{Version} & \textbf{Size} & \textbf{Blobs} & \textbf{Floats} & \textbf{When} & \textbf{Record (byte offsets in the blob)} \\
\midrule
\endfirsthead
\toprule
\textbf{Type} & \textbf{Version} & \textbf{Size} & \textbf{Blobs} & \textbf{Floats} & \textbf{When} & \textbf{Record (byte offsets in the blob)} \\
\midrule
\endhead
\bottomrule
\endlastfoot
0xCA & 0x20 & 17 & 2,821 & 160 & 95\% pre-mission & event: u24 code @5 (0x004002 in 83\%; also 0x000100, 0x004000-0x004003, 0x002000, 0x000600, 0x001000), u32 time @8, f32 value @12 \\
0xD1 & 0x10 & 22 & 2,809 & 155 & 95\% pre-mission & u16 counter @5, u32 time @7, u16 count @11, u32 time @13, u32 value @17 (0 in nearly all) \\
0xAE & 0x11 (0x10) & 13 (12) & 483 & 117 & pre-mission only & u16 0 @5, u16 float serial number @7, then 3 (v0x10: 2) configuration bytes \\
0xEC & 0x05, 0x08 & 15, 10 & 89 & 4 & in mission (2 pre-mission) & not decoded, e.g. \code{05 01 0e 00 01 07 01 0b 00 01} (11151) \\
0xFA & 0x0A, 0x02 & 281 & 24 & 10 & 23 pre-mission & text: a processor fault dump, or all NUL \\
0xDB & — & 12-328 & 36 & 4 & in mission & CTD exchange: u16 command length @4, command, u16 reply length, reply \\
\end{longtable}
\endgroup

\begin{itemize}
\item \textbf{Times} are UNIX seconds (e.g. 0x62BF7945 = 2022-07-01T22:46:29Z, 23 minutes before that packet's Iridium session), or seconds since the controller started when they are below 2000-01-01 (37\% of the 0xCA times, e.g. 402). The decoder lists the latter as \code{+402s}.
\item \textbf{0xCA} records come in pairs with a 0xD1 record 8-17 s earlier, several pairs per startup packet. The value is a few tenths of a unit around zero (-0.32 to +0.10 in the examples), most likely the surface pressure in dBar measured at the event; code 0x004002 is the most common event. The 0xD1 counter increases by one per pair across a float's life (4 on 9166 in 2022; 71-73 on 9304 in 2026), so it is probably a controller start (reset) count, with the start time and, in the second time field, an uptime in seconds (\code{+207s}). What \code{count} counts is unknown.
\item \textbf{0xAE} carries the float's serial number (0x23FB = 9211, 0x245B = 9307, 0x2D6D = 11629) and 2-3 small configuration values, sent during the startup tests. They may be the pump and battery types that the BIT format 11 “traveler” errors (B.2.6, MissingPumpType, MissingBatteryType) refer to.
\item \textbf{0xFA} is an ARM Cortex-M fault report: \code{r7}, \code{sp}, \code{excrt} (EXC\_\allowbreak{}RETURN), \code{psr}, the fault status registers \code{hfsr}, \code{cfsr}, \code{mmfar}, \code{bfar}, \code{shcsr}, and the stacked \code{s.\allowbreak{}r0}-\code{s.\allowbreak{}r3}, \code{s.\allowbreak{}r12}, \code{s.\allowbreak{}lr}, \code{s.\allowbreak{}pc}, \code{s.\allowbreak{}psr}, one per line (CR LF), in hex. It records a controller crash, for example 9304 dive 0 (firmware 10.12.1.31) with \code{cfsr 00000082} (a data access violation at a precise address). Empty reports are all NUL.
\item \textbf{0xDB} records the RBR CTD exchanges of 9280 (10.6.2.1) and three other floats (11148, 11151, 11330): the command text, e.g. \code{ctdrw fetch\textbackslash{}\allowbreak{}r|\allowbreak{}5} (a literal backslash-r; the \code{|\allowbreak{}} digit is perhaps a timeout), and the CTD's reply, e.g. \code{2024-03-15 22:\allowbreak{}57:\allowbreak{}12.\allowbreak{}000,\allowbreak{} 25.\allowbreak{}1567,\allowbreak{} 10.\allowbreak{}1491,\allowbreak{} 117.\allowbreak{}6462,\allowbreak{} 25.\allowbreak{}6469\textbackslash{}\allowbreak{}r\textbackslash{}\allowbreak{}n} or a calibration (\code{calibration 3 label =\allowbreak{} pressure\_\allowbreak{}00,\allowbreak{} .\allowbreak{}.\allowbreak{}.\allowbreak{}}) or \code{id model =\allowbreak{} RBRargo3,\allowbreak{} .\allowbreak{}.\allowbreak{}.\allowbreak{}} line. The lengths account for every byte of the blob.
\end{itemize}

\textbf{Listing.} Under the section banner “Extended Record Messages (0xFF)”, one line per record: \code{XFF SN dive ridx 0x\allowbreak{}TT 0x\allowbreak{}VV record}, where \code{record} is the decoded fields of the record types above, the text of 0xDB and 0xFA records (CR, LF and other non-printable bytes escaped as \code{\textbackslash{}\allowbreak{}r}, \code{\textbackslash{}\allowbreak{}n} and \code{\textbackslash{}\allowbreak{}x\allowbreak{}NN}), or \code{hex} and the record bytes for other types:

\begin{Verbatim}
#      SN dive ridx type  ver  record
XFF  9166   -1    0 0xCA 0x20  event 0x004002  time 2022-07-01T22:46:29Z  value -0.0100
XFF  9166   -1    1 0xD1 0x10  counter 4  time 2022-07-01T22:46:12Z  count 1  time +62s  value 0
XFF  9211  -15    0 0xAE 0x11  serial 9211  config 0x01 0x02 0x01
XFF  9280    2    5 0xDB 0x00  command "ctdrw fetch\r|5"  reply "2024-03-15 22:57:12.000, 25.1567, 10.1491, 117.6462, 25.6469\r\n"
XFF  9304    0   35 0xFA 0x0A  text "\nr7    20016E60\r\nsp    20016E30\r\nexcrt FFFFFFFD\r\n..."
\end{Verbatim}

Open questions: the meaning of the 0xCA event codes and value, the 0xD1 \code{count} and last field, the 0xAE configuration bytes, the 0xEC records, and the \code{|\allowbreak{}} digit of the CTD commands.

\subsection{\texorpdfstring{C.7 WIMS configuration parameters (0xC5)}{C.7 WIMS configuration parameters (0xC5)}}\label{sec:c7-wims-configuration-parameters-0xc5}

{}[inferred from production] No format document describes this message. A WIMS float sends the configuration of its WIMS wave sensor in it:

\begin{itemize}
\item once in its diagnostic dive (dive 0), and
\item in the next dive after a \code{WCG -1} upload (WIMS Command Get, all parameters).
\end{itemize}

The production RAW files (2014–2026) hold 54, from 21 floats:

\begin{itemize}
\item dive 0 of 9246, 9248, 9249, 9265–9271, 9273, 9291–9293, 9296, 9297, 9299, 9303 and 9304;
\item after \code{WCG -1} uploads to 9239 and 9245 (2020), and 9292, 9296 and 9297 (2024).
\end{itemize}

9310 echoed a \code{WCG -1} upload in dive 201 (2026-09-24) but sent no 0xC5. Before build 3336 the decoder logged 0xC5 as an unknown X message.

\begingroup\footnotesize\setlength{\tabcolsep}{3pt}
\begin{longtable}{@{}>{\raggedleft\arraybackslash}p{0.0625\linewidth}>{\raggedleft\arraybackslash}p{0.0446\linewidth}>{\raggedright\arraybackslash}p{0.0714\linewidth}>{\raggedright\arraybackslash}p{0.1071\linewidth}>{\raggedright\arraybackslash}p{0.6429\linewidth}@{}}
\toprule
\textbf{Offset} & \textbf{Size} & \textbf{Type} & \textbf{Field} & \textbf{Notes} \\
\midrule
\endfirsthead
\toprule
\textbf{Offset} & \textbf{Size} & \textbf{Type} & \textbf{Field} & \textbf{Notes} \\
\midrule
\endhead
\bottomrule
\endlastfoot
0 & 1 & U1 & ID & 0xC5. \\
1 & 2 & U2 & \code{jj} & 12-bit size. \\
3 & 1 & U1 & Version & 1: 42 parameters (2020 floats, 90 bytes). 2: 43 parameters (2024 floats, 92 bytes). \\
4 & 1 & U1 & Part index & 0. One dive of 9239 had three copies with parts 0, 1 and 2, otherwise identical except parameters 36 and 37. \\
5 + 2i & 2 & U2 / I2 & Parameter i & One value per WIMS parameter, by the float's parameter index (the index of the \code{WCG} and \code{WCS} commands). \\
jj − 1 & 1 & char & \code{;\allowbreak{}} &  \\
\end{longtable}
\endgroup

The parameter index is the float's, not the WIMS user manual's: \code{WCS} uploads to 9296 and 9297 (\code{WCS 24 45;\allowbreak{} 28 10;\allowbreak{} 35 5;\allowbreak{} 39 17;\allowbreak{} 42 32}) appear at exactly those positions in the next 0xC5.

\textbf{Uploads.} \code{WCG <index>} (WIMS Command Get; 0-41, or -1 for all) and \code{WCV} (WIMS Command saVe) appear to relay the WIMS firmware's \code{get} and \code{save} commands; \code{WCS <index> <value>} relays \code{set}, which changes a parameter in RAM until \code{save} writes it to the sensor's EEPROM. The Deployment and Piloting Guide v1.0 lists all three as Reserved (Auxiliary Sensor Systems). Every upload in \code{/\allowbreak{}data1/\allowbreak{}alamo/\allowbreak{}data/\allowbreak{}sbd/\allowbreak{}outgoing} that uses them, but the first one to 9239, is zero or more \code{WCS} commands followed by \code{WCG -1;\allowbreak{} WCV;\allowbreak{}}:

\begin{itemize}
\item 2020-07 to 09, 9239 and 9245 (about 20 uploads): \code{wcs 22} 2 or 4, \code{wcs 30} 7 or 14, \code{wcs 34 180}, \code{wcs 36} 0, 65, 98, 131 or 164, \code{wcs 39} 9, 19 or 39, \code{wcs 40} 64, 128 or 256.
\item 2020-09-16, 9105, 9246-9249 and 9265-9273: \code{wcs 22 2;\allowbreak{} 30 14;\allowbreak{} 34 180;\allowbreak{} 36 131;\allowbreak{} 40 256} (9105 also \code{wcs 20 0}, then \code{wcs 31 6}).
\item 2024-09-26, 9292, 9296 and 9297: \code{WCS 24 45;\allowbreak{} 28 10;\allowbreak{} 35 5;\allowbreak{} 39 17;\allowbreak{} 42 32}, then \code{WCS 31 2}.
\item 2026-09-24, 9310: \code{WCS 39 17;\allowbreak{} 42 32}.
\end{itemize}

\begingroup\footnotesize\setlength{\tabcolsep}{3pt}
\begin{longtable}{@{}>{\raggedleft\arraybackslash}p{0.0536\linewidth}>{\raggedright\arraybackslash}p{0.2857\linewidth}>{\raggedright\arraybackslash}p{0.6161\linewidth}@{}}
\toprule
\textbf{Index} & \textbf{Name} & \textbf{Confirmed by} \\
\midrule
\endfirsthead
\toprule
\textbf{Index} & \textbf{Name} & \textbf{Confirmed by} \\
\midrule
\endhead
\bottomrule
\endlastfoot
0 & \code{brd.\allowbreak{}leds}? & always 0 \\
1 & unknown & always 0 \\
2–4 & \code{cal.\allowbreak{}acofx}, \code{cal.\allowbreak{}acofy}, \code{cal.\allowbreak{}acofz} & accelerometer offsets, e.g. -213, -352, -1193 (signed) \\
5–7 & \code{cal.\allowbreak{}acscx} … \code{cal.\allowbreak{}acscz} & accelerometer scales, 0 in every message (signed) \\
8–10 & \code{cal.\allowbreak{}gcofx} … \code{cal.\allowbreak{}gcofz} & gyroscope offsets (signed) \\
11–13 & \code{cal.\allowbreak{}gcscx} … \code{cal.\allowbreak{}gcscz} & gyroscope scales, 0 in every message (signed) \\
14–16 & \code{cal.\allowbreak{}mcofx} … \code{cal.\allowbreak{}mcofz} & magnetometer offsets (signed) \\
17–19 & \code{cal.\allowbreak{}mcscx}? … \code{cal.\allowbreak{}mcscz}? & per-float values such as 1300, 0, 181 (signed) \\
20 & unknown & 1 (v1), 0 (v2); \code{wcs 20 0} \\
21 & \code{imu.\allowbreak{}dlpf}? & 20 or 24 \\
22 & \code{imu.\allowbreak{}fsra} & 2 or 4 g; matches the \code{WVS} ASSF (8192 at 2 g, 2048 at 4 g); \code{wcs 22} 2 and 4 \\
23 & \code{imu.\allowbreak{}fsrg} & 500 °/s \\
24 & \code{imu.\allowbreak{}srate} & 40, 45 or 125 Hz; matches the \code{WVS} \code{fbase}, and \code{WCS 24 45} \\
25–26 & unknown & 0x0A62, 0xEE57 in every message \\
27 & unknown & per float and dive, e.g. 5, 8, 25, 57, 66 \\
28 & \code{wav.\allowbreak{}afilt}? & 2 or 10; \code{WCS 28 10} \\
29 & \code{wav.\allowbreak{}mxsam}? & 32768 \\
30 & \code{wav.\allowbreak{}mxseg}? & 7, 12 or 14; \code{wcs 30} 7 and 14 \\
31 & unknown & 0, 2 or 6; \code{wcs 31} 2 and 6 \\
32 & \code{wav.\allowbreak{}nfwax}? & 0 \\
33 & \code{wav.\allowbreak{}obeta}? & 1343, the manual's default \\
34 & \code{wav.\allowbreak{}ostsc}? & 90 or 180 s; \code{wcs 34 180} \\
35 & \code{wav.\allowbreak{}omode}? & 5; \code{WCS 35 5} \\
36 & \code{wav.\allowbreak{}ozeta}? & 65, 98, 131 or 164; changes between messages of the same float because it was tuned with \code{wcs 36} 0, 65, 98, 131 and 164 \\
37 & unknown & changes between messages of the same float, 1–16 \\
38 & \code{wav.\allowbreak{}smode}? & 2 \\
39 & \code{wav.\allowbreak{}srdiv} & 4, 9 or 17; \code{WVS} \code{fdiv} = value + 1 (doubled for every-other-bin firmware, C.3.8), and \code{wcs 39} 9, 17, 19 and 39 \\
40 & \code{wav.\allowbreak{}slen} & 64 or 256; matches the \code{WVS} \code{slen}; \code{wcs 40} 64, 128 and 256 \\
41 & \code{wav.\allowbreak{}twalp}? & 64, the value the ESSF (C.3.6) assumes \\
42 & \code{wav.\allowbreak{}gaina} & v2 only; 32 = 4.0, matching the \code{WVS} \code{again}, and \code{WCS 42 32} \\
\end{longtable}
\endgroup

Names ending in “?” are inferred from their values (defaults in the WIMS user manual, rev. 3.15) and position; the order is close to, but not the same as, the manual's alphabetical order. Values are listed signed for the calibration coefficients and unsigned otherwise.

\textbf{Listing.} Under the section banner “WIMS Configuration Parameters (0xC5)”, one line per parameter: \code{WIP SN dive midx ver part pidx name value hex}, with \code{midx} the message's index in the dive and \code{-} for an unknown name:

\begin{Verbatim}
#      SN dive midx  ver part pidx  name         value     hex
WIP  9297   15    0    2    0    2  cal.acofx     -213  0xFF2B
WIP  9297   15    0    2    0   24  imu.srate       45  0x002D
WIP  9297   15    0    2    0   25  -             2658  0x0A62
WIP  9297   15    0    2    0   39  wav.srdiv       17  0x0011
WIP  9297   15    0    2    0   42  wav.gaina       32  0x0020
\end{Verbatim}

Open questions: parameters 1, 20, 25–27, 31 and 37, the inferred names, and why 9310 sent no reply to \code{WCG -1}.

\subsection{\texorpdfstring{Appendix C.A: Offset maps of every MFA layout (decoder)}{Appendix C.A: Offset maps of every MFA layout (decoder)}}\label{sec:appendix-ca-offset-maps-of-every-mfa-layout-decoder}

\code{NAME@offset} gives the decoder's read offset, counted from the ID byte. Fields are U2 unless a suffix says otherwise (\code{:\allowbreak{}U1}, \code{:\allowbreak{}U4}, \code{:\allowbreak{}I2}, \code{:\allowbreak{}R4}). \code{T\_\allowbreak{}DESCENT}, \code{T\_\allowbreak{}PARK}, \code{T\_\allowbreak{}ASCENT} and \code{T\_\allowbreak{}SURFACE} are U4 state start times in seconds since 2000-01-01. Gaps between fields are bytes that exist but aren't decoded (see C.1.9). In layouts from 4.2.62 on, offset 8 is the parameters version and offset 9 the CTD type (shown as \code{(\allowbreak{}params ver.\allowbreak{})} and \code{(\allowbreak{}CTD type)}).

\begin{itemize}
\item \textbf{Fw02006200 D0}: MSN@8, DTD@10, MDS@12, APT@14:U1, XAS@15:U1, CPS@16:U1, CST@17:U1, CSP@18:U4, CPB1@22, CPB2@24, CPB3@26, CPN1@28, CPN2@30, CPN3@32, CPP1@34, CPP2@36, CPP3@38, MTD@40, MTT@42:U4, MAM@46:U1, MAT@47, ACA@49, ACG@51, ACR@53, PNB@55, PNC@57, PXC@59, PBL@61, PBH@63, PP0@65, PP1@67, PVL@69, PVH@71, PV0@73, PV5@75, PC0@77, PC1@79, T\_\allowbreak{}DESCENT@99:U4, DNS@103, DRA@105, DSD@107, DST@109, DDF@111, DDC@113, DSI@115:U4, DDD@119, DDT@121, DFT@123, DCM@125:U1, DAL@126:U1, DCT@127, DCA@129:U1, DRC@130, DRT@132, DCX@134, DCP@136, DCF@138, DCD@140, DCR@142, DSW@144, DUC@146, DPT@148, DUT@150, DPD@152, DVT@154, PFR@156, PFC@158, DDP@160, DDS@162, T\_\allowbreak{}PARK@179:U4, PNS@183, PTO@185:U4, PTH@189:U4, PSA@193:U1, PSI@194:U4, PFM@198:U1, PFT@199, PFD@201, PFV@203, PDT@205, PSM@207:U1, PSD@208, PST@210, PSV@212, PSW@214, T\_\allowbreak{}ASCENT@227:U4, ASN@231, ARR@233, APM@235, APN@237, APF@239, ASI@241:U4, ASD@245, ADD@247, ADT@249, ARC@251, AHR@253, AHD@255, T\_\allowbreak{}SURFACE@271:U4, SFN@275, SDC@277:I2, SGT@279, SET@281, SBM@283:U1, SBT@284, SDT@286:U4, SSA@290, SSI@292:U4, SDS@296, SAM@298:U1, SAX@299, SAN@301, SAO@303, SAP@305, SPT@307
\item \textbf{Fw02006900 D0}: MSN@8, DTD@10, MDS@12, APT@14:U1, XAS@15:U1, CPS@16:U1, CST@17:U1, CTF@18, CTL@20, CTU@22, CSP@24:U4, CPB1@28, CPB2@30, CPB3@32, CPN1@34, CPN2@36, CPN3@38, CPP1@40, CPP2@42, CPP3@44, MTD@46, MTT@48:U4, MAM@52:U1, MAT@53, ACA@55, ACG@57, ACR@59, PNB@61, PNC@63, PXC@65, PBL@67, PBH@69, PP0@71, PP1@73, PVL@75, PVH@77, PV0@79, PV5@81, PC0@83, PC1@85, ICM@87:U1, IMM@88, ICA0@90, ICA1@92, ICB0@94, ICB1@96, IMT0@98, IMT1@100, ILT0@102, ILT1@104, ICN@106, ICX@108, ICT@110:I2, ICS@112, ISS0@114:U1, ISS1@115:U1, IST@116, ISO@118, ISP@120:U1, IUE0@121:U1, IUE1@122:U1, ICD0@123, ICD1@125, ICR0@127:U4, ICR1@131:U4, ICG0@135:U1, ICG1@136:U1, ICO0@137, ICO1@139, ICF@141:U1, ICW@145:U1, T\_\allowbreak{}DESCENT@149:U4, DNS@153, DRA@155, DSD@157, DST@159, DDF@161, DDC@163, DSI@165:U4, DDD@169, DDT@171, DFT@173, DCM@175:U1, DAL@176:U1, DCT@177, DCA@179:U1, DRC@180, DRT@182, DCX@184, DCP@186, DCF@188, DCD@190, DCR@192, DSW@194, DUC@196, DPT@198, DUT@200, DPD@202, DVT@204, PFR@206, PFC@208, DDP@210, DDS@212, T\_\allowbreak{}PARK@214:U4, PNS@218, PTO@220:U4, PTH@224:U4, PSA@228:U1, PSI@229:U4, PFM@233:U1, PFT@234, PFD@236, PFV@238, PDT@240, PSM@242:U1, PSD@243, PST@245, PSV@247, PSW@249, T\_\allowbreak{}ASCENT@251:U4, ASN@255, ARR@257, APM@259, APN@261, APF@263, ASI@265:U4, ASD@269, ADD@271, ADT@273, ARC@275, AHR@277, AHD@279, T\_\allowbreak{}SURFACE@281:U4, SFN@285, SDC@287:I2, SGT@289, SET@291, SBM@293:U1, SBT@294, SDT@296:U4, SSA@300, SSI@302:U4, SDS@306, SAM@308:U1, SAX@309, SAN@311, SAO@313, SAP@315, SPT@317
\item \textbf{Fw03006500 D0}: MSN@8, DTD@10, MDS@12, APT@14:U1, XAS@15:U1, CPS@16:U1, CST@17:U1, CTF@18, CTL@20, CTU@22, CSP@24:U4, CPB1@28, CPB2@30, CPB3@32, CPN1@34, CPN2@36, CPN3@38, CPP1@40, CPP2@42, CPP3@44, MTD@46, MTT@48:U4, MAM@52:U1, MAT@53, ACA@55, ACG@57, ACR@59, PNB@61, PNC@63, PXC@65, PBL@67, PBH@69, PP0@71, PP1@73, PVL@75, PVH@77, PV0@79, PV5@81, PC0@83, PC1@85, ICM@87:U1, IMM@88, ICA0@90, ICA1@92, ICB0@94, ICB1@96, IMT0@98, IMT1@100, ILT0@102, ILT1@104, ICN@106, ICX@108, ICT@110:I2, ICS@112, ISS0@114:U1, ISS1@115:U1, IST@116, ISO@118, ISP@120:U1, IUE0@121:U1, IUE1@122:U1, ICD0@123, ICD1@125, ICR0@127:U4, ICR1@131:U4, ICG0@135:U1, ICG1@136:U1, ICO0@146, ICO1@148, ICF@150:U1, ICW@154:U1, T\_\allowbreak{}DESCENT@158:U4, DNS@162, DRA@164, DSD@166, DST@168, DDF@170, DDC@172, DSI@174:U4, DDD@178, DDT@180, DFT@182, DCM@184:U1, DAL@185:U1, DCT@186, DCA@188:U1, DRC@189, DRT@191, DCX@193, DCP@195, DCF@197, DCD@199, DCR@201, DSW@203, DUC@205, DPT@207, DUT@209, DPD@211, DVT@213, PFR@215, PFC@217, DDP@219, DDS@221, T\_\allowbreak{}PARK@223:U4, PNS@227, PTO@229:U4, PTH@233:U4, PSA@237:U1, PSI@238:U4, PFM@242:U1, PFT@243, PFD@245, PFV@247, PDT@249, PSM@251:U1, PSD@252, PST@254, PSV@256, PSW@258, T\_\allowbreak{}ASCENT@260:U4, ASN@264, ARR@266, APM@268, APN@270, APF@272, ASI@274:U4, ASD@278, ADD@280, ADT@282, ARC@284, AHR@286, AHD@288, T\_\allowbreak{}SURFACE@290:U4, SFN@294, SDC@296:I2, SGT@298, SET@300, SBM@302:U1, SBT@303, SDT@305:U4, SSA@309, SSI@311:U4, SDS@315, SAM@317:U1, SAX@318, SAN@320, SAO@322, SAP@324
\item \textbf{Fw03016700 D0}: MSN@8, DTD@10, MDS@12, APT@14:U1, XAS@15:U1, CPS@16:U1, CST@17:U1, CTF@18, CTL@20, CTU@22, CSP@24:U4, CPB1@28, CPB2@30, CPB3@32, CPN1@34, CPN2@36, CPN3@38, CPP1@40, CPP2@42, CPP3@44, MTD@46, MTT@48:U4, MAM@52:U1, MAT@53, ACA@55, ACG@57, ACR@59, PNB@61, PNC@63, PXC@65, PBL@67, PBH@69, PP0@71, PP1@73, PVL@75, PVH@77, PV0@79, PV5@81, PC0@83, PC1@85, ICM@87:U1, IMM@88, ICA0@90, ICA1@92, ICB0@94, ICB1@96, IMT0@98, IMT1@100, ILT0@102, ILT1@104, ICN@106, ICX@108, ICT@110:I2, ICS@112, ISS0@114:U1, ISS1@115:U1, IST@117, ISO@119, ISP@121:U1, IUE0@122:U1, IUE1@123:U1, ICD0@124, ICD1@126, ICR0@128:U4, ICR1@132:U4, ICG0@136:U1, ICG1@137:U1, ICO0@147, ICO1@149, ICF@151:U1, ICW@155:U1, T\_\allowbreak{}DESCENT@159:U4, DNS@163, DRA@165, DSD@167, DST@169, DDF@171, DDC@173, DSI@175:U4, DDD@179, DDT@181, DFT@183, DCM@185:U1, DAL@186:U1, DCT@187, DCA@189:U1, DRC@190, DRT@192, DCX@194, DCP@196, DCF@198, DCD@200, DCR@202, DSW@204, DUC@206, DPT@208, DUT@210, DPD@212, DVT@214, PFR@216, PFC@218, DDP@220, DDS@222, T\_\allowbreak{}PARK@224:U4, PNS@228, PTO@230:U4, PTH@234:U4, PSA@238:U1, PSI@239:U4, PFM@243:U1, PFT@244, PFD@246, PFV@248, PDT@250, PSM@252:U1, PSD@253, PST@255, PSV@257, PSW@259, T\_\allowbreak{}ASCENT@261:U4, ASN@265, ARR@267, APM@269, APN@271, APF@273, ASI@275:U4, ASD@279, ADD@281, ADT@283, ARC@285, AHR@287, AHD@289, T\_\allowbreak{}SURFACE@291:U4, SFN@295, SDC@297:I2, SGT@299, SET@301, SBM@303:U1, SBT@304, SDT@306:U4, SSA@310, SSI@312:U4, SDS@316, SAM@318:U1, SAX@319, SAN@321, SAO@323, SAP@325
\item \textbf{Fw03026100 D0}: MSN@8, DTD@10, MDS@12, APT@14:U1, BTM@15:U1, XAS@16:U1, CPS@17:U1, CST@18:U1, CTF@19, CTL@21, CTU@23, CSP@25:U4, CPB1@29, CPB2@31, CPB3@33, CPN1@35, CPN2@37, CPN3@39, CPP1@41, CPP2@43, CPP3@45, MTD@47, MTT@49:U4, MAM@53:U1, MAT@54, ACA@56, ACG@58, ACR@60, PNB@62, PNC@64, PXC@66, PBL@68, PBH@70, PP0@72, PP1@74, PVL@76, PVH@78, PV0@80, PV5@82, PC0@84, PC1@86, ICM@88:U1, IMM@89, ICA0@91, ICA1@93, ICB0@95, ICB1@97, IMT0@99, IMT1@101, ILT0@103, ILT1@105, ICN@107, ICX@109, ICT@111:I2, ICS@113, ISS0@115:U1, ISS1@116:U1, IST@118, ISO@120, ISP@122:U1, IUE0@123:U1, IUE1@124:U1, ICD0@125, ICD1@127, ICR0@129:U4, ICR1@133:U4, ICG0@137:U1, ICG1@138:U1, ICO0@148, ICO1@150, ICF@152:U1, ICW@156:U1, T\_\allowbreak{}DESCENT@160:U4, DNS@164, DRA@166, DSD@168, DST@170, DDF@172, DDC@174, DSI@176:U4, DDD@180, DDT@182, DFT@184, DCM@186:U1, DAL@187:U1, DCT@188, DCA@190:U1, DRC@191, DRT@193, DCX@195, DCP@197, DCF@199, DCD@201, DCR@203, DSW@205, DUC@207, DPT@209, DUT@211, DPD@213, DVT@215, PFR@217, PFC@219, DDP@221, DDS@223, T\_\allowbreak{}PARK@225:U4, PNS@229, PTO@231:U4, PTH@235:U4, PSA@239:U1, PSI@240:U4, PFM@244:U1, PFT@245, PFD@247, PFV@249, PDT@251, PSM@253:U1, PSD@254, PST@256, PSV@258, PSW@260, T\_\allowbreak{}ASCENT@262:U4, ASN@266, ARR@268, APM@270, APN@272, APF@274, ASI@276:U4, ASD@280, ADD@282, ADT@284, ARC@286, AHR@288, AHD@290, T\_\allowbreak{}SURFACE@292:U4, SFN@296, SDC@298:I2, SGT@300, SET@302, SBM@304:U1, SBT@305, SDT@307:U4, SSA@311, SSI@313:U4, SDS@317, SAM@319:U1, SAX@320, SAN@322, SAO@324
\item \textbf{Fw04016100 D0}: MSN@8, DTD@10, MDS@12, APT@14:U1, BTM@15:U1, XAS@16:U1, CPS@17:U1, CST@18:U1, CTF@19, CTL@21, CTU@23, CSP@25:U4, CPB1@29, CPB2@31, CPB3@33, CPN1@35, CPN2@37, CPN3@39, CPP1@41, CPP2@43, CPP3@45, MTD@47, MTT@49:U4, MAM@53:U1, MAT@54, ACA@56, ACG@58, ACR@60, PNB@62, PNC@64, PXC@66, PBL@68, PBH@70, PP0@72, PP1@74, PVL@76, PVH@78, PV0@80, PV5@82, PC0@84, PC1@86, ICM@88:U1, IMM@89, ICA0@91, ICA1@93, ICB0@95, ICB1@97, IMT0@99, IMT1@101, ILT0@103, ILT1@105, ICN@107, ICX@109, ICT@111:I2, ICS@113, ISS0@115:U1, ISS1@116:U1, IST@118, ISO@120, ISP@122:U1, IUE0@123:U1, IUE1@124:U1, ICD0@125, ICD1@127, ICR0@129:U4, ICR1@133:U4, ICG0@137:U1, ICG1@138:U1, ICO0@148, ICO1@150, ICF@152:U1, ICW@156:U1, T\_\allowbreak{}DESCENT@160:U4, DNS@164, DRA@166, DSD@168, DST@170, DDF@172, DDC@174, DSI@176:U4, DDD@180, DDT@182, DFT@184, DCM@186:U1, DAL@187:U1, DCT@188, DCA@190:U1, DRC@191, DRT@193, DCX@195, DCP@197, DCF@199, DCD@201, DCR@203, DSW@205, DUC@207, DPT@209, DUT@211, DPD@213, DVT@215, PFR@217, PFC@219, DDP@221, DDS@223, T\_\allowbreak{}PARK@225:U4, PNS@229, PTO@231:U4, PTH@235:U4, PSA@239:U1, PSI@240:U4, PFM@244:U1, PFT@245, PFD@247, PFV@249, PDT@251, PSM@253:U1, PSD@254, PST@256, PSV@258, PSW@260, T\_\allowbreak{}ASCENT@262:U4, ASN@266, ARR@268, APM@270, APN@272, APF@274, ASI@276:U4, ASD@280, ADD@282, ADT@284, ARC@286, AHR@288, AHD@290, T\_\allowbreak{}SURFACE@292:U4, SFN@296, SDC@298:I2, SGT@300, SET@302, SBM@304:U1, SBT@305, SDT@307:U4, SSA@311, SSI@313:U4, SDS@317, SAM@319:U1, SAX@320, SAN@322, SAO@324
\item \textbf{Fw04026200 D0}: (params ver.)@8:U1, (CTD type)@9:U1, MSN@10, DTD@12, MDS@14, APT@16:U1, BTM@17:U1, XAS@18:U1, CPS@19:U1, CST@20:U1, CTF@21, CTL@23, CTU@25, CSP@27:U4, CPB1@31, CPB2@33, CPB3@35, CPN1@37, CPN2@39, CPN3@41, CPP1@43, CPP2@45, CPP3@47, MTD@49, MTT@51:U4, MAM@55:U1, MAT@56, ACA@58, ACG@60, ACR@62, PNB@64, PNC@66, PXC@68, PBL@70, PBH@72, PP0@74, PP1@76, PVL@78, PVH@80, PV0@82, PV5@84, PC0@86, PC1@88, ICM@90:U1, IMM@91, ICA0@93, ICA1@95, ICB0@97, ICB1@99, IMT0@101, IMT1@103, ILT0@105, ILT1@107, ICN@109, ICX@111, ICT@113:I2, ICS@115, ISS0@117:U1, ISS1@118:U1, IST@120, ISO@122, ISP@124:U1, IUE0@125:U1, IUE1@126:U1, ICD0@127, ICD1@129, ICR0@131:U4, ICR1@135:U4, ICG0@139:U1, ICG1@140:U1, ISA@141:U1, ISD@142, IAT@144, ISV@146, ISW@148, ICO0@150, ICO1@152, ICF@154:U4, ICW@158:U4, DNS@162, DRA@164, DSD@166, DST@168, DDF@170, DDC@172, DSI@174:U4, DDD@178, DDT@180, DFT@182, DCM@184:U1, DAL@185:U1, DCT@186, DCA@188:U1, DRC@189, DRT@191, DCX@193, DCP@195, DCF@197, DCD@199, DCR@201, DSW@203, DUC@205, DPT@207, DUT@209, DPD@211, DVT@213, DVD@215, DVW@217, PFR@219, PFC@221, DDP@223, DDS@225, PNS@227, PTO@229:U4, PTH@233:U4, PSA@237:U1, PSI@238:U4, PFM@242:U1, PFT@243, PFD@245, PFV@247, PDT@249, PSM@251:U1, PSD@252, PST@254, PSV@256, PSW@258, ASN@260, ARR@262, APM@264, APN@266, APF@268, ASI@270:U4, ASD@274, ADD@276, ADT@278, ARC@280, AHR@282, AHD@284, SFN@286, SDC@288:I2, SGT@290, SET@292, SBM@294:U1, SBT@295, SDT@297:U4, SSA@301, SSI@303:U4, SDS@307, SAM@309:U1, SAX@310, SAN@312, SAO@314, SAP@316, SPT@318, SSP@320
\item \textbf{Fw04036100 D0}: (params ver.)@8:U1, (CTD type)@9:U1, MSN@10, DTD@12, MDS@14, APT@16:U1, BTM@17:U1, XAS@18:U1, CPS@19:U1, CST@20:U1, CTF@21, CTL@23, CTU@25, CSP@27:U4, CPB1@31, CPB2@33, CPB3@35, CPN1@37, CPN2@39, CPN3@41, CPP1@43, CPP2@45, CPP3@47, MTD@53, MTT@55:U4, MAM@59:U1, MAT@60, ACA@62, ACG@64, ACR@66, PNB@68, PNC@70, PXC@72, PBL@74, PBH@76, PP0@78, PP1@80, PVL@82, PVH@84, PV0@86, PV5@88, PC0@90, PC1@92, ICM@94:U1, IMM@95, ICA0@97, ICA1@99, ICB0@101, ICB1@103, IMT0@105, IMT1@107, ILT0@109, ILT1@111, ICN@113, ICX@115, ICT@117:I2, ICS@119, ISS0@121:U1, ISS1@122:U1, IST@124, ISO@126, ISP@128:U1, IUE0@129:U1, IUE1@130:U1, ICD0@131, ICD1@133, ICR0@135:U4, ICR1@139:U4, ICG0@143:U1, ICG1@144:U1, ISA@145:U1, ISD@146, IAT@148, ISV@150, ISW@152, ICO0@154, ICO1@156, ICF@158:U4, ICW@162:U4, DNS@166, DRA@168, DSD@170, DST@172, DDF@174, DDC@176, DSI@178:U4, DDD@182, DDT@184, DFT@186, DCM@188:U1, DAL@189:U1, DCT@190, DCA@192:U1, DRC@193, DRT@195, DCX@197, DCP@199, DCF@201, DCD@203, DCR@205, DSW@207, DUC@209, DPT@211, DUT@213, DPD@215, DVT@217, DVD@219, DVW@221, PFR@223, PFC@225, DDP@227, DDS@229, PNS@231, PTO@233:U4, PTH@237:U4, PSA@241:U1, PSI@242:U4, PFM@246:U1, PFT@247, PFD@249, PFV@251, PDT@253, PSM@255:U1, PSD@256, PST@258, PSV@260, PSW@262, ASN@264, ARR@266, APM@268, APN@270, APF@272, ASI@274:U4, ASD@278, ADD@280, ADT@282, ARC@284, AHR@286, AHD@288, SFN@290, SDC@292:I2, SGT@294, SET@296, SBM@298:U1, SBT@299, SDT@301:U4, SSA@305, SSI@307:U4, SDS@311, SAM@313:U1, SAX@314, SAN@316, SAO@318, SAP@320, SPT@322, SSP@324
\item \textbf{Fw04046100 D0}: (params ver.)@8:U1, (CTD type)@9:U1, MSN@10, DTD@12, MDS@14, APT@16:U1, BTM@17:U1, XAS@18:U1, CPS@19:U1, CST@20:U1, CTF@21, CTL@23, CTU@25, CSP@28:U4, CPB1@32, CPB2@34, CPB3@36, CPN1@38, CPN2@40, CPN3@42, CPP1@44, CPP2@46, CPP3@48, MTD@54, MTT@56:U4, MAM@60:U1, MAT@61, ACA@63, ACG@65, ACR@67, PNB@69, PNC@71, PXC@73, PBL@75, PBH@77, PP0@79, PP1@81, PVL@83, PVH@85, PV0@87, PV5@89, PC0@91, PC1@93, ICM@95:U1, IMM@96, ICA0@98, ICA1@100, ICB0@102, ICB1@104, IMT0@106, IMT1@108, ILT0@110, ILT1@112, ICN@114, ICX@116, ICT@118:I2, ICS@120, ISS0@122:U1, ISS1@123:U1, IST@125, ISO@127, ISP@129:U1, IUE0@130:U1, IUE1@131:U1, ICD0@132, ICD1@134, ICR0@136:U4, ICR1@140:U4, ICG0@144:U1, ICG1@145:U1, ISA@146:U1, ISD@147, IAT@149, ISV@151, ISW@153, ICO0@155, ICO1@157, ICF@159:U4, ICW@163:U4, DNS@167, DRA@169, DSD@171, DST@173, DDF@175, DDC@177, DSI@179:U4, DDD@183, DDT@185, DFT@187, DCM@189:U1, DAL@190:U1, DCT@191, DCA@193:U1, DRC@194, DRT@196, DCX@198, DCP@200, DCF@202, DCD@204, DCR@206, DSW@208, DUC@210, DPT@212, DUT@214, DPD@216, DVT@218, DVD@220, DVW@222, PFR@224, PFC@226, DDP@228, DDS@230, PNS@232, PTO@234:U4, PTH@238:U4, PSA@242:U1, PSI@243:U4, PFM@247:U1, PFT@248, PFD@250, PFV@252, PDT@254, PSM@256:U1, PSD@257, PST@259, PSV@261, PSW@263, ASN@265, ARR@267, APM@269, APN@271, APF@273, ASI@275:U4, ASD@279, ADD@281, ADT@283, ARC@285, AHR@287, AHD@289, SFN@291, SDC@293:I2, SGT@295, SET@297, SBM@299:U1, SBT@300, SDT@302:U4, SSA@306, SSI@308:U4, SDS@312, SAM@314:U1, SAX@315, SAN@317, SAO@319, SAP@321, SPT@323, SSP@325
\item \textbf{Fw04056300 D0}: (params ver.)@8:U1, (CTD type)@9:U1, MSN@10, DTD@12, MDS@14, APT@16:U1, BTM@17:U1, XAS@18:U1, CPS@19:U1, CST@20:U1, CTF@21, CTL@23, CTU@25, CTC@27:U1, CSP@28:U4, CPB1@32, CPB2@34, CPB3@36, CPN1@38, CPN2@40, CPN3@42, CPP1@44, CPP2@46, CPP3@48, CRD@50, CRS@52, MTD@54, MTT@56:U4, MAM@60:U1, ACA@67, ACG@69, ACR@71, PNB@73, PNC@75, PXC@77, PBL@79, PBH@81, PP0@83, PP1@85, PVL@87, PVH@89, PV0@91, PV5@93, PC0@95, PC1@97, ICM@99:U1, IMM@100, ICA0@102, ICA1@104, ICB0@106, ICB1@108, IMT0@110, IMT1@112, ILT0@114, ILT1@116, ICN@118, ICX@120, ICT@122:I2, ICS@124, ISS0@126:U1, ISS1@127:U1, IST@129, ISO@131, ISP@133:U1, IUE0@134:U1, IUE1@135:U1, ICD0@136, ICD1@138, ICR0@140:U4, ICR1@144:U4, ICG0@148:U1, ICG1@149:U1, ISA@150:U1, ISD@151, IAT@153, ISV@155, ISW@157, ICO0@159, ICO1@161, ICF@163:U4, ICW@167:U4, DNS@171, DRA@173, DSD@175, DST@177, DDF@179, DDC@181, DSI@183:U4, DDD@187, DDT@189, DFT@191, DCM@193:U1, DAL@194:U1, DCT@195, DCA@197:U1, DRC@198, DRT@200, DCX@202, DCP@204, DCF@206, DCD@208, DCR@210, DSW@212, DUC@214, DPT@216, DUT@218, DPD@220, DVT@222, DVD@224, DVW@226, PFR@228, PFC@230, DDP@232, DDS@234, PNS@236, PTO@238:U4, PTH@242:U4, PSA@250:U1, PSI@251:U4, PFM@255:U1, PFT@256, PFD@258, PFV@260, PDT@262, PSM@264:U1, PSD@265, PST@267, PSV@269, PSW@271, ASN@273, ARR@275, APM@277, APN@279, APF@281, ASI@283:U4, ASD@287, ADD@289, ADT@291, ARC@293, AHR@295, AHD@297, SFN@299, SDC@301:I2, SGT@303, SET@305, SBM@307:U1, SBT@308, SDT@310:U4, SSA@314, SSI@316:U4, SDS@320, SAM@322:U1, SAX@323, SAN@325
\item \textbf{Fw04056300 D1}: SAO@3, SAP@5, SPT@7, SSP@9
\item \textbf{Fw04066100 D0}: (params ver.)@8:U1, (CTD type)@9:U1, MSN@10, DTD@12, MDS@14, APT@16:U1, BTM@17:U1, XAS@18:U1, CPS@19:U1, CST@20:U1, CTF@21, CTL@23, CTU@25, CSP@28:U4, CPB1@32, CPB2@34, CPB3@36, CPN1@38, CPN2@40, CPN3@42, CPP1@44, CPP2@46, CPP3@48, MTD@54, MTT@56:U4, MAM@60:U1, ACA@67, ACG@69, ACR@71, PNB@73, PNC@75, PXC@77, PBL@79, PBH@81, PP0@83, PP1@85, PVL@87, PVH@89, PV0@91, PV5@93, PC0@95, PC1@97, ICM@99:U1, IMM@100, ICA0@102, ICA1@104, ICB0@106, ICB1@108, IMT0@110, IMT1@112, ILT0@114, ILT1@116, ICN@118, ICX@120, ICT@122:I2, ICS@124, ISS0@126:U1, ISS1@127:U1, IST@129, ISO@131, ISP@133:U1, IUE0@134:U1, IUE1@135:U1, ICD0@136, ICD1@138, ICR0@140:U4, ICR1@144:U4, ICG0@148:U1, ICG1@149:U1, ISA@150:U1, ISD@151, IAT@153, ISV@155, ISW@157, ICO0@159, ICO1@161, ICF@163:U4, ICW@167:U4, DNS@171, DRA@173, DSD@175, DST@177, DDF@179, DDC@181, DSI@183:U4, DDD@187, DDT@189, DFT@191, DCM@193:U1, DAL@194:U1, DCT@195, DCA@197:U1, DRC@198, DRT@200, DCX@202, DCP@204, DCF@206, DCD@208, DCR@210, DSW@212, DUC@214, DPT@216, DUT@218, DPD@220, DVT@222, DVD@224, DVW@226, PFR@228, PFC@230, DDP@232, DDS@234, PNS@236, PTO@238:U4, PTH@242:U4, PSA@250:U1, PSI@251:U4, PFM@255:U1, PFT@256, PFD@258, PFV@260, PDT@262, PSM@264:U1, PSD@265, PST@267, PSV@269, PSW@271, ASN@273, ARR@275, APM@277, APN@279, APF@281, ASI@283:U4, ASD@287, ADD@289, ADT@291, ARC@293, AHR@295, AHD@297, SFN@299, SDC@301:I2, SGT@303, SET@305, SBM@307:U1, SBT@308, SDT@310:U4, SSA@314, SSI@316:U4, SDS@320, SAM@322:U1, SAX@323, SAN@325
\item \textbf{Fw05006100 D0}: (params ver.)@8:U1, (CTD type)@9:U1, MSN@10, MPS@12:U1, MPC@13, MNM0@15, MNM1@17, MNM2@19, MNS0@21:U1, MNS1@22:U1, MNS2@23:U1, DTD0@24, DTD1@26, DTD2@28, MDS@30, APT@32:U1, BTM@33:U1, DDB@34:U1, XAS@35:U1, CPS@36:U1, CST@37:U1, CTF@38, CTL@40, CTU@42, CTC@44:U1, CSP@45:U4, CPB1@49, CPB2@51, CPB3@53, CPN1@55, CPN2@57, CPN3@59, CPP1@61, CPP2@63, CPP3@65, CRD@67, CRS@69, MTD@71, MTT@73:U4, MTS@77:U4, MAM@81:U1, MAT1@82, MAT2@84, MAP@86, ACA@88, ACG@90, ACR@92, PNB@94, PNC@96, PXC@98, PBL@100, PBH@102, PP0@104, PP1@106, PVL@108, PVH@110, PV0@112, PV5@114, PC0@116, PC1@118, ICM@120:U1, IMM@121, ICA0@123, ICA1@125, ICB0@127, ICB1@129, IMT0@131, IMT1@133, ILT0@135, ILT1@137, ICN@139, ICX@141, ICT@143:I2, ICS@145, ISS0@147:U1, ISS1@148:U1, ISM@149:U1, IST@150, ISO@152, ISP@154:U1, IUE0@155:U1, IUE1@156:U1, ICD0@157, ICD1@159, ICR0@161:U4, ICR1@165:U4, ICG0@169:U1, ICG1@170:U1, ISA@171:U1, ISD@172, IAT@174, ISV@176, ISW@178, ICO0@180, ICO1@182, ICF@184:U4, ICW@188:U4, DNS@192, DRA@194, DSD0@196, DSD1@198, DSD2@200, DST0@202, DST1@204, DST2@206, DDF@208, DDC@210, DSI@212:U4, DDD@216, DDT@218, DFT0@220, DFT1@222, DFT2@224, DCM@226:U1, DAL@227:U1, DCT0@228, DCT1@230, DCT2@232, DCA@234:U1, DRC@235, DRT@237, DCX@239, DCP@241, DCF@243, DCD@245, DCR@247, DSW@249, DUC@251, DPT@253, DUT@255, DPD@257, DVT@259, DVD@261, DVW@263, PFR@265, PFC@267, DDP@269, DDS@271, PNS@273, PTO0@275:U4, PTO1@279:U4, PTO2@283:U4, PTH0@287:U4, PTH1@291:U4, PTH2@295:U4, PTD@299:U4, PSA@303:U1, PSI@304:U4, PFM0@308:U1, PFM1@309:U1, PFM2@310:U1, PFT0@311, PFT1@313, PFT2@315, PFD0@317, PFD1@319, PFD2@321, PFV0@323, PFV1@325
\item \textbf{Fw05006100 D1}: PFV2@3, PDT@5, PSM@7:U1, PSD@8, PST@10, PSV@12, PSW@14, ASN@16, ARR@18, APM@20, APN@22, APF@24, ASI@26:U4, ASD@30, ADD@32, ADT@34, ARC@36, AHR@38, AHD@40, SFN@42, SDC@44:I2, SGT@46, SET@48, SBM@50:U1, SBT@51, SDT@53:U4, SSA@57, SSI@59:U4, SDS@63, SAM@65:U1, SAX@66, SAN@68, SAO@70, SAP@72, SPT@74, SSP@76
\item \textbf{Fw05006101 D0}: (params ver.)@8:U1, (CTD type)@9:U1, MSN@10, MPS@12:U1, MPC@13, MNM0@15, MNM1@17, MNM2@19, MNS0@21:U1, MNS1@22:U1, MNS2@23:U1, DTD0@24, DTD1@26, DTD2@28, MDS@30, APT@32:U1, BTM@33:U1, XAS@34:U1, CPS@35:U1, CST@36:U1, CTF@37, CTL@39, CTU@41, CTC@43:U1, CSP@44:U4, CPB1@48, CPB2@50, CPB3@52, CPN1@54, CPN2@56, CPN3@58, CPP1@60, CPP2@62, CPP3@64, CRD@66, CRS@68, MTD@70, MTT@72:U4, MTS@76:U4, MAM@80:U1, MAT2@81, ACA@83, ACG@85, ACR@87, PNB@89, PNC@91, PXC@93, PBL@95, PBH@97, PP0@99, PP1@101, PVL@103, PVH@105, PV0@107, PV5@109, PC0@111, PC1@113, ICM@115:U1, IMM@116, ICA0@118, ICA1@120, ICB0@122, ICB1@124, IMT0@126, IMT1@128, ILT0@130, ILT1@132, ICN@134, ICX@136, ICT@138:I2, ICS@140, ISS0@142:U1, ISS1@143:U1, ISM@144:U1, IST@145, ISO@147, ISP@149:U1, IUE0@150:U1, IUE1@151:U1, ICD0@152, ICD1@154, ICR0@156:U4, ICR1@160:U4, ICG0@164:U1, ICG1@165:U1, ISA@166:U1, ISD@167, IAT@169, ISV@171, ISW@173, ICO0@175, ICO1@177, ICF@179:U4, ICW@183:U4, DNS@187, DRA@189, DSD0@191, DSD1@193, DSD2@195, DST0@197, DST1@199, DST2@201, DDF@203, DDC@205, DSI@207:U4, DDD@211, DDT@213, DFT0@215, DFT1@217, DFT2@219, DCM@221:U1, DAL@222:U1, DCT0@223, DCT1@225, DCT2@227, DCA@229:U1, DRC@230, DRT@232, DCX@234, DCP@236, DCF@238, DCD@240, DCR@242, DSW@244, DUC@246, DPT@248, DUT@250, DPD@252, DVT@254, DVD@256, DVW@258, PFR@260, PFC@262, DDP@264, DDS@266, PNS@268, PTO0@270:U4, PTO1@274:U4, PTO2@278:U4, PTH0@282:U4, PTH1@286:U4, PTH2@290:U4, PTD@294:U4, PSA@298:U1, PSI@299:U4, PFM0@303:U1, PFM1@304:U1, PFM2@305:U1, PFT0@306, PFT1@308, PFT2@310, PFD0@312, PFD1@314, PFD2@316, PFV0@318, PFV1@320, PFV2@322, PDT@324, PSM@326:U1
\item \textbf{Fw05006101 D1}: PSD@3, PST@5, PSV@7, PSW@9, ASN@11, ARR@13, APM@15, APN@17, APF@19, ASI@21:U4, ASD@25, ADD@27, ADT@29, ARC@31, AHR@33, AHD@35, SFN@37, SDC@39:I2, SGT@41, SET@43, SBM@45:U1, SBT@46, SDT@48:U4, SSA@52, SSI@54:U4, SDS@58, SAM@60:U1, SAX@61, SAN@63, SAO@65, SAP@67, SPT@69, SSP@71
\item \textbf{Fw06000000 D0}: (params ver.)@8:U1, (CTD type)@9:U1, MSN@10, MPS@12:U1, MPC@13, MNM0@15, MNM1@17, MNM2@19, MNS0@21:U1, MNS1@22:U1, MNS2@23:U1, DTD0@24, DTD1@26, DTD2@28, MDS@30, APT@32:U1, BTM@33:U1, DDB@34:U1, XAS@35:U1, CPS@36:U1, CST@37:U1, CTF@38, CTL@40, CTU@42, CTC@44:U1, CSP@45:U4, CPB1@49, CPB2@51, CPB3@53, CPN1@55, CPN2@57, CPN3@59, CPP1@61, CPP2@63, CPP3@65, CRD@67, CRS@69, MTD@71, MTT@73:U4, MTS@77:U4, MAM@81:U1, MAT1@82, MAT2@84, MAP@86, ACA@88, ACG@90, ACR@92, PNB@94, PNC@96, PXC@98, PBL@100, PBH@102, PP0@104, PP1@106, PVL@108, PVH@110, PV0@112, PV5@114, PC0@116, PC1@118, ICM@120:U1, IMM@121, ICA0@123, ICA1@125, ICB0@127, ICB1@129, IMT0@131, IMT1@133, ILT0@135, ILT1@137, ICN@139, ICX@141, ICT@143:I2, ICS@145, ISS0@147:U1, ISS1@148:U1, ISM@149:U1, IST@150, ISO@152, ISP@154:U1, IUE0@155:U1, IUE1@156:U1, ICD0@157, ICD1@159, ICR0@161:U4, ICR1@165:U4, ICG0@169:U1, ICG1@170:U1, ISA@171:U1, ISD@172, IAT@174, ISV@176, ISW@178, ICO0@180, ICO1@182, ICF@184:U4, ICW@188:U4, DNS@192, DRA@194, DSD0@196, DSD1@198, DSD2@200, DST0@202, DST1@204, DST2@206, DDF@208, DDC@210, DSI@212:U4, DDD@216, DDT@218, DFT0@220, DFT1@222, DFT2@224, DCM@226:U1, DAL@227:U1, DCT0@228, DCT1@230, DCT2@232, DCA@234:U1, DRC@235, DRT@237, DCX@239, DCP@241, DCF@243, DCD@245, DCR@247, DSW@249, DUC@251, DPT@253, DUT@255, DPD@257, DVT@259, DVD@261, DVW@263, PFR@265, PFC@267, DDP@269, DDS@271, PNS@273, PTO0@275:U4, PTO1@279:U4, PTO2@283:U4, PTH0@287:U4, PTH1@291:U4, PTH2@295:U4, PTD@299:U4, PSA@303:U1, PSI@304:U4, PFM0@308:U1, PFM1@309:U1, PFM2@310:U1, PFT0@311, PFT1@313, PFT2@315, PFD0@317, PFD1@319, PFD2@321, PFV0@323, PFV1@325
\item \textbf{Fw06000000 D1}: PFV2@3, PDT@5, PSM@7:U1, PSD@8, PST@10, PSV@12, PSW@14, ASN@16, ARR@18, APM@20, APS@22, APN@24, APF@26, ASI@28:U4, ASD@32, ADD@34, ADT@36, ARC@38, AHR@40, AHD@42, SFN@44, SDC@46:I2, SGT@48, SET@50, SBM@52:U1, SBT@53, SDT@55:U4, SSA@59, SSI@61:U4, SDS@65, SAM@67:U1, SAX@68, SAN@70, SAO@72, SAP@74, SPT@76, SSP@78
\item \textbf{Fw06010300 D0}: (params ver.)@8:U1, (CTD type)@9:U1, MSN@10, MPS@12:U1, MPC@13, MNM0@15, MNM1@17, MNM2@19, MNS0@21:U1, MNS1@22:U1, MNS2@23:U1, DTD0@24, DTD1@26, DTD2@28, MDS@30, APT@32:U1, BTM@33:U1, DDB@34:U1, XAS@35:U1, CPS@36:U1, CST@37:U1, CTF@38, CTL@40, CTU@42, CTC@44:U1, CSP@45:U4, CPB1@49, CPB2@51, CPB3@53, CPN1@55, CPN2@57, CPN3@59, CPP1@61, CPP2@63, CPP3@65, CRD@67, CRS@69, MTD@71, MTT@73:U4, MTS@77:U4, MAM@81:U1, MAT1@82, MAT2@84, MAP@86, ACA@88, ACG@90, ACR@92, PNB@94, PNC@96, PXC@98, PBL@100, PBH@102, PP0@104, PP1@106, PVL@108, PVH@110, PV0@112, PV5@114, PC0@116, PC1@118, ICM@120:U1, IMM@121, ICA0@123, ICA1@125, ICB0@127, ICB1@129, IMT0@131, IMT1@133, ILT0@135, ILT1@137, ICN@139, ICX@141, ICT@143:I2, ICS@145, ISS0@147:U1, ISS1@148:U1, ISM@149:U1, IST@150, ISO@152, ISP@154:U1, IUE0@155:U1, IUE1@156:U1, ICD0@157, ICD1@159, ICR0@161:U4, ICR1@165:U4, ICG0@169:U1, ICG1@170:U1, ISA@171:U1, ISD@172, IAT@174, ISV@176, ISW@178, ICO0@180, ICO1@182, ICF@184:U4, ICW@188:U4, DNS@192, DRA@194, DSD0@196, DSD1@198, DSD2@200, DST0@202, DST1@204, DST2@206, DDF@208, DDC@210, DSI@212:U4, DDD@216, DDT@218, DFT0@220, DFT1@222, DFT2@224, DCM@226:U1, DAL@227:U1, DCT0@228, DCT1@230, DCT2@232, DCA@234:U1, DRC@235, DRT@237, DCX@239, DCP@241, DCF@243, DCD@245, DCR@247, DSW@249, DUC@251, DPT@253, DUT@255, DPD@257, DVT@259, DVD@261, DVW@263, PFR@265, PFC@267, DDP@269, DDS@271, PNS@273, PTO0@275:U4, PTO1@279:U4, PTO2@283:U4, PTH0@287:U4, PTH1@291:U4, PTH2@295:U4, PTD@299:U4, PSA@303:U1, PSI@304:U4, PFM0@308:U1, PFM1@309:U1, PFM2@310:U1, PFT0@311, PFT1@313, PFT2@315, PFD0@317, PFD1@319, PFD2@321, PFV0@323, PFV1@325
\item \textbf{Fw06010300 D1}: PFV2@3, PDT@5, PSM@7:U1, PSD@8, PST@10, PSV@12, PSW@14, ASN@16, ARR@18, APM@20, APS@22, APN@24, APF@26, ASI@28:U4, ASD@32, ADD@34, ADT@36, ARC@38, AHR@40, AHD@42, SFN@44, SDC@46:I2, SGT@48, SET@50, SBM@52:U1, SBT@53, SDT@55:U4, SSA@59, SSI@61:U4, SDS@65, SAM@67:U1, SAX@68, SAN@70, SAO@72, SAP@74, SPT@76, SSP@78
\item \textbf{Fw06000100 D0}: (params ver.)@8:U1, (CTD type)@9:U1, MSN@10, MPS@12:U1, MPC@13, MNM0@15, MNM1@17, MNM2@19, MNS0@21:U1, MNS1@22:U1, MNS2@23:U1, DTD0@24, DTD1@26, DTD2@28, MDS@30, APT@32:U1, BTM@33:U1, DDB@34:U1, XAS@35:U1, CPS@36:U1, CST@37:U1, CTF@38, CTL@40, CTU@42, CTC@44:U1, CSP@45:U4, CPB1@49, CPB2@51, CPB3@53, CPN1@55, CPN2@57, CPN3@59, CPP1@61, CPP2@63, CPP3@65, CRD@67, CRS@69, MTD@71, MTT@73:U4, MTS@77:U4, MAM@81:U1, MAT1@82, MAT2@84, MAP@86, ACA@88, ACG@90, ACR@92, PNB@94, PNC@96, PXC@98, PBL@100, PBH@102, PP0@104, PP1@106, PVL@108, PVH@110, PV0@112, PV5@114, PC0@116, PC1@118, ICM@120:U1, IMM@121, ICA0@123, ICA1@125, ICB0@127, ICB1@129, IMT0@131, IMT1@133, ILT0@135, ILT1@137, ICN@139, ICX@141, ICT@143:I2, ICS@145, ISS0@147:U1, ISS1@148:U1, ISM@149:U1, IST@150, ISO@152, ISP@154:U1, IUE0@155:U1, IUE1@156:U1, ICD0@157, ICD1@159, ICR0@161:U4, ICR1@165:U4, ICG0@169:U1, ICG1@170:U1, ISA@171:U1, ISD@172, IAT@174, ISV@176, ISW@178, ICO0@180, ICO1@182, ICF@184:U4, ICW@188:U4, DNS@192, DRA@194, DSD0@196, DSD1@198, DSD2@200, DST0@202, DST1@204, DST2@206, DDF@208, DDC@210, DSI@212:U4, DDD@216, DDT@218, DFT0@220, DFT1@222, DFT2@224, DCM@226:U1, DAL@227:U1, DCT0@228, DCT1@230, DCT2@232, DCA@234:U1, DRC@235, DRT@237, DCX@239, DCP@241, DCF@243, DCD@245, DCR@247, DSW@249, DUC@251, DPT@253, DUT@255, DPD@257, DVT@259, DVD@261, DVW@263, PFR@265, PFC@267, DDP@269, DDS@271, PNS@273, PTO0@275:U4, PTO1@279:U4, PTO2@283:U4, PTH0@287:U4, PTH1@291:U4, PTH2@295:U4, PTD@299:U4, PSA@303:U1, PSI@304:U4, PFM0@308:U1, PFM1@309:U1, PFM2@310:U1, PFT0@311, PFT1@313, PFT2@315, PFD0@317, PFD1@319, PFD2@321, PFV0@323, PFV1@325
\item \textbf{Fw06000100 D1}: PFV2@3, PDT@5, PSM@7:U1, PSD@8, PST@10, PSV@12, PSW@14, ASN@16, ARR@18, APM@20, APS@22, APN@24, APF@26, ASI@28:U4, ASD@32, ADD@34, ADT@36, ARC@38, AHR@40, AHD@42, SFN@44, SDC@46:I2, SGT@48, SET@50, SBM@52:U1, SBT@53, SDT@55:U4, SSA@59, SSI@61:U4, SDS@65, SAM@67:U1, SAX@68, SAN@70, SAO@72, SAP@74, SPT@76, SSP@78
\item \textbf{Fw06040900 D0}: (params ver.)@8:U1, (CTD type)@9:U1, MSN@10, MPS@12:U1, MPC@13, MNM0@15, MNM1@17, MNM2@19, MNS0@21:U1, MNS1@22:U1, MNS2@23:U1, DTD0@24, DTD1@26, DTD2@28, MDS@30, APT@32:U1, BTM@33:U1, DDB@34:U1, XAS@35:U1, CPS@36:U1, CST@37:U1, CTF@38, CTL@40, CTU@42, CTC@44:U1, CSP@45:U4, CPB1@49, CPB2@51, CPB3@53, CPN1@55, CPN2@57, CPN3@59, CPP1@61, CPP2@63, CPP3@65, CRD@67, CRS@69, MTD@71, MTT@73:U4, MTS@77:U4, MAM@81:U1, MAT1@82, MAT2@84, MAP@86, ACA@88, ACG@90, ACR@92, PNB@94, PNC@96, PXC@98, PBL@100, PBH@102, PP0@104, PP1@106, PVL@108, PVH@110, PV0@112, PV5@114, PC0@116, PC1@118, ICM@120:U1, IMM@121, ICA0@123, ICA1@125, ICB0@127, ICB1@129, IMT0@131, IMT1@133, ILT0@135, ILT1@137, ICN@139, ICX@141, ICT@143:I2, ICS@145, ISS0@147:U1, ISS1@148:U1, ISM@149:U1, IST@150, ISO@152, ISP@154:U1, IUE0@155:U1, IUE1@156:U1, ICD0@157, ICD1@159, ICR0@161:U4, ICR1@165:U4, ICG0@169:U1, ICG1@170:U1, ISA@171:U1, ISD@172, IAT@174, ISV@176, ISW@178, ICO0@180, ICO1@182, ICF@184:U4, ICW@188:U4, DNS@192, DRA@194, DSD0@196, DSD1@198, DSD2@200, DST0@202, DST1@204, DST2@206, DDF@208, DDC@210, DSI@212:U4, DDD@216, DDT@218, DFT0@220, DFT1@222, DFT2@224, DCM@226:U1, DAL@227:U1, DCT0@228, DCT1@230, DCT2@232, DCA@234:U1, DRC@235, DRT@237, DCX@239, DCP@241, DCF@243, DCD@245, DCR@247, DSW@249, DUC@251, DPT@253, DUT@255, DPD@257, DVT@259, DVD@261, DVW@263, PFR@265, PFC@267, DDP@269, DDS@271, PNS@273, PTO0@275:U4, PTO1@279:U4, PTO2@283:U4, PTH0@287:U4, PTH1@291:U4, PTH2@295:U4, PTD@299:U4, PSA@303:U1, PSI@304:U4, PFM0@308:U1, PFM1@309:U1, PFM2@310:U1, PFT0@311, PFT1@313, PFT2@315, PFD0@317, PFD1@319, PFD2@321, PFV0@323, PFV1@325
\item \textbf{Fw06040900 D1}: PFV2@3, PDT@5, PSM@7:U1, PSD@8, PST@10, PSV@12, PSW@14, ASN@16, ARR@18, APM@20, APS@22, APN@24, APF@26, ASI@28:U4, ASD@32, ADD@34, ADT@36, ARC@38, AHR@40, AHD@42, SFN@44, SDC@46:I2, SGT@48, SET@50, SBM@52:U1, SBT@53, SDT@55:U4, SSA@59:U1, SSI@60, SAM@63:U1, SAX@64, SAN@66, SAO@68, SAP@70, SPT@72, SSP@74
\item \textbf{Fw06050100 D0}: (params ver.)@8:U1, (CTD type)@9:U1, MSN@10, MPS@12:U1, MPC@13, MNM0@15, MNM1@17, MNM2@19, MNS0@21:U1, MNS1@22:U1, MNS2@23:U1, DTD0@24, DTD1@26, DTD2@28, MDS@30, APT@32:U1, BTM@33:U1, DDB@34:U1, XAS@35:U1, CPS@36:U1, CST@37:U1, CTF@38, CTL@40, CTU@42, CTC@44:U1, CSP@45:U4, CPB1@49, CPB2@51, CPB3@53, CPN1@55, CPN2@57, CPN3@59, CPP1@61, CPP2@63, CPP3@65, CRD@67, CRS@69, MTD@71, MTT@73:U4, MTS@77:U4, MAM@81:U1, MAT1@82, MAT2@84, MAP@86, ACA@88, ACG@90, ACR@92, PNB@94, PNC@96, PXC@98, PBL@100, PBH@102, PP0@104, PP1@106, PVL@108, PVH@110, PV0@112, PV5@114, PC0@116, PC1@118, ICM@120:U1, IMM@121, ICA0@123, ICA1@125, ICB0@127, ICB1@129, IMT0@131, IMT1@133, ILT0@135, ILT1@137, ICN@139, ICX@141, ICT@143:I2, ICS@145, ISS0@147:U1, ISS1@148:U1, ISM@149:U1, IST@150, ISO@152, ISP@154:U1, IUE0@155:U1, IUE1@156:U1, ICD0@157, ICD1@159, ICR0@161:U4, ICR1@165:U4, ICG0@169:U1, ICG1@170:U1, ISA@171:U1, ISD@172, IAT@174, ISV@176, ISW@178, ICO0@180, ICO1@182, ICF@184:U4, ICW@188:U4, DNS@192, DRA@194, DSD0@196, DSD1@198, DSD2@200, DST0@202, DST1@204, DST2@206, DDF@208, DDC@210, DSI@212:U4, DDD@216, DDT@218, DFT0@220, DFT1@222, DFT2@224, DCM@226:U1, DAL@227:U1, DCT0@228, DCT1@230, DCT2@232, DCA@234:U1, DRC@235, DRT@237, DCX@239, DCP@241, DCF@243, DCD@245, DCR@247, DSW@249, DUC@251, DPT@253, DUT@255, DPD@257, DVT@259, DVD@261, DVW@263, PFR@265, PFC@267, DDP@269, DDS@271, PNS@273, PTO0@275:U4, PTO1@279:U4, PTO2@283:U4, PTH0@287:U4, PTH1@291:U4, PTH2@295:U4, PTD@299:U4, PSA@303:U1, PSI@304:U4, PFM0@308:U1, PFM1@309:U1, PFM2@310:U1, PFT0@311, PFT1@313, PFT2@315, PFD0@317, PFD1@319, PFD2@321, PFV0@323, PFV1@325
\item \textbf{Fw06050100 D1}: PFV2@3, PDT@5, PSM@7:U1, PSD@8, PST@10, PSV@12, PSW@14, ASN@16, ARR@18, APM@20, APS@22, APN@24, APF@26, ASI@28:U4, ASD@32, ADD@34, ADT@36, ARC@38, AHR@40, AHD@42, SFN@44, SDC@46:I2, SGT@48, SET@50, SBM@52:U1, SBT@53, SSA@55, SSI@57:U4, SDS@61, SAM@63:U1, SAX@64, SAN@66, SAO@68, SAP@70, SPT@72, SSP@74
\item \textbf{Fw06057700 D0}: (params ver.)@8:U1, (CTD type)@9:U1, MSN@10, MPS@12:U1, MPC@13, MNM0@15, MNM1@17, MNM2@19, MNS0@21:U1, MNS1@22:U1, MNS2@23:U1, DTD0@24, DTD1@26, DTD2@28, MDS@30, APT@32:U1, BTM@33:U1, DDB@34:U1, XAS@35:U1, CPS@36:U1, CST@37:U1, CTF@38, CTL@40, CTU@42, CTC@44:U1, CSP@45:U4, CPB1@49, CPB2@51, CPB3@53, CPN1@55, CPN2@57, CPN3@59, CPP1@61, CPP2@63, CPP3@65, CRD@67, CRS@69, MTD@71, MTT@73:U4, MTS@77:U4, MAM@81:U1, MAT1@82, MAT2@84, MAP@86, ACA@88, ACG@90, ACR@92, PNB@94, PNC@96, PXC@98, PBL@100, PBH@102, PP0@104, PP1@106, PVL@108, PVH@110, PV0@112, PV5@114, PC0@116, PC1@118, ICM@120:U1, IMM@121, ICA0@123, ICA1@125, ICB0@127, ICB1@129, IMT0@131, IMT1@133, ILT0@135, ILT1@137, ICN@139, ICX@141, ICT@143:I2, ICS@145, ISS0@147:U1, ISS1@148:U1, ISM@149:U1, IST@150, ISO@152, ISP@154:U1, IUE0@155:U1, IUE1@156:U1, ICD0@157, ICD1@159, ICR0@161:U4, ICR1@165:U4, ICG0@169:U1, ICG1@170:U1, ISA@171:U1, ISD@172, IAT@174, ISV@176, ISW@178, ICO0@180, ICO1@182, ICF@184:U4, ICW@188:U4, DNS@192, DRA@194, DSD0@196, DSD1@198, DSD2@200, DST0@202, DST1@204, DST2@206, DDF@208, DDC@210, DSI@212:U4, DDD@216, DDT@218, DFT0@220, DFT1@222, DFT2@224, DCM@226:U1, DAL@227:U1, DCT0@228, DCT1@230, DCT2@232, DCA@234:U1, DRC@235, DRT@237, DCX@239, DCP@241, DCF@243, DCD@245, DCR@247, DSW@249, DUC@251, DPT@253, DUT@255, DPD@257, DVT@259, DVD@261, DVW@263, PFR@265, PFC@267, DDP@269, DDS@271, PNS@273, PTO0@275:U4, PTO1@279:U4, PTO2@283:U4, PTH0@287:U4, PTH1@291:U4, PTH2@295:U4, PTD@299:U4, PSA@303:U1, PSI@304:U4, PFM0@308:U1, PFM1@309:U1, PFM2@310:U1, PFT0@311, PFT1@313, PFT2@315, PFD0@317, PFD1@319, PFD2@321, PFV0@323
\item \textbf{Fw06057700 D1}: PFV1@3, PFV2@5, PDT@7, PSM@9:U1, PSD@10, PST@12, PSV@14, PSW@16, ASN@18, ARR@20, APM@22, APS@24, APN@26, APF@28, ASI@30:U4, ASD@34, ADD@36, ADT@38, ARC@40, AHR@42, AHD@44, SFN@46, SDC@48:I2, SGT@50, SET@52, SBM@54:U1, SBT@55, SSA@57, SSI@59:U4, SDS@63, SAM@65:U1, SAX@66, SAN@68, SAO@70, SAP@72, SPT@74, SSP@76
\item \textbf{Fw06080500 D0}: (params ver.)@8:U1, (CTD type)@9:U1, MSN@10, MPS@12:U1, MPC@13, MNM0@15, MNM1@17, MNM2@19, MNS0@21:U1, MNS1@22:U1, MNS2@23:U1, DTD0@24, DTD1@26, DTD2@28, MDS@30, APT@32:U1, BTM@33:U1, DDB@34:U1, XAS@35:U1, CPS@36:U1, CST@37:U1, CTF@38, CTL@40, CTU@42, CTC@44:U1, CSP@45:U4, CPB1@49, CPB2@51, CPB3@53, CPN1@55, CPN2@57, CPN3@59, CPP1@61, CPP2@63, CPP3@65, CRD@67, CRS@69, MTD@71, MTT@73:U4, MTS@77:U4, MAM@81:U1, MAT1@82, MAT2@84, MAP@86, ACA@88, ACG@90, ACR@92, PNB@96, PNC@98, PXC@100, PBL@102, PBH@104, PP0@106, PP1@108, PVL@110, PVH@112, PV0@114, PV5@116, PC0@118, PC1@120, ICM@122:U1, IMM@123, ICA0@125, ICA1@127, ICB0@129, ICB1@131, IMT0@133, IMT1@135, ILT0@137, ILT1@139, ICN@141, ICX@143, ICT@145:I2, ICS@147, ISS0@149:U1, ISS1@150:U1, ISM@151:U1, IST@152, ISO@154, ISP@156:U1, IUE0@157:U1, IUE1@158:U1, ICD0@159, ICD1@161, ICR0@163:U4, ICR1@167:U4, ICG0@171:U1, ICG1@172:U1, ISA@173:U1, ISD@174, IAT@176, ISV@178, ISW@180, ICO0@182, ICO1@184, ICF@186:U4, ICW@190:U4, DNS@194, DRA@196, DSD0@198, DSD1@200, DSD2@202, DST0@204, DST1@206, DST2@208, DDF@210, DDC@212, DSI@214:U4, DDD@218, DDT@220, DFT0@222, DFT1@224, DFT2@226, DCM@228:U1, DAL@229:U1, DCT0@230, DCT1@232, DCT2@234, DCA@236:U1, DRC@237, DRT@239, DCX@241, DCP@243, DCF@245, DCD@247, DCR@249, DSW@251, DUC@253, DPT@255, DUT@257, DPD@259, DVT@261, DVD@263, DVW@265, PFR@267, PFC@269, DDP@271, DDS@273, PNS@275, PTO0@277:U4, PTO1@281:U4, PTO2@285:U4, PTH0@289:U4, PTH1@293:U4, PTH2@297:U4, PTD@301:U4, PSA@305:U1, PSI@306:U4, PFM0@310:U1, PFM1@311:U1, PFM2@312:U1, PFT0@313, PFT1@315, PFT2@317, PFD0@319, PFD1@321, PFD2@323, PFV0@325
\item \textbf{Fw06080500 D1}: PFV1@3, PFV2@5, PDT@7, PSM@9:U1, PSD@10, PST@12, PSV@14, PSW@16, ASN@18, ARR@20, APM@22, APS@24, APN@26, APF@28, ASI@30:U4, ASD@34, ADD@36, ADT@38, ARC@40, AHR@42, AHD@44, SFN@46, SDC@48:I2, SGT@50, SET@52, SBM@54:U1, SBT@55, SSA@57, SSI@59:U4, SDS@63, SAM@65:U1, SAX@66, SAN@68, SAO@70, SAP@72, SPT@74, SSP@76
\item \textbf{Fw07020200 D0}: (params ver.)@8:U1, (CTD type)@9:U1, MSN@10, MPS@12:U1, MPC@13, MNM0@15, MNM1@17, MNM2@19, MNS0@21:U1, MNS1@22:U1, MNS2@23:U1, DTD0@24, DTD1@26, DTD2@28, MDS@30, APT@32:U1, BTM@33:U1, DDB@34:U1, XAS@35:U1, CPS@36:U1, CST@37:U1, CTF@38, CTL@40, CTU@42, CTC@44:U1, CSP@45:U4, CPB1@49, CPB2@51, CPB3@53, CPN1@55, CPN2@57, CPN3@59, CPP1@61, CPP2@63, CPP3@65, CRD@67, CRS@69, MTD@71, MTT@73:U4, MTS@77:U4, MAM@81:U1, MAT1@82, MAT2@84, MAP@86, ACA@88, ACG@90, ACR@92, PNB@94, PNC@96, PXC@98, PBL@100, PBH@102, PP0@104, PP1@106, PVL@108, PVH@110, PV0@112, PV5@114, PC0@116, PC1@118, ICM@120:U1, IMM@121, ICA0@123, ICA1@125, ICB0@127, ICB1@129, IMT0@131, IMT1@133, ILT0@135, ILT1@137, ICN@139, ICX@141, ICT@143:I2, ICS@145, ISS0@147:U1, ISS1@148:U1, ISM@149:U1, IST@150, ISO@152, ISP@154:U1, IUE0@155:U1, IUE1@156:U1, ICD0@157, ICD1@159, ICR0@161:U4, ICR1@165:U4, ICG0@169:U1, ICG1@170:U1, ISA@171:U1, ISD@172, IAT@174, ISV@176, ISW@178, ICO0@180, ICO1@182, ICF@184:U4, ICW@188:U4, DNS@192, DRA@194, DSD0@196, DSD1@198, DSD2@200, DST0@202, DST1@204, DST2@206, DDF@208, DDC@210, DSI@212:U4, DDD@216, DDT@218, DFT0@220, DFT1@222, DFT2@224, DCM@226:U1, DAL@227:U1, DCT0@228, DCT1@230, DCT2@232, DCA@234:U1, DRC@235, DRT@237, DCX@239, DCP@241, DCF@243, DCD@245, DCR@247, DSW@249, DUC@251, DPT@253, DUT@255, DPD@257, DVT@259, DVD@261, DVW@263, PFR@265, PFC@267, DDP@269, DDS@271, PNS@273, PTO0@275:U4, PTO1@279:U4, PTO2@283:U4, PTH0@287:U4, PTH1@291:U4, PTH2@295:U4, PTD@299:U4, PSA@303:U1, PSI@304:U4, PFM0@308:U1, PFM1@309:U1, PFM2@310:U1, PFT0@311, PFT1@313, PFT2@315, PFD0@317, PFD1@319, PFD2@321, PFV0@323, PFV1@325
\item \textbf{Fw07020200 D1}: PFV2@3, PDT@5, PSM@7:U1, PSD@8, PST@10, PSV@12, PSW@14, PSS@16:U1, ASN@17, ARR@19, APM@21, APS@23, APN@25, APF@27, ASI@29:U4, ASD@33, ADD@35, ADT@37, ARC@39, AHR@41, AHD@43, SFN@45, SDC@47:I2, SGT@49, SET@51, SBM@53:U1, SBT@54, SSA@56, SSI@58:U4, SDS@62, SAM@64:U1, SAX@65, SAN@67, SAO@69, SPT@71, SSP@73
\item \textbf{Fw08000000 D0}: (params ver.)@8:U1, (CTD type)@9:U1, MSN@10, MPS@12:U1, MPC@13, MNM0@15, MNM1@17, MNM2@19, MNS0@21:U1, MNS1@22:U1, MNS2@23:U1, DTD0@24, DTD1@26, DTD2@28, MDS@30, APT@32:U1, BTM@33:U1, DDB@34:U1, XAS@35:U1, CPS@36:U1, CST@37:U1, CTF@38, CTL@40, CTU@42, CTC@44:U1, CSP@45:U4, CPB1@49, CPB2@51, CPB3@53, CPN1@55, CPN2@57, CPN3@59, CPP1@61, CPP2@63, CPP3@65, CRD@67, CRS@69, MTD@71, MTT@73:U4, MTS@77:U4, MAM@81:U1, MAT1@82, MAT2@84, MAP@86, ACA@88, ACG@90, ACR@92, PNB@94, PNC@96, PXC@98, PBL@100, PBH@102, PP0@104, PP1@106, PVL@108, PVH@110, PV0@112, PV5@114, PC0@116, PC1@118, ICM@120:U1, IMM@121, ICA0@123, ICA1@125, ICB0@127, ICB1@129, IMT0@131, IMT1@133, ILT0@135, ILT1@137, ICN@139, ICX@141, ICT@143:I2, ICS@145, ISS0@147:U1, ISS1@148:U1, ISM@149:U1, IST@150, ISO@152, ISP@154:U1, IUE0@155:U1, IUE1@156:U1, ICD0@157, ICD1@159, ICR0@161:U4, ICR1@165:U4, ICG0@169:U1, ICG1@170:U1, ISA@171:U1, ISD@172, IAT@174, ISV@176, ISW@178, ICO0@180, ICO1@182, ICF@184:U4, ICW@188:U4, DNS@192, DRA@194, DSD0@196, DSD1@198, DSD2@200, DST0@202, DST1@204, DST2@206, DDF@208, DDC@210, DSI@212:U4, DDD@216, DDT@218, DFT0@220, DFT1@222, DFT2@224, DCM@226:U1, DAL@227:U1, DCT0@228, DCT1@230, DCT2@232, DCA@234:U1, DRC@235, DRT@237, DCX@239, DCP@241, DCF@243, DCD@245, DCR@247, DSW@249, DUC@251, DPT@253, DUT@255, DPD@257, DVT@259, DVD@261, DVW@263, PFR@265, PFC@267, DDP@269, DDS@271, PNS@273, PTO0@275:U4, PTO1@279:U4, PTO2@283:U4, PTH0@287:U4, PTH1@291:U4, PTH2@295:U4, PTD@299:U4, PSA@303:U1, PSI@304:U4, PFM0@308:U1, PFM1@309:U1, PFM2@310:U1, PFT0@311, PFT1@313, PFT2@315, PFD0@317, PFD1@319, PFD2@321, PFV0@323, PFV1@325
\item \textbf{Fw08000000 D1}: PFV2@3, PDT@5, PSM@7:U1, PSD@8, PST@10, PSV@12, PSW@14, PSS@16:U1, ASN@17, ARR@19, APM@21, APS@23, APN@25, APF@27, ASI@29:U4, ASD@33, ADD@35, ADT@37, ARC@39, AHR@41, AHD@43, SFN@45, SDC@47:I2, SGT@49, SET@51, SBM@53:U1, SBT@54, SSA@56, SSI@58:U4, SDS@62, SAM@64:U1, SAX@65, SAN@67, SAO@69, SPT@71, SSP@73, WMT@75
\item \textbf{Fw08010000 D0}: (params ver.)@8:U1, (CTD type)@9:U1, MSN@10, MPS@12:U1, MPC@13, MNM0@15, MNM1@17, MNM2@19, MNS0@21:U1, MNS1@22:U1, MNS2@23:U1, DTD0@24, DTD1@26, DTD2@28, MDS@30, APT@32:U1, BTM@33:U1, DDB@34:U1, XAS@35:U1, CPS@36:U1, CST@37:U1, CTF@38, CTL@40, CTU@42, CTC@44:U1, CSP@45:U4, CPB1@49, CPB2@51, CPB3@53, CPN1@55, CPN2@57, CPN3@59, CPP1@61, CPP2@63, CPP3@65, CRD@67, CRS@69, MTD@71, MTT@73:U4, MTS@77:U4, MAM@81:U1, MAT1@82, MAT2@84, MAP@86, ACA@88, ACG@90, ACR@92, PNB@94, PNC@96, PXC@98, PBL@100, PBH@102, PP0@104, PP1@106, PVL@108, PVH@110, PV0@112, PV5@114, PC0@116, PC1@118, ICM@120:U1, IMM@121, ICA0@123, ICA1@125, ICB0@127, ICB1@129, IMT0@131, IMT1@133, ILT0@135, ILT1@137, ICN@139, ICX@141, ICT@143:I2, ICS@145, ISS0@147:U1, ISS1@148:U1, ISM@149:U1, IST@150, ISO@152, ISP@154:U1, IUE0@155:U1, IUE1@156:U1, ICD0@157, ICD1@159, ICR0@161:U4, ICR1@165:U4, ICG0@169:U1, ICG1@170:U1, ISA@171:U1, ISD@172, IAT@174, ISV@176, ISW@178, ICO0@180, ICO1@182, ICF@184:U4, ICW@188:U4, DNS@192, DRA@194, DSD0@196, DSD1@198, DSD2@200, DST0@202, DST1@204, DST2@206, DDF@208, DDC@210, DSI@212:U4, DDD@216, DDT@218, DFT0@220, DFT1@222, DFT2@224, DCM@226:U1, DAL@227:U1, DCT0@228, DCT1@230, DCT2@232, DCA@234:U1, DRC@235, DRT@237, DCX@239, DCP@241, DCF@243, DCD@245, DCR@247, DSW@249, DUC@251, DPT@253, DUT@255, DPD@257, DVT@259, DVD@261, DVW@263, PFR@265, PFC@267, DDP@269, DDS@271, PNS@273, PTO0@275:U4, PTO1@279:U4, PTO2@283:U4, PTH0@287:U4, PTH1@291:U4, PTH2@295:U4, PTD@299:U4, PSA@303:U1, PSI@304:U4, PFM0@308:U1, PFM1@309:U1, PFM2@310:U1, PFT0@311, PFT1@313, PFT2@315, PFD0@317, PFD1@319, PFD2@321, PFV0@323, PFV1@325
\item \textbf{Fw08010000 D1}: PFV2@3, PDT@5, PSM@7:U1, PSD@8, PST@10, PSV@12, PSW@14, PSS@16:U1, ASN@17, ARR@19, APM@21, APS@23, APN@25, APF@27, ASI@29:U4, ASD@33, ADD@35, ADT@37, ARC@39, AHR@41, AHD@43, SFN@45, SDC@47:I2, SGT@49, SET@51, SBM@53:U1, SBT@54, SSA@56, SSI@58:U4, SDS@62, SAM@64:U1, SAX@65, SAN@67, SAO@69, SPT@71, SSP@73
\item \textbf{Fw0A050000 D0}: (params ver.)@8:U1, (CTD type)@9:U1, MSN@10, MPS@12:U1, MPC@13, MNM0@15, MNM1@17, MNM2@19, MNS0@21:U1, MNS1@22:U1, MNS2@23:U1, DTD0@24, DTD1@26, DTD2@28, MDS@30, APT@32:U1, BTM@33:U1, DDB@34:U1, XAS@35:U1, CPS@36:U1, CST@37:U1, CTF@38, CTL@40, CTU@42, CTC@44:U1, CSP@45:U4, CPB1@49, CPB2@51, CPB3@53, CPN1@55, CPN2@57, CPN3@59, CPP1@61, CPP2@63, CPP3@65, CRD@67, CRS@69, MTD@71, MTT@73:U4, MTS@77:U4, MAM@81:U1, MAT1@82, MAT2@84, MAP@86, ACA@88, ACG@90, ACR@92, PNB@94, PNC@96, PXC@98, PBL@100, PBH@102, PP0@104, PP1@106, PVL@108, PVH@110, PV0@112, PV5@114, PC0@116, PC1@118, ICM@120:U1, IMM@121, ICA0@123, ICA1@125, ICB0@127, ICB1@129, IMT0@131, IMT1@133, ILT0@135, ILT1@137, ICN@139, ICX@141, ICT@143:I2, ICS@145, ISS0@147:U1, ISS1@148:U1, ISM@149:U1, IST@150, ISO@152, ISP@154:U1, IUE0@155:U1, IUE1@156:U1, ICD0@157, ICD1@159, ICR0@161:U4, ICR1@165:U4, ICG0@169:U1, ICG1@170:U1, ISA@171:U1, ISD@172, IAT@174, ISV@176, ISW@178, ICO0@180, ICO1@182, ICF@184:U4, ICW@188:U4, DNS@192, DRA@194, DSD0@196, DSD1@198, DSD2@200, DST0@202, DST1@204, DST2@206, DDF@208, DDC@210, DSI@212:U4, DDD@216, DDT@218, DFT0@220, DFT1@222, DFT2@224, DCM@226:U1, DAL@227:U1, DCT0@228, DCT1@230, DCT2@232, DCA@234:U1, DRC@235, DRT@237, DCX@239, DCP@241, DCF@243, DCD@245, DCR@247, DSW@249, DUC@251, DPT@253, DUT@255, DPD@257, DVT@259, DVD@261, DVW@263, PFR@265, PFC@267, DDP@269, DDS@271, PNS@273, PTO0@275:U4, PTO1@279:U4, PTO2@283:U4, PTH0@287:U4, PTH1@291:U4, PTH2@295:U4, PTD@299:U4, PSA@303:U1, PSI@304:U4, PFM0@308:U1, PFM1@309:U1, PFM2@310:U1, PFT0@311, PFT1@313, PFT2@315, PFD0@317, PFD1@319, PFD2@321, PFV0@323, PFV1@325
\item \textbf{Fw0A050000 D1}: PFV2@3, PDT@5, PSM@7:U1, PSD@8, PST@10, PSV@12, PSW@14, PSS@16:U1, ASN@17, ARR@19, APM@21, APS@23, APN@25, APF@27, ASI@29:U4, ASD@33, ADD@35, ADT@37, ARC@39, AHR@41, AHD@43, SFN@45, SDC@47:I2, SGT@49, SET@51, SBM@53:U1, SBT@54, SSA@56, SSI@58:U4, SDS@62, SAM@64:U1, SAX@65, SAN@67, SAO@69, SPT@71, SSP@73, PPC0@75:R4, PPC1@79:R4, PPC2@83:R4, PPC3@87:R4, PPE@91:U1
\item \textbf{Fw0A060231 D0}: (params ver.)@8:U1, (CTD type)@9:U1, MSN@10, MPS@12:U1, MPC@13, MNM0@15, MNM1@17, MNM2@19, MNS0@21:U1, MNS1@22:U1, MNS2@23:U1, DTD0@24, DTD1@26, DTD2@28, MDS@30, APT@32:U1, BTM@33:U1, DDB@34:U1, XAS@35:U1, CPS@36:U1, CST@37:U1, CTF@38, CTL@40, CTU@42, CTC@44:U1, CSP@45:U4, CPB1@49, CPB2@51, CPB3@53, CPN1@55, CPN2@57, CPN3@59, CPP1@61, CPP2@63, CPP3@65, CRD@67, CRS@69, MTD@71, MTT@73:U4, MTS@77:U4, MAM@81:U1, MAT1@82, MAT2@84, MAP@86, ACA@88, ACG@90, ACR@92, PNB@94, PNC@96, PXC@98, PBL@100, PBH@102, PP0@104, PP1@106, PVL@108, PVH@110, PV0@112, PV5@114, PC0@116, PC1@118, ICM@120:U1, IMM@121, ICA0@123, ICA1@125, ICB0@127, ICB1@129, IMT0@131, IMT1@133, ILT0@135, ILT1@137, ICN@139, ICX@141, ICT@143:I2, ICS@145, ISS0@147:U1, ISS1@148:U1, ISM@149:U1, IST@150, ISO@152, ISP@154:U1, IUE0@155:U1, IUE1@156:U1, ICD0@157, ICD1@159, ICR0@161:U4, ICR1@165:U4, ICG0@169:U1, ICG1@170:U1, ISA@171:U1, ISD@172, IAT@174, ISV@176, ISW@178, ICO0@180, ICO1@182, ICF@184:U4, ICW@188:U4, DNS@192, DRA@194, DSD0@196, DSD1@198, DSD2@200, DST0@202, DST1@204, DST2@206, DDF@208, DDC@210, DSI@212:U4, DDD@216, DDT@218, DFT0@220, DFT1@222, DFT2@224, DCM@226:U1, DAL@227:U1, DCT0@228, DCT1@230, DCT2@232, DCA@234:U1, DRC@235, DRT@237, DCX@239, DCP@241, DCF@243, DCD@245, DCR@247, DSW@249, DUC@251, DPT@253, DUT@255, DPD@257, DVT@259, DVD@261, DVW@263, PFR@265, PFC@267, DDP@269, DDS@271, PNS@273, PTO0@275:U4, PTO1@279:U4, PTO2@283:U4, PTH0@287:U4, PTH1@291:U4, PTH2@295:U4, PTD@299:U4, PSA@303:U1, PSI@304:U4, PFM0@308:U1, PFM1@309:U1, PFM2@310:U1, PFT0@311, PFT1@313, PFT2@315, PFD0@317, PFD1@319, PFD2@321, PFV0@323, PFV1@325
\item \textbf{Fw0A060231 D1}: PFV2@3, PDT@5, PSM@7:U1, PSD@8, PST@10, PSV@12, PSW@14, PSS@16:U1, ASN@17, ARR@19, APM@21, APS@23, APN@25, APF@27, ASI@29:U4, ASD@33, ADD@35, ADT@37, ARC@39, AHR@41, AHD@43, SFN@45, SDC@47:I2, SGT@49, SET@51, SBM@53:U1, SBT@54, SSA@56, SSI@58:U4, SDS@62, SAM@64:U1, SAX@65, SAN@67, SAO@69, SPT@71, SSP@73, PPC0@75:R4, PPC1@79:R4, PPC2@83:R4, PPC3@87:R4, PPE@91:U1, WMT@92
\item \textbf{Fw0A060200 D0}: (params ver.)@8:U1, (CTD type)@9:U1, MSN@10, MPS@12:U1, MPC@13, MNM0@15, MNM1@17, MNM2@19, MNS0@21:U1, MNS1@22:U1, MNS2@23:U1, DTD0@24, DTD1@26, DTD2@28, MDS@30, APT@32:U1, BTM@33:U1, DDB@34:U1, XAS@35:U1, CPS@36:U1, CST@37:U1, CTF@38, CTL@40, CTU@42, CTC@44:U1, CSP@45:U4, CPB1@49, CPB2@51, CPB3@53, CPN1@55, CPN2@57, CPN3@59, CPP1@61, CPP2@63, CPP3@65, CRD@67, CRS@69, MTD@71, MTT@73:U4, MTS@77:U4, MAM@81:U1, MAT1@82, MAT2@84, MAP@86, ACA@88, ACG@90, ACR@92, PNB@94, PNC@96, PXC@98, PBL@100, PBH@102, PP0@104, PP1@106, PVL@108, PVH@110, PV0@112, PV5@114, PC0@116, PC1@118, ICM@120:U1, IMM@121, ICA0@123, ICA1@125, ICB0@127, ICB1@129, IMT0@131, IMT1@133, ILT0@135, ILT1@137, ICN@139, ICX@141, ICT@143:I2, ICS@145, ISS0@147:U1, ISS1@148:U1, ISM@149:U1, IST@150, ISO@152, ISP@154:U1, IUE0@155:U1, IUE1@156:U1, ICD0@157, ICD1@159, ICR0@161:U4, ICR1@165:U4, ICG0@169:U1, ICG1@170:U1, ISA@171:U1, ISD@172, IAT@174, ISV@176, ISW@178, ICO0@180, ICO1@182, ICF@184:U4, ICW@188:U4, DNS@192, DRA@194, DSD0@196, DSD1@198, DSD2@200, DST0@202, DST1@204, DST2@206, DDF@208, DDC@210, DSI@212:U4, DDD@216, DDT@218, DFT0@220, DFT1@222, DFT2@224, DCM@226:U1, DAL@227:U1, DCT0@228, DCT1@230, DCT2@232, DCA@234:U1, DRC@235, DRT@237, DCX@239, DCP@241, DCF@243, DCD@245, DCR@247, DSW@249, DUC@251, DPT@253, DUT@255, DPD@257, DVT@259, DVD@261, DVW@263, PFR@265, PFC@267, DDP@269, DDS@271, PNS@273, PTO0@275:U4, PTO1@279:U4, PTO2@283:U4, PTH0@287:U4, PTH1@291:U4, PTH2@295:U4, PTD@299:U4, PSA@303:U1, PSI@304:U4, PFM0@308:U1, PFM1@309:U1, PFM2@310:U1, PFT0@311, PFT1@313, PFT2@315, PFD0@317, PFD1@319, PFD2@321, PFV0@323, PFV1@325
\item \textbf{Fw0A060200 D1}: PFV2@3, PDT@5, PSM@7:U1, PSD@8, PST@10, PSV@12, PSW@14, PSS@16:U1, ASN@17, ARR@19, APM@21, APS@23, APN@25, APF@27, ASI@29:U4, ASD@33, ADD@35, ADT@37, ARC@39, AHR@41, AHD@43, SFN@45, SDC@47:I2, SGT@49, SET@51, SBM@53:U1, SBT@54, SSA@56, SSI@58:U4, SDS@62, SAM@64:U1, SAX@65, SAN@67, SAO@69, SPT@71, SSP@73, PPC0@75:R4, PPC1@79:R4, PPC2@83:R4, PPC3@87:R4, PPE@91:U1
\item \textbf{Fw0A070101 D0}: (params ver.)@8:U1, (CTD type)@9:U1, MSN@10, MPS@12:U1, MPC@13, MNM0@15, MNM1@17, MNM2@19, MNS0@21:U1, MNS1@22:U1, MNS2@23:U1, DTD0@24, DTD1@26, DTD2@28, MDS@30, APT@32:U1, BTM@33:U1, DDB@34:U1, XAS@35:U1, CPS@36:U1, CST@37:U1, CTF@38, CTL@40, CTU@42, CTC@44:U1, CSP@45:U4, CPB1@49, CPB2@51, CPB3@53, CPN1@55, CPN2@57, CPN3@59, CPP1@61, CPP2@63, CPP3@65, CRD@67, CRS@69, MTD@71, MTT@73:U4, MTS@77:U4, MAM@81:U1, MAT1@82, MAT2@84, MAP@86, ACA@88, ACG@90, ACR@92, PBL@94, PBH@96, PP0@98, PP1@100, PVL@102, PVH@104, PV0@106, PV5@108, PC0@110, PC1@112, ICM@114:U1, IMM@115, ICA0@117, ICA1@119, ICB0@121, ICB1@123, IMT0@125, IMT1@127, ILT0@129, ILT1@131, ICN@133, ICX@135, ICT@137:I2, ICS@139, ISS0@141:U1, ISS1@142:U1, ISM@143:U1, IST@144, ISO@146, ISP@148:U1, IUE0@149:U1, IUE1@150:U1, ICD0@151, ICD1@153, ICR0@155:U4, ICR1@159:U4, ICG0@163:U1, ICG1@164:U1, ISA@165:U1, ISD@166, IAT@168, ISV@170, ISW@172, ICO0@174, ICO1@176, ICF@178:U4, ICW@182:U4, DNS@186, DRA@188, DSD0@190, DSD1@192, DSD2@194, DST0@196, DST1@198, DST2@200, DDF@202, DDC@204, DSI@206:U4, DDD@210, DDT@212, DFT0@214, DFT1@216, DFT2@218, DCM@220:U1, DAL@221:U1, DCT0@222, DCT1@224, DCT2@226, DCA@228:U1, DRC@229, DRT@231, DCX@233, DCP@235, DCF@237, DCD@239, DCR@241, DSW@243, DUC@245, DPT@247, DUT@249, DPD@251, DVT@253, DVD@255, DVW@257, PFR@259, PFC@261, DDP@263, DDS@265, PNS@267, PTO0@269:U4, PTO1@273:U4, PTO2@277:U4, PTH0@281:U4, PTH1@285:U4, PTH2@289:U4, PTD@293:U4, PSA@297:U1, PSI@298:U4, PFM0@302:U1, PFM1@303:U1, PFM2@304:U1, PFT0@305, PFT1@307, PFT2@309, PFD0@311, PFD1@313, PFD2@315, PFV0@317, PFV1@319, PFV2@321, PDT@323, PSM@325:U1
\item \textbf{Fw0A070101 D1}: PSD@3, PST@5, PSV@7, PSW@9, PSS@11:U1, ASN@12, ARR@14, APM@16, APS@18, APN@20, APF@22, ASI@24:U4, ASD@28, ADD@30, ADT@32, ARC@34, AHR@36, AHD@38, SFN@40, SDC@42:I2, SGT@44, SET@46, SBM@48:U1, SBT@49, SSA@51, SSI@53:U4, SDS@57, SAM@59:U1, SAX@60, SAN@62, SAO@64, SPT@66, SSP@68
\item \textbf{Fw0A080004 D0}: (params ver.)@8:U1, (CTD type)@9:U1, MSN@10, MPS@12:U1, MPC@13, MNM0@15, MNM1@17, MNM2@19, MNS0@21:U1, MNS1@22:U1, MNS2@23:U1, DTD0@24, DTD1@26, DTD2@28, MDS@30, APT@32:U1, BTM@33:U1, DDB@34:U1, XAS@35:U1, CPS@36:U1, CST@37:U1, CTF@38, CTL@40, CTU@42, CTC@44:U1, CSP@45:U4, CPB1@49, CPB2@51, CPB3@53, CPN1@55, CPN2@57, CPN3@59, CPP1@61, CPP2@63, CPP3@65, CRD@67, CRS@69, MTD@71, MTT@73:U4, MTS@77:U4, MAM@81:U1, MAT1@82, MAT2@84, MAP@86, ACA@88, ACG@90, ACR@92, PBL@94, PBH@96, PP0@98, PP1@100, PVL@102, PVH@104, PV0@106, PV5@108, PC0@110, PC1@112, ICM@114:U1, IMM@115, ICA0@117, ICA1@119, ICB0@121, ICB1@123, IMT0@125, IMT1@127, ILT0@129, ILT1@131, ICN@133, ICX@135, ICT@137:I2, ICS@139, ISS0@141:U1, ISS1@142:U1, ISM@143:U1, IST@144, ISO@146, ISP@148:U1, IUE0@149:U1, IUE1@150:U1, ICD0@151, ICD1@153, ICR0@155:U4, ICR1@159:U4, ICG0@163:U1, ICG1@164:U1, ISA@165:U1, ISD@166, IAT@168, ISV@170, ISW@172, ICO0@174, ICO1@176, ICF@178:U4, ICW@182:U4, DNS@186, DRA@188, DSD0@190, DSD1@192, DSD2@194, DST0@196, DST1@198, DST2@200, DDF@202, DDC@204, DSI@206:U4, DDD@210, DDT@212, DFT0@214, DFT1@216, DFT2@218, DCM@220:U1, DAL@221:U1, DCT0@222, DCT1@224, DCT2@226, DCA@228:U1, DRC@229, DRT@231, DCX@233, DCP@235, DCF@237, DCD@239, DCR@241, DSW@243, DUC@245, DPT@247, DUT@249, DPD@251, DVT@253, DVD@255, DVW@257, PFR@259, PFC@261, DDP@263, DDS@265, PNS@267, PTO0@269:U4, PTO1@273:U4, PTO2@277:U4, PTH0@281:U4, PTH1@285:U4, PTH2@289:U4, PTD@293:U4, PSA@297:U1, PSI@298:U4, PFM0@302:U1, PFM1@303:U1, PFM2@304:U1, PFT0@305, PFT1@307, PFT2@309, PFD0@311, PFD1@313, PFD2@315, PFV0@317, PFV1@319, PFV2@321, PDT@323, PSM@325:U1
\item \textbf{Fw0A080004 D1}: PSD@3, PST@5, PSV@7, PSW@9, PSS@11:U1, ASN@12, ARR@14, APM@16, APS@18, APN@20, APF@22, ASI@24:U4, ASD@28, ADD@30, ADT@32, ARC@34, AHR@36, AHD@38, SFN@40, SDC@42:I2, SGT@44, SET@46, SBM@48:U1, SBT@49, SSA@51, SSI@53:U4, SDS@57, SAM@59:U1, SAX@60, SAN@62, SAO@64, SPT@66, SSP@68, PPC0@70:R4, PPC1@74:R4, PPC2@78:R4, PPC3@82:R4, PPE@86:U1
\item \textbf{Fw0A080031 D0}: (params ver.)@8:U1, (CTD type)@9:U1, MSN@10, MPS@12:U1, MPC@13, MNM0@15, MNM1@17, MNM2@19, MNS0@21:U1, MNS1@22:U1, MNS2@23:U1, DTD0@24, DTD1@26, DTD2@28, MDS@30, APT@32:U1, BTM@33:U1, DDB@34:U1, XAS@35:U1, CPS@36:U1, CST@37:U1, CTF@38, CTL@40, CTU@42, CTC@44:U1, CSP@45:U4, CPB1@49, CPB2@51, CPB3@53, CPN1@55, CPN2@57, CPN3@59, CPP1@61, CPP2@63, CPP3@65, CRD@67, CRS@69, MTD@71, MTT@73:U4, MTS@77:U4, MAM@81:U1, MAT1@82, MAT2@84, MAP@86, ACA@88, ACG@90, ACR@92, PBL@94, PBH@96, PP0@98, PP1@100, PVL@102, PVH@104, PV0@106, PV5@108, PC0@110, PC1@112, ICM@114:U1, IMM@115, ICA0@117, ICA1@119, ICB0@121, ICB1@123, IMT0@125, IMT1@127, ILT0@129, ILT1@131, ICN@133, ICX@135, ICT@137:I2, ICS@139, ISS0@141:U1, ISS1@142:U1, ISM@143:U1, IST@144, ISO@146, ISP@148:U1, IUE0@149:U1, IUE1@150:U1, ICD0@151, ICD1@153, ICR0@155:U4, ICR1@159:U4, ICG0@163:U1, ICG1@164:U1, ISA@165:U1, ISD@166, IAT@168, ISV@170, ISW@172, ICO0@174, ICO1@176, ICF@178:U4, ICW@182:U4, DNS@186, DRA@188, DSD0@190, DSD1@192, DSD2@194, DST0@196, DST1@198, DST2@200, DDF@202, DDC@204, DSI@206:U4, DDD@210, DDT@212, DFT0@214, DFT1@216, DFT2@218, DCM@220:U1, DAL@221:U1, DCT0@222, DCT1@224, DCT2@226, DCA@228:U1, DRC@229, DRT@231, DCX@233, DCP@235, DCF@237, DCD@239, DCR@241, DSW@243, DUC@245, DPT@247, DUT@249, DPD@251, DVT@253, DVD@255, DVW@257, PFR@259, PFC@261, DDP@263, DDS@265, PNS@267, PTO0@269:U4, PTO1@273:U4, PTO2@277:U4, PTH0@281:U4, PTH1@285:U4, PTH2@289:U4, PTD@293:U4, PSA@297:U1, PSI@298:U4, PFM0@302:U1, PFM1@303:U1, PFM2@304:U1, PFT0@305, PFT1@307, PFT2@309, PFD0@311, PFD1@313, PFD2@315, PFV0@317, PFV1@319, PFV2@321, PDT@323, PSM@325:U1
\item \textbf{Fw0A080031 D1}: PSD@3, PST@5, PSV@7, PSW@9, PSS@11:U1, ASN@12, ARR@14, APM@16, APS@18, APN@20, APF@22, ASI@24:U4, ASD@28, ADD@30, ADT@32, ARC@34, AHR@36, AHD@38, SFN@40, SDC@42:I2, SGT@44, SET@46, SBM@48:U1, SBT@49, SSA@51, SSI@53:U4, SDS@57, SAM@59:U1, SAX@60, SAN@62, SAO@64, SPT@66, SSP@68, WMT@70, PPC0@72:R4, PPC1@76:R4, PPC2@80:R4, PPC3@84:R4, PPE@88:U1
\end{itemize}

\end{document}
